%*******************************************************
% Chapter 1
%*******************************************************
% STRUCTURE : version révisée, 11 sections, 27 avril 2026
% §1.1 Avant l'électricité (avant 1794) — RÉDIGÉ
% §1.2 Le télégraphe (1794–1880) — RÉDIGÉ (v5)
% §1.3 Le calcul mécanisé (1820–1914) — RÉDIGÉ (v4d)
% §1.4 Le système mécanographique (1890–1940) — RÉDIGÉ
% §1.5 Le téléphone (~1876–1895) — RÉDIGÉ (v2, 15 études)
% §1.6 La communication sans fil et la radiodiffusion (~1895–1930) — RÉDIGÉ (v1, wireless)
% §1.7 L'électrification et le substrat commun (~1880–1930) — ouverture rédigée v3
%      stubs : broadcasting de masse, IBM tab, Shannon, Enigma
% §1.8 Le calcul programmable et la révolution électronique (1945–1975) — stub
% §1.9 La convergence sur le réseau (1975–2007) — stub
% §1.10 L'information qui consomme (2007–) — stub
% §1.11 Des cas aux archétypes — réduction typologique
%
% DÉCOUPAGE en fichiers séparés pour contourner la limite d'édition Overleaf.
% Chaque \input{} charge une section complète.
% Le fichier source de vérité reste le repo Git (branche main).
%
% Estimation : ~100 pages
%*******************************************************

% =============================================================================
% PRÉAMBULE : titre du chapitre, méthode, tableaux des dimensions
% =============================================================================
\makeatletter
\@ifundefined{dimmark}{%
  \newcommand{\dimmark}[2]{\marginpar{\footnotesize\color{teal}\textbf{#1}\,#2}}%
}{}
\makeatother

\chapter[Transmettre, stocker, calculer]{Transmettre, stocker, calculer~: l'émergence des techniques de l'information au regard de l'économie et de l'énergie}\label{ch:history}

Ce chapitre retrace comment les techniques de la transmission, du stockage et du calcul se sont inscrites dans les structures de l'économie~--- comment elles ont pesé sur le capital, le travail, l'organisation des marchés ou le financement des infrastructures, entre autres~--- et à quel coût matériel et énergétique elles l'ont fait. Le parcours est chronologique~--- des tablettes d'argile sumériennes au numérique contemporain~--- mais le geste est analytique~: chaque cas historique est observé au prisme de l'économie et du couple énergie-matière, et les questions qu'il soulève émergent du récit lui-même. Ces observations sont rassemblées en fin de chapitre par réduction typologique. À notre connaissance, aucune des traditions intellectuelles mobilisées~--- économie de la croissance, analyse évolutionniste, histoire des systèmes techniques, économie de l'information, économie de l'énergie\footnote{Les cinq traditions et leurs apports respectifs sont développés dans l'introduction générale de la thèse.}~--- n'a constitué le croisement entre structures économiques, infrastructure informationnelle et substrat énergétique en objet d'étude dans la longue durée. Nous pensons pourtant que c'est à leurs croisements que se logent des éléments de structure qu'aucune, prise isolément, ne peut voir.



% =============================================================================
% PRÉAMBULE MÉTHODOLOGIQUE
% =============================================================================

\subsection*{Méthode et protocole}\label{subsec:grammaire}

Le chapitre traverse une dizaine de cas historiques dans l'ordre chronologique. Le mode dominant est l'induction comparative~: chaque cas est observé pour lui-même, les catégories analytiques émergent de la comparaison entre les cas, et les patterns provisoires qu'un cas fait apparaître sont modifiés ou infirmés par le suivant. C'est ce que \textcite{znaniecki_method_1934} appelle l'induction analytique~: non pas la généralisation à partir d'un échantillon, mais la construction progressive de catégories par confrontation séquentielle du concept au cas. Le regroupement final en dimensions économiques n'intervient qu'en fin de parcours, sous forme d'idéaux-types\footnote{\og Il [l'idéal-type] est obtenu par l'accentuation unilatérale d'un ou de plusieurs points de vue et par l'enchaînement d'une multitude de phénomènes individuels, diffus et discrets, présents tantôt en grand nombre, tantôt en petit nombre, par endroits pas du tout, qui s'ordonnent selon ces points de vue unilatéralement choisis, pour former un tableau de pensée homogène\fg{} \parencite[p.~190--191]{weber_objektivitat_1904}.}.

L'induction analytique au sens de Znaniecki n'est pas une protection contre le biais de confirmation~: elle expose au contraire la construction des catégories à la critique en rendant visible le geste sélectif. Le protocole adopté ici en assume la conséquence~: une catégorie n'est conservée que si elle survit au cas suivant~; un cas qui contredit le pattern n'est jamais écarté comme exception mais traité comme révision de la catégorie. Les trois moments abductifs identifiés plus bas sont précisément des points où cette révision a eu lieu.

L'information, dans ses trois modes~--- la transmission, le stockage et le calcul~--- constitue le niveau d'analyse principal~; selon les cas, la technologie est étudiée en elle-même ou dans le système technique plus vaste où elle s'insère et dont elle tire ses effets.

Mais l'induction seule ne produit pas de résultat théorique nouveau~: elle accumule des régularités. Les moments les plus productifs du parcours sont ceux où un cas \emph{contredit} le pattern établi par les précédents. Lorsque sept pays convergent vers le monopole télégraphique, l'induction identifie un pattern~; lorsque la calculatrice mécanique, technologie comparable, produit un marché fragmenté au lieu d'un monopole, le pattern casse. C'est cette rupture qui déclenche un geste inférentiel distinct~: l'abduction, au sens de Peirce\footnote{Charles Sanders Peirce (1839--1914), fondateur du pragmatisme américain. L'abduction, ou \emph{inference to the best explanation}, est formalisée dans ses \emph{Collected Papers} (1931--1958, vol.~5, §~189). Peirce la distingue de l'induction (généralisation à partir de cas) et de la déduction (conséquence d'une règle). Voir \textcite{douven_abduction_2021} pour une présentation contemporaine.}. Un fait surprenant au regard de ce qui précède engendre une hypothèse qui, si elle était vraie, rendrait le fait non surprenant~; l'hypothèse est ensuite mise à l'épreuve dans le cas suivant. \textcite{timmermans_theory_2012} formalisent cette articulation~: l'induction identifie les régularités, l'abduction génère les hypothèses quand les régularités cassent, et la prédiction qui en résulte est testée déductivement dans le cas d'après. Le chapitre suit ce cycle.

Les moments abductifs du parcours sont les suivants~: le contraste entre la calculatrice (marché fragmenté) et la tabulatrice (monopole à 85,7\,\%), qui engendre l'hypothèse que c'est le mode d'appropriation du capital, et non la technologie, qui détermine la structure de marché~; la mécanisation du \emph{Nautical Almanac}, qui coûte plus cher et emploie plus de monde que le calcul humain, ce qui engendre l'hypothèse que la mécanisation réorganise le travail au lieu de le remplacer~; l'écart de quarante-sept ans entre l'invention du \emph{block system} et son adoption obligatoire, qui engendre l'hypothèse que la technique ne se déploie pas sans un véhicule de financement adapté à la structure des coûts et des bénéfices. Ces hypothèses ne sont pas des conclusions~: elles sont les résultats provisoires que le chapitre~2 devra formaliser et que le chapitre~4 devra tester.

Le parcours met également en évidence une direction causale que le plan du chapitre pourrait occulter. Si chaque cas est présenté comme l'irruption d'une technique dans un contexte économique, le matériau montre aussi, de façon récurrente, que la technique est appelée par une demande de coordination que les moyens existants ne satisfont plus~: la coordination territoriale en temps de guerre pour le télégraphe optique, la sécurité ferroviaire sur voie unique pour le \emph{block system}, la gestion d'un recensement devenu ingérable pour la tabulatrice. Cette bidirectionnalité~--- la technique transforme l'économie, mais l'économie appelle la technique~--- n'est pas traitée séparément dans le chapitre~1~; elle deviendra un objet d'analyse à part entière dans le chapitre~2, où les mécanismes de demande induite devront être formalisés au même titre que les mécanismes d'offre technologique\footnote{Cette posture critique de l'innovation-centrisme et l'attention portée aux conditions d'usage et aux configurations institutionnelles de la diffusion partagent avec un courant historiographique récent une orientation commune. \textcite{edgerton_innovation_1998}, dans le numéro \og Histoire des techniques\fg{} dirigé par Y.~Cohen et D.~Pestre pour les \emph{Annales HSS}, en formule la version qui a fait référence en France~; \textcite{winston_media_1998} en propose la modélisation par le couple \og \emph{supervening necessity}\fg{} / \og \emph{suppression of radical potential}\fg{}~; \textcite{gardey_ecrire_2008} en prolonge la version annaliste en étendant l'\og outillage mental\fg{} de \textcite{febvre_outillage_1938} vers la matérialité des inscriptions, des supports et des gestes. Notre démarche s'en distingue par son ancrage CIRED~--- qui suit en parallèle les grandeurs physiques et monétaires~--- et par son objectif final, qui est de produire des mécanismes testables dans un modèle énergie-économie.}.

Comparer des objets aussi éloignés dans le temps et dans la forme suppose une décision qu'il faut expliciter. Ces dispositifs ne partagent ni le substrat matériel, ni le contexte institutionnel, ni même ce que leurs contemporains appelaient \emph{économie}~--- le mot n'a pas le même sens dans la redistribution des palais sumériens, dans le capitalisme industriel victorien et dans les plateformes numériques actuelles. La comparaison ne tient donc pas à une supposée équivalence économique des objets, mais à un geste assumé~: interroger le matériau historique avec des questions que notre époque a les moyens de formuler, en sachant que ces questions sont elles-mêmes datées. C'est ce que Loraux nommait \emph{anachronisme contrôlé}~--- non un défaut à neutraliser, mais une opération productive qui rend visible, par contraste, ce que les régimes antérieurs ne pouvaient pas voir d'eux-mêmes\footnote{\textcite{loraux_elogeanachronisme_1993}~; pour une discussion contemporaine en histoire économique, \textcite{monnet_perspectives_2025}.}. Le contrôle de l'opération tient à deux disciplines~: que l'anachronisme soit nommé comme tel à chaque mobilisation~--- jamais présenté comme évidence intrinsèque du matériau~--- et qu'il soit borné par les sources, c'est-à-dire qu'il ne fasse dire aux acteurs que ce que leurs effets produisaient, jamais ce que leur conscience aurait dû formuler.

Encore faut-il distinguer ce que l'on rétroprojette. Les concepts économiques que mobilise le chapitre~--- externalité de réseau, coûts de transaction, demande dérivée, frontière firme-marché~--- sont des outils du capitalisme industriel et post-industriel~; appliqués à des dispositifs antérieurs, ils ne décrivent pas ce que leurs acteurs faisaient mais ce que leurs effets \emph{produisaient} au sens que notre vocabulaire permet de formuler. La tripartition qui structure ce chapitre~--- \emph{transmettre, stocker, calculer}~--- est dans une situation analogue, mais d'une autre profondeur~: elle ne nous vient pas d'une économie ancienne dont nous aurions changé le vocabulaire, elle naît avec la cybernétique, à la fin de la période même que ce chapitre traite. Babbage en~1834 sépare le \emph{store} et le \emph{mill} dans l'\emph{Analytical Engine}~; \textcite{vonneumann_first_1945} formalise la séparation entre mémoire, unité arithmétique et entrées-sorties~; \textcite{wiener_cybernetics_1948} définit le contrôle comme communication augmentée d'une rétroaction, laquelle exige stockage de l'état passé et calcul de la correction~; \textcite{kittler_grammophon_1986} pose la tripartition \emph{Speicherung~/ \"Ubertragung~/ Verarbeitung} comme structure médiologique des techniques modernes~; \textcite{hilbert_worlds_2011} sont les premiers à les quantifier comme trois grandeurs mondiales distinctes, dont les rythmes de croissance entre~1986 et~2007 (23, 28 et 58~\% par an respectivement) suggèrent qu'il s'agit bien de fonctions \emph{séparables}\footnote{La distinction prolonge la partition \emph{communication technology~/ information technology} de \textcite{vanderkooij_invention_2015} en explicitant le stockage comme fonction propre. Elle est par ailleurs \emph{orthogonale} à la hiérarchie donnée~$\rightarrow$~information~$\rightarrow$~connaissance~$\rightarrow$~sagesse de \textcite{ackoff_data_1989}, reprise par \textcite{fouquet_digitalisation_2023} pour ses séries longues énergie-information~: chacun de ces niveaux peut être transmis, stocké ou calculé. Sur le plan mathématique, \textcite{bailly_mathematics_2011} et \textcite{longo_computing_2011} fournissent le fondement de la distinction discret/continu qui sous-tend l'argument que la transition analogique/numérique est une différence de nature, non de degré~; \textcite{spreng_interdependency_2015} théorise par ailleurs la substituabilité entre énergie, temps et information comme trois intrants fondamentaux de la production~--- un cadre complémentaire au nôtre, qui décompose non les intrants mais les fonctions du système informationnel.}. La rétroprojection de cette grille sur le réseau Chappe ou sur la mécanographie hollerithienne n'est donc pas neutre~: c'est une lecture que la cybernétique rend possible, et qui dévoile dans le matériau ancien des structurations dont les acteurs ne disposaient pas pour se penser eux-mêmes.

C'est sur la question de l'énergie-matière, cependant, que la rétroprojection devient elle-même un objet de la thèse. La catégorie d'\emph{énergie} comme grandeur unifiée n'existe pas avant l'unification thermodynamique du milieu du XIX\textsuperscript{e}~siècle (Mayer, Joule, Helmholtz)~; son importation dans la pensée économique sous la forme d'un processus productif converti en équivalent énergétique reste marginale, depuis les tentatives de Podolinsky en~1880 jusqu'aux fondateurs de l'économie écologique au XX\textsuperscript{e}~siècle \parencite{vianna_franco_history_2023}. Avant 1845, ce n'est pas que les acteurs ou les économistes \emph{évacuaient} l'énergie~: la catégorie qui aurait permis de la nommer comme telle n'était pas disponible. Le silence des comptes hanovriens du Calenberg ou des budgets du \emph{Nautical Almanac} sur un poste \emph{énergie} n'est pas un oubli mais une impossibilité conceptuelle, et cette impossibilité est elle-même historique. Ce que la thèse identifiera plus loin comme un changement de régime du couplage énergie-information se laisse, à cette aune, formuler autrement~: le moment où la coordination énergie-information cesse d'appartenir au seul historien qui la lit \emph{après coup}, pour entrer dans le calcul des acteurs eux-mêmes. La transversale \textbf{0-TER} n'est donc pas une grille ahistorique appliquée uniformément à toutes les périodes~: c'est une catégorie dont l'apparition même fait partie de la chronologie que ce chapitre cherche à reconstituer. La seconde transversale, la \emph{gouvernance} (\textbf{0-GOV}), est dans un statut différent~: elle désigne l'ensemble des modes de coordination autoritaire qui prescrivent et contraignent le système informationnel~--- l'État souverain, l'État régulateur, l'organisme international, la normalisation technique, la corégulation public-privé, l'autorégulation industrielle. Ces modes traversent toutes les périodes mais sous des formes historiquement contrastées, et l'un des résultats que le chapitre cherche à établir est qu'ils ne se substituent pas mais se déplacent~: la modalité souveraine-westphalienne pour le télégraphe Chappe, la régulation antitrust pour le téléphone Bell et la mécanographie hollerithienne, l'allocation internationale du spectre pour le wireless, la gouvernance multi-stakeholder pour Internet, le retour de la modalité souveraine à la période contemporaine à travers les politiques de Cloud souverain, l'\emph{AI Act} et le contrôle des semi-conducteurs. Nommer cette transversale \emph{gouvernance} plutôt que \emph{souveraineté} signale d'emblée cette pluralité des modes et permet, dans la suite de la thèse, l'articulation avec les leviers politiques modélisables (taxes, subsides, régulations, standards techniques, commande publique) qui structureront la cartographie empirique du chapitre~3 et les scénarios du chapitre~4\footnote{Les promesses qui mobilisent le financement (capital privé) et la commande étatique ont, plus largement, une dimension performative que Jasanoff (2015) théorise sous le nom de \emph{sociotechnical imaginaries}~--- visions collectivement tenues, institutionnellement stabilisées et publiquement performées de futurs désirables, qui orientent le déploiement des technologies. La promesse de dématérialisation qui accompagne aujourd'hui le numérique, comme la promesse de souveraineté territoriale qui accompagnait le télégraphe en~1794, relèvent de ce registre. Nous ne mobilisons pas ce cadre interprétatif dans le corps de l'analyse, dont la finalité est de produire des mécanismes testables dans un modèle énergie-économie, mais nous reconnaissons sa pertinence pour qui voudrait prolonger la lecture sur le versant des imaginaires.}.

Le découpage chronologique suit les changements de substrat technique~--- du signal optique au signal électrique, du calcul mécanique au calcul électronique, de l'infrastructure filaire au réseau numérique. Ces bornes ne sont pas des ruptures nettes~; elles marquent des basculements dans les conditions matérielles de la transmission, du stockage et du calcul. Chaque section correspond à un régime technique distinct~--- cette correspondance est une discipline de lecture, pas un postulat~: elle s'ajustera au fil de la thèse.

La sélection des cas suit le principe de la traversée représentative des substrats techniques, pas celui de l'exhaustivité géographique~: chaque section retient le ou les cas qui font apparaître le plus nettement le basculement technique en cause, et les contre-exemples connus (Suisse, Belgique pour le télégraphe~; calculatrice mécanique américaine pour la mécanographie) sont mobilisés là où ils contredisent le pattern qu'établissent les cas dominants. La régularité que le chapitre identifie n'est donc jamais affirmée sur l'absence de contre-exemples mais sur leur intégration argumentée.

Chaque cas est abordé sous cinq aspects~: \emph{technique} (l'artefact et sa trajectoire d'amélioration), \emph{socio-économique} (les marchés, les acteurs, la régulation), \emph{conceptuel} (les cadres théoriques existants, mobilisés quand le cas les appelle), \emph{matériel} (les flux physiques, les extractions, l'empreinte énergétique), et \emph{conditions d'émergence}~: qui demande le déploiement de la technologie, qui le finance et pour quelle promesse, quel contexte institutionnel le rend possible ou le bloque, et pourquoi le déploiement prend cette forme ici et pas ailleurs\footnote{La promesse qui mobilise le financement peut être commerciale (la rente informationnelle du câble transatlantique), militaire (la coordination territoriale du réseau Chappe), scientifique (la précision des éphémérides au \emph{Nautical Almanac}), ou politique (la souveraineté nationale dans le cas Morse au Congrès). Sur la dimension performative de ces promesses, voir la note précédente sur \emph{sociotechnical imaginaries}.}. Les cas de non-émergence~--- la \emph{Difference Engine} de Babbage, financée mais sans demande~; le télégraphe français, bloqué pendant douze ans par le lobby du sémaphore~--- sont traités avec la même attention que les cas d'émergence réussie, parce qu'ils révèlent les conditions nécessaires par leur absence.

Les observations qui émergent du récit sont étiquetées selon trois niveaux de certitude~: \textbf{*}~(fait documenté indépendamment par des sources distinctes), \textbf{+}~(mécanisme proposé comme explication, compatible avec le matériau mais non testé causalement), \textbf{?}~(piste ouverte, pas encore d'argument construit). Ces marqueurs accompagnent la traversée dans les marges et structurent les bilans de fin de section. Le style du chapitre~--- phrases longues, connecteurs logiques, prose continue~--- porte un risque que ces marqueurs sont conçus pour contenir~: le risque de présenter comme nécessaire ce qui est contingent, ou comme démontré ce qui est proposé. Un pattern observé dans sept pays n'est pas un mécanisme~; un mécanisme compatible avec les données n'est pas un résultat causalement identifié~; et un résultat causalement identifié ne vaut que pour le cas et la période qu'il couvre.

Ces marqueurs prolongent un débat méthodologique sur les rapports entre histoire et économie que \textcite{trannoy_economistes_2025} ont récemment cartographié. \textcite{grenier_causalite_2025} y distingue le paradigme de Laplace, pour lequel la régression causale est possible et l'on peut remonter des effets aux causes, et celui de Cournot, pour lequel l'histoire est ordonnée mais contingente et la régression causale périlleuse~; l'historiographie post-laplacienne reconnaît l'existence de relations causales mais doute de la possibilité de les isoler sur un cas singulier. \textcite{trannoy_causalite_2025} prolonge cette prudence en notant qu'un raisonnement par cas ne peut réfuter qu'un seul type de causalité à la fois, nécessaire ou suffisante, jamais les deux. \textcite{monnet_perspectives_2025} identifie trois types de causalité que la thèse mêle sans toujours les nommer~: la causalité \emph{structurelle}, qui articule un modèle théorique et un choc exogène, et que la traversée mobilise quand un mécanisme économique générique est convoqué~; la causalité \emph{interventionniste}, qui repose sur des quasi-expériences, et sous laquelle se rangent les travaux de \textcite{steinwender_real_2018} et \textcite{juhasz_spinning_2018} mobilisés dans les sections suivantes~; et la causalité \emph{processuelle}, qui reconstruit la chaîne d'actions et où le récit est lui-même le véhicule de l'explication, registre dominant du présent chapitre. Une seconde précaution s'ajoute aux marqueurs et porte non sur la régularité observée mais sur le mécanisme qui la produirait. \textcite{sewell_logics_2005}, cité par \textcite{monnet_perspectives_2025}, rappelle que les structures causales sous-jacentes subissent elles-mêmes des mutations dans le temps~: lorsque le chapitre identifie un pattern qui se reproduit sur cent soixante ans, il n'en suit pas que le mécanisme générateur soit stable. La stabilité est une hypothèse à tester dans les chapitres formels, pas un postulat de la traversée historique.

Le regroupement typologique des questions soulevées par le récit intervient en synthèse de fin de chapitre (\S\ref{sec:reduction})~; sa formalisation comme grille de questionnement relève du chapitre~suivant.

% =============================================================================
% TABLEAUX DÉSACTIVÉS — déplacés vers le chapitre 2 (formalisation des dimensions).
% Présents ici en \iffalse...\fi pour archive ; le préambule du chapitre 1 ne les
% mobilise plus, par cohérence avec la posture inductive (les catégories émergent
% du récit, le regroupement typologique intervient en synthèse de fin de chapitre).
% =============================================================================
\iffalse
\begin{table}[htbp]
\caption{Les dix canaux économiques du couplage ICT-économie~--- présentés ici pour le contexte de lecture~; le déploiement systématique relève du chapitre~2}\label{tab:dimensions}
\centering\small
\begin{tabularx}{\textwidth}{@{}cp{3.2cm}Xp{3.8cm}@{}}
\toprule
 & \textbf{Dimension} & \textbf{Questions soulevées par le récit historique} & \textbf{Exemples de travaux} \\
\midrule
D1 & Propriétés du bien informationnel & Quelles propriétés intrinsèques~--- rivalité, codifiabilité, excluabilité, périssabilité, entre autres~--- déterminent la portée des effets économiques de l'information~? & Samuelson 1954, Arrow 1962, Romer 1990, Shapiro \& Varian 1999, Jones \& Tonetti 2020 \\[4pt]
D2 & Rendements d'échelle et externalités de réseau & Le réseau de communication tend-il vers le monopole~? Le coût de production de l'information diminue-t-il ou augmente-t-il avec l'ambition du service~? & Katz \& Shapiro 1985, Arthur 1989, Tirole 1988 \\[4pt]
D3 & Nature du capital informationnel & La technique informationnelle remplace-t-elle du travail par du capital~? Quel type de capital~--- intangible, physique, hybride~? & Jorgenson 2001, Haskel \& Westlake 2018, Field 1992 \\[4pt]
D4 & Travail, qualification et division cognitive & L'information qualifie-t-elle ou déqualifie-t-elle la main-d'\oe{}uvre~? Crée-t-elle de nouvelles tâches ou en détruit-elle~? & Autor 2003, Acemoglu \& Restrepo 2019, Daston 2018 \\[4pt]
D5 & Demande, tarification et préférences & La baisse du coût de l'information engendre-t-elle un rebond de consommation~? La numérisation transforme-t-elle les préférences~? & Schmookler 1966, Jevons 1865, Allcott et al.\ 2020 \\[4pt]
D6 & Signal de prix et intégration des marchés & L'information synchrone réduit-elle les frictions de marché ou crée-t-elle de nouveaux marchés~? L'effet dépend-il de la nature de ce qui est transmis, stocké ou calculé~? & Hayek 1945, Steinwender 2018, Akerlof 1970 \\[4pt]
D7 & Coûts de transaction et formes organisationnelles & Quand le système de prix ne coordonne pas, quelles formes d'organisation émergent~? La technique informationnelle déplace-t-elle la frontière entre la firme et le marché~? & Coase 1937, Williamson 1985, Chandler 1977 \\[4pt]
D8 & Financement et cycles d'infrastructure & Comment les réseaux d'information sont-ils financés~? La séquence investissement massif, faillite, recapitalisation est-elle récurrente~? & Perez 2002, Gordon 2002, Nonnenmacher 2001 \\[4pt]
D9 & Géographie de la production et commerce & Qui produit les composants, les matériaux et l'énergie nécessaires~? La localisation des infrastructures informationnelles suit-elle la géographie de l'énergie~? & Krugman 1991, Headrick 1988, Kander et al.\ 2013 \\[4pt]
D10 & Trajectoires d'innovation et dépendance de sentier & Les choix techniques passés contraignent-ils les choix futurs~? L'innovation incrémentale crée-t-elle des irréversibilités~? & Nelson \& Winter 1982, Rosenberg 1982, Dosi 1982 \\
\bottomrule
\end{tabularx}
\end{table}

\begin{table}[htbp]
\caption{Deux forces transversales~--- l'énergie-matière et la souveraineté~--- qui traversent plusieurs des canaux précédents sans être elles-mêmes des canaux économiques}\label{tab:transversales}
\centering\small
\begin{tabularx}{\textwidth}{@{}cp{3.2cm}Xp{3.8cm}@{}}
\toprule
 & \textbf{Transversale} & \textbf{Questions soulevées par le récit historique} & \textbf{Exemples de travaux} \\
\midrule
0-TER & Énergie, matière et empreinte physique & Quels flux de matière et d'énergie le système informationnel mobilise-t-il~? Quelles sont les extractions, les infrastructures, les externalités environnementales qui le rendent possible~? Ces flux croissent-ils avec l'ambition fonctionnelle, ou peuvent-ils en être découplés~? & Wrigley 2010, Smil 2017, Pottier 2014, Ensmenger 2018 \\[4pt]
0-GOV & Souveraineté, commande étatique et usage militaire-stratégique & Par qui le système informationnel est-il commandité, financé, contraint~? Quelle est la part de l'État-régulateur, de l'État-militaire, de la géopolitique des infrastructures dans la trajectoire technique~? L'information sert-elle le marché, le pouvoir, ou les deux inextricablement~? & Hughes 1983, Headrick 1991, Edwards 1996, Mazzucato 2013 \\
\bottomrule
\end{tabularx}
\end{table}

\begin{table}[htbp]
\caption{Travaux de référence en économie de l'ICT pour chacune des dix dimensions~--- auteurs ayant spécifiquement étudié les effets économiques des techniques de l'information}\label{tab:dimensions-ict}
\centering\small
\begin{tabularx}{\textwidth}{@{}cp{2.8cm}Xp{4.2cm}@{}}
\toprule
 & \textbf{Dimension} & \textbf{Résultat clé en économie de l'ICT} & \textbf{Travaux spécifiques ICT} \\
\midrule
D1 & Propriétés du bien informationnel & Le bien numérique a un coût marginal nul~; les données sont non rivales mais leur valeur dépend de l'agrégation & Shapiro \& Varian 1999, Jones \& Tonetti 2020, Goldfarb \& Tucker 2019 \\[4pt]
D2 & Rendements et externalités de réseau & Les réseaux de communication convergent vers le monopole par un cycle ouverture-fermeture récurrent & Katz \& Shapiro 1985, Economides 1996, Wu 2010 \\[4pt]
D3 & Nature du capital & L'investissement intangible dépasse le tangible aux É.-U. vers 2000~; le capital informationnel se substitue au capital physique avec un levier de 1:7 à 1:13 & Corrado et al.\ 2005, Haskel \& Westlake 2018, Field 1992 \\[4pt]
D4 & Travail et division cognitive & L'ICT polarise l'emploi par substitution des tâches routinières~; l'automatisation du \emph{switching} téléphonique a pris 60~ans & Autor et al.\ 2003, Acemoglu \& Restrepo 2019, Feigenbaum \& Gross 2024 \\[4pt]
D5 & Demande et préférences & Le surplus du consommateur numérique est sous-estimé par le PIB~; les biens gratuits masquent le coût de production & Brynjolfsson et al.\ 2019, Allcott et al.\ 2020, Varian 2018 \\[4pt]
D6 & Signal de prix et marchés & Le télégraphe, le téléphone mobile et Internet créent du commerce qui n'aurait pas eu lieu sans eux~--- l'information est constitutive, pas lubrifiant & Steinwender 2018, Jensen 2007, Aker 2010 \\[4pt]
D7 & Coûts de transaction et organisations & L'IT exige des transformations organisationnelles complémentaires~; la baisse des coûts de transaction est asymétrique selon les dimensions & Brynjolfsson \& Hitt 2000, Bloom et al.\ 2012, Goldfarb \& Tucker 2019 \\[4pt]
D8 & Financement et cycles & Les vagues d'infrastructure ICT suivent le cycle irruption-frenzy-crash-synergy~; la séquence se reproduit du câble transatlantique à l'IA & Perez 2002, Nonnenmacher 2001, Janeway 2012 \\[4pt]
D9 & Géographie et commerce & La distance pèse encore dans le commerce numérique~; la diffusion de l'Internet est géographiquement inégale & Blum \& Goldfarb 2006, Forman et al.\ 2012, Headrick 1988 \\[4pt]
D10 & Innovation et dépendance de sentier & L'ICT peut être analysée comme GPT~; le cadre GPT-énergie n'existe pas encore & Bresnahan \& Trajtenberg 1995, Lipsey et al.\ 2005, Nuvolari 2020 \\
\bottomrule
\end{tabularx}
\end{table}
\fi


% =============================================================================

% =============================================================================
% §1.1 Avant l'électricité (avant 1794)
% =============================================================================
\clearpage
\section{Avant l'électricité (avant 1794)}\label{sec:era0}
% =============================================================================

\subsection{L'écriture comptable et les supports de mémoire}\label{subsec:ecriture}

Les techniques de l'information ne commencent pas avec l'électricité. Elles commencent avec l'argile. Dans le sud de la Mésopotamie, entre l'Euphrate et le Tigre, l'État sumérien qui se forme au troisième millénaire avant notre ère développe un système d'écriture cunéiforme dont la fonction première n'est pas littéraire mais comptable~: enregistrer les stocks de grain, les surfaces irriguées, les journées de travail, les rations distribuées \parencite{nissen_archaic_1993}. Comme l'ont établi \textcite{nissen_archaic_1993} et \textcite{englund_hard_1991}, ce système informationnel naît des besoins d'une économie redistributive et donne à l'administration un instrument de contrôle sur des projets hydrauliques de grande échelle~--- canaux, digues, bassins de rétention. \dimmark{D7}{?}%
L'écriture rend l'irrigation gouvernable, et l'irrigation rend l'État viable. \dimmark{0-TER}{?}%
Mais les mêmes canaux, gérés sur plusieurs siècles, provoquent la salinisation des sols du sud mésopotamien \parencite{dickson_circumscription_1987}~: l'un des premiers épisodes documentés de destruction environnementale à l'échelle d'un écosystème est indissociable de la première technologie de l'information.

Parallèlement à cette trajectoire du stockage, le calcul a sa propre histoire. L'abaque apparaît sous des formes indépendantes dans des civilisations qui ne communiquent pas entre elles~: jetons en argile à Sumer ($\sim$2500~av.~J.-C.), tables à rainures en Grèce et à Rome, bouliers en Chine (dynastie Han) et au Japon (\emph{soroban}). \dimmark{D5}{?}%
Cette convergence suggère que la demande de calcul naît partout où la complexité économique~--- commerce, fiscalité, gestion des stocks~--- dépasse la capacité du calcul mental. Le système de numération positionnelle sexagésimale des Babyloniens, transmis par l'astronomie arabe puis repris par les astronomes européens, est encore celui des heures, des minutes et des degrés~: \dimmark{D10}{?}%
c'est un cas précoce de verrouillage de standard, où le coût de basculement maintient une convention en place sur quatre millénaires. Le mécanisme d'Anticythère ($\sim$200~av.~J.-C.), retrouvé dans une épave en 1901, est un calculateur astronomique à engrenages capable de prédire les éclipses et les positions planétaires \parencite{freeth_decoding_2006}. Mais il n'aura pas de descendance~: en l'absence d'une demande économique ou administrative suffisante pour justifier sa reproduction, le calcul mécanique sera réinventé au \textsc{xvii}\ieme{} siècle, pas hérité de l'Antiquité.

\subsection{Le transport de l'information}\label{subsec:transport}

La transmission d'information à distance repose sur le porteur humain et animal. Le \emph{cursus publicus} romain, organisé par Auguste, relie Rome aux provinces par des relais espacés de dix à quinze miles~; sa vitesse moyenne ne dépasse pas cinquante miles par jour \parencite{kolb_transport_2000}. Mille ans plus tard, la poste de la famille Taxis, à partir de 1490, couvre les Pays-Bas, l'Italie et l'Empire. La vitesse maximale de transmission de l'information reste celle du cheval au galop~--- environ trente kilomètres par jour en moyenne, relais compris. Le message ne voyage pas plus vite que le messager. \dimmark{D6}{?}%
La portée de la coordination économique~--- marchés, administrations, empires~--- était ainsi bornée par la vitesse du transport physique\footnote{Le principe est posé dès 1776 par Adam Smith~: la division du travail est limitée par l'étendue du marché (\emph{Wealth of Nations}, livre~I, chap.~3). La vitesse de transmission de l'information est l'une des contraintes sur cette étendue.}.

\subsection{L'imprimerie et la reproduction du savoir}\label{subsec:imprimerie}

La transformation du stockage intervient avec l'imprimerie à caractères mobiles de Gutenberg ($\sim$1440), \dimmark{D2}{?}%
qui réduit le coût de reproduction du texte de plusieurs ordres de grandeur\footnote{\textcite{eisenstein_printing_1979} a montré que cette chute du coût marginal de reproduction transforma non seulement la diffusion du savoir mais aussi les conditions de l'activité scientifique elle-même, en rendant possible la comparaison systématique des sources et la cumulativité des résultats.} sans modifier la nature du support~: le papier, introduit en Europe depuis la Chine via le monde arabe au \textsc{xii}\ieme{} siècle, reste un support inerte qui ne consomme pas d'énergie pour exister. La diffusion du \emph{Liber Abaci} de Fibonacci (1202), qui introduit la numération indo-arabe en Europe, et celle des tables de logarithmes de Napier (1614), sont des effets de cette infrastructure de stockage. L'imprimerie est aussi, comme le soutient \textcite{mokyr_gifts_2002}, l'infrastructure de la \og \emph{knowledge economy}\fg{} qui rend possible la révolution industrielle.

\subsection{Pascal, Leibniz et les premières machines à calculer}\label{subsec:pascaline}

Le calcul mécanique apparaît au \textsc{xvii}\ieme{} siècle avec deux prototypes sans descendance industrielle. La Pascaline (1642) de Blaise Pascal, conçue pour aider son père, commissaire aux impôts en Haute-Normandie, exécute l'addition et la soustraction par un mécanisme de roues à dix dents et de report de retenue. Le \emph{Stepped Reckoner} de Leibniz (1673) étend le mécanisme aux quatre opérations en introduisant le tambour cannelé (\emph{Leibniz wheel}), un cylindre à neuf dents de longueur croissante. Ni Pascal ni Leibniz ne résoudront le problème de la production en série~: la Pascaline sera fabriquée à une cinquantaine d'exemplaires, le \emph{Stepped Reckoner} restera à l'état de prototype. \dimmark{D5}{?}%
Comme pour le mécanisme d'Anticythère, la capacité technique précède la demande~: le calcul mécanique ne passera à l'échelle industrielle qu'avec l'Arithmometer de Thomas de Colmar, breveté en 1820, lorsque les compagnies d'assurance et les administrations fiscales fourniront un marché suffisant.

\subsection{Prony, le Nautical Almanac et la division du travail intellectuel}\label{subsec:prony}

Le calcul humain organisé, en revanche, précède de loin la machine et détermine les conditions de son apparition. En France, Gaspard de Prony organise en 1791 le calcul de tables logarithmiques et trigonométriques en mobilisant soixante-dix à quatre-vingts personnes réparties en trois niveaux hiérarchiques~: les mathématiciens qui conçoivent les formules, les calculateurs qualifiés qui les décomposent en opérations élémentaires, les calculateurs non qualifiés qui exécutent ces opérations. \dimmark{D4}{?}%
Babbage, qui visite le projet, en tirera dans \emph{On the Economy of Machinery and Manufactures} (1832) la conclusion que la division du travail de Smith s'applique aussi au traitement de l'information. En Angleterre, le Nautical Almanac Office, créé en 1767 par Nevil Maskelyne, emploie en permanence des \emph{computers}~--- le mot désigne alors une personne, pas une machine~--- pour calculer les éphémérides astronomiques nécessaires à la navigation\footnote{Les éphémérides ne sont pas un bien de consommation~: elles sont un intrant de la flotte marchande britannique, et leur production est financée par l'État pour des raisons de puissance commerciale et militaire. La demande de calcul est ici dérivée de la demande de commerce maritime.}.

\subsection{Le métier Jacquard et le contrôle par l'information stockée}\label{subsec:jacquard}

Le contrôle automatisé par l'information apparaît sous deux formes distinctes à la fin du \textsc{xviii}\ieme{} siècle. Le régulateur centrifuge de Watt (1788), qui coupe l'alimentation en vapeur quand la vitesse du moteur dépasse un seuil, est un mécanisme de rétroaction physique~--- du \emph{feedback} mécanique, pas du contrôle informationnel. Le métier Jacquard (1804), en revanche, est le premier cas de contrôle d'un processus physique par une séquence d'information stockée~: les cartes perforées encodent le motif du tissu, et le métier l'exécute sans intervention humaine. \dimmark{D3/D5}{?}%
L'innovation répond à une demande précise~: la complexité croissante des motifs exigés par le marché de la soie lyonnaise dépassait la capacité de coordination manuelle des tireurs de lacs\footnote{Avant le métier Jacquard, la réalisation d'un motif complexe nécessitait un \emph{tireur de lacs}, ouvrier spécialisé qui actionnait manuellement les fils de chaîne selon les instructions du dessinateur. Le métier Jacquard substituait du capital (les cartes, le mécanisme) à ce travail qualifié.}.

Ces trajectoires~--- transmettre, calculer, stocker, contrôler~--- coexistent sans converger. Les gens qui calculent les tables de Prony n'ont aucun rapport avec le réseau postal des Taxis, ni avec les imprimeurs de Gutenberg, ni avec les tisserands lyonnais qui travaillent sur les métiers Jacquard. L'histoire longue que cette section a esquissée n'est pas un prologue décoratif~: elle pose les conditions initiales à partir desquelles le télégraphe optique de Chappe (1794), puis le télégraphe électrique de Morse et Cooke \& Wheatstone (1837), vont pour la première fois mécaniser l'une de ces quatre fonctions~--- la transmission~--- et ouvrir la trajectoire que le reste du chapitre va parcourir.




% =============================================================================

% =============================================================================
% §1.2 Le télégraphe (1794–1880)
% =============================================================================
\clearpage
\section{Le télégraphe (1794--1880)}\label{sec:telegraph}
% =============================================================================

\subsection{Le télégraphe optique et électrique}\label{subsec:telegraph}

La transmission d'information à distance n'attendit pas l'électricité, mais les raisons pour lesquelles elle devint, au cours du dix-neuvième siècle, un enjeu économique croissant doivent être posées d'entrée. Entre 1840 et 1914, le commerce international change d'échelle. Le commerce extérieur de la France passe de 2,5~milliards de francs courants en 1847 à 15~milliards en 1913~; celui de l'Angleterre de 13 à 35~milliards entre 1870 et 1913, celui de l'Allemagne de 5 à 25~milliards \parencite{mitchell_historical_2007}\footnote{Valeurs exprimées en francs-or courants par \textcite{mitchell_historical_2007}. Les ordres de grandeur de la croissance sont comparables en valeur réelle.}. Des navires chargent du coton à New York et le déchargent à Liverpool, mais l'information sur les cours, les commandes et les stocks voyage au même rythme qu'eux, quand elle ne prend pas de retard. La poste, qui repose sur le transfert physique du document, ne peut qu'aligner la vitesse de l'information sur celle des personnes\footnote{\textcite{juhasz_spinning_2018} le formulent avec précision~: avant le télégraphe, l'information ne pouvait voyager qu'à la vitesse du mode de transport physique le plus rapide. \textcite{carey_communication_1989} fait de cette indissociabilité entre communication et transport le point de départ de son analyse~: le télégraphe est la première technologie à séparer les deux.}. \dimmark{D6}{?}%
Or la croissance du volume des échanges accroît le nombre de décisions~--- acheter, vendre, acheminer, assurer~--- qui dépendent d'une information distante~; chaque décision prise avec un retard de dix jours sur les cours réels est une décision prise dans l'obscurité. Le problème n'est pas que l'information soit lente~; c'est que le système qu'elle doit coordonner grossit plus vite qu'elle.

C'est dans cet écart que le télégraphe va s'insérer. Le principe de la transmission électrique d'un signal était connu depuis les expériences de Lesage à Genève (1774) et de Sömmering à Munich (1809), mais ces systèmes étaient des impasses~: celui de Lesage exigeait autant de fils que de lettres dans l'alphabet, celui de Sömmering reposait sur un bullage électrochimique trop lent pour un usage opérationnel \parencite{beauchamp2001telegraphy}. C'est la découverte de l'électromagnétisme par Oersted à Copenhague en 1820 qui ouvrit les deux trajectoires viables~: d'un côté le galvanomètre, dont Schilling à Saint-Pétersbourg puis Gauss et Weber à Göttingen tirèrent un télégraphe à aiguilles~; de l'autre l'électro-aimant, dont Sturgeon à Londres puis Henry à Albany firent un actionneur capable de commander un bras métallique à distance \parencite{kooij_telegraph_2015}\footnote{Le relais électromagnétique de Henry est le composant clé~: il permettait de régénérer le signal sur de longues distances, rendant le réseau possible. C'est le même principe~--- un circuit électrique qui en commande un autre~--- qui servira dans les réseaux de distribution électrique de la génération suivante.}. Les deux lignées aboutirent indépendamment au télégraphe électrique viable en 1837, la première avec Cooke et Wheatstone en Angleterre, la seconde avec Morse en Amérique. La capacité technique existait donc, mais elle ne suffit pas à expliquer le déploiement~: il fallait aussi une demande solvable, un acteur prêt à financer le réseau, et un cadre institutionnel qui en autorise ou en commande la construction\footnote{On distingue ici, par commodité analytique, la capacité technique (l'existence d'un artefact fonctionnel) des conditions de son déploiement effectif (demande, financement, cadre réglementaire). Cette distinction est un choix de cadrage~: la sociologie des techniques a montré que la capacité technique est elle-même socialement construite~--- un appareil existe quand un réseau d'acteurs le fait exister. Nous la maintenons parce qu'elle organise la comparaison entre cas, en rendant visible le contraste entre une capacité partagée et des trajectoires divergentes.}. Ces conditions se manifestèrent de manière hétérogène~: demande militaire en France (la Convention avait besoin de coordonner un territoire en guerre), demande ferroviaire en Angleterre (la sécurité des trains sur voie unique), demande commerciale aux États-Unis (l'arbitrage sur les cours du coton), commande étatique en Prusse (le suivi du Parlement de Francfort depuis Berlin). Dans chaque cas, la forme institutionnelle du réseau~--- monopole d'État, concession privée, initiative ferroviaire~--- refléta l'intérêt de l'acteur qui le finançait. La question est alors de savoir si ces conditions d'émergence, radicalement différentes, produisent des formes organisationnelles tout aussi différentes, ou si les propriétés du réseau lui-même finissent par dominer les conditions initiales. La réponse n'est pas donnée d'avance~; les cas qui suivent la documenteront. Ce que l'on peut déjà noter, c'est que les composantes de la demande n'ont pas la même temporalité~: la coordination militaire en temps de guerre, immédiate et étatique~; la coordination commerciale en temps de paix, progressive et privée~; la coordination ferroviaire, tardive et organisationnelle.

Que la communication rapide à distance fût possible, le réseau de sémaphores de Claude Chappe l'avait déjà montré en France. Le réseau fut déployé à partir de 1793 non par curiosité technique mais parce que la Convention nationale, en guerre contre les puissances européennes, avait un besoin urgent de coordonner les mouvements militaires sur un territoire immense~; il couvrait 534~stations en 1844 et plus de cinq mille kilomètres de lignes \parencite{holzmann_early_1995, field_french_1994}. La ligne Paris-Lille transmettait le premier symbole d'un message en neuf minutes~; la ligne Paris-Strasbourg mettait six heures pour une dépêche de 25~mots \parencite[p.~87]{holzmann_early_1995}. Mais le principe restait mécanique et optique --- des bras articulés observés à la longue-vue, reproduits de tour en tour, inutilisables la nuit et par temps de brouillard. En 1838, l'administration comptait près de mille stationnaires pour une cinquantaine d'inspecteurs \parencite{ducamp_telegraphe_1869, field_french_1994}. \dimmark{D3}{?}%
Le poste dominant n'était pas le capital matériel~--- les tours étaient des constructions modestes~--- mais le travail humain, irréductiblement attaché à chaque n\oe{}ud~: doubler la longueur du réseau doublait le nombre d'opérateurs.

C'est ce rapport entre capital et travail que le télégraphe électrique allait modifier~--- et l'observation la plus frappante, documentée indépendamment dans chaque contexte national, est qu'il le fit partout de la même manière. Le télégraphe électrique substituait du capital physique~--- le fil métallique, les poteaux, les isolateurs, les batteries~--- au travail humain des opérateurs de tour\footnote{Le fait de cette substitution est documenté indépendamment pour chaque pays par \textcite{holzmann_early_1995}, \textcite{kieve_electric_1973}, \textcite{hubbard_cooke_1965} et \textcite{headrick_invisible_1991}. Aucun réseau télégraphique électrique ne conserva le principe des relais humains visuels.}. Il fonctionnait la nuit. Il fonctionnait par temps de brouillard. Sa disponibilité ne dépendait plus du climat mais de l'intégrité matérielle de la ligne. En contrepartie, il exigeait un investissement initial considérable en infrastructure linéaire~: les séries du \emph{National Bureau of Economic Research}\footnote{\emph{NBER} --- \emph{National Bureau of Economic Research}, organisme américain de recherche économique fondé en 1920, principal producteur de séries statistiques longues pour l'économie des États-Unis.} compilées par \textcite{gordon_capital_nineteenth} montrent que le réseau américain passa de 12\,000~miles de fil en 1850 à plus d'un million en 1900. La compétition entre les deux systèmes rivaux~--- le télégraphe à aiguilles de Cooke et Wheatstone\footnote{William Fothergill Cooke (1806--1879) et Charles Wheatstone (1802--1875), inventeurs britanniques du télégraphe électrique à aiguilles, breveté en 1837. Leur système utilisait initialement cinq aiguilles aimantées dont les déviations indiquaient les lettres, puis fut simplifié à deux aiguilles.}, qui nécessitait cinq fils puis deux \parencite{hubbard_cooke_1965}, et le système de Morse\footnote{Samuel Finley Breese Morse (1791--1872), peintre américain sans formation scientifique, conceptualisa le télégraphe en octobre~1832, sur le paquebot \emph{Sully} qui le ramenait de trois ans en Europe. Pendant les six semaines de traversée, le Dr.~Charles T.~Jackson de Boston lui montra un électro-aimant~; Morse esquissa son système dans son carnet de croquis. En débarquant à New York, il dit au capitaine Pell~: \og\emph{should you ever hear of the telegraph one of these days as the wonder of the world, remember the discovery was made on the good ship Sully}\fg{} \parencite[p.~251--257]{prime_life_1875}. Il passa cinq ans de misère à construire un prototype avec l'aide du mécanicien Alfred Vail (1807--1859), qui améliora considérablement le matériel et contribua à l'élaboration du code. Le système fut breveté en 1840~; la première dépêche publique fut transmise le 24~mai 1844 sur la ligne Washington-Baltimore~: \og\emph{What hath God wrought}\fg{} \parencite[p.~17]{bertho_histoire_1984}. Morse n'était pas électricien~; comme Cooke (ancien militaire), comme Bell (enseignant pour sourds), comme Marconi (autodidacte), il pensait en termes de fonction~--- communiquer à distance~--- pas en termes de substrat électrique \parencite{kooij_telegraph_2015}.}, qui n'en requérait qu'un~--- illustre la manière dont la structure de coûts orientait les choix techniques, sans les déterminer seule~: le code Morse économisait du cuivre, mais son succès tenait aussi à la qualité du matériel amélioré par Vail, au financement fédéral de la première ligne, et à la familiarité des opérateurs avec le système une fois installé~; d'autres trajectoires étaient possibles, et le télégraphe à aiguilles resta en usage en Grande-Bretagne jusqu'au \textsc{xx}\ieme{} siècle. Les opérateurs Morse devinrent, selon l'expression de \textcite{gabler_american_1988}, les premiers \og travailleurs de l'information\fg{} au sens propre~--- là où Wheatstone avait minimisé le besoin de compétence cognitive au prix d'un capital physique accru, \dimmark{D3/D4}{*}%
Morse fit l'inverse.

\dimmark{D2/D3}{*}%
La substitution du capital au travail dans la transmission d'information est documentée dans chaque pays qui déploie un réseau télégraphique~--- un fait d'observation, pas encore un mécanisme. Ce qui suit examine si cette régularité technique s'accompagne d'une régularité institutionnelle, ou si les conditions d'émergence~--- radicalement différentes d'un pays à l'autre~--- produisent des formes organisationnelles tout aussi différentes.

\dimmark{D2}{*}%
En Angleterre, la demande initiale ne vint ni de l'État ni du commerce, mais des compagnies ferroviaires, qui avaient besoin de signaler les trains sur voie unique. William Fothergill Cooke (1806--1879), ancien officier de l'armée des Indes reconverti en étudiant en médecine, assista le 22~avril 1836 à Heidelberg à une démonstration du télégraphe de Schilling~; il abandonna la médecine sur-le-champ et se consacra au télégraphe, qu'il identifiait d'emblée comme un instrument ferroviaire, non comme un moyen de communication publique \parencite{kooij_telegraph_2015, hubbard_cooke_1965}. Charles Wheatstone (1802--1875), facteur d'instruments de musique devenu expérimentateur en électricité, apportait la compétence technique qui manquait à Cooke~; ensemble, ils installèrent en 1839 une liaison d'une vingtaine de kilomètres pour la \emph{Great Western Railway} \parencite{hubbard_cooke_1965}. Le réseau se développa ensuite par l'initiative privée, sans intervention étatique\footnote{Le médecin Edward Davy (1806--1885) avait, dès 1838, démarché huit compagnies ferroviaires et créé une société, la \emph{Voltaic Telegraph Company}, capitalisée à 500\,000~livres~; il émigra en Australie avant d'aboutir, et son brevet fut racheté pour 600~livres \parencite{kooij_telegraph_2015}.}. Mais l'initiative privée ne produisit pas la concurrence~: elle produisit la concentration. \textcite{kieve_electric_1973} retrace la dynamique~: en 1852, l'\emph{Electric Telegraph Company} dominait avec environ quatre mille miles en service~; la croissance s'accompagna d'un service médiocre dans les zones peu rentables et de tarifs élevés~; la presse, particulièrement déçue, milita pour la nationalisation, qui intervint en 1869\footnote{Le \emph{Telegraph Act} de 1863 avait tenté de corriger les abus sans succès.}. L'initiative privée avait construit le réseau~; la concentration l'avait transformé en monopole~; et c'est l'État, sous la pression de la presse et non du marché, qui finit par en prendre le contrôle.

\dimmark{D2}{*}%
Aux États-Unis, le contexte était entièrement différent. Aucune compagnie privée ne portait la demande~; c'est l'État fédéral qui finança la première ligne, sur une promesse de souveraineté~--- le Congrès débattit pendant six ans de la question de savoir si la communication instantanée relevait de l'utilité générale et justifiait un financement fédéral\footnote{Le débat portait sur la souveraineté et l'utilité publique au sens politique du terme, non sur le concept de \og bien public\fg{} au sens technique de \textcite{samuelson_pure_1954} (non rival, non excluable), qui est postérieur d'un siècle.}. \textcite{thompson_wiring_1947} détaille les six années d'obstruction parlementaire~--- les représentants du Tennessee et de l'Alabama firent obstacle au vote de la subvention~--- avant que la première ligne Washington-Baltimore (64~km) n'ouvre en mai 1844. Le réseau se développa ensuite dans une concurrence anarchique~--- vingt compagnies en 1850, quatorze dans le seul Ohio en 1852, profits insignifiants, comptes mal tenus, tarifs arbitraires \parencite{bertho_histoire_1984, nonnenmacher_telegraph_2001}. La concentration capitaliste l'emporta~: le 12~juin 1866, l'\emph{American Telegraph Company} et la \emph{Western Union} fusionnèrent pour former le premier monopole privé américain \parencite{thompson_wiring_1947, nonnenmacher_telegraph_2001}. Le monopole n'avait été ni voulu ni prévu par le Congrès qui avait financé la première ligne~; il était le produit d'une concurrence que personne n'avait su réguler.

\dimmark{D2/D8}{*}%
En France, un réseau de télégraphie optique couvrait déjà le territoire, et son administration~--- militaire dans son principe, rattachée au ministère de l'Intérieur~--- n'avait aucune raison de céder la place. \textcite{ducamp_telegraphe_1869} rapporte en détail la fraude boursière qui déclencha le changement~: entre 1834 et 1836, les frères Blanc faisaient passer les cours de la rente d'État sur la ligne Paris-Bordeaux aux spéculateurs provinciaux avant l'arrivée de la malle-poste\footnote{Le procès (Tours, mars 1837) se termina par un acquittement faute de base légale. Alexandre Dumas s'en inspire dans \emph{Le Comte de Monte-Cristo} (1844), chapitres LX--LXI.}. Aucune loi n'interdisant la fraude, le ministère fit voter la loi du 2~mai 1837, dont l'article unique établissait le monopole d'État~: \og \emph{Quiconque transmettra sans autorisation des signaux d'un lieu à un autre, soit à l'aide de machines télégraphiques, soit par tout autre moyen, sera puni}\fg{}\footnote{Loi du 2~mai 1837, article unique (\emph{Bulletin des Lois}, 1837).}. Le débat avait vu s'affronter trois thèses~: liberté totale~--- les libéraux soutenaient que \og \emph{le monopole des télégraphes ne doit pas être plus concédé au gouvernement que le monopole de la presse}\fg{}\footnote{Les trois positions sont rapportées par \textcite[pp.~22--23]{bertho_histoire_1984}, qui s'appuie sur les comptes rendus de la Chambre.}~---, accès régulé au réseau d'État (les hommes d'affaires), monopole strict (le ministre). Le comte de Gasparin, ministre de l'Intérieur, l'emporta en invoquant un argument économique~--- un réseau de communication tendait par nature vers le monopole~--- et un argument policier, que le \emph{Moniteur universel} du 6~janvier 1837 rapporte en ces termes~: \og \emph{Que serait-il arrivé si les canuts insurgés de Lyon avaient pu transmettre le signal de la révolte sur tout le territoire~?}\fg{} Et Gasparin d'ajouter~: \og \emph{Les poisons et les poudres ne sont débités que par l'autorisation de l'État et certes la télégraphie entre des mains malveillantes pourrait devenir une arme dangereuse.}\fg{} Comme le documentent \textcite{ducamp_telegraphe_1869} et \textcite{belloc_telegraphie_1894}, l'administration résista ensuite au télégraphe électrique pendant douze ans~: elle commanda aux ateliers Bréguet un appareil qui reproduisait les mouvements du sémaphore optique~--- pour ne pas avoir à former les opérateurs à un nouveau code~--- et le lobby du télégraphe optique fit voter en 1842 des crédits pour un système de sémaphores nocturnes. Il fallut un changement de régime~--- l'arrivée de Louis Napoléon Bonaparte~--- pour que le réseau s'ouvre au public en 1851\footnote{Loi du 29~novembre 1850. Le ministre de l'Intérieur déclara devant les députés, en des termes qui marquaient le renversement complet de la position de 1837~: \og \emph{Il est impossible que le gouvernement français refuse de faire participer le commerce et l'industrie de notre pays aux facilités qu'offre le télégraphe électrique}\fg{} \parencite[p.~25]{bertho_histoire_1984}. Plus de six millions de francs-or furent alloués à la construction du réseau pour la seule année 1852 \parencite{ducamp_telegraphe_1869}.}. Le cas français est le seul où le monopole d'État sur la communication précéda la technologie qu'il allait gouverner~; c'est aussi le seul où une technique~--- le sémaphore optique~--- bloqua activement, par ses défenseurs institutionnels, la technique appelée à la remplacer.

\dimmark{0-GOV}{*}%
En Allemagne, le réseau fut d'emblée public. \textcite{siemens_lebenserinnerungen_1892} raconte comment, en 1842, quelques appareils de Wheatstone lui parvinrent alors qu'il était jeune officier d'artillerie~; il s'associa avec le mécanicien Halske pour présenter en 1846 un modèle dérivé du télégraphe anglais. La première ligne, Francfort-Berlin, fut construite en 1848 pour permettre à la capitale de la Prusse de suivre les travaux du Parlement fédéral \parencite[ch.~2]{headrick_invisible_1991, bertho_histoire_1984}. La ligne fut enterrée le long des voies ferrées pour la protéger des saboteurs. Le réseau se construisit ensuite sous tutelle de la poste prussienne, en même temps que l'unité politique allemande. En Allemagne, pas de résistance institutionnelle, pas de concurrence anarchique, pas de nationalisation tardive~; le réseau fut public d'emblée, et le monopole n'eut pas besoin d'émerger~--- il fut posé.

\dimmark{D2/D9/0-GOV}{+}%
Le schéma ne s'arrêtait pas à l'Europe et à l'Amérique du Nord. En Inde, l'administration coloniale britannique inaugura dès 1851 la ligne Calcutta-Diamond Harbour~; en 1865 le réseau couvrait plus de 28\,000~miles, au service du contrôle impérial \parencite[ch.~3]{headrick_invisible_1991}. Au Japon, le gouvernement Meiji fit construire la ligne Tokyo-Yokohama en 1869, dans le cadre d'un programme de modernisation étatique accéléré~: le télégraphe y fut d'emblée un monopole public au service de l'unité nationale \parencite[ch.~4]{headrick_invisible_1991}. En Chine, le réseau fut déployé à partir des années~1870, d'abord sous concession étrangère~--- la \emph{Great Northern Telegraph Company} danoise opéra les premières lignes~--- puis sous administration impériale chinoise. \textcite{hao_wiring_2022} et \textcite{gao_communication_2021} ont montré que les effets sur l'intégration des marchés du riz furent comparables à ceux observés sur le coton transatlantique. Dans ces trois cas, les conditions initiales différaient radicalement~--- colonisation, modernisation autoritaire, concession étrangère~--- mais le résultat institutionnel fut le même~: le monopole, public ou colonial.

\dimmark{D2/0-GOV}{*}%
Que sept conditions d'émergence radicalement différentes produisent le même résultat institutionnel~--- une forme ou une autre de monopole~--- n'est pas, en soi, surprenant~: la théorie économique des monopoles naturels prédit ce résultat pour toute infrastructure de réseau à coûts fixes élevés et rendements croissants. La convergence est documentée indépendamment par des sources distinctes~: \textcite{nonnenmacher_telegraph_2001} pour les États-Unis, \textcite{kieve_electric_1973} pour l'Angleterre, \textcite{bertho_histoire_1984} pour la France, \textcite{headrick_invisible_1991} pour les réseaux coloniaux. Ce que le détour par les sept cas ajoute à la théorie, c'est autre chose~: la diversité des véhicules politiques de la concentration (contrôle policier en France, souveraineté fédérale aux États-Unis, commande impériale en Inde), le fait que le contrôle politique et les rendements croissants jouent en parallèle plutôt qu'en séquence, et surtout le cas français de non-émergence~--- douze ans de blocage par l'administration du sémaphore, de 1837 à 1851~--- qui révèle que la pression vers le monopole peut retarder la transition technique autant qu'elle l'accélère.

Le mécanisme proposé pour rendre compte de cette convergence est celui que \textcite{wu2010master} a synthétisé sous le nom de \emph{master switch}~: la combinaison des rendements croissants du réseau, du désir de contrôle politique et de la logique capitalistique de consolidation exercerait une pression structurelle vers la concentration, quelle que soit la trajectoire institutionnelle initiale. La proposition de Wu est un essai d'historien-juriste, non un article académique formalisé~; elle opère par compilation interprétative de trois mécanismes économiques chacun bien établis par ailleurs dans la littérature~: les rendements croissants de réseau \parencite{katz_network_1985, economides_economics_1996}, le contrôle politique sur les infrastructures stratégiques \parencite{mazzucato_entrepreneurial_2013, headrick_invisible_1991}, et la consolidation capitalistique des infrastructures lourdes \parencite{chandler_visible_1977, perez_technological_2002}. Ce n'est donc pas l'invocation d'un théorème mais l'usage d'un cadre interprétatif qui rassemble trois résultats économiques solides en une régularité descriptive applicable aux réseaux de communication. Cette explication est compatible avec les sept trajectoires observées~; elle relève toutefois de ce que l'on pourrait appeler, en suivant la distinction de \textcite[ch.~6]{trannoy_economistes_2025}, une causalité structurelle~: un modèle théorique qui propose un mécanisme, non un résultat isolé par une variation exogène. Ce que les sept cas permettent d'établir, c'est la régularité du résultat~; ce que Wu propose, c'est un mécanisme qui en rend compte~; ce qui reste ouvert, c'est de savoir si cette pression est propre aux réseaux de communication ou commune à toutes les infrastructures lourdes.

\emergence{D2~--- Sept trajectoires institutionnelles, un seul résultat~: le monopole. La convergence est une régularité documentée, attendue par la théorie des monopoles naturels~; ce que les cas révèlent en propre, c'est la diversité des véhicules politiques, le jeu parallèle rendements croissants / contrôle étatique, et la non-émergence française comme révélateur des conditions de blocage. \questionch{2}{D2~--- Cette convergence se reproduit-elle pour le téléphone, la radio, Internet~? Est-elle un trait des réseaux de communication en tant que tels, ou des infrastructures lourdes en général~?}}

\emergence{D8~--- La loi de 1837, votée pour un réseau de sémaphores, s'applique encore au télégraphe électrique grâce à la formule \og soit par tout autre moyen\fg{}. C'est un cas où le cadre réglementaire survit à la technologie pour laquelle il a été conçu. \questionch{2}{D8~--- Ce phénomène de persistance réglementaire se retrouve-t-il aux vagues technologiques ultérieures~?}}

\emergence{D9/0-GOV~--- La géographie du déploiement télégraphique reflète les rapports de puissance~: réseau colonial en Inde, concession étrangère en Chine, domination britannique sur les câbles sous-marins. \questionch{2}{D9/0-GOV~--- La localisation des infrastructures informationnelles suit-elle la géographie du pouvoir économique et politique~?}}

La tarification du télégramme~--- au mot dans la plupart des systèmes~--- imposait une discipline de compression du langage et donnait à l'information un prix unitaire explicite. Sur les réseaux nationaux, les prix baissèrent considérablement, mais selon des mécanismes différents~: par la concurrence aux États-Unis~--- 1,55~dollar pour dix mots entre New York et Chicago en 1850, 40~cents en 1890 selon les données de \textcite{nonnenmacher_telegraph_2001}~---, par la politique tarifaire volontariste en France, où une dépêche de vingt mots de Paris à Marseille coûtait 18,44~francs en 1851 et 1~franc en 1878 \parencite{ducamp_telegraphe_1869}\footnote{Le tarif de 1851 représentait, selon la formule de \textcite{bertho_histoire_1984}, pp.~33--35, \og le prix d'un vêtement du dimanche qu'une ouvrière portera toute l'année\fg{}.}. Le câble transatlantique, en revanche, restait réservé au trafic d'affaires~: \dimmark{D5}{*}%
10~dollars par mot en 1866\footnote{Le bureau de Londres de l'\emph{Associated Press} envoyait 35\,000~mots par jour en 1869, pour des dépenses atteignant 2\,000~dollars par jour \parencite{frederix_siecle_1959, bertho_histoire_1984}.}.

\paragraph{Le câble transatlantique et la formation des marchés}\mbox{}\\
La coordination marchande, elle, dispose d'un résultat d'une force probante rare dans cette section, fondé sur une stratégie d'identification économétrique. Avant le câble transatlantique de 1866, le prix du coton à Liverpool et celui de New York pouvaient diverger pendant les dix jours de traversée~; la bourse du coton de New York fut créée en 1870, quatre ans après la pose du câble \parencite{standage1998victorian, gordon_thread_2002}. \textcite{steinwender_real_2018}, à partir de données quotidiennes du \emph{New York Times} et du \emph{Liverpool Mercury} et d'une stratégie d'identification fondée sur la rupture temporelle du câble, a montré que la connexion transatlantique réduisit l'écart de prix moyen de 35\,\%, l'équivalent de la suppression d'un tarif \emph{ad valorem} de 7\,\%, et que les gains d'efficience représentaient environ 8\,\% de la valeur annuelle des exportations. L'effet n'était pas symétrique~: les prix à New York réagissaient fortement aux nouvelles de Liverpool, tandis que Liverpool réagissait peu aux nouvelles de New York~--- parce que le marché britannique était le marché dominant du coton mondial, et que New York était preneur de prix, non faiseur. Surtout, l'effet n'était pas seulement redistributif~: les exportations elles-mêmes répondaient aux nouvelles de prix, ce qui signifie que le câble ne se contentait pas de réallouer les profits entre acteurs mieux et moins bien informés~--- \dimmark{D6}{*}%
il créait du commerce qui n'aurait pas eu lieu sans lui. Un effet comparable avait déjà été observé sur les marchés intérieurs américains~: en 1846, les prix du blé et du maïs à Buffalo avaient quatre jours de retard sur ceux de New York~; en 1848, après la connexion télégraphique entre les deux villes, les prix étaient synchronisés \parencite{nonnenmacher_telegraph_2001}.

\textcite{juhasz_spinning_2018} ont montré, sur les textiles de coton britanniques exportés vers 75~pays entre 1845 et 1880, que l'effet du télégraphe dépendait de la \emph{codifiabilité} du produit~: le fil (\emph{yarn}), dont la variété se résumait à un chiffre standardisé, vit ses importations augmenter de 10,6\,\% après la connexion télégraphique~; le tissu uni (\emph{plain cloth}), dont les caractéristiques pouvaient être décrites en quelques mots, de 5,0\,\%~; le tissu imprimé (\emph{finished cloth}), dont les motifs ne pouvaient être spécifiés que par l'examen d'un échantillon physique, ne bénéficia d'aucun effet statistiquement significatif\footnote{Les élasticités estimées par variables instrumentales (la rugosité du fond marin comme prédicteur de la date de connexion) sont de $-0{,}27$ pour le fil, $-0{,}13$ pour le tissu uni, et $-0{,}03$ (non significatif) pour le tissu fini. \textcite{juhasz_spinning_2018} montrent également que le télégraphe facilita la diffusion technologique internationale, les pays connectés important davantage de machines et d'intrants intermédiaires~--- l'ICT ne facilitait pas seulement le commerce, elle accélérait le transfert de savoir technique.}. \dimmark{D1}{?}%
Le télégraphe accélérait les mots, pas les choses~: pour les produits codifiables, il ne réduisait pas un coût de transaction préexistant mais \emph{rendait possible la spécification à distance}, et donc la commande, qui n'aurait pas eu lieu sans lui~; pour les produits non codifiables, l'échantillon continuait de voyager par bateau, et rien ne changeait. On notera l'asymétrie~: la ligne télégraphique est un bien rival~--- un message occupe le circuit, et le suivant doit attendre~---, mais l'information qu'elle transporte ne l'est pas~--- une fois le cours du coton connu à Liverpool, rien n'empêche sa diffusion immédiate à tous les acteurs du marché.

\emergence{D6/D9/D10~--- Le résultat de \textcite{juhasz_spinning_2018} soulève trois questions structurelles. Premièrement, l'effet de l'ICT sur le commerce n'est pas uniforme~: il dépend de la correspondance entre ce que la technologie sait transmettre et ce que le produit exige comme spécification. Si ce résultat se transpose, chaque saut technologique~--- du télégraphe (mots) au téléphone (voix) à Internet (images, vidéo)~--- redéfinirait la frontière entre les biens spécifiables à distance et ceux qui exigent la co-présence physique. \questionch{2}{Cette frontière est-elle la même que celle de Rauch (1999) entre biens \og organisés\fg{} et biens \og différenciés\fg{}~? L'ICT la déplace-t-elle continûment~?} Deuxièmement, le télégraphe a davantage augmenté le commerce d'intrants intermédiaires (le fil, $+10{,}6\,\%$) que de produits finis (le tissu imprimé, effet nul)~: le télégraphe ne se contente pas d'augmenter le volume du commerce, il en modifie la \emph{composition}, en favorisant la fragmentation de la chaîne de valeur. Un pays connecté importe du fil au lieu du tissu fini~--- et se met à tisser localement, important aussi des machines. Le télégraphe facilite le commerce, mais il réorganise aussi la division internationale du travail et accélère le transfert technologique. \questionch{2}{Ce mécanisme de fragmentation~--- que la littérature contemporaine documente pour Internet et l'offshoring \parencite{grossman_trading_2008}~--- opérait-il déjà au \textsc{xix}\ieme{} siècle avec le télégraphe~? Si oui, la codifiabilité est moins une dimension parmi d'autres qu'un filtre qui module l'effet de toutes les dimensions.} Troisièmement, la codifiabilité n'est pas une propriété fixe des produits~--- elle est relative à la technologie disponible. Le tissu imprimé qui résistait au télégraphe (il fallait voir l'échantillon) deviendrait potentiellement codifiable dès qu'une technologie saurait transmettre des images. Si cette logique se poursuit, chaque saut technologique ne se contenterait pas d'accélérer la transmission de ce qui est déjà codifiable~: il \emph{rendrait codifiable ce qui ne l'était pas}, et déplacerait la frontière entre les transactions réalisables à distance et celles qui exigent la co-présence. \questionch{2}{Si l'IA générative rend codifiable presque tout (image, voix, intention, spécification), est-elle une ICT de plus dans la série, ou une ICT qui élimine le filtre même qui limitait toutes les précédentes~?} \questionch{3}{Si la codifiabilité croissante fragmente les chaînes de valeur de services comme le télégraphe a fragmenté celle des textiles, chaque fragment échangé à distance consomme de la bande passante, du calcul et du stockage~--- c'est-à-dire de l'énergie. La fragmentation informationnelle est-elle un vecteur de rebond énergétique~?}}

Le résultat de \textcite{steinwender_real_2018} n'est pas isolé~: il a été répliqué sur d'autres produits, dans d'autres géographies, par des méthodes différentes. \textcite{ejrnaes_gains_2010} montrent sur le blé Chicago-Liverpool que la demi-vie des chocs à la loi du prix unique passe de 3,5 à 2,6~semaines après le câble, et que les gains d'efficacité du marché sont d'un ordre de grandeur comparable aux gains liés à la baisse des coûts de transport~--- l'information n'est pas un facteur secondaire. \textcite{hao_wiring_2022} répliquent le résultat dans une économie radicalement différente~--- le riz chinois, neuf provinces, 1870-1911 ~--- et observent en outre que la volatilité des prix \emph{augmente} dans les marchés auparavant isolés, exactement comme les flux transatlantiques de Steinwender, parce que la connexion informationnelle expose les marchés locaux à des chocs distants qu'ils ne percevaient pas. \textcite{andrabi_information_2021}, sur le grain indien (192~districts, 1862-1920), isolent l'effet informationnel de l'effet transport en montrant que le télégraphe réduit la dispersion des prix \emph{avant} l'arrivée du chemin de fer ~--- mais seulement là où un moyen de transport physique existe déjà (la Gange, la Grand Trunk Road). L'information sans transport ne produit pas d'arbitrage\footnote{\textcite{gao_communication_2021}, dans un résultat complémentaire sur le réseau chinois~--- où les chemins de fer sont quasi absents ---, montrent que le télégraphe réduit l'incidence des prix extrêmes et atténue la réponse des prix aux chocs météo locaux, mais augmente la réactivité aux chocs dans les régions connectées à distance. C'est un résultat de \emph{risk-sharing} spatial~: la volatilité converge~--- les marchés stables devenant plus volatils, les marchés instables se stabilisant. L'effet dépend de la densité de voies navigables, ce qui confirme la conditionnalité au transport.}. \dimmark{D6}{+}%
L'ensemble dessine une régularité convergente~: le télégraphe modifie la coordination marchande partout où il est déployé, mais cet effet est conditionnel à l'existence d'un moyen d'arbitrage physique, et il a un coût~--- l'exposition des marchés locaux à la contagion des chocs distants.

\emergence{D6~--- Le télégraphe ne réduit pas un coût de transaction préexistant~: il rend possible la transaction elle-même. Le câble transatlantique crée du commerce qui n'aurait pas eu lieu sans lui \parencite{steinwender_real_2018}, et cet effet est répliqué sur le blé, le riz chinois et le grain indien. Mais il est conditionnel~: sans moyen de transport physique, l'information ne produit pas d'arbitrage~; et il a un coût~: l'exposition des marchés locaux à la contagion des chocs distants. \questionch{2}{D6~--- Ce rôle constitutif de l'information synchrone se retrouve-t-il dans les marchés coordonnés par le téléphone, puis par Internet~?}}

L'augmentation du trafic se heurta aux capacités des lignes. Selon \textcite{beauchamp2001telegraphy}, les appareils Morse transmettaient 25~mots à la minute~; les appareils Hughes, déployés en France à partir de 1862, portèrent ce débit à 45~mots~; le système Baudot (1874), fondé sur le principe du temps partagé, atteignit 60~mots. Mais ces gains de débit étaient conditionnels à la qualité matérielle du réseau. \textcite{brooks_report_1875}, ingénieur télégraphique américain envoyé à l'Exposition de Vienne en 1873, observait que le Hughes français tournait sous la pluie à pleine vitesse sur Marseille-Paris ou Lyon-Paris~; aux États-Unis, sur des lignes mal isolées et hâtivement posées, la même technologie était inutilisable, et les opérateurs revinrent à la lecture par sondeur du Morse, plus rustique mais robuste à l'instabilité électrique. \dimmark{0-TER/D9}{*}%
La matérialité du substrat ne conditionnait pas seulement l'extraction des inputs~; elle filtrait quels artefacts pouvaient effectivement fonctionner sur le réseau, et donc l'horizon des trajectoires techniques nationales. Parallèlement, le multiplexage progressait~: duplex de Stearn (1870), quadruplex d'Edison (1874)\footnote{Thomas Alva Edison (1847--1931), ancien opérateur de télégraphe itinérant sur les lignes du Midwest, avait passé ses nuits à expérimenter avec les circuits. Le quadruplex~--- quatre messages simultanés sur un seul fil~--- fut sa première grande invention et lui rapporta 100\,000~dollars de Jay Gould et de Western Union \parencite{kooij_telegraph_2015}. C'est cette fortune et cette maîtrise du circuit électrique, acquises dans le domaine de la communication, qu'il transposa dans le domaine de la puissance~: le quadruplex est le pont entre le monde du télégraphe et celui de l'électrification, via un homme qui était dans les deux. Les capacités transatlantiques passèrent de 15--17~mots par minute vers 1870 à 90~mots par minute après l'introduction de câbles de gros diamètre en 1894.}. Mais c'est en cherchant à réaliser le multiplexage par différenciation des fréquences\footnote{Le multiplexage par différenciation des fréquences consiste à transmettre simultanément plusieurs signaux sur un même fil en attribuant à chacun une fréquence porteuse distincte. Le récepteur sépare les signaux en filtrant les fréquences. C'est le principe qui sera repris, un siècle plus tard, par la radio FM et les réseaux cellulaires.}~--- le \og télégraphe harmonique\fg{}~--- que Graham Bell et Elisha Gray découvrirent le téléphone. Alexander Graham Bell (1847--1922), enseignant pour sourds à Boston, fils d'Alexander Melville Bell (inventeur du \emph{Visible Speech}), expérimentait le soir dans son atelier avec des électro-aimants, des bobines et des diapasons. Il était financé par Gardiner Hubbard, père d'une de ses élèves sourdes~--- Mabel, qu'il épousera~--- et par George Sanders, père d'un autre. Hubbard voulait briser le monopole de Western Union~; il finançait Bell pour trouver une alternative au quadruplex d'Edison. C'est en travaillant sur ce problème commercial~--- le multiplexage~--- que Bell découvrit la transmission de la voix \parencite{kooij_telephone_2016}. Elisha Gray (1835--1901), ingénieur à Western Electric, déposa son \emph{caveat} pour le téléphone le même jour que Bell~--- le 14~février 1876. Gray est l'homme qui ne vit pas la valeur du téléphone parce que sa métrique était celle du télégraphe~: messages par fil, coût par mot. Le téléphone était un sous-produit de la compétition télégraphique, engendré par la recherche d'économies de cuivre. \dimmark{D2/D10}{+}%
La trajectoire du multiplexage donna ainsi naissance à deux inventions qui la dépassaient~: le téléphone de Bell et, par Edison, l'électrification.

\emergence{Le multiplexage est D2 par excellence~: le coût de la ligne est fixe, toute technique qui augmente le nombre de messages par fil augmente le rendement. \questionch{2}{D10~--- la course au débit est-elle une constante structurelle, de Baudot à la 5G~?}}

La guerre de Crimée (1853--1856) illustre un mécanisme d'émergence distinct de tous les précédents~: la demande militaire, financée sans contrainte budgétaire par l'État en guerre, créa en quelques mois une ligne Paris-Vienne-Bucarest-Crimée que le commerce n'aurait pas justifiée. Le télégraphe transforma les conditions de commandement~--- au point que le général Pélissier demanda à être relevé d'un commandement qu'il jugeait \og \emph{\dimmark{D7/0-GOV}{*}%
impossible à exercer à l'extrémité parfois paralysante d'un fil électrique}\fg{}\footnote{Échange entre Napoléon~III et Pélissier rapporté dans \textcite{frederix_siecle_1959}, \textcite{ducamp_telegraphe_1869} et \textcite[p.~44]{bertho_histoire_1984}. L'expression de Pélissier (1855) est peut-être la première formulation du problème de la micro-gestion à distance par les télécommunications. L'agence Havas avait célébré à tort la prise de Sébastopol en septembre~1853, ce qui faillit ruiner son crédit \parencite{frederix_siecle_1959}.}.

Ce que ces usages militaires, commerciaux et financiers ne comptabilisaient pas, c'est ce que le signal présupposait en amont de lui-même. Le signal télégraphique consommait une quantité d'énergie négligeable~--- les piles de Daniell fournissaient quelques milliampères dont le coût n'apparaît nulle part dans les comptes d'exploitation\footnote{Argument \emph{ex silentio}. Western Union dégageait 30 à 40\,\% de marge nette \parencite{nonnenmacher_telegraph_2001}~; les comptes d'exploitation publiés ne détaillent pas les postes de consommation énergétique. \textcite{kieve_electric_1973} fournit des bilans comptables pour le réseau britannique, mais sans décomposition des coûts par nature de ressource.}. \dimmark{0-TER}{+}%
Mais le système reposait sur une chaîne d'extraction dont l'ampleur n'a été documentée que rétrospectivement.

\dimmark{D9/0-TER}{+}%
Aux États-Unis, le réseau terrestre exigeait des centaines de milliers de poteaux en bois dont l'approvisionnement impliquait, comme l'a montré \textcite{fitzmaurice_material_2023}, une déforestation géographiquement délocalisée. Western Union qualifiait ses poteaux de \emph{perishable property}~; les poteaux non traités duraient environ dix ans \parencite{gordon_capital_nineteenth}. Le réseau terrestre laissait entrevoir ce que les câbles sous-marins allaient confirmer à plus grande échelle~: l'infrastructure de transmission avait une empreinte matérielle dont l'ampleur n'apparaissait dans aucun bilan.

\dimmark{D9/0-GOV}{*}%
Sur les câbles sous-marins, en effet, la dépendance matérielle était plus spectaculaire encore. Entre 1850 et 1880, comme le documentent \textcite{headrick_invisible_1991}, \textcite{hunt_imperial_2021} et \textcite{bertho_histoire_1984}, la télégraphie sous-marine internationale était entièrement dominée par l'Angleterre~: les financiers de Londres contrôlaient l'approvisionnement en cuivre et en \emph{gutta-percha}, fixaient les prix et les quantités\footnote{La concentration s'étendait à la fabrication elle-même~: \textcite{muller_telegraph_2015} documentent que Telcon (\emph{Telegraph Construction and Maintenance Company}) manufacturait sur les rives de la Tamise la quasi-totalité des câbles sous-marins mondiaux jusque dans les années~1930.}. Le câble de 1865 pesait 980~kg/km, dont 73~kg/km de cuivre et 98~kg/km de gutta-percha~--- un latex naturel extrait d'arbres tropicaux du genre \emph{Palaquium}, endémiques de Malaisie et Bornéo, dont un arbre mature de trente ans produisait environ 600~grammes utilisables \parencite{greene_dark_2005, sanzo_around_2023}. \textcite{tully_victorian_2009} a estimé, par une reconstruction rétrospective fondée sur les tonnages d'exportation, que les seuls câbles atlantiques de 1865--66 avaient nécessité l'abattage d'au moins 1,2~million d'arbres\footnote{Estimation reconstruite, non décompte exhaustif.}. Telcon estime avoir consommé 47~millions d'arbres sur un siècle \parencite{greene_dark_2005}. Les prix grimpèrent de 8 à 60~dollars espagnols par picul\footnote{Le picul (équivalent à environ 60~kg) était l'unité de mesure standard du commerce de la gutta-percha à Singapour. Le dollar espagnol, monnaie de référence du commerce asiatique au \textsc{xix}\ieme{} siècle, valait environ cinq francs-or ou un dollar américain. La hausse de 8 à 60~dollars par picul en moins de dix ans représente donc une multiplication par sept et demi du prix de la matière première.} entre 1844 et 1853~; les arbres furent épuisés à Singapour dès 1847 \parencite{sanzo_around_2023, buckley_anecdotal_1902}.

\dimmark{D8/D9}{*}%
Le cas du câble transatlantique révèle, au-delà de la matérialité, la structure du financement et la promesse qui le porte. Cyrus Field, homme d'affaires new-yorkais sans compétence technique, finança l'entreprise sur une promesse de rente informationnelle~: celui qui contrôlerait le câble contrôlerait l'arbitrage entre les marchés du coton, du blé et des titres. La promesse mobilisa des capitaux britanniques considérables~--- mais l'accès aux matières premières (cuivre, gutta-percha) constituait une contrainte que personne ne quantifiait, et le financement lui-même était une aventure. Le premier câble (1857--1858), qui ne fonctionna que trois semaines, avait coûté environ 300\,000~livres sterling à l'\emph{Atlantic Telegraph Company}~; la compagnie fit faillite et ne fut reconstituée qu'en 1864, avec les capitaux de Pender et Brassey, pour un câble réussi en 1866 dont le coût atteignit 600\,000 à 700\,000~livres \parencite{gordon_thread_2002}\footnote{Le pattern investissement massif en infrastructure ICT $\rightarrow$ faillite $\rightarrow$ recapitalisation $\rightarrow$ succès se reproduit avec la \emph{railway mania} des années~1840, la \emph{corporate mania} électrique des années~1880--1890 et la bulle dot-com de~2000.}. Les faillites des compagnies françaises de câbles sous-marins~--- la \emph{Compagnie Erlanger} rachetée en 1873, la \emph{Compagnie française du câble Paris-New York} en faillite en 1895~--- montrèrent les limites du développement par concessions dans un secteur à capitaux massifs et risque immense \parencite{bata_cables_1982, headrick_invisible_1991}. Aux faillites par concession s'ajoute, à partir des années~1890, un autre type d'échec~: l'infrastructure posée \emph{avant} la demande. \textcite{muller_telegraph_2015} documentent comment les câbles transpacifiques de 1902 (canadien) et 1904 (américain), reliant l'Amérique du Nord à l'Asie via Hawaï et Midway, échouèrent commercialement. Le trafic ne s'approcha jamais des coûts d'infrastructure. L'Atlantique de 1866 avait suivi une route commerciale dense préexistante~; les câbles pacifiques, eux, furent posés pour ouvrir un marché que l'on espérait voir naître~: la rivalité sino-américaine, les ambitions canadiennes vers l'Asie, la doctrine de la porte ouverte fournissaient les motifs~; la demande commerciale, elle, n'existait pas. \dimmark{D6/D8}{*}%
C'est la position symétrique du blocage français de 1837-1851~: là, la demande existait mais l'infrastructure ne se construisait pas~; ici, l'infrastructure se construit mais la demande ne se manifeste pas. Le câble pacifique fonctionne, le réseau est complet, et les comptes ne rentrent pas. Le réseau mondial reposait sur du cuivre des Cornouailles et du Michigan, du bois du Wisconsin, du latex de Bornéo, du charbon des navires câbliers. \textcite{belloc_telegraphie_1894}, qui écrivait au c\oe{}ur de cette période, constatait dans sa préface que la télégraphie électrique avait été \og surtout remarquable par ses conséquences économiques\fg{}~--- sans éprouver le besoin de mentionner ses conséquences matérielles. La coïncidence est au moins suggestive~: en 1865, la même année que le premier câble transatlantique, \textcite{jevons_coal_1865} publiait \emph{The Coal Question}. Les préoccupations énergétiques de l'Angleterre victorienne portaient sur l'épuisement du charbon et l'avenir de la puissance industrielle~--- pas sur le cuivre, le latex ou le bois que le réseau de communication consommait. On ne peut pas en conclure que l'occultation était structurelle, mais on peut constater que les catégories de l'époque séparaient l'énergie~--- question de puissance~--- et l'information~--- question de commerce et de contrôle. Cette séparation, qu'elle fût délibérée ou simplement un effet des ordres de grandeur, n'a pas été corrigée depuis~: ni les historiens de l'énergie \parencite{smil2017energy, wrigley2010energy}, ni les historiens de l'ICT \parencite{headrick_invisible_1991, standage1998victorian}, ni les \emph{environmental historians} n'ont fait de la relation entre infrastructure informationnelle et substrat matériel un objet d'étude dans la longue durée\footnote{Les seules exceptions sont \textcite{tully_victorian_2009} sur la gutta-percha, \textcite{fitzmaurice_material_2023} sur la matérialité du réseau américain, et l'article programmatique de \textcite{ensmenger_environmental_2018} sur l'histoire environnementale du \emph{computing}. \textcite{magalhaes_accumuler_2024} a produit pour les infrastructures de transport françaises (routes, autoroutes, béton) le type d'histoire environnementale qui manque pour les réseaux de communication. Il n'existe à ce jour aucune analyse du cycle de vie du réseau télégraphique du \textsc{xix}\ieme{} siècle au sens ISO~14040/44.}.

\questionch{3}{Matérialité occultée~: le pattern gutta-percha/poteaux/cuivre préfigure-t-il les terres rares, le lithium et l'eau de refroidissement des datacenters~? La domination britannique sur les matériaux du câble préfigure-t-elle la domination américaine sur les semi-conducteurs ou chinoise sur les terres rares~? \textcite{pirson_environmental_2022} montrent que l'empreinte environnementale par cm\textsuperscript{2} de fabrication de circuits intégrés, qui baissait depuis les années~1990, remonte à partir du n\oe{}ud de 7~nm~: la miniaturisation exige un bouquet élémentaire toujours plus large et plus pur, dont les impacts croissent avec l'ambition fonctionnelle~--- le même pattern, à un siècle et demi de distance.}

\emergence{D1/D6/0-TER~--- Steinwender montre que le stockage physique de coton atténue les frictions informationnelles mais ne les élimine pas~: les négociants utilisent leurs stocks comme tampon face à l'incertitude sur les prix distants. L'analogie avec le stockage par batteries dans la transition énergétique est à analyser~: la batterie atténue l'intermittence des renouvelables mais ne l'élimine pas. Dans les deux cas, l'information en temps réel (le câble) se \emph{substitue} partiellement au stock physique~; si l'analogie tient, le smart grid pourrait jouer un rôle comparable vis-à-vis des batteries~--- le couple information + stockage serait plus puissant que chacun des deux seuls. \questionch{3}{La complémentarité entre fonctions ICT (transmit = câble, store = stock) et substrat matériel est-elle une constante structurelle du couplage ICT-énergie~?}}

Avant de poursuivre, un point de pause~--- non sur le contenu mais sur le statut de ce qui vient d'être avancé. Le récit qui précède a procédé par trois voies que les historiens et les économistes ne confondent pas toujours, et qu'il importe de distinguer\footnote{La distinction entre ces trois registres s'inspire des catégories proposées par \textcite{monnet_perspectives_2025} dans sa réflexion sur les discordes méthodologiques entre économistes et historiens~: causalité processuelle (le récit de l'enchaînement des actions), causalité structurelle (un modèle théorique qui propose un mécanisme), causalité interventionniste (une quasi-expérience qui estime un effet). Voir \textcite[ch.~6--8]{trannoy_economistes_2025} pour un traitement approfondi de ces questions.}.

La reconstruction narrative des trajectoires nationales~--- qui a déployer, pour quelle demande, dans quel cadre institutionnel, et avec quel résultat~--- relève de ce que l'on peut appeler une évidence processuelle~: un enchaînement documenté d'actions, de décisions et de conséquences, attesté par des sources indépendantes. La substitution du capital au travail dans sept contextes institutionnels \parencite{gordon_capital_nineteenth, holzmann_early_1995, gabler_american_1988}, la convergence vers le monopole \parencite{wu2010master, nonnenmacher_telegraph_2001, kieve_electric_1973, bertho_histoire_1984}, le pattern investissement-faillite-recapitalisation du câble \parencite{gordon_thread_2002, headrick_invisible_1991}, la matérialité extractive du réseau \parencite{fitzmaurice_material_2023, tully_victorian_2009}~: ces régularités sont documentées indépendamment par des sources distinctes. Elles ne dépendent pas d'un cadre théorique pour être observées.

Les mécanismes proposés pour en rendre compte relèvent d'un autre registre. Le \emph{master switch} de Wu (rendements croissants + contrôle politique + consolidation), la codifiabilité de Juhász comme filtre des effets du commerce, le seuil de McCallum au-delà duquel la coordination hiérarchique devient nécessaire~: chacun de ces mécanismes est compatible avec le matériau observé, mais aucun n'est isolé par une variation exogène. Ce sont des propositions explicatives, non des résultats testés.

Le résultat de \textcite{steinwender_real_2018} relève d'un troisième registre encore~: une stratégie d'identification économétrique fondée sur la rupture temporelle du câble transatlantique, qui permet d'estimer l'effet de l'information synchrone sur l'intégration des marchés. Ce résultat a été répliqué dans quatre autres géographies par des méthodes différentes \parencite{ejrnaes_gains_2010, hao_wiring_2022, andrabi_information_2021, gao_communication_2021}, ce qui renforce sa portée~; il reste toutefois conditionnel aux hypothèses de l'instrument et au contexte spatio-temporel dans lequel il a été estimé, et ne se transforme pas en loi universelle.

Parmi les pistes qui restent au stade de l'hypothèse~: la tension entre efficacité informationnelle et empreinte matérielle, la substitution partielle entre information et stockage physique, et la question de savoir si l'information synchrone modifie les préférences ou seulement les contraintes.

Mais le télégraphe ne fonctionnait jamais seul. Il s'insérait dans des systèmes techniques plus vastes, et c'est dans cette insertion que l'on peut observer comment une infrastructure de transmission interagit avec un substrat physique et énergétique concret. Le cas qui suit~--- le télégraphe dans le chemin de fer~--- n'est pas un septième cas d'émergence nationale~; c'est un changement de registre. Les sept trajectoires qui précèdent examinaient le télégraphe comme réseau autonome~: qui le déploie, sous quelle forme institutionnelle, avec quel effet sur les marchés. Le cas ferroviaire examine le télégraphe comme composant d'un système plus large~: comment une infrastructure de transmission s'insère dans un système de transport, quelles formes organisationnelles cette insertion engendre, et quel rapport le circuit informationnel entretient avec le substrat énergétique du système qu'il coordonne.


\subsection*{Le télégraphe dans le chemin de fer}\label{subsec:railroad}

\paragraph{La collision du Western Railroad (1841)}\mbox{}\\
C'est à partir du réseau ferroviaire américain et britannique, entre 1840 et 1907, que \textcite{chandler_visible_1977} a construit sa théorie de la firme multidivisionnelle et \textcite{beniger_control_1986} sa théorie de la société de l'information. Le phénomène n'était pas limité aux pays anglo-saxons~--- les réseaux prussiens, le PLM français\footnote{La Compagnie des chemins de fer de Paris à Lyon et à la Méditerranée, fondée en 1857, exploitait le plus grand réseau français avec plus de 10\,000~km de lignes et utilisait le télégraphe pour la coordination du trafic dès les années~1860.}, le réseau colonial indien connaissaient des problèmes de coordination comparables~--- mais c'est le cas américain que Chandler et Beniger ont analysé avec une profondeur qui n'a pas d'équivalent ailleurs, et c'est sur cette analyse que nous nous appuyons. Le revisiter, avec les corrections que lui a apportées \textcite{schwantes_train_2019}, permet de voir non seulement ce que le télégraphe faisait dans ce système, mais aussi ce qu'il ne faisait pas~--- et surtout ce que le système tout entier impliquait comme rapport à l'énergie sans que personne, à l'époque ni depuis, ne le comptabilise.

Le problème se formulait simplement. L'application de la vapeur au transport à partir des années~1830 avait accéléré la vitesse et le volume des flux matériels dans des proportions que les dispositifs de coordination existants ne pouvaient pas absorber. \textcite{beniger_control_1986} donne à ce phénomène le nom de \emph{crise de contrôle}~: pour la première fois dans l'histoire, \dimmark{D6/D7}{*}%
des biens se déplaçaient plus vite que l'information les concernant. Le cas fondateur est celui du \emph{Western Railroad}, Massachusetts. Le 5~octobre~1841, sur 156~miles de voie unique entre Worcester et Albany, deux trains de passagers entrèrent en collision frontale~: deux morts, dix-sept blessés. L'enquête, reconstituée par \textcite[p.~183]{salsbury_state_1967}, identifia la cause~: six trains circulaient quotidiennement, nécessitant neuf croisements sur voie unique, sans signaux standardisés ni communication en temps réel. La programmation du trafic reposait sur des horaires imprimés dont l'exécution supposait une ponctualité parfaite~--- condition que la vapeur, par sa puissance même, rendait impossible à garantir\footnote{La ponctualité parfaite est une hypothèse de l'horaire~; la vapeur est une réalité physique dont la variabilité défie l'horaire. La contradiction est structurelle, pas accidentelle.}.

La crise de contrôle n'était pas un accident local. Elle se reproduisait partout où le réseau s'étendait au-delà de la portée de la supervision directe. Un surintendant de ligne, écrivait Daniel McCallum\footnote{Daniel Craig McCallum (1815--1878), ingénieur et administrateur d'origine écossaise. L'\emph{Erie Railroad}, achevé en 1851, reliait l'Hudson au lac Érié sur 483~miles à travers le sud de l'État de New York~--- à l'époque le plus long chemin de fer du monde.} en 1856 dans son rapport au conseil d'administration de l'\emph{Erie Railroad}, pouvait donner son attention personnelle à une ligne de cinquante miles~; mais à cinq cents miles, les mêmes méthodes devenaient \og \emph{entirely inadequate}\fg{} \parencite[p.~21--22]{chandler_strategy_1962}. Ce n'était pas une question de degré~: c'était un seuil qualitatif, au-delà duquel le fonctionnement du système exigeait non pas plus de la même chose, mais autre chose. Aucune organisation antérieure, notait \textcite[p.~94]{chandler_visible_1977}, n'avait nécessité une coordination aussi constante sur des centaines de miles~--- ni les armées, ni les manufactures, ni les administrations d'État.

\paragraph{McCallum, les six formes informationnelles et l'Erie Railroad}

La réponse à la crise ne fut pas le télégraphe. Elle fut un ensemble ordonné de formes informationnelles, déployées séquentiellement, dont le télégraphe n'était que la quatrième. En reprenant la chronologie établie par \textcite[ch.~5--6]{beniger_control_1986} et \textcite[p.~30, p.~148]{chandler_henry_1956}, on distingue six formes successives. Les \emph{rule books} et les horaires, d'abord, codifiaient les règles d'exploitation dans un document distribué à chaque agent~: information figée, exécutée mécaniquement, qui fonctionnait tant que les conditions réelles ne déviaient pas du plan. La hiérarchie managériale de McCallum, ensuite, instaurée sur l'\emph{Erie Railroad} à partir de 1854, créait le premier système hiérarchique formel de collecte, de traitement et de communication de l'information~: un organigramme en arbre, des uniformes codés par subdivision et par grade, une architecture de traitement~--- pas un simple signal\footnote{\textcite{yates_control_1989} montre que ce \og software organisationnel\fg{} comprenait aussi des formulaires standardisés, un classement vertical des documents, et des technologies de duplication qui rendaient l'information reproductible à l'intérieur de la firme.}. Le reporting systématique, troisièmement, rendait le réseau \emph{observable} depuis le centre~: rapports horaires, quotidiens, mensuels, compilés par huit clercs à temps plein \parencite[p.~148]{chandler_henry_1956}. Le télégraphe, quatrièmement, ajoutait la communication en temps réel~: le premier usage opérationnel documenté, sur l'\emph{Erie}, date de 1851, lorsque le surintendant Minot\footnote{Charles Minot (1810--1866), surintendant de la division de l'\emph{Erie Railroad}. L'épisode de 1851 est le premier usage documenté du télégraphe pour le contrôle opérationnel du trafic ferroviaire.} retarda un train à quatorze miles de distance par un message télégraphique \parencite{schwantes_train_2019}. La comptabilité analytique, cinquièmement, transformait les données brutes en indicateurs décisionnels~: Thomson\footnote{John Edgar Thomson (1808--1874), président du \emph{Pennsylvania Railroad} de 1852 à 1874. Le \emph{Pennsylvania} était le principal concurrent de l'\emph{Erie} et devint, sous sa direction, la plus grande entreprise privée du monde.}, sur le \emph{Pennsylvania Railroad}, introduisit le coût par tonne-mile~; Fink\footnote{Albert Fink (1827--1897), ingénieur d'origine allemande, vice-président du \emph{Louisville and Nashville Railroad}. Sa comptabilité analytique décomposait les coûts d'exploitation par segment de ligne et par type de service.} développa une comptabilité analytique complète \parencite[p.~231--232]{beniger_control_1986}\footnote{Mais \textcite{thompson_myth_1991} montre que les hypothèses de coûts fixes sur lesquelles reposait cette comptabilité étaient erronées~--- preuve que l'on peut innover dans la forme organisationnelle tout en se trompant sur le contenu analytique.}. Enfin, la standardisation infrastructurelle inscrivait l'information dans la matière elle-même~: écartement national en 1886, wagons standardisés en 1867, signaux et \emph{interlocking}\footnote{L'\emph{interlocking} est un dispositif mécanique qui rend physiquement impossible l'activation simultanée de signaux contradictoires~--- par exemple, un signal de voie libre et un aiguillage dirigeant un train vers une voie occupée. L'information de sécurité est inscrite dans le mécanisme lui-même.} en 1875, fuseaux horaires en 1883 \parencite[p.~235--236]{beniger_control_1986}.

\dimmark{D3/D10}{+}%
Ces six formes ne sont pas du même ordre. Les unes~--- le \emph{rule book}, la hiérarchie managériale, le reporting, la comptabilité analytique~--- sont des technologies organisationnelles dont le coût de diffusion est faible~: du \emph{software} social, reproductible au prix de l'impression et de la formation. Les autres~--- le télégraphe et la standardisation infrastructurelle~--- sont du \emph{hardware} matériel dont les coûts fixes sont élevés. \textcite{kooij_telegraph_2015}, dans son analyse micro-économique de l'innovation dans les réseaux de communication, formalise cette distinction sous le nom de \emph{communication engines}~: des artefacts techniques dont l'adoption exige un investissement en capital que toutes les compagnies ne pouvaient pas consentir. L'installation de la ligne télégraphique sur le \emph{Great Western Railway} coûtait 3\,270~livres sterling pour treize miles~; un \emph{rule book} ne coûtait que le prix de l'impression. La capitalisation de la compagnie devenait ainsi une variable médiatrice de l'adoption~: les moins capitalisées pouvaient s'organiser mieux, pas s'équiper. \textcite{schwantes_train_2019} l'a confirmé en corrigeant le récit de Beniger~: même en 1861, seules les grandes lignes utilisaient le télégraphe~; l'adoption massive vint après la Guerre civile, portée par des chocs exogènes~; et le \emph{Hours of Service Act} de 1907 accéléra \dimmark{D7/D4}{*}%
le passage au téléphone.

\paragraph{Le block system britannique}

Le cas britannique renversait la perspective. En Amérique, le problème était la distance~--- cinq cents miles de voie unique en territoire continental~--- et le télégraphe avait été un \og mariage immobilier\fg{} avec les chemins de fer, les lignes longeant les voies pour des raisons de commodité, pas de conception. En Grande-Bretagne, le problème était la densité~--- davantage de trains par mile de voie~--- et le couplage entre télégraphe et chemin de fer était intentionnel dès l'origine. \textcite{hubbard_cooke_1965} rappelle que William Fothergill Cooke avait conçu son télégraphe \emph{pour} les chemins de fer dès 1836~; en 1842, le pamphlet \emph{Telegraphic Railways} proposait formellement le \emph{block system}~--- le principe selon lequel un tronçon de voie ne pouvait être occupé que par un seul train à la fois, la vérification étant assurée par le télégraphe\footnote{Le premier \emph{block working} effectif date probablement de 1839 au \emph{Clay Cross Tunnel}, mais la datation exacte est disputée.}. Le \emph{Railway Regulation Act} de 1840 créa un organe de contrôle public, le \emph{Board of Trade}, dont les inspecteurs recommandèrent le \emph{block system} pendant quatre décennies avant qu'il ne devînt obligatoire. L'écart de quarante-sept ans entre l'invention du \emph{block system} et son adoption obligatoire en 1889~--- après la catastrophe d'Armagh, quatre-vingts morts~--- était le résultat d'une défaillance de marché au sens classique~: \dimmark{D2/D6}{+}%
le coût du télégraphe était privé, mais le bénéfice de la sécurité était en grande partie public~: les compagnies capturaient certes une part du bénéfice sous forme de réputation et d'évitement de procès, mais le risque de collision affectait d'abord les passagers, et la compagnie qui investissait dans la sécurité ne pouvait pas en facturer le prix à ses concurrentes qui ne le faisaient pas\footnote{Quinze mille instruments de télégraphe à aiguilles étaient en service sur le réseau britannique à la fin du siècle, certains jusqu'aux années~1930.}.

Ce que le cas ferroviaire faisait apparaître avec une netteté particulière, c'était la nature du problème de coordination. On peut le reformuler, avec un anachronisme assumé, comme un \emph{problème hayekien}\footnote{McCallum ne connaissait évidemment pas Hayek, dont l'article sur l'usage de la connaissance dispersée date de 1945. La reformulation est analytique~: elle éclaire la structure du problème, pas son histoire.}~: de l'information dispersée parmi des centaines d'agents répartis sur des centaines de miles devait être collectée, traitée et transformée en décisions en temps quasi réel \parencite{hayek_use_1945}. Mais le mécanisme de coordination que Hayek avait identifié comme le plus efficace~--- le système de prix~--- ne semblait pas adapté. Trois raisons structurelles l'en empêchaient~: la temporalité d'abord~--- le prix synthétise une information sur des jours ou des semaines, là où la coordination ferroviaire exigeait des décisions à la minute~; la rivalité physique absolue ensuite~--- deux trains sur la même voie signifient collision, pas compétition~; l'interdépendance séquentielle enfin~--- les agents ferroviaires sont enchaînés dans une séquence, pas en concurrence parallèle sur un marché. La réponse historique, telle que la documentent \textcite{chandler_visible_1977} et \textcite{beniger_control_1986}, ne fut pas le prix mais la hiérarchie informationnelle~--- ce que \textcite{coase_nature_1937} allait théoriser comme l'émergence de la firme là où la coordination marchande coûte plus cher que la coordination hiérarchique\footnote{La taxonomie de \textcite{hayek_use_1945}~--- la distinction entre savoir général et savoir des circonstances particulières de temps et de lieu~--- se projetait néanmoins avec exactitude sur les six formes informationnelles~: le \emph{rule book} et la standardisation codifient du savoir général~; le reporting et le télégraphe transmettent le savoir des circonstances~; la hiérarchie de McCallum est le dispositif qui articule les deux~; et la comptabilité analytique de Thomson est une tentative de construire l'équivalent du prix hayekien à l'intérieur de la firme.}. Autrement dit, le \emph{problème hayekien}~--- l'agrégation d'information dispersée~--- admettait des solutions non hayekiennes, et ces solutions étaient elles-mêmes des techniques de l'information dont le coût et la disponibilité déterminaient quelle forme organisationnelle était praticable. Le cas ferroviaire forçait ainsi la question qui allait traverser toute l'histoire des techniques de l'information~: quand le système de prix ne suffit pas à coordonner des flux matériels complexes, \dimmark{D6/D7}{*}%
quelles formes organisationnelles émergent~--- et à quel coût~?

\paragraph{Le bilan énergétique du réseau}

Or ce circuit informationnel avait un rapport au substrat énergétique que personne, à l'époque, ne pensait à mesurer~--- parce que l'asymétrie des ordres de grandeur le rendait inposable comme question. Le dispositif de \emph{process} de la division~--- les huit clercs de McCallum, les piles Daniell du télégraphe, l'éclairage du bureau~--- mobilisait au plus l'équivalent de quelques centaines de watts à quelques kilowatts de puissance dissipée en continu. En face, chaque locomotive en service développait plusieurs centaines de kilowatts, et le parc d'une division en comptait plusieurs dizaines. \dimmark{0-TER}{+}%
Quelle que soit la base de comparaison retenue, le rapport entre la puissance du circuit informationnel et la puissance du système qu'il pilotait restait inférieur de plusieurs ordres de grandeur. L'asymétrie était telle que la question du coût énergétique du circuit ne pouvait pas se poser~: l'information était, énergétiquement, gratuite.

Ce n'est pas un résultat anodin. À notre connaissance, il n'existe pas d'étude dédiée quantifiant les conséquences énergétiques des innovations informationnelles dans les chemins de fer\footnote{Résultat d'une recherche systématique sur la base Elicit (mars 2026), portant sur cinquante articles, dont dix retenus pour leur pertinence. Aucun ne traite du bilan énergétique des dispositifs informationnels ferroviaires en tant que tel.}. \textcite{field_magnetic_1992} a montré que le télégraphe contribuait à la productivité ferroviaire par la gestion serrée du capital, pas par la transmission des prix~; \textcite{fishlow_productivity_1966} a estimé que l'utilisation des capacités expliquait un sixième à un tiers de la croissance de la productivité du secteur~; \textcite{hawke_railways_1970}, cité par \textcite[p.~123--124]{miller_railway_2003}, a évalué la réduction des inventaires permise par les chemins de fer à un montant négligeable~--- entre zéro et un dixième de pour cent du revenu national en 1865. Aucun de ces travaux ne fait le dernier pas vers le bilan énergétique. Le lien entre les innovations informationnelles et la consommation d'énergie du réseau reste un trou dans la littérature~--- un trou qui s'explique par la spécialisation disciplinaire~: les historiens de la technologie ne comptent pas l'énergie, les économistes de l'énergie ne voient pas l'information, et les économistes de l'innovation traitent l'information comme analytiquement distincte de la matière.

\dimmark{D3}{+}%
\textcite{field_magnetic_1992} propose un cadre qui modifie la hiérarchie des effets. Il distingue les données de \emph{prix} --- cours du coton, actions, nouvelles politiques --- des données de \emph{quantité} --- localisation, direction et vitesse d'actifs physiques, état des stocks. Les premières conservent une valeur économique même avec un retard de quelques jours~; les secondes ne valent rien une fois que le train est arrivé ou que la viande a pourri. Le contrôle logistique en temps réel exige un dispositif à faible latence~; pour la diffusion des prix, le courrier postal eût presque suffi. Cette distinction conduit Field à réévaluer la contribution du télégraphe~: sa valeur sociale tient moins à la convergence des prix sur les bourses de matières premières qu'à l'économie de capital physique qu'il rend possible dans les systèmes de transport. Le réseau ferroviaire américain en 1890 était à 73\,\% en voie unique~; sans le télégraphe, il eût fallu doubler l'ensemble pour éviter les collisions, à un coût estimé de 957~millions de dollars --- contre 70 à 147~millions pour l'ensemble du capital télégraphique, soit un rapport de substitution de un pour sept à un pour treize. Le même raisonnement s'applique à l'industrie de la viande réfrigérée, où la périssabilité du produit interdit le stockage tampon et rend le routage en temps réel indispensable --- tandis que la sidérurgie, qui stocke cinq à trente mois de minerai, n'avait pas besoin du télégraphe pour fonctionner. La variable discriminante n'est donc pas la taille de la firme ni la complexité du réseau, mais le coût du stockage alternatif~: quand ce coût est faible, l'information optimise~; quand il est prohibitif, elle est une condition de possibilité du système.

\dimmark{D5/0-TER}{?}%
On peut en revanche conjecturer la structure du lien entre information et énergie dans ce système. Un réseau mieux coordonné consomme moins de charbon par tonne-mile~--- mais il rend aussi possible l'expansion du réseau lui-même, de trois mille miles en 1840 à plus de deux cent cinquante mille en 1900, et donc l'augmentation massive de la consommation agrégée. Si les deux effets coexistent~--- efficacité unitaire accrue et consommation totale en hausse~---, le schéma relève de ce que \textcite{jevons_coal_1865} avait identifié pour le charbon. L'hypothèse d'un \emph{effet de Jevons ferroviaire} reste une conjecture~; elle est formulée ici pour être testée au chapitre~3.

Ce cas n'est pas un septième cas de la trajectoire informationnelle~; il est le premier test du cadre analytique sur un système complet. Les dimensions que la section précédente avait identifiées isolément~--- le monopole naturel, la structure du capital, la coordination marchande, le changement technique~--- apparaissent ici en interaction. Le monopole du réseau crée un problème de coordination que le prix ne résout pas~; la coordination hiérarchique qui émerge en réponse exige un capital informationnel dont la structure diffère du capital matériel~; et le rapport de ce circuit au substrat énergétique est masqué par l'asymétrie des ordres de grandeur. Ce qui frappe, rétrospectivement, c'est la stabilité apparente du schéma~: au dix-neuvième siècle, l'énergie crée le problème~--- la vapeur dépasse la capacité de coordination~--- et l'information le résout, à un coût énergétique invisible. Le chapitre~3 examinera si, au vingt et unième siècle, ce rapport commence à s'inverser.

\emergence{D3/D7/0-TER~--- Le cas ferroviaire fournit trois résultats propres. Premièrement, l'information ne se réduit pas au télégraphe~: elle circule dans un circuit de six formes séquentielles dont le télégraphe n'est que la quatrième. Deuxièmement, il existe un seuil qualitatif au-delà duquel l'information cesse d'être un facilitateur pour devenir une condition de fonctionnement du système~--- ce que l'on pourrait nommer un \emph{seuil de McCallum}. Troisièmement, le rapport entre la puissance du circuit informationnel et celle du système coordonné est de plusieurs ordres de grandeur en faveur du second, ce qui rend le couplage analytiquement invisible. \questionch{2}{D3/0-TER~--- Cet écart se réduit-il avec les phases successives de la trajectoire~?}}

\emergence{D6/D7~--- Le cas ferroviaire force la question hayekienne~: quand le prix ne coordonne pas, quelles formes alternatives émergent~? \questionch{2}{D6/D7~--- la frontière firme-marché est-elle déterminée par les techniques d'information disponibles~?}}

\emergence{D5/0-TER~--- Le \og Jevons ferroviaire\fg{} est conjecture, pas preuve. L'information améliore l'efficacité unitaire ET permet l'expansion qui augmente la consommation agrégée. \questionch{3}{D5/0-TER~--- Le même schéma opère-t-il pour le compute AI~?}}



%%% --- BILAN §1.2 ---

\subsection*{Bilan de la section~: le télégraphe au prisme des dix dimensions}\label{subsec:bilan_telegraph}

Le tableau~\ref{tab:telegraphe-dimensions} récapitule ce que la traversée du cas télégraphe~--- réseau autonome puis cas imbriqué ferroviaire~--- a permis d'établir pour chacune des dix dimensions de la grille de lecture. Trois niveaux de certitude sont distingués~: \textbf{*}~(régularité documentée indépendamment par des sources distinctes), \textbf{+}~(mécanisme proposé, compatible avec les observations mais non testé), \textbf{?}~(piste ouverte, pas encore d'argument construit).

\begin{table}[p]
\caption{Le télégraphe au prisme des dix dimensions~: acquis, incertitudes et questions ouvertes}\label{tab:telegraphe-dimensions}
\centering\footnotesize
\renewcommand{\arraystretch}{1.15}
\begin{tabularx}{\textwidth}{@{}cXcp{4.8cm}@{}}
\toprule
\textbf{Dim.} & \textbf{Acquis du cas télégraphe} & & \textbf{Question pour la suite} \\
\midrule
D1 & Signal rival (un message occupe le circuit), information non rivale (le cours du coton est diffusable). & ? & Le numérique supprime-t-il la rivalité du canal~? \\[3pt]
D2 & Sept modèles (UK, US, FR, DE, Inde, Japon, Chine) convergent vers le monopole. Documenté indépendamment (Nonnenmacher, Kieve, Bertho, Headrick, Wu). & * & Trait des réseaux de communication, ou des infrastructures lourdes en général~? \\[3pt]
D3 & Substitution du capital linéaire (fil, poteaux) au travail (opérateurs Chappe). Quatre pays. Morse économise du cuivre au prix d'un opérateur qualifié. Le cas ferroviaire révèle deux registres~: un \emph{software} organisationnel (\emph{rule book}, hiérarchie, comptabilité) à coût faible, un \emph{hardware} matériel (télégraphe, standardisation) à coûts fixes élevés. & * & La substitution K/L de l'information est-elle du même type que celle des machines~? \\[3pt]
D4 & Opérateurs Morse~: premiers travailleurs de l'information. Huit clercs chez McCallum. Qualification non analysée. & + & L'information qualifie-t-elle ou déqualifie-t-elle~? \\[3pt]
D5 & Prix documentés (18,44~F $\rightarrow$ 1~F~; 1,55\$ $\rightarrow$ 0,40\$). Pas d'analyse d'élasticité. & ? & Effet de rebond par les tarifs~? \\[3pt]
D6 & L'information synchrone \emph{crée} le marché (Steinwender, 4~réplications). La codifiabilité filtre l'effet (Juhász). & * & La codifiabilité module-t-elle \emph{toutes} les dimensions~? \\[3pt]
D7 & Le système de prix ne coordonne pas les flux ferroviaires~: un \emph{problème hayekien} qui appelle des solutions non hayekiennes. La hiérarchie informationnelle émerge (McCallum, Chandler). Six formes séquentielles, le télégraphe n'est que la quatrième. & * & La frontière firme-marché dépend-elle des techniques d'information disponibles~? \\[3pt]
D8 & Câble transatlantique~: investissement $\rightarrow$ faillite $\rightarrow$ recapitalisation. Non formalisé. & + & Même structure que \emph{railway mania} et dot-com~? \\[3pt]
D9 & La matérialité conditionne le réseau à deux niveaux~: extraction (cuivre, gutta-percha, bois~; domination britannique) et déploiement (le réseau US trop instable ne peut faire tourner Hughes ou Baudot, qui restent cantonnés à l'Europe). & + & Même pattern pour semi-conducteurs et terres rares~? \\[3pt]
D10 & Épisodes isolés (sémaphore électrique, \emph{needle telegraph}). Pas d'argument général. & ? & Constante des réseaux de communication~? \\[3pt]
\midrule
0-TER & Asymétrie de plusieurs ordres de grandeur entre la puissance du circuit informationnel et celle du système coordonné. Matérialité extractive occultée. Aucune quantification dans la littérature. Effet de Jevons ferroviaire conjecturé. & +/? & Cet écart se réduit-il avec les phases successives de la trajectoire~? À quel seuil le couplage information-énergie devient-il visible~? \\
\bottomrule
\end{tabularx}
\renewcommand{\arraystretch}{1.0}
\end{table}

Quatre dimensions ressortent du cas télégraphe avec une évidence processuelle forte, documentée par des sources indépendantes~: la convergence vers le monopole (D2), la substitution capital-travail et la dualité \emph{software}/\emph{hardware} du capital informationnel (D3), le rôle constitutif de l'information synchrone dans la formation des marchés (D6) et l'émergence de la hiérarchie informationnelle comme réponse à la défaillance du système de prix dans le cas ferroviaire (D7). Le résultat de Steinwender sur D6 s'appuie en outre sur une stratégie d'identification économétrique répliquée dans quatre géographies. Trois dimensions sont présentes mais incomplètes~: le travail (D4), le financement (D8) et la géographie matérielle (D9)~--- les observations existent, la formalisation manque. Deux restent au stade de la piste~: la rivalité du signal (D1) et l'élasticité de la demande aux tarifs (D5). La dimension énergie (0-TER) est transversale~: elle apparaît sur trois registres~--- extraction (gutta-percha, poteaux, cuivre), invisibilité par les ordres de grandeur (ratio 1:2\,000), conjecture de Jevons ferroviaire~--- sans constituer un résultat mesurable. C'est ce silence empirique, et l'explication structurelle qu'en donne le ratio, que le chapitre~3 devra instruire.


% =============================================================================

% =============================================================================
% §1.3 Le calcul mécanisé (1820–1914)
% =============================================================================
\clearpage
\section{Le calcul mécanisé~: l'outil isolé (1820--1914)}\label{sec:calculating}
% =============================================================================

Pendant que le télégraphe s'électrifie et se déploie le long des voies ferrées, une autre fonction informationnelle croît dans l'ombre~--- mais elle ne s'électrifie pas, ne sort pas de l'organisation qui la commande, et ne crée aucun marché. Les mêmes décennies qui voient Morse breveter son télégraphe (1837), Cooke installer sa ligne sur le \emph{Great Western Railway} (1839) et le câble transatlantique unifier les cours du coton (1866) voient aussi Thomas de Colmar breveter l'Arithmomètre (1820), Babbage concevoir l'\emph{Analytical Engine} (1834) et le \emph{Census Bureau} de Washington crouler sous le dépouillement du recensement de~1880. Les ateliers de mécanique de précision sont parfois les mêmes~: Bréguet, à Paris, fabrique des appareils télégraphiques pour l'administration des lignes et des instruments scientifiques pour les observatoires\footnote{La maison Bréguet, fondée par Abraham-Louis Bréguet en~1775 comme atelier d'horlogerie, fabrique sous son petit-fils Louis-François-Clément les appareils télégraphiques commandés par l'administration française à partir des années~1840 \parencite{bertho_histoire_1984}. Le même atelier produit des chronomètres de marine, des instruments de mesure électrique et des pièces d'horlogerie de haute précision. La mécanique fine du télégraphe et celle de la calculatrice sortent du même monde artisanal.}. Mais les trajectoires divergent~: le télégraphe s'électrifie dès~1837 parce que sa fonction~--- transmettre un signal à distance~--- exige un vecteur que seule l'électricité peut fournir~; le calcul reste mécanique, puis électromécanique, jusqu'à l'arrivée de l'électronique en~1946~: sa fonction~--- combiner des nombres~--- peut être accomplie par des engrenages, des leviers et des doigts humains, puis par des relais électromagnétiques, sans appeler le passage à un substrat radicalement différent. Cette asymétrie~--- la transmission s'électrifie un siècle avant le calcul~--- n'est pas un accident de calendrier~; elle tient à ce que la transmission a une demande externe (le commerce, l'armée, les chemins de fer) qui finance le réseau, tandis que le calcul a une demande interne (la comptabilité, l'actuariat, l'astronomie) qui ne justifie que des outils isolés.

Le calcul est, au dix-neuvième siècle, l'un des travaux les plus répandus et les moins visibles de l'économie industrielle. Les compagnies d'assurance, les administrations fiscales, les chemins de fer emploient des commis par centaines dont le métier est d'additionner des colonnes, de vérifier des sommes, de recopier des résultats d'un registre à l'autre~: un travail répétitif, qualifié, et entièrement interne, dont les chiffres ne sortent pas de l'organisation qui les a commandés. Au moment où les premières machines apparaissent pour l'accélérer, cette activité n'a pas de nom\footnote{Ce n'est qu'en~1962 que \textcite{machlup_production_1962} constitue le traitement de l'information en objet d'analyse économique, sous le nom de \og production de la connaissance\fg{}. Le calcul de bureau du dix-neuvième siècle n'apparaît dans aucune des catégories comptables ou statistiques de l'époque.}. La tension, pourtant, est ancienne. En~1757, le mathématicien Alexis-Claude Clairaut\footnote{Alexis-Claude Clairaut (1713--1765), mathématicien français, membre de l'Académie des sciences à dix-huit ans. Il avait déjà résolu le problème de la figure de la Terre (1743) et clarifié la théorie de la Lune. La prédiction du retour de Halley était pour lui une confirmation expérimentale de la gravitation universelle.} voulut prédire le retour de la comète de Halley en résolvant numériquement le problème à trois corps~: le calcul des perturbations gravitationnelles de Jupiter et de Saturne sur l'orbite de la comète, un test décisif de la mécanique newtonienne que l'astronomie attendait depuis un demi-siècle. Le travail dépassait les forces d'un seul homme~; Clairaut recruta l'astronome Jérôme Lalande\footnote{Jérôme Lalande (1732--1807), astronome français, professeur au Collège de France et directeur de l'Observatoire de Paris. Il sera l'un des astronomes les plus célèbres de son temps~; en~1757, à vingt-cinq ans, il est encore un jeune collaborateur de Clairaut.} et la mathématicienne Nicole-Reine Lepaute\footnote{Nicole-Reine Lepaute (1723--1788), née Étable, épouse de l'horloger Jean-André Lepaute. Mathématicienne et astronome, elle calcule des éphémérides, des tables du Soleil et contribue à la prédiction d'éclipses. Son rôle dans le calcul de la comète sera effacé par Clairaut lui-même (voir la note suivante).}, et les trois s'installèrent au Palais du Luxembourg pour décomposer le problème en pas élémentaires calculés à la plume. Lalande écrira~: \og \emph{Pendant plus de six mois, nous calculâmes depuis le matin jusqu'au soir, quelquefois même à table}\fg{} \parencite[p.~678]{lalande_bibliographie_1803}\footnote{C'est le premier cas documenté de division du travail de calcul entre des personnes recrutées spécifiquement pour cette tâche \parencite[ch.~1]{grier_when_2005}. Clairaut, en publiant sa \emph{Théorie des comètes} (1760), omit de mentionner Lepaute parmi les calculateurs. Lalande lui-même, dans la \emph{Bibliographie astronomique} qu'il fit paraître quarante ans plus tard, rétablit le rôle de Lepaute et critiqua publiquement l'omission~; la querelle avait brisé son amitié avec Clairaut \parencite{lalande_bibliographie_1803}. L'épisode anticipe un trait qui traversera toute la section~: le travail de calcul est le premier à être effacé, et l'effacement est un acte~--- non un oubli~--- qui rend invisibles, en priorité, les contributrices.}. En novembre~1758, Clairaut annonça le retour de la comète pour le 13~avril de l'année suivante. La comète passa à son périhélie le 13~mars~1759, un mois avant la date prédite~: trente et un jours d'écart sur une orbite de soixante-seize ans, soit une erreur inférieure à cinq pour cent. Le calcul avait fonctionné. Il prouvait que les lois de Newton suffisaient à prédire le mouvement d'un objet perturbé par deux planètes, et que six mois de labeur arithmétique pouvaient produire un résultat que l'observation confirmait. D'Alembert, coéditeur de l'\emph{Encyclopédie} et rival de Clairaut à l'Académie des sciences, s'empara de l'erreur résiduelle pour dénoncer ce qu'il appelait \og l'esprit de calcul\fg{} et juger l'entreprise \og plus laborieuse que profonde\fg{}\footnote{La source primaire est \textcite[t.~1, douzième mémoire]{alembert_opuscules_1761}. La querelle est rapportée par \textcite[ch.~1]{grier_when_2005}. Cette dépréciation du calcul \emph{numérique} (résolution itérative, parfois en équipe) au profit du calcul \emph{analytique} (équations différentielles closes, attribuables à un seul auteur) n'est pas une opinion isolée de d'Alembert~: elle structure une part de la culture mathématique académique du XIX\textsuperscript{e}~siècle, jusqu'à la réhabilitation des méthodes numériques au XX\textsuperscript{e}~--- précisément quand la mécanisation du calcul leur donnera un substrat industriel.}. La formule est révélatrice par ce qu'elle refuse de voir~: Clairaut avait divisé le travail intellectuel entre trois personnes, conçu un protocole de calcul séquentiel, et obtenu un résultat que la nature avait confirmé. Mais pour d'Alembert, le labeur numérique ne pouvait pas constituer une méthode. C'est cette tension entre le calcul comme travail et le calcul comme pensée qui traverse la présente section~: qui calcule, à quel coût, et selon quelle organisation~?


\subsection{La demande~: un travail en expansion}\label{subsec:demande_calcul}

En~1774, à Hanovre, un certain Anton Dies, registraire du fonds de pension pour veuves du Calenberg, écrivit au comité de la Trésorerie pour demander une augmentation. Il tenait les comptes du fonds depuis vingt ans, calculait les intérêts composés année par année, vérifiait les tables actuarielles, gérait la correspondance avec les souscripteurs. Lorsqu'une crise des pensions éclata en~1780, le volume de calcul explosa~: le fonds comptait alors plus de 3\,700~couples et 700~veuves, et la réforme exigeait des projections constamment révisées. Dies souffrait d'une paralysie du bras droit. Son successeur Eisendecker estimait sa charge à trois fois celle de son prédécesseur. Un troisième commis, Rath, secondé pour absorber le surplus, sombra dans ce que sa femme qualifia de \og \emph{Gemüthskrankheit}\fg{}, un désordre mental. L'administration lui accorda une gratification de onze thalers\footnote{\textcite{rosenhaft_hands_2003} a reconstitué ce cas à partir des archives de la Landschaft de Hanovre (HStAHann, Dep.~7B et Hann.~93). Les détails biographiques et les montants proviennent de la correspondance administrative entre le comité de la Trésorerie et la Couronne, 1766--1785.}. \dimmark{D5}{?}%
Le cas du Calenberg est modeste, mais il rend visible ce que les comptes d'exploitation du siècle suivant ne montreront jamais~: le calcul est un travail qui use le corps et l'esprit, et dont le volume croît avec la complexité de l'organisation qu'il sert. Reformulée dans le vocabulaire marshallien~--- que les administrateurs hanovriens ne mobilisaient évidemment pas~---, la structure du cas peut être résumée ainsi~: la demande de calcul est dérivée (on calcule pour assurer, pour dénombrer, pour acheminer, jamais pour calculer), peu sensible au prix de la machine~--- un commis coûte le même salaire que la machine existe ou non, et l'investissement dans une machine ne réduit pas le calcul que l'organisation doit produire~---, et croissante avec la complexité de l'organisation servie~: plus le nombre de souscripteurs, de polices, de wagons, de courses croît, plus le volume de calcul à exécuter croît. Ces propriétés sont qualitatives, reconstituées à partir d'un cas individuel~; elles ne prétendent pas à l'estimation d'élasticités, qui supposerait des séries de prix et de quantités dont les archives administratives hanovriennes ne disposaient pas. \dimmark{0-TER}{?}%
Ce que les sources documentent, c'est l'usure du corps et de l'esprit~: le bras paralysé de Dies, la \og \emph{Gemüthskrankheit}\fg{} de Rath, l'épuisement d'Eisendecker. Aucun poste budgétaire n'enregistre ces coûts, et cette absence n'est pas seulement empirique~: elle est constitutive. La catégorie même d'\emph{énergie} comme grandeur unifiée n'existait pas en~1780~--- elle ne se stabilisera qu'avec l'unification thermodynamique du milieu du XIX\textsuperscript{e}~siècle (Mayer, Joule, Helmholtz)~---, et son importation dans la pensée économique sous la forme d'un travail humain converti en équivalent énergétique restera marginale, depuis les tentatives de Sergei Podolinsky (1880) jusqu'à leur reprise tardive par Nicholas Georgescu-Roegen un siècle plus tard\footnote{Sur la généalogie de l'\emph{énergie} comme catégorie économique~--- de Podolinsky, Sacher, Bogdanov, Ostwald et Neurath jusqu'aux fondateurs de l'économie écologique au XX\textsuperscript{e}~siècle~---, on renverra à Martinez-Alier, \emph{Ecological Economics. Energy, Environment and Society} (Oxford, Blackwell, 1987), et à la reconstruction récente de \textcite{vianna_franco_history_2023}. La présente section mobilise quelques jalons de cette histoire sans en faire un objet pour lui-même.}. Que le calcul du Calenberg n'apparaisse pas comme \og dépense énergétique\fg{} ne tient donc pas à un oubli des comptables hanovriens~: le concept qui aurait permis de l'enregistrer comme tel n'était pas encore disponible. À l'invisibilité d'ordre de grandeur~--- le coût métabolique d'un commis est dérisoire devant le système qu'il sert~--- s'ajoute donc une invisibilité catégorielle plus profonde~: un coût qui n'a pas de nom n'a pas de poste.

\dimmark{D4}{?}%
Cette invisibilité n'est pas accidentelle. \textcite{downey_virtual_2001} a montré que les discussions sur la \og société de l'information\fg{}, contemporaines comme rétrospectives, effacent systématiquement le travail humain dans les réseaux d'information. Le travail de calcul en bureau n'a pas laissé de traces comparables aux registres commerciaux ou aux brevets d'invention~; les budgets des entreprises du dix-neuvième siècle, quand ils sont conservés, ne distinguent pas la ligne \og calcul\fg{} de la masse salariale administrative. \textcite{rosenhaft_hands_2003} suggère que cette invisibilité tient en partie au caractère \og \emph{low-tech}\fg{} du travail~: à la différence des télégraphistes ou des opératrices téléphoniques, dont le travail était visiblement lié à un appareil, les commis qui additionnaient des colonnes et recopiaient des totaux ne semblaient pas exercer une activité qui fût en prise avec la technique.

On peut néanmoins en mesurer la croissance. Le nombre de commis aux États-Unis passe de un à deux pour cent de la population active en~1870 à plus de dix pour cent en~1940 \parencite{cortada_before_1993}\footnote{La catégorie \og clerk\fg{} dans les recensements américains a connu plusieurs reclassifications entre~1880 et~1930. \textcite{cortada_before_1993} construit ses séries en harmonisant autant que possible les nomenclatures successives, mais les comparaisons longues sur sept décennies restent sensibles à ces reclassifications et sont à manier avec précaution.}. \textcite{rotella_transformation_1981} a établi, à partir des données du Census Bureau, la correspondance temporelle précise entre l'adoption des machines de bureau et le changement de composition par genre de l'emploi clérical entre~1870 et~1930. En Angleterre, \textcite{boot_salaries_1991} a montré que le nombre de clercs de la Bank of Scotland plus que tripla entre~1730 et~1880, alors que le nombre de guichetiers ne fit que doubler. \dimmark{D7}{+}%
L'écart de croissance entre le \og \emph{front office}\fg{} et le \og \emph{back office}\fg{} suggère un effet multiplicateur interne~: chaque augmentation du volume d'affaires engendrait un accroissement plus que proportionnel du travail de tenue des registres, de vérification et de correspondance. C'est dans cet écart que les premières machines vont s'insérer. Mais le fait remarquable est la lenteur de cette insertion~: la première machine capable d'effectuer les quatre opérations arithmétiques est brevetée en~1820, et il faudra attendre un demi-siècle pour que la demande soit suffisamment concentrée pour en absorber la production.

\emergence{D5/D6/0-TER~: le calcul mécanisé diffère structurellement du télégraphe sur trois dimensions. Premièrement, la demande de calcul est dérivée et peu sensible au prix de la machine (D5)~: là où la baisse du prix du télégramme augmente le trafic, la baisse du prix de la calculatrice ne crée pas de demande nouvelle de calcul~; c'est la complexité de l'organisation servie qui détermine le volume. Deuxièmement, les effets du calcul restent internes à l'organisation~: la calculatrice ne crée pas de marchés, ne synchronise pas de prix, ne rend pas possible des transactions qui n'auraient pas eu lieu sans elle (D6 négatif, par contraste avec le résultat de Steinwender pour le télégraphe). Troisièmement, l'énergie du calcul est métabolique~: elle passe par le corps du commis, et à ce titre elle est doublement invisible~: invisible dans les comptes parce qu'elle ne constitue pas un poste identifié, et invisible dans les ordres de grandeur parce qu'elle est dérisoire comparée au système qu'elle sert (0-TER). \questionch{2}{D5/D6~: l'effet des techniques de l'information sur les marchés dépend-il de la fonction exercée (transmission, stockage ou calcul)~? Si oui, la codifiabilité de Juhász n'est pas le seul filtre~: la nature même de la fonction ICT en est un autre.}}


\subsection{De l'Arithmomètre aux calculatrices industrielles}\label{subsec:arithmometre}

En~1820, Charles Xavier Thomas de Colmar\footnote{Thomas de Colmar (1785--1870), né Charles Xavier Thomas, directeur de la compagnie d'assurance Le~Phénix à Paris. Son brevet d'invention (n\textsuperscript{o}~1420, 18~novembre 1820, archives INPI) décrit un \og arithmomètre ou machine à calculer\fg{}. L'appareil est décrit dans le catalogue officiel de la \emph{Great Exhibition} de Londres en~1851, où il reçut une médaille, et dans la notice technique de \textcite{sebert_machines_1879}.}, directeur de la compagnie d'assurance Le~Phénix à Paris, déposa un brevet pour un appareil qu'il appelait \og arithmomètre\fg{}~--- une machine capable d'effectuer les quatre opérations arithmétiques par un mécanisme de tambours cannelés\footnote{Le tambour cannelé (\emph{Staffelwalze} ou \emph{Leibniz wheel}) est un cylindre portant neuf dents de longueur croissante, conçu par Leibniz pour son \emph{Stepped Reckoner} (1694). Le principe permet de réaliser la multiplication par rotation différentielle. Leibniz en donna une description dans les \emph{Miscellanea Berolinensia} \parencite{leibniz_machina_1710}. Le mécanisme resta sans descendance industrielle pendant plus d'un siècle, jusqu'à ce que Thomas de Colmar le reprenne dans un boîtier de dimensions compatibles avec un bureau.} dérivé du \emph{Stepped Reckoner} de Leibniz. Que la première machine à calculer produite en série fût l'\oe{}uvre d'un assureur n'est pas un hasard~: le calcul actuariel~--- tables de mortalité, primes, réserves~--- était l'un des rares domaines où le volume de calcul justifiait l'investissement dans un appareil coûteux et lent. L'Arithmomètre fonctionnait~; il fonctionnait assez bien pour que des observatoires et quelques administrations l'achètent à leur tour, et assez longtemps pour rester en service jusqu'à la Première Guerre mondiale \parencite{cortada_before_1993, sebert_machines_1879}\footnote{La machine n'était pas équipée d'imprimante~; pour pallier cette absence et garantir la fiabilité de la transcription des résultats, \textcite{sebert_machines_1879} recommandait dans son rapport à la Société d'encouragement pour l'industrie nationale d'avoir recours à \og deux employés, un qui tourne la manivelle et l'autre qui transcrit\fg{}. La calculatrice ne supprimait pas le travail de calcul~: elle imposait une division du travail à deux personnes autour d'elle, ce qui constitue l'un des premiers cas documentés de réorganisation du travail induite par une machine de bureau~--- pattern que l'on retrouvera, amplifié, autour des tabulatrices Hollerith au \S\,\ref{sec:mechanography}.}. Mais la diffusion fut d'une lenteur remarquable. \textcite{cortada_before_1993} estime que Thomas vendit \og \emph{possibly several hundred}\fg{} machines sur trois à quatre décennies~; \textcite{nordhaus_two_2007} retient un stock cumulé d'environ cinq cents machines avant~1865. La capacité technique précédait la demande~: la machine existait, mais le marché qui aurait pu l'absorber~--- des organisations dont les besoins de calcul fussent suffisamment volumineux et répétitifs pour justifier l'investissement~--- ne se constituerait qu'avec la croissance des compagnies d'assurance et des administrations statistiques dans la seconde moitié du siècle.

En~1834, Charles Babbage conçut l'\emph{Analytical Engine}, une machine dont l'architecture se distinguait radicalement de l'Arithmomètre\footnote{La machine à différences elle-même, sur laquelle Babbage avait travaillé depuis les années~1820, n'était pas conceptuellement nouvelle~: la même proposition avait été faite par un officier hessois, J.~H.~Müller, en~1786, sans que Babbage en eût connaissance \parencite[p.~42]{augarten_bit_1984}. L'idée d'une mécanisation du calcul polynomial par différences finies était disponible dans la culture mathématique européenne depuis près d'un demi-siècle~; ce qui manquait n'était pas le concept mais les conditions économiques de sa réalisation.}. L'appareil séparait deux organes~: le \emph{store}, un ensemble de colonnes de roues capables de retenir cinquante nombres de cinquante chiffres, et le \emph{mill}, une unité de calcul capable de les combiner selon quatre opérations. Les instructions devaient être codées sur des cartes perforées empruntées au métier Jacquard~; le même mécanisme servait à sélectionner les nombres stockés et à commander la séquence des opérations. L'\emph{Analytical Engine} ne fut jamais construit~; mais sa conception posait, pour la première fois dans un artefact technique, la distinction entre le stockage et le traitement de l'information comme deux fonctions architecturalement séparées\footnote{Ada Lovelace, dans ses notes de~1843 sur le mémoire de Menabrea, en décrivit les possibilités avec une précision qui en fait le premier document de programmation de l'histoire. \textcite{schaffer_babbage_1994} souligne un paradoxe~: attribuer l'intelligence à la machine occultait le travail humain qui restait nécessaire pour la concevoir, la programmer et en vérifier les résultats.}.

La machine à différences, en revanche, fut effectivement construite~--- mais en Suède, par un autre que Babbage. Georg Scheutz, juriste et éditeur de journal stockholmois, lut un compte rendu des travaux de Babbage dans l'\emph{Edinburgh Review} de~1834 et entreprit, avec son fils Edvard et l'appui de l'Académie royale des sciences de Suède, de réaliser l'appareil. Une première machine fonctionnait en~1844~; une version industrielle, capable de produire des stéréotypes d'imprimerie pour la composition automatique des tables, fut achevée en~1855 et exposée à l'Exposition universelle de Paris la même année. L'une de ces machines fut achetée par le \emph{General Register Office} britannique et y fonctionna entre~1859 et~1864 pour calculer les tables actuarielles du \emph{English Life Table}~\parencite[ch.~1]{goldstine_computer_1972}\footnote{La machine Scheutz est aujourd'hui conservée au Smithsonian. Elle constitue, à notre connaissance, le seul cas documenté d'utilisation industrielle d'une machine babbagéenne au XIX\textsuperscript{e}~siècle. \textcite{winston_media_1998} en fait l'illustration de ce qu'il appelle un \og prototype accepté en niche\fg{}~: un dispositif techniquement abouti et économiquement utile dans une demande étroite, mais qui n'entraîne pas la généralisation faute d'une nécessité sociale survenante de plus large portée. Les performances effectives de la machine, soulignent toutefois \textcite[p.~328]{warwick_laboratory_1995}, restèrent faibles, en raison notamment de défauts d'impression~--- ce qui n'invalide pas l'argument de niche réalisée mais en nuance la portée.}. \dimmark{D5}{?}%
Le cas Scheutz importe pour notre argument~: il établit que la machine à différences était techniquement réalisable et économiquement utile dans une niche, et que la non-émergence de l'industrie babbagéenne ne tient ni à un verrou technique ni à une impossibilité économique en soi, mais à l'étroitesse de la niche. Le \emph{General Register Office} avait besoin de tables de mortalité~; il ne lui en fallait pas plusieurs centaines par an. La machine Scheutz desservait une demande qui se satisfaisait d'un exemplaire unique. Babbage, lui, projetait un calculateur universel pour une demande qui n'existait pas encore. La distinction entre les deux projets n'est pas seulement technique~--- la machine à différences calcule des polynômes par différences finies tandis que l'\emph{Analytical Engine} aurait pu, par sa programmabilité altérable, traiter n'importe quel problème calculable\footnote{C'est précisément l'intuition de Babbage et Lovelace que reprendra Turing un siècle plus tard~: la programmabilité altérable~--- la capacité d'une machine à modifier son propre programme selon les résultats intermédiaires~--- distingue le calculateur de l'ordinateur \parencite[ch.~8]{winston_media_1998}. Cette distinction reste \og indistincte\fg{} (\emph{indistinct conception}) dans les carnets tardifs de Babbage, ce qui le maintient en deçà du seuil conceptuel franchi en~1936 par \textcite{turing_computable_1937}.}~--- elle est aussi économique~: une niche ne nourrit pas une industrie, et seule une demande répétée en volume élevé pouvait justifier l'investissement qu'exigeait une production sérielle de machines à calculer.

\dimmark{D8}{?}%
L'échec de Babbage n'est pas seulement technique~; il est économique~--- et le paradoxe est que Babbage lui-même avait formulé, dans son \emph{On the Economy of Machinery and Manufactures} (Londres, Charles Knight, 1832), l'une des premières théories systématiques de la mécanisation comme division du travail\footnote{Le \og principe Babbage\fg{}~--- décomposer une tâche complexe en sous-tâches de qualifications hétérogènes, afin de ne payer chaque qualification qu'au prix qu'elle exige~--- est exposé au chapitre~19 de l'\emph{Economy of Machinery}. Marx en reprendra l'argument dans \emph{Le Capital} (livre~I, ch.~14). Babbage disposait donc du cadre analytique qui aurait pu rendre intelligible son propre projet~--- mais ce cadre supposait précisément le marché qu'il n'avait pas.}. Le \emph{Difference Engine} n\textsuperscript{o}~1, conçu pour calculer des tables polynomiales par la méthode des différences finies, avait reçu environ 17\,470~livres sterling du gouvernement britannique entre~1823 et~1842, une somme comparable au coût d'un navire de guerre de l'époque \parencite{swade_difference_2001}. Le projet fut abandonné après que Babbage et son ingénieur Joseph Clement se brouillèrent sur la propriété des outils, et le gouvernement refusa de financer l'\emph{Analytical Engine}. Le contraste avec le câble transatlantique est instructif~: le câble suivait la séquence investissement, faillite, recapitalisation, succès~; la machine de Babbage suivait la séquence investissement, pas de produit, pas de recapitalisation, abandon. La différence tient à la demande~: en~1840, aucune organisation n'avait besoin de calculer en volume au point de justifier un investissement de cette ampleur~; cinquante ans plus tard, le Census Bureau américain en aurait besoin.

Ce que l'échec de Babbage confirmait, c'est ce que l'Arithmomètre suggérait par son succès modeste~: la machine à calculer ne pouvait se déployer que là où la demande de calcul était déjà organisée en flux régulier, et cette condition ne serait remplie, à grande échelle, que dans le dernier quart du siècle.

Le passage de l'Arithmomètre artisanal à la calculatrice industrielle s'accomplit entre~1875 et~1914 par une convergence technique qui mit fin au monopole de fait du tambour de Leibniz. En~1875, Frank Baldwin aux États-Unis et Willgodt Odhner en Russie développèrent indépendamment un mécanisme à disque à broches rétractables~--- plus compact, plus fiable, moins coûteux à fabriquer que le tambour cannelé. L'ajout du clavier, emprunté à la machine à écrire, transforma la calculatrice en un instrument de bureau ergonomique. En une génération, quatre \emph{dominant designs} s'imposèrent~: le Comptometer de Dorr Felt (1886, à clavier direct, sans manivelle), la machine Burroughs (1886, à clavier avec impression du résultat), les machines Odhner et Brunsviga (à manivelle, dominantes en Europe) \parencite{cortada_before_1993}.


\subsection{La diffusion des calculatrices et la structure du marché}\label{subsec:burroughs}

La croissance des ventes, documentée par les archives de la \emph{Burroughs Adding Machine Company} exploitées par \textcite{cortada_before_1993}, fut spectaculaire. L'\emph{American Arithmometer Company} (future Burroughs) vendit 286~machines en~1895 à un prix moyen de 223~dollars~; 1\,399 en~1900~; 5\,008 en~1904~; 15\,763 en~1909~--- le double de ce que William Burroughs pensait, en~1885, pouvoir vendre dans le monde entier. Les ventes annuelles passèrent de 4,9~millions de dollars en~1905 à 39,9~millions en~1916, soit un facteur huit en onze ans~--- une croissance largement supérieure à celle du produit national brut américain sur la même période. En Europe, Brunsviga vendit plus de vingt mille machines entre~1885 et~1912 et obtint 130~brevets allemands et 300~brevets étrangers~; le Millionaire suisse, capable de multiplication directe par lecture d'une table de multiplication mécaniquement stockée dans la machine, plaça 4\,600~machines entre~1894 et~1935 \parencite{cortada_before_1993}\footnote{La table de multiplication du Millionaire, gravée sur un bloc métallique mobile, anticipe matériellement la fonction de mémoire morte (\emph{ROM}) des calculateurs électroniques~: un savoir codifié, fixé dans la matière de la machine, et lu sans recalcul à chaque opération. Ce dispositif illustre que la fonction de stockage~--- ce que la présente thèse appelle STORE~--- n'est pas née avec l'électronique~; elle est constitutive de toute machine à calculer dès qu'elle dépasse l'opération atomique \parencite[ch.~8]{winston_media_1998}. Le principe des \og barèmes matérialisés\fg{}, c'est-à-dire l'inscription mécanique d'une table de multiplication dans la machine elle-même sous forme de chevilles en saillie sur des plaques, avait été développé indépendamment par Léon Bollée au Mans dès~1889~--- précédant Steiger qui ne breveta le Millionaire qu'en~1892~--- et présenté à la Société d'encouragement pour l'industrie nationale par \textcite{sebert_bollee_1890}.}. \textcite{nordhaus_two_2007}, à partir d'une reconstruction patiente de la productivité du calcul sur deux siècles, estime que le stock cumulé mondial de calculatrices mécaniques passa de cinq cents avant~1865 à environ 130\,000 en~1910 et 900\,000 en~1920, et que le coût par opération arithmétique baissa d'environ trois pour cent par an entre~1850 et~1945.

\dimmark{D2}{+}%
Le marché des calculatrices mécaniques ne convergeait pas vers le monopole. Une centaine de firmes fabriquaient des calculatrices aux États-Unis en~1904~; treize fabricants restaient actifs dans les années~1920 \parencite{cortada_before_1993}. Les barrières à l'entrée étaient faibles~: la mécanique de précision ne requérait pas d'infrastructure de réseau, la vente directe ou par agents ne supposait pas de couverture nationale, et la machine ne créait aucune dépendance au fournisseur une fois achetée. \textcite{leffingwell_office_1926}, dans le \emph{Office Appliance Manual} publié par l'association professionnelle du secteur, recensait des dizaines de modèles concurrents vendus à des prix comparables par des canaux similaires~--- un marché fragmenté, disputé, sans concentration visible.

\dimmark{D6}{--}%
L'effet de ces machines sur l'économie restait confiné à l'intérieur des organisations. La calculatrice n'agissait pas sur les prix entre entreprises, ne créait pas de marchés nouveaux, ne modifiait pas les échanges~: elle agissait sur les coûts de fonctionnement, sur la division interne du travail, sur la productivité du commis. \textcite{mills_white_1951}, dans \emph{White Collar}, observait que les machines de bureau ne furent pas utilisées de manière extensive avant~1910~--- une date que \textcite{cortada_before_1993} corrige à la hausse en montrant que les volumes de vente étaient déjà substantiels dans les années~1890. Mais l'un comme l'autre constatent que les effets de ces machines n'apparaissaient pas sur les marchés~: ils restaient dans les budgets internes des organisations, et les budgets des entreprises du dix-neuvième siècle ne sont, pour l'essentiel, pas conservés.


\subsection{La restructuration du travail de bureau}\label{subsec:calcul_travail}

\dimmark{D4}{*}%
Le déclin du coût unitaire du calcul n'entraîna pas de déclin de l'emploi~: au contraire, le nombre de commis aux États-Unis passa de un à deux pour cent de la population active en~1870 à plus de dix pour cent en~1940 \parencite{cortada_before_1993}. \textcite{reinstaller_market_2001}, dans le seul travail qui construise le mécanisme causal complet, ont avancé que le biais technique américain~--- capital-intensif, \emph{labour-saving}~--- transféra à l'intérieur du bureau la même logique de substitution qui avait transformé l'usine~: la corrélation entre le ratio capital/travail global de l'économie américaine et le ratio capital/travail des employés de bureau atteint~0,921 entre~1889 et~1937\footnote{Une corrélation aussi élevée sur deux séries également tendancielles ne suffit pas, en toute rigueur, à établir un mécanisme causal commun~: elle est compatible avec l'effet d'une tendance partagée. Reinstaller \& Hölzl mobilisent cependant l'argument dans une chaîne plus large (biais technique, baisse du salaire réel, féminisation), où la corrélation joue le rôle d'une régularité statistique parmi d'autres pièces.}. La part des femmes dans l'emploi de bureau passa de 2,4\,\% en~1870 à 48,4\,\% en~1920~; les salaires réels des commis baissèrent de 1\,920 à 1\,385~dollars (en prix de~1929) entre~1899 et~1919, les femmes étant payées 25 à 45\,\% de moins que les hommes \parencite{reinstaller_market_2001}. \textcite{davies_womans_1982} a documenté, à partir d'archives d'entreprise et de données du Census Bureau, la dynamique de ce basculement~: les employeurs embauchaient des femmes non parce qu'elles étaient formées aux machines, mais parce que les machines rendaient possible le recrutement d'une main-d'\oe{}uvre moins payée et la réorganisation du travail de bureau selon un modèle plus proche de la manufacture. \textcite{strom_beyond_1992} a étendu l'analyse en montrant que la féminisation ne se réduisait pas à un effet de prix~: elle s'accompagnait d'une transformation des hiérarchies, des espaces et des rythmes de travail dans le bureau américain entre~1900 et~1930. La calculatrice ne supprimait pas le travail~: elle le restructurait, le féminisait, et en comprimait le coût~--- non pas en éliminant des postes mais en substituant à des commis qualifiés une main-d'\oe{}uvre moins qualifiée, moins payée, assistée par la machine. \textcite{wootton_making_2007} confirme le résultat par un autre angle~: l'introduction des machines de bureau dans la comptabilité britannique ne réduisit pas le nombre de \emph{bookkeepers} mais modifia la nature de leur tâche, la partie mécanique du calcul étant transférée à la machine et la partie cognitive~--- vérification, interprétation, décision~--- restant humaine.

L'énergie, dans cette histoire, était presque inexistante. La seule donnée quantitative que nous ayons trouvée provient des archives du \emph{Nautical Almanac} exploitées par \textcite{daston_calculation_2018}~: en~1929, l'électricité nécessaire au fonctionnement des machines Hollerith et Burroughs représentait 9~livres sterling sur un budget total de 607~livres, soit 1,5\,\% du coût. Cette donnée est isolée~--- un point dans une institution une année~--- et il convient de la prendre comme illustrative plutôt que comme l'établissement d'une régularité~; nous n'avons pas trouvé de séries comparables sur d'autres budgets de l'époque\footnote{La part électrique des budgets administratifs et bancaires des années~1920--1930 reste un point aveugle de l'historiographie comptable. \textcite{cortada_before_1993} et \textcite{wootton_making_2007} ne discutent pas spécifiquement le poste \og énergie\fg{} dans les budgets qu'ils exploitent. Le résultat de \textcite{daston_calculation_2018} ne peut donc pas être généralisé sans réserve, mais il indique un ordre de grandeur cohérent avec ce que l'on attend d'un calcul électromécanique d'avant l'électronique.}. \dimmark{0-TER}{+}%
Sous cette réserve, le couplage entre calcul et énergie présentait deux propriétés qui le rendaient invisible~: d'abord par les ordres de grandeur~--- l'électricité, lorsqu'elle existait, restait un poste négligeable dans les budgets disponibles~---, ensuite par le fait que le calcul ne pilotait pas directement un système physique. La calculatrice produisait des chiffres qui restaient sur le papier~; elle ne contrôlait pas un réseau, ne coordonnait pas des flux matériels, ne consommait pas de matières premières au-delà du papier et de l'encre. L'industrie entière des machines de bureau~--- calculatrices, tabulatrices, machines à écrire, caisses enregistreuses~--- pesait 217,8~millions de dollars en~1929, soit deux dixièmes de pour cent du produit national brut américain \parencite{cortada_before_1993}. La question énergétique du calcul ne se poserait qu'au moment où le calcul changerait de substrat~--- du mécanique à l'électronique~--- et où son volume croîtrait de plusieurs ordres de grandeur.

\emergence{D2~--- Le marché des calculatrices mécaniques ne converge pas vers le monopole~: une centaine de firmes en~1904, treize fabricants encore actifs dans les années~1920. Les barrières à l'entrée sont faibles, le produit ne crée pas de dépendance au fournisseur, et aucun rendement croissant de réseau ne joue. \questionch{2}{D2~--- La convergence vers le monopole est-elle liée aux propriétés du réseau de communication, ou peut-elle procéder d'un autre mécanisme dans le calcul~?}}

\emergence{D4~--- La calculatrice ne remplace pas le travail~: elle le restructure. L'emploi de bureau croît, se féminise, et le salaire réel baisse. \textcite{davies_womans_1982} et \textcite{strom_beyond_1992} documentent la transformation des hiérarchies et des espaces de travail qui accompagne la mécanisation. \questionch{2}{D4~--- Cette restructuration sans suppression est-elle propre à la calculatrice, ou se retrouve-t-elle avec la tabulatrice et l'ordinateur~?}}

Au-delà des résultats par dimension, ce parcours traverse en filigrane une transformation plus large de l'économie matérielle de l'écrit administratif, dont \textcite{gardey_ecrire_2008} a proposé la périodisation et le cadre interprétatif. Entre~1890 et~1940, l'organisation du travail de bureau bascule d'un régime où la sécurité, la conservation et la traçabilité sont les valeurs cardinales~--- celles qui justifiaient le commis qualifié, l'écriture lente, le registre relié~--- vers un régime où la vitesse, la malléabilité et la mobilité de l'information prennent le pas. La confiance que les pratiques anciennes plaçaient dans les humains et dans leur jugement est progressivement déléguée aux artefacts et aux dispositifs qui en organisent le repérage. Cette bascule normative, que Gardey nomme une \og économie morale\fg{} de l'inscription, n'est pas le simple effet de la mécanisation~: elle en est la condition. La calculatrice, dont les effets restaient internes à l'organisation, accompagne ce déplacement sans le saturer~; la mécanographie qui suit, en couplant le stockage industrialisé et le calcul automatisé, viendra l'inscrire dans un dispositif d'une autre nature. C'est l'objet de la section suivante.


% =============================================================================

% =============================================================================
% §1.4 Le système mécanographique (1890–1940)
% =============================================================================
\clearpage
\section{Le système mécanographique~: le couplage stockage-calcul (1890--1940)}\label{sec:mechanography}
% =============================================================================

La calculatrice mécanique accélérait un geste~: un opérateur entrait des nombres, la machine produisait un résultat, rien ne persistait. Le résultat apparaissait sur un cadran ou sur une bande de papier que l'opérateur recopiait dans un registre, et la machine n'en gardait aucune trace. Chaque calcul était un acte isolé, sans mémoire, sans lien avec le calcul précédent ni avec le suivant. C'est la raison pour laquelle cent firmes pouvaient coexister sur ce marché~: une machine Burroughs et une machine Brunsviga faisaient la même chose, et aucune des deux ne créait de dépendance au-delà de l'habitude de l'opérateur.

Ce que Herman Hollerith construisit à la fin des années~1880 était d'une nature différente. La carte perforée~--- un rectangle de carton portant des trous disposés selon un code~--- était un support de données séparé de l'opérateur et de la machine, réutilisable, triable, stockable. Elle pouvait être perforée par une perforatrice, triée par une trieuse, tabulée par une tabulatrice, reproduite par une reproductrice~: le même jeu de cartes passait d'une machine à l'autre dans un flux intégré. Hollerith réalise ainsi, pour la première fois dans un système opérationnel, le couplage entre un stockage sur support \emph{machine-lisible} et un traitement \emph{en volume} de l'information qui y est inscrite. La précision n'est pas verbale. Le métier Jacquard (1804) couplait certes une carte perforée à un dispositif mécanique, mais la chaîne de cartes y constituait \emph{un programme unique} exécuté de façon répétée sur un substrat continu (le tissu)~--- non un ensemble de \emph{données distinctes} traitées par un opérateur agrégateur. L'\emph{Analytical Engine} de Babbage, conçue dès les années~1830, séparait le \emph{store} et le \emph{mill} dans une architecture analogue à celle de Hollerith, mais n'avait jamais été réalisée. Les tables de Prony, enfin, couplaient un calcul humain divisé à un stockage papier réutilisable, mais le support n'était pas lisible par une machine. Aucun de ces dispositifs ne satisfait la combinaison \emph{stockage machine-lisible / traitement en volume de données distinctes / système opérationnel} qui définit le geste hollerithien~--- et c'est cette combinaison, non l'un quelconque de ses composants, qui restructure ce qui suit. Ce n'était plus un outil qui accélérait un geste~; c'était un système qui restructurait un processus.

\textcite{cortada_before_1993} l'a formulé avec netteté~: \og \emph{unlike typewriters or calculating machines, which were single units of equipment that handled small amounts of information in a restricted fashion, tabulating gear operated on thousands, even millions of pieces of data}\fg{}. La distinction entre l'outil isolé et le système intégré n'est pas une question de degré mais de nature. Elle conditionne tout ce qui suit~--- le modèle économique, la forme du monopole, le rapport au travail, et la trajectoire qui mène à l'ordinateur.


\subsection{La fiche comme infrastructure~: le stockage avant les machines}\label{subsec:fiche_infrastructure}

Avant que la tabulatrice ne couple le stockage au calcul, le stockage existait déjà comme dispositif technico-organisationnel autonome. Dans la seconde moitié du XIX\textsuperscript{e}~siècle, les bibliothèques, les administrations, les bureaux de banque et les services médicaux développent un appareil~--- la fiche, le classeur, le meuble à tiroirs~--- qui constitue, indépendamment de toute mécanisation, une infrastructure de l'information. \textcite{gardey_ecrire_2008} a montré que cet appareil ne se réduit ni à ses objets matériels, ni à ses modes d'usage individuels~: il est d'emblée un \emph{système}, dont les pièces ne prennent leur sens que dans l'ensemble qui les ordonne. \og Une, la fiche n'est rien~; elle ne prend son sens que conçue dans cet ensemble (non fini) dont elle peut être extraite\fg{}\footnote{\textcite[ch.~4]{gardey_ecrire_2008}, qui en fait la \og lecture forte\fg{} du dispositif. La formule restitue ce que les manuels d'organisation du travail de bureau de l'époque~--- \textcite{leffingwell_office_1926} aux États-Unis, \textcite{bolle_utilisation_1929} en France~--- donnent comme évidence pratique~: une fiche utilisée seule est un fragment, et l'investissement n'a de rendement que s'il porte sur le système complet, c'est-à-dire l'index, les meubles, les protocoles d'inscription et les opératrices formées.}.

\dimmark{D8}{*}%
Cette transformation, datée par Gardey de la période~1890--1920, reconfigure ce que \textcite{febvre_outillage_1938} nommera dans les années~1930 l'\og outillage mental\fg{} d'une époque~: non plus seulement les concepts et les mots disponibles à la pensée, mais les supports, les instruments et les gestes qui rendent certaines opérations cognitives possibles\footnote{La filiation est explicitement revendiquée par \textcite[introduction]{gardey_ecrire_2008}, qui propose son ouvrage comme \og prolongation matérielle\fg{} de la notion d'outillage mental introduite par Febvre dans les premiers \emph{Annales}. La présente section retient cette extension matérielle de l'outillage mental comme cadre opératoire pour l'analyse de la fonction STORE~: ce ne sont pas seulement des objets que l'époque produit, ce sont des dispositifs distribués qui rendent disponibles des opérations~--- classer, retrouver, croiser, comparer~--- que les supports manuscrits antérieurs ne supportaient qu'au prix d'un travail considérable.}. La fiche, dans cette perspective, n'est pas un perfectionnement marginal du registre manuscrit~: elle inaugure une économie matérielle de l'écrit dans laquelle la mobilité, la malléabilité et la combinatoire prennent le pas sur la sécurité et la conservation. La confiance que les pratiques anciennes plaçaient dans les humains qualifiés~--- le commis qui retrouve un dossier, le bibliothécaire qui connaît son fonds~--- est progressivement déléguée aux artefacts et aux protocoles qui en organisent le repérage.

L'autre vertu théorique de la fiche, dont la sociologie des sciences a tiré des conséquences, est sa plasticité interprétative. \textcite{star_griesemer_1989} ont qualifié d'\og objets-frontières\fg{} (\emph{boundary objects}) les artefacts qui circulent entre des communautés de pratiques distinctes en gardant assez de stabilité pour permettre la coordination, mais en restant assez plastiques pour servir des usages différents. La fiche médicale, la fiche bibliographique, la fiche actuarielle, la fiche de stocks~: ces objets matériellement comparables soutiennent des opérations cognitives radicalement différentes selon le système qui les ordonne. Ce qui les rend \emph{infrastructure} n'est pas leur identité matérielle mais leur capacité à devenir le substrat partagé de coordinations hétérogènes\footnote{La théorie des objets-frontières de Star et Griesemer est postérieure aux faits que cette section décrit~; elle est ici employée rétrospectivement comme grille de lecture, dans la logique d'anachronisme contrôlé exposée au préambule du chapitre.}.

\dimmark{D7}{+}%
Cette infrastructure-fiche fonctionne déjà à grande échelle bien avant que la mécanisation ne s'en saisisse. Au tournant du siècle, le bibliographe belge Paul~Otlet et son associé Henri~La~Fontaine entreprennent de constituer, à Bruxelles, un \emph{Répertoire bibliographique universel} dont l'ambition est de fournir un index unifié de l'ensemble du savoir publié~; ils accumuleront, jusqu'à la Première Guerre mondiale, plusieurs millions de fiches selon une classification décimale qu'ils ont eux-mêmes affinée\footnote{Sur le projet d'Otlet et de La~Fontaine et leur \emph{Mundaneum} bruxellois, la référence biographique et archivistique canonique reste \textcite{rayward_universe_1975}.}. Otlet, dans les textes de cette période repris dans son \emph{Traité de documentation}~\parencite{otlet_traite_1934}, emploie déjà l'expression de \og société de l'information\fg{} pour désigner le régime que ce dispositif rend possible\footnote{Sur cette occurrence pré-1914 du syntagme et sur le projet plus large d'une économie politique de la documentation, on renverra à l'enquête d'Anne~Rasmussen reprise par \textcite[ch.~7]{gardey_ecrire_2008}~: l'expression circule dans les milieux internationaux de la documentation scientifique et technique avant d'être réinventée, sous une autre généalogie, par \textcite{machlup_production_1962}, qui en méconnait la profondeur historique.}. Que le concept précède de plus d'un demi-siècle l'industrie qui s'en réclamera ne tient pas à une anticipation visionnaire~: il tient à ce que l'infrastructure-fiche, autonome, suffisait déjà à poser comme objet le traitement systématique de l'information.

C'est sur cette infrastructure préexistante que vient se greffer la mécanographie. La carte perforée d'Hollerith n'invente ni la fiche, ni le système qui l'ordonne, ni l'opératrice qui la manipule~: elle couple ces éléments à un dispositif de calcul électromécanique. Ce que la sous-section suivante examine n'est donc pas l'irruption d'une fonction nouvelle~--- la fonction STORE existe avant elle~--- mais la transformation que produit son couplage à la fonction COMPUTE. La singularité du moment Hollerith-IBM tient précisément à ce couplage~: le stockage, déjà infrastructure, devient \emph{lisible par une machine}, et le calcul, déjà industriel, devient capable de \emph{traiter le stocké en volume}.

\emergence{STORE~--- La fonction de stockage de l'information n'apparaît pas avec la mécanographie~: elle est constituée comme infrastructure technico-organisationnelle dès la seconde moitié du XIX\textsuperscript{e}~siècle, autour de la fiche, du classeur et des protocoles administratifs et bibliographiques qui les ordonnent. Ce que la mécanographie introduit n'est pas la fonction mais son couplage avec la fonction de calcul. Ce déplacement~--- du stockage comme infrastructure manuelle à son intégration dans une chaîne automatisée~--- est l'objet propre de la présente section. \questionch{2}{D7/D8~--- La distinction entre l'infrastructure-fiche pré-mécanographique et son couplage avec le calcul électromécanique permet-elle de définir un seuil structurel~? Si oui, ce seuil est-il l'analogue, pour STORE, de ce qu'est l'unification de marché par le télégraphe pour TRANSMIT~?}}


\subsection{Hollerith et le recensement de~1890}\label{subsec:hollerith}

Herman Hollerith (1860--1929) était employé au \emph{Bureau of the Census} quand, en~1881, John Shaw Billings, directeur de la division des statistiques démographiques, lui fit la remarque qui allait orienter sa carrière~: il devait bien exister un moyen mécanique de tabuler les données du recensement, quelque chose sur le principe du métier Jacquard \parencite{cortada_before_1993}. Le recensement de~1880, portant sur cinquante millions d'habitants, avait pris huit ans à dépouiller~; celui de~1890, avec une population estimée à plus de soixante millions et un questionnaire considérablement élargi, menaçait de ne pas être achevé avant le suivant.

Hollerith construisit, entre~1884 et~1889, un système complet~: une perforatrice manuelle pour inscrire les données sur des cartes, un trieur électrique pour regrouper les cartes par catégories, et une tabulatrice pour compter les trous. La tabulatrice de~1890 ne faisait qu'une chose~--- compter~--- mais elle le faisait vite~: 415~cartes par minute, soit 80\,000 à 90\,000 par jour \parencite{cortada_before_1993}. En~1890, ses machines tabulèrent les 62\,622\,250~habitants des États-Unis et permirent de publier les premiers résultats en six semaines, contre des années par les méthodes antérieures. Le \emph{Bureau of the Census} estima l'économie à cinq millions de dollars par rapport au dépouillement manuel. Pour le recensement de~1900, le Bureau utilisa 311~tabulatrices, 20~trieuses automatiques et 1\,021~perforatrices, pour un coût total de 428\,239~dollars incluant les machines, le service et les cartes \parencite{cortada_before_1993}.

Le succès du recensement américain ouvrit les portes internationales~: l'Autriche utilisa les machines Hollerith pour son recensement de~1891, le Canada et la Norvège suivirent, puis la Russie en~1897 pour ses 129~millions de citoyens \parencite{cortada_before_1993}. Mais le recensement était un événement décennal~; ce dont Hollerith avait besoin pour pérenniser son entreprise, c'était d'une base de clientèle continue. Il la trouva dans les chemins de fer~--- le \emph{New York Central Railroad} signa en~1895 un contrat pour traiter près de quatre millions de lettres de voiture par an~--- puis dans les assureurs, les sidérurgistes et les grands distributeurs. Avant~1914, les applications couvraient la paie, la gestion des stocks, la comptabilité analytique, le suivi des coûts de production~--- toutes les fonctions que les calculatrices mécaniques ne pouvaient pas traiter en volume.


\subsection{Le modèle économique~: location, cartes, brevets}\label{subsec:ibm_model}

\dimmark{D2}{*}%
Le modèle économique d'Hollerith, puis de C-T-R et d'IBM, était structurellement différent de celui des fabricants de calculatrices. Les machines n'étaient pas vendues mais louées~: une trieuse coûtait environ 10~dollars par mois, une tabulatrice 40~dollars \parencite{cortada_before_1993}. L'amortissement intervenait en vingt mois pour les tabulatrices, trente pour les trieuses. Mais la vraie source de profit n'était pas la location des machines~--- c'étaient les cartes. Le contrat de location stipulait que le client devait acheter ses cartes exclusivement auprès d'Hollerith (puis d'IBM). Les cartes coûtaient entre 85~cents et un dollar le millier~; leur production revenait à environ 30~cents, soit une marge de 60 à 80\,\% \parencite{cortada_before_1993}. En~1929, les ventes de cartes représentaient un tiers des revenus d'IBM~; pendant la dépression, cette proportion monta à 39\,\%, parce que les clients qui gardaient leurs machines devaient continuer à acheter des cartes. Au pire de la crise, IBM vendait cent millions de cartes par an à 1,40~dollar le millier. En~1938, les cartes rapportaient 5~millions de dollars sur 34,7~millions de revenus totaux et étaient le produit le plus rentable de la gamme \parencite{cortada_before_1993}.

Le mécanisme de monopole qui en résultait procédait du verrouillage du consommable~: la machine était un véhicule pour la vente de cartes, et le client qui avait perforé ses données sur des cartes IBM ne pouvait pas les faire lire par des machines Powers\footnote{Les cartes IBM à~80~colonnes (1928) utilisaient des perforations rectangulaires, délibérément différentes des perforations rondes des cartes Hollerith antérieures et des cartes Powers. \textcite{heide_punched_2009} note que les perforations rectangulaires furent choisies non seulement pour des raisons techniques (lecture électrique plus fiable, carton plus solide) mais aussi parce qu'elles étaient brevetables.}. \textcite{austrian_hollerith_1982} montre que Hollerith avait imposé cette clause d'exclusivité dès les premiers contrats de location au \emph{Bureau of the Census}, et que le modèle fut maintenu systématiquement par Watson après la création de C-T-R. À cela s'ajoutait un portefeuille de brevets systématiquement élargi par des rachats~: IBM acquit les brevets de Peirce en~1922 pour un montant équivalant à 25\,\% de ses bénéfices nets annuels, et racheta en~1930 l'ensemble des brevets de l'Autrichien Gustav Tauschek~--- non pour utiliser sa technologie, mais pour empêcher un concurrent de contrôler la multiplication par cartes perforées \parencite{heide_punched_2009}.

Les chiffres du \emph{Department of Justice} américain, à la fin de~1935, quantifient le résultat~: IBM contrôlait 85,7\,\% des tabulatrices installées aux États-Unis, 86,1\,\% des trieuses, 81,6\,\% des perforatrices \parencite{cortada_before_1993}. L'antitrust américain attaqua IBM et Remington Rand en~1933, et la Cour suprême confirma en~1936 l'interdiction d'obliger les clients à acheter les cartes auprès du loueur. Mais le monopole survécut au jugement~: la dépendance technique créée par les formats de cartes propriétaires et la supériorité de la force de vente IBM suffirent à maintenir la position dominante.

\dimmark{D2}{*}%
\textcite{heide_punched_2009} éclaire un aspect que les données américaines seules ne font pas voir. La licence que Watson accorda à Powers en~1914, donnant accès aux brevets Hollerith contre une redevance de 20\,\% du chiffre d'affaires sur les tabulatrices et trieuses, était un instrument de contrôle autant qu'un geste de prudence antitrust. En Allemagne, un tribunal avait imposé une redevance de 5\,\% seulement pour une licence équivalente \parencite{heide_punched_2009}. Le différentiel~--- 20\,\% aux États-Unis, 5\,\% en Allemagne~--- illustre une différence institutionnelle que Heide formule explicitement~: le système américain de brevets renforçait le premier entrant en lui permettant de fixer les conditions de la licence, tandis que le système allemand, avec ses licences imposées par voie judiciaire, facilitait la concurrence.

\textcite{mounier-kuhn_punched_2015}, dans sa recension de Heide, propose une périodisation en quatre cycles de clôture (\emph{closure}) du système des cartes perforées, chacun fondé sur un brevet-clé renouvelé~: le brevet Hollerith initial (lecture électrique), le brevet sur le groupe automatique (1914, accordé en~1931 seulement), le brevet sur la carte à~80~colonnes (1928), et les brevets de multiplication. À chaque cycle, le renouvellement du brevet-clé permettait de prolonger le verrouillage du système et de bloquer ou de retarder les challengers.

\emergence{D2~--- Le monopole IBM procède d'un mécanisme différent de celui du télégraphe~: pas les rendements de réseau mais le triptyque location + cartes exclusives + brevets. Deux monopoles, deux mécanismes, un même résultat~--- la concentration du marché dans les mains d'un opérateur dominant. \questionch{2}{D2~--- Le monopole par verrouillage du consommable (IBM/cartes) est-il structurellement plus stable que le monopole par rendements de réseau (télégraphe)~?}}

\dimmark{D3}{+}%
Les revenus d'IBM passèrent de 2,5~millions de dollars en~1919 à 20,3~millions en~1931, fléchirent à 17,6~millions en~1933, puis remontèrent à 45,3~millions en~1940 grâce aux contrats du New Deal \parencite{cortada_before_1993}. Le profit net passa de 1,4~million en~1922 à 9,1~millions en~1939, et IBM paya des dividendes sans interruption pendant toute la dépression. En~1928, IBM n'était que le quatrième fournisseur de l'industrie des machines de bureau par le chiffre d'affaires (19,7~millions, derrière Remington Rand à 59,6, NCR à 49, et Burroughs à 32,1), mais le premier par la marge~: 26,9\,\% de profit sur le chiffre d'affaires, contre 25,9\,\% pour Burroughs, 15,9\,\% pour NCR et 10,1\,\% pour Remington Rand \parencite{cortada_before_1993}. En~1939, les positions s'étaient inversées~: IBM était premier en profit (9,1~millions), devant NCR (3,1) et Burroughs (2,9), alors que Remington Rand, avec un chiffre d'affaires encore supérieur (43,4~millions), ne dégageait que 1,6~million de bénéfice. C'est la marge qui distinguait IBM, pas la taille~--- et cette marge venait des cartes.

Pour autant, la taille relative d'IBM dans l'économie américaine restait modeste. En~1930, IBM représentait 12\,\% de la valeur de l'industrie des machines de bureau (20,3~millions sur 165,3) et moins de trois centièmes de pour cent du produit national brut. L'industrie entière~--- calculatrices, tabulatrices, machines à écrire, caisses enregistreuses~--- pesait 217,8~millions de dollars en~1929, soit un peu plus de deux dixièmes de pour cent du PNB \parencite{cortada_before_1993}. Watson lui-même reconnaissait en~1929 que son équipement n'était installé que dans un cinquième des organisations qui pourraient l'utiliser~--- le marché potentiel n'était servi qu'à la marge. La raison en était le seuil de rentabilité~: un client devait effectuer au moins 1\,500~opérations par jour pour justifier l'investissement \parencite{cortada_before_1993}, ce qui excluait de fait les petites organisations.


\subsection{Une trajectoire alternative~: les classi-compteurs de Lucien March (1901--1940)}\label{subsec:march_classicompteurs}

En~1896, lorsque Lucien March (1859--1933) prend la direction de la Statistique générale de France, le bureau parisien dépouille le recensement quinquennal au moyen des machines Hollerith\footnote{Polytechnicien de formation, Lucien March dirige la Statistique générale de France de~1896 à~1920. Sur sa carrière administrative et son rôle dans la modernisation statistique française, voir \textcite{cinquanteans_insee_1996} et, pour le contexte plus large de l'épistémologie statistique française, \textcite{desrosieres_politique_1993}.}~; en~1901, il ne le fait plus. Le bureau s'en décharge à l'occasion d'un dispositif que March a conçu lui-même et que la Statistique utilisera, sans interruption majeure, pour tous les recensements jusqu'à celui de~1936. March avait perçu, dès l'expérience du recensement de~1896, deux limites au système américain~: son coût~--- les machines Hollerith étaient louées pour des sommes importantes, et leur usage continu n'était pas justifié par le rythme décennal du recensement~--- et la séparation qu'il imposait entre la lecture des bulletins et le comptage statistique, en obligeant à une perforation préalable comme étape distincte du dépouillement.

Le classi-compteur, breveté par March à la fin du XIX\textsuperscript{e}~siècle et progressivement perfectionné, est un clavier compté de soixante touches, chacune correspondant à une donnée portée sur le formulaire de recensement~--- sexe, date et lieu de naissance, profession, situation matrimoniale, et ainsi de suite. L'opératrice~--- les sources administratives parlent des \og dames\fg{}~--- enfonce successivement les touches en lisant chaque bulletin, et les compteurs internes additionnent par catégorie. Une fois la série traitée, l'abaissement d'un châssis imprime sur une bande les chiffres enregistrés, le compteur est remis à zéro et un nouveau lot peut être dépouillé. Le rythme de fonctionnement est de huit mille bulletins par jour \parencite[ch.~7]{gardey_ecrire_2008}.

\dimmark{D10}{*}%
L'innovation analytique du dispositif tient à un geste~: March élimine la carte perforée comme support intermédiaire. Là où le système Hollerith requiert une perforation préalable de la donnée sur un support indépendant, qui doit ensuite être trié et tabulé en flux successifs, le classi-compteur fait coïncider la lecture du bulletin avec son comptage. Le bulletin papier reste comme document source, mais l'objet-fiche-perforée disparaît du processus~--- ce qui supprime à la fois le poste budgétaire le plus coûteux du système américain (la location des perforatrices, l'achat des cartes) et l'étape de travail la plus exposée à l'erreur de transcription. L'inconvénient inverse~--- l'absence d'un support intermédiaire matériel qui permette une vérification à distance et a~posteriori~--- pose en contrepartie des problèmes en matière de détection d'erreurs, mais sans remettre en cause le choix du dispositif pour l'usage ciblé~: le dépouillement statistique d'État.

\dimmark{D2}{*}%
La Statistique générale de France utilise les classi-compteurs pour les recensements de~1901, 1906, 1911, 1921, 1926, 1931 et~1936, avec des perfectionnements successifs portant sur la mécanique du clavier et l'extension du jeu de catégories codables. L'avantage économique est documenté~: moins cher que les locations Hollerith, le système permet par ailleurs aux opératrices de manipuler directement les bulletins sans étape de perforation séparée. Au-delà du cas français, le système est adopté par les administrations statistiques d'Italie, du Portugal et du Brésil~\parencite[ch.~7]{gardey_ecrire_2008}~; \textcite{peaucelle_mecanographie_2004}, dans une étude économique rétrospective, juge que le choix de la mécanographie classi-compteurs en France relève d'une rationalité économique défendable au regard des coûts d'opportunité de l'époque, et non d'un retard ou d'une singularité culturelle.

Ce cas mérite d'être inscrit avec soin dans la lecture comparative que cette section conduit. Il ne s'agit pas d'une variante secondaire du système Hollerith~: c'est un dispositif structurellement différent~--- le clavier-compteur sans support intermédiaire~--- qui a fonctionné quarante ans dans plusieurs pays, et qui n'a pas été supplanté par les cartes perforées dans son espace d'usage. Ce que ce cas révèle, c'est que la trajectoire Hollerith-IBM, qui domine l'historiographie de la mécanographie, n'était ni la seule disponible ni la seule économiquement viable au tournant du siècle. La concentration de l'historiographie sur Hollerith tient en partie à la croissance ultérieure du conglomérat américain, mais cette croissance ne peut être rétrojetée comme nécessité technique~: à l'époque où le système IBM se déployait aux États-Unis et conquérait progressivement les administrations statistiques européennes, en France, le classi-compteur dépouillait les recensements à un coût inférieur, et continuerait à le faire pendant quatre décennies. La non-émergence d'IBM en France pour le dépouillement statistique d'État n'est donc pas un retard~; c'est une trajectoire alternative documentée. \textcite{mounier_kuhn_informatique_1999} insiste sur ce point dans son enquête sur l'informatique française d'avant le \emph{Plan calcul}~: les institutions statistiques françaises ne se sont rangées au modèle des cartes perforées qu'au sortir de la Seconde Guerre mondiale, et seulement parce que le volume et la complexité du dépouillement~--- introduction de variables nouvelles, croisements multivariés~--- débordaient désormais les capacités du clavier-compteur.

\emergence{D2/D10~--- Le cas March établit empiriquement qu'à demande équivalente (le dépouillement décennal d'un recensement national), plusieurs dispositifs technico-organisationnels coexistent durablement et occupent la même fonction par des architectures matérielles distinctes. La domination ultérieure d'IBM ne relève donc pas d'une supériorité technique intrinsèque mais d'un choix d'écosystème institutionnel et économique. \questionch{2}{D2/D10~--- Quel mécanisme institutionnel a permis à la trajectoire alternative française de se maintenir quarante ans, et qu'est-ce qui a fait qu'elle ne s'est pas exportée au-delà de quelques pays latins~? Le facteur est-il l'absence d'industrie de fabrication concurrente, l'absence de stratégie d'exportation comparable à celle de Watson, ou la spécificité du dépouillement administratif comme niche fonctionnelle relativement à d'autres usages~?}}


\subsection{La tabulatrice dans les organisations~: le Nautical Almanac et le chemin de fer}\label{subsec:zoom_tabulating}

Le Census de 1890 avait démontré que la tabulatrice pouvait traiter des millions d'enregistrements en un temps que le calcul humain ne pouvait pas atteindre. Mais le recensement est un événement décennal~; ce qui allait révéler les effets de la mécanisation sur l'organisation du travail, c'est l'usage quotidien de ces machines dans des institutions qui calculent en continu~--- les observatoires, les compagnies d'assurance, les administrations ferroviaires. C'est dans ces usages que la tabulatrice cesse d'être un exploit technique et devient un fait organisationnel, dont les conséquences sur le travail, le coût et la supervision n'avaient pas été anticipées par ses promoteurs.

\paragraph{Le Nautical Almanac~: la mécanisation et ses coûts}\mbox{}\\
Le cas le mieux documenté est celui du \emph{Nautical Almanac} britannique, dont les archives administratives ont été exploitées par \textcite{daston_calculation_2018} à partir du fonds RGO~16 de la Cambridge University Library. Le \emph{Nautical Almanac}, publié depuis 1767 sous la direction de l'\emph{Astronomer Royal}, fournissait les éphémérides nécessaires à la navigation de la flotte marchande et militaire britannique. Sa production exigeait des dizaines de milliers de calculs répétitifs~--- positions lunaires, prédictions de marées, tables de réfraction~--- dont l'exécution avait été organisée, depuis Maskelyne, en un réseau de calculateurs à domicile dispersés dans toute l'Angleterre \parencite{croarken_human_2009}. Chaque calculateur recevait des \og \emph{precepts}\fg{}~--- des algorithmes décomposés en étapes sur des formulaires pré-imprimés~--- et rendait un mois complet de tables, vérifié par un \emph{anti-computer} et un \emph{comparer}. Le système relevait du \emph{putting-out}, comme le tissage à domicile~: du travail qualifié, décentralisé, faiblement supervisé.

En décembre~1928, l'Amirauté autorisa l'achat d'une machine Burroughs et la location d'une machine Hollerith pour une durée de six mois. Le surintendant Comrie justifia la dépense par la promesse d'une réduction de la fatigue mentale et d'une augmentation de la précision. Les archives permettent de reconstituer le bilan. La location de la Hollerith coûtait environ 264~livres~; les dix mille cartes perforées nécessaires, 100~livres~; les salaires de quatre opératrices pendant six mois, puis de deux opératrices pendant six mois supplémentaires, 234~livres~; l'électricité, 9~livres. Le total atteignait 607~livres pour un calcul que les méthodes anciennes accomplissaient pour 500~livres par an. La mécanisation coûtait plus cher, pas moins. Et elle exigeait davantage de travail, pas moins~: au lieu de dix mille sommes tirées de sept tables différentes, il fallait désormais perforer douze millions de chiffres sur trois cent mille cartes pour les faire passer dans la machine.

\dimmark{D4}{*}%
Ce qui avait changé, ce n'était pas le volume de travail mais sa nature et sa répartition. Les anciens calculateurs~--- des hommes qualifiés, souvent retraités du service ou enseignants complétant leurs revenus, travaillant chez eux à leur rythme~--- furent remplacés par des opératrices recrutées sur concours dans les catégories les moins bien payées de la fonction publique. L'Amirauté exigeait des candidates qu'elles eussent passé un examen compétitif portant sur \og \emph{English, Arithmetic, General Knowledge and Mathematics}\fg{}~; la réglementation de la \emph{Civil Service} interdisait l'emploi de femmes mariées. Le surintendant Comrie, dans une lettre de 1931, estimait que des femmes pourraient occuper les postes d'assistante aussi bien que des hommes, mais recommandait que les fonctions de surintendant et de surintendant adjoint \og \emph{be reserved for men, especially in view of the fact that the greater part of the calculation is now performed by mechanical means}\fg{}~--- logique remarquable où la mécanisation du calcul, loin de rendre le genre du calculateur indifférent, justifiait au contraire la réservation de l'autorité sur les machines aux hommes qui ne les opéraient pas.

La centralisation devint obligatoire~: les machines étaient coûteuses et ne pouvaient pas quitter le bureau. Le travail à domicile disparut. Et le travail de supervision explosa. Là où l'ancien système consistait à dire à un calculateur~--- un certain W.~F.~Doaken, \emph{M.A.}~--- \og \emph{Do the Moon Transit}\fg{}, et à recevoir cinq mois plus tard la copie prête pour l'imprimeur, le nouveau système distribuait le travail entre six ou sept personnes, auxquelles le surintendant devait fournir \og \emph{perhaps 100 to 120 different sets of instructions}\fg{}. \dimmark{D7}{+}%
Comrie estimait que la préparation du travail pour les machines et leurs opératrices représentait désormais vingt à trente pour cent du calcul total~--- un cinquième à un tiers du temps du surintendant consacré non pas à calculer, ni à superviser au sens classique, mais à \emph{concevoir la séquence} dans laquelle humains et machines devaient s'articuler. C'est ce que \textcite{daston_calculation_2018} nomme l'\emph{analytical intelligence}~: l'intelligence qui organise le calcul, distincte de celle qui l'exécute et de celle qui conçoit les formules. Elle occupait le deuxième étage de la pyramide de Prony~--- les sept ou huit algébristes qui traduisaient les mathématiques en procédures élémentaires~--- mais la mécanisation l'avait considérablement alourdie, parce que les machines ne calculaient pas de la même manière que les humains, et que chaque machine calculait différemment des autres.

Les frictions étaient concrètes. Les frères Daniels, calculateurs anciens et irremplaçables pour la relecture d'épreuves, étaient jugés \og \emph{temperamentally unfitted to supervise subordinate staff}\fg{} et \og \emph{too stereotyped in their habits to adapt themselves to the use of machines}\fg{}. Miss~Stocks et Miss~Burroughs, affectées à la transformation des coordonnées héliocentriques en coordonnées géocentriques sur machines Brunsviga, nécessitèrent chacune trois mois de formation individuelle dispensée par le surintendant lui-même. Le nouveau système était, à terme, environ vingt pour cent moins cher que l'ancien~--- mais seulement parce que les opératrices étaient payées moins que les anciens calculateurs, pas parce que les machines avaient réduit le travail.

\paragraph{Le PLM~: \og d'abord organiser, ensuite mécaniser\fg{}}\mbox{}\\
Le même schéma se reproduisait dans un contexte radicalement différent. En 1929, Georges Bolle, polytechnicien, chef de la comptabilité de la Compagnie des chemins de fer Paris-Lyon-Méditerranée, publiait dans la \emph{Revue générale des chemins de fer} un article détaillé sur l'introduction des machines Hollerith pour le suivi du fret \parencite{bolle_utilisation_1929}. Bolle ne cachait pas son enthousiasme~: les avantages économiques des nouvelles machines étaient, à ses yeux, \og vraiment incalculables\fg{}. Mais l'essentiel de l'article ne portait pas sur les machines~; il portait sur le travail nécessaire \emph{avant} que les machines ne puissent fonctionner. Chaque détail du flux de travail devait être méticuleusement pensé à l'avance~: comment maximiser l'usage des quarante-cinq colonnes d'une carte Hollerith, comment numéroter les catégories de fret pour accélérer le codage, quels pictogrammes créer pour les marchandises les plus fréquemment transportées. \og \emph{Chaque détail doit être minutieusement examiné, discuté, pesé dans un travail de cette nature}\fg{}, écrivait Bolle, avant de formuler la devise qui résumait l'expérience du PLM~: \og \emph{D'abord organiser, ensuite mécaniser}\fg{}\footnote{Cité également par \textcite{couffignal_machines_1933}, p.~79, qui en fait un principe général de la mécanisation du calcul.}.

Comme au \emph{Nautical Almanac}, l'introduction des machines avait entraîné la centralisation du travail (à Paris), le recrutement d'une main-d'œuvre féminine bon marché, et un effort de supervision sans proportion avec celui de l'ancien système. Les opératrices perforaient trois cents cartes de quarante-cinq colonnes par jour, pendant six heures et demie au maximum, et pas plus de quatorze jours consécutifs par mois \parencite{bolle_utilisation_1929}. Une étude psychophysique conduite peu après par \textcite{lahy_selection_1931} sur les opératrices de machines Elliot-Fischer dans les compagnies ferroviaires françaises confirma que chaque calcul impliquait seize étapes distinctes~--- de l'insertion du papier dans la machine au réglage des compteurs en passant par la vérification des noms et des montants~--- et que l'attention requise était \og \emph{continue et concentrée}\fg{}, impossible à soutenir longtemps \og \emph{sans repos}\fg{}. L'effort que les machines avaient promis de supprimer~--- la fatigue mentale du calcul répétitif~--- n'avait pas disparu~; il s'était déplacé des calculateurs vers les opératrices (qui devaient maintenir une vigilance constante sur la machine) et vers les superviseurs (qui devaient concevoir les séquences d'opérations). \dimmark{D4}{*}%

\paragraph{Le résultat~: la mécanisation réorganise, elle ne remplace pas}\mbox{}\\
Les deux cas~--- l'un scientifique et britannique, l'autre industriel et français~--- convergent sur un résultat que ni Comrie ni Bolle n'avaient anticipé et que la littérature ultérieure n'a pas suffisamment relevé. La mécanisation du calcul ne réduit pas le nombre de travailleurs~: elle le maintient ou l'augmente, et elle en transforme le profil. Des calculateurs qualifiés, autonomes, travaillant à domicile, sont remplacés par des opératrices moins qualifiées, moins payées, supervisées de près, auxquelles il faut ajouter un surintendant dont le temps est désormais absorbé, à hauteur de vingt à trente pour cent, par la conception des procédures. Le coût total peut baisser~--- il baisse de vingt pour cent au \emph{Nautical Almanac}~--- mais cette baisse tient à la compression salariale, pas à l'économie de travail.

\dimmark{D4}{*}%
La devise de Bolle~--- \og d'abord organiser, ensuite mécaniser\fg{}~--- fait écho, sans que Bolle pût le savoir, aux six formes informationnelles de McCallum sur l'\emph{Erie Railroad}. Dans les deux cas, la technique matérielle (le télégraphe, la tabulatrice) n'est efficace qu'insérée dans un dispositif organisationnel qui la précède et la conditionne. Dans les deux cas, le coût de ce dispositif est sous-estimé par les promoteurs de la technique. Et dans les deux cas, ce coût est essentiellement un coût de travail qualifié~--- l'\emph{analytical intelligence} de Daston, la \og supervision\fg{} de Comrie, l'\og effort cérébral très laborieux\fg{} de Bolle~--- qui ne disparaît pas avec la machine mais se concentre sur un nombre restreint de personnes dont la charge cognitive augmente.

\emergence{D4~--- La mécanisation du calcul ne remplace pas le travail humain~: elle le réorganise. Le nombre de travailleurs se maintient ou augmente~; leur profil change (calculateurs qualifiés $\rightarrow$ opératrices moins payées + superviseurs plus chargés). Le coût de supervision (20--30\,\% du travail total) est une charge nouvelle créée par la mécanisation elle-même. \questionch{2}{D4~--- Cette restructuration sans suppression se conserve-t-elle avec le computing électronique~? L'ordinateur remplace-t-il enfin les calculateurs, ou réorganise-t-il le travail une deuxième fois~?}}

\emergence{D7~--- La devise de Bolle (\og d'abord organiser, ensuite mécaniser\fg{}) confirme, dans un contexte entièrement différent, le résultat du cas ferroviaire~: la technique informationnelle n'est efficace qu'insérée dans un dispositif organisationnel qui la précède. Le coût de ce dispositif (20--30\,\% du travail total au \emph{Nautical Almanac}) est un coût de coordination au sens de \textcite{coase_nature_1937}~--- mais un coût que la technique informationnelle \emph{crée} au lieu de le réduire. \questionch{2}{D7~--- La mécanisation du calcul augmente-t-elle les coûts de coordination à l'intérieur de la firme avant de les réduire~? Ce schéma se reproduit-il avec l'informatisation des années~1960--1980~?}}


%%% --- BILAN §1.3--§1.4 ---

\subsection*{Bilan~: deux mécanisations, deux logiques}\label{subsec:bilan_compute}

Les deux sections qui précèdent ont traversé, dans l'ordre chronologique, deux phénomènes que le mot \og mécanisation du calcul\fg{} tend à confondre. La calculatrice mécanique (§\,\ref{sec:calculating}) est un outil individuel qui accélère un geste élémentaire~--- l'addition, la multiplication~--- sans modifier le processus dans lequel ce geste s'insère, sans créer de mémoire externe, sans lier l'opérateur à un fournisseur particulier. Le système mécanographique (§\,\ref{sec:mechanography}) est une architecture intégrée~--- perforatrices, trieuses, tabulatrices, reproductrices~--- qui couple pour la première fois, dans un système opérationnel, le stockage de données sur un support lisible par la machine et le traitement en volume de ces données. La différence n'est pas de degré~; elle est de nature, et elle produit des résultats économiques opposés~: un marché concurrentiel dans un cas, un monopole par verrouillage dans l'autre~; un instrument invisible à l'extérieur de l'organisation dans un cas, un système qui restructure les organisations et leurs rapports avec leurs fournisseurs dans l'autre.

L'un et l'autre partagent cependant un trait commun, que le zoom sur le \emph{Nautical Almanac} et le PLM a fait apparaître~: la mécanisation du calcul ne supprime pas le travail humain~; elle le réorganise. Le nombre de travailleurs se maintient ou augmente~; leur profil change~--- des calculateurs qualifiés et autonomes sont remplacés par des opératrices moins payées et plus étroitement supervisées~; et le coût de supervision~--- ce que \textcite{daston_calculation_2018} nomme l'\emph{analytical intelligence}~--- croît au lieu de décroître. La mécanisation améliore la productivité du calcul, mais elle le fait en restructurant le travail, pas en le supprimant~--- un résultat que \textcite{reinstaller_market_2001}, \textcite{davies_womans_1982} et \textcite{strom_beyond_1992} documentent dans trois registres différents (quantitatif, genre, organisation). \textcite{yates_control_1989} montre que cette restructuration ne s'explique pas par la machine seule~: les systèmes de classement, les formulaires standardisés et les technologies de duplication interne précèdent les machines de bureau et conditionnent leur adoption.

\textcite{kooij_information_2022} est, dans la littérature sur les technologies à usage général, l'auteur qui accorde le plus d'attention à la trajectoire du calcul mécanique. Il mentionne l'Arithmomètre, les tabulatrices Hollerith, Babbage, et parle d'une \og \emph{First Mechanical Calculator Revolution}\fg{}. Mais il condense cent trente ans en un seul paragraphe de contexte, en route vers le transistor~; son cadre analytique complet (phases de vie de la technologie, cycles d'affaires, identification des brevets) n'est pas appliqué à cette période. Hollerith reçoit une demi-phrase~; Morse, Bell et Marconi reçoivent chacun un volume entier. Ce n'est pas un oubli~: c'est le signe que la mécanisation du calcul était trop petite, relativement à l'économie, pour que le cadre des technologies à usage général la capte comme un fait de premier ordre. L'industrie entière des machines de bureau pesait, à son pic de~1929, deux dixièmes de pour cent du produit national brut~; IBM seul en représentait un centième. À cette première relativisation, il faut en superposer une seconde~: l'industrie des machines à cartes perforées elle-même, dont IBM contrôlait~85,7\,\%, ne représentait jusqu'aux années~1930 qu'une part minime de la mécanographie totale. Les tâches comptables et bancaires courantes étaient assurées principalement par les machines comptables traditionnelles~--- Burroughs, Comptometer, Underwood Bookkeeping~---, pas par les tabulatrices Hollerith \parencite[p.~158]{cortada_before_1993}. \textcite{heide_punched_2009} confirme la lecture pour les marchés européens~: le système hollerithien ne s'imposa pas comme architecture universelle de la mécanographie de bureau, mais comme spécialisation d'un sous-segment~--- celui du dépouillement statistique en grand volume et de la gestion des flux dans les administrations et les grandes entreprises. Le monopole d'IBM portait donc, jusqu'à la veille de la Seconde Guerre mondiale, sur un sous-segment d'une industrie elle-même marginale dans l'économie américaine. L'énergie consommée par le calcul mécanique était, comme on l'a vu, 1,5\,\% du coût total d'une opération de mécanisation~--- un rapport qui rendait la question énergétique littéralement invisible aux acteurs comme aux observateurs\footnote{L'invisibilité énergétique ne signifie pas l'absence d'empreinte matérielle. \textcite{little_toxic_2014} a documenté que le site d'Endicott (New York), ville-usine d'IBM depuis les années~1920, est aujourd'hui contaminé par un panache de trichloroéthylène sur 300~acres, avec plus de 200~millions de dollars dépensés en décontamination. La continuité du site~--- de la mécanographie des années~1930 à la fabrication de composants électroniques de l'après-guerre~--- montre que la dimension matérielle traverse la trajectoire d'IBM bien au-delà du carton perforé. Voir également \textcite{lecuyer_clean_2017} sur les pollutions de la Silicon Valley et \textcite{chiu_dark_2011} sur celles de Taïwan.}.

Ce n'est qu'au moment où le calcul changera de substrat~--- du mécanique à l'électronique~--- et où son volume croîtra de plusieurs ordres de grandeur que les ratios commenceront à se modifier. Le tableau~\ref{tab:compute-dimensions} récapitule ce que la traversée des deux cas~--- la calculatrice mécanique et le système mécanographique~--- a permis d'établir pour chacune des dix dimensions de la grille de lecture.

\begin{table}[p]
\caption{La mécanisation du calcul au prisme des dix dimensions~: observations, incertitudes et questions ouvertes}\label{tab:compute-dimensions}
\centering\footnotesize
\renewcommand{\arraystretch}{1.15}
\begin{tabularx}{\textwidth}{@{}cXcp{4.5cm}@{}}
\toprule
\textbf{Dim.} & \textbf{Observations des cas \u00a7\,\ref{sec:calculating}--\ref{sec:mechanography}} & & \textbf{Question pour la suite} \\
\midrule
D1 & \emph{Calculatrice}~: traite des données internes, pas de support persistant. \emph{Mécanographie}~: la carte perforée est le premier support de données lisible par la machine, réutilisable et stockable. Le carton ne consomme pas d'énergie pour conserver l'information. & + & Le passage au support électronique modifie-t-il cette propriété~? \\[3pt]
D2 & \emph{Calculatrice}~: marché fragmenté ($\sim$100~firmes US, 1904). \emph{Mécanographie}~: monopole IBM à~85,7\,\% (DoJ, 1935) par le triptyque location + cartes exclusives + brevets (Peirce~1922, Tauschek~1930). Différentiel institutionnel~: licence à Powers~: 20\,\% US vs 5\,\% Allemagne \parencite{heide_punched_2009}. Quatre cycles de clôture identifiés par \textcite{mounier-kuhn_punched_2015}. & * & Ce mécanisme de verrouillage par le consommable se retrouve-t-il dans d'autres secteurs~? \\[3pt]
D3 & \emph{Calculatrice}~: bien d'équipement vendu ($\sim$\$220). \emph{Mécanographie}~: location + flux de consommables. Cartes = 33\,\% des revenus IBM (1929), 39\,\% pendant la dépression \parencite{cortada_before_1993}. & + & Le modèle location + consommable constitue-t-il un type de capital informationnel distinct~? \\[3pt]
D4 & La mécanisation ne semble pas réduire l'emploi de bureau~: il passe de 1--2\,\% (1870) à 10\,\% (1940) de la population active \parencite{cortada_before_1993}. Corrélation K/L~: 0,921 \parencite{reinstaller_market_2001}. Au \emph{Nautical Almanac}~: supervision = 20--30\,\% du travail \parencite{daston_calculation_2018}. Féminisation et compression salariale documentées par \textcite{davies_womans_1982} et \textcite{strom_beyond_1992}. & + & Le computing électronique modifie-t-il ce rapport~? \\[3pt]
D5 & Coût par opération en baisse de $\sim$3\,\%/an, 1850--1945 \parencite{nordhaus_two_2007}. Seuil client~: 1\,500~opérations/jour \parencite{cortada_before_1993}. \emph{Pas d'analyse d'élasticité disponible.} & ? & La baisse du coût engendre-t-elle un rebond de la demande de calcul~? \\[3pt]
D6 & Pas d'effet documenté sur les prix ou l'intégration des marchés. Le calcul agit à l'intérieur des organisations~; ses effets restent dans les budgets internes, qui ne sont généralement pas conservés. & ? & L'absence d'effet observable reflète-t-elle une différence de nature, ou un défaut de sources~? \\[3pt]
D7 & \og D'abord organiser, ensuite mécaniser\fg{} \parencite{bolle_utilisation_1929}. 16~étapes par calcul \parencite{lahy_selection_1931}. \textcite{yates_control_1989}~: les systèmes de classement précèdent les machines et conditionnent leur adoption. & + & La mécanisation augmente-t-elle les coûts de coordination avant de les réduire~? \\[3pt]
D8 & IBM revenus~: \$2,5M (1919) $\rightarrow$ \$45,3M (1940). Dividendes payés sans interruption pendant la dépression. & ? & Le financement du calcul suit-il un cycle propre~? \\[3pt]
D9 & Pas de chaîne d'extraction identifiée pour les machines de bureau. La carte est du carton~; la machine est de la mécanique de précision. & -- & La matérialité apparaît-elle avec le changement de substrat~? \\[3pt]
D10 & Powers disposait d'une technologie jugée supérieure en~1912--14 mais de capacités de fabrication et de commercialisation inférieures \parencite{cortada_before_1993, heide_punched_2009}. Carte 80~colonnes~: perforations rectangulaires délibérément brevetables \parencite{heide_punched_2009}. & + & L'avantage du premier entrant tient-il à la technique ou au modèle d'affaires~? \\[3pt]
\midrule
0-TER & £9 sur £607 \parencite{daston_calculation_2018}~: 1,5\,\% du coût total. L'énergie est un poste négligeable~; la question énergétique ne se pose pas à cette échelle. Industrie entière = 0,2\,\% du PNB (1929). & + & À quel moment et par quel mécanisme le coût énergétique du calcul change-t-il d'ordre de grandeur~? \\[3pt]
\bottomrule
\end{tabularx}
\renewcommand{\arraystretch}{1.0}
\end{table}


% =============================================================================

% =============================================================================
% §1.5 Le téléphone (1876–1940)
% =============================================================================
\clearpage
\section{Le téléphone (1876--1940)}\label{sec:telephone}
% =============================================================================


% §1.5 — DRAFT v1 — 26 avril 2026
% Structure en quatre temps :
%   (1) Ouverture : le non-codifiable et le déplacement du canal
%   (2) Bell, Gray et le brevet : focusing device inversé
%   (3) Le monopole et les modèles institutionnels (Vail, 4 pays)
%   (4) La commutation et les demoiselles (Feigenbaum & Gross)
% Bilan : tableau D1–D10 + 0-TER
% Sources : VdK Telephone, Bertho 1984 (pp. 51–79), Feigenbaum & Gross
%   (QJE 2024, Management Science 2024), Mueller 1993 (via VdK)

Au milieu des années~1870, l'industrie du télégraphe est la plus puissante des industries de communication. Western Union, constituée en~1856, est l'une des premières grandes entreprises américaines~; elle possède des dizaines de milliers de kilomètres de lignes, emploie des milliers d'opérateurs, et génère des profits considérables. En Europe, les administrations d'État (Prusse, France, Angleterre) contrôlent des réseaux tout aussi étendus. C'est dans ce terreau industriel~--- dans ses ateliers, avec ses ingénieurs, contre ses problèmes~--- que naît le téléphone.

Le problème qui mobilise les meilleurs esprits de cette industrie est celui de la saturation des lignes. Les conducteurs de cuivre coûtent cher, et chaque fil ne porte qu'un message à la fois. Faire passer plusieurs messages simultanément sur un même fil~--- le multiplexage~--- est le graal technique des années~1870. Edison y travaille avec le quadruplex (1874), qui permet quatre transmissions simultanées. Elisha Gray, ingénieur chez Western Union depuis~1870, y travaille avec la télégraphie harmonique, qui exploite la superposition de fréquences. Ce sont des hommes du télégraphe, formés par le télégraphe, qui raisonnent dans les catégories du télégraphe~: la vitesse, le débit, le nombre de messages par fil.

Mais l'idée de transmettre la voix humaine par un circuit électrique est plus ancienne que cette course, et ne vient pas seulement du monde anglo-saxon. En~1854, Charles Bourseul, employé de l'administration française des télégraphes, avait décrit dans un article de \emph{L'Illustration} le principe d'un appareil capable de transmettre la parole, sans jamais parvenir à le réaliser\footnote{Bourseul écrivait~: \og Supposez que l'on parle près d'une plaque mobile assez flexible pour ne perdre aucune des vibrations produites par la voix, que cette plaque établisse et interrompe successivement la communication avec une pile~: vous pourrez avoir à distance une autre plaque qui exécutera exactement les mêmes vibrations.\fg{} Cité par \textcite{bertho_histoire_1984}, p.~51.}. En~1861, l'Allemand Philipp Reis construisit un appareil qu'il appela \emph{Telephon}, capable de transmettre des sons musicaux et, dans des conditions favorables, des fragments de parole~; il le présenta devant la Société de physique de Francfort, mais ne parvint pas à intéresser les industriels du télégraphe, qui n'y virent qu'un jouet acoustique\footnote{Reis mourut en~1874, deux ans avant le brevet de Bell, sans avoir obtenu de reconnaissance pour son travail. L'appareil de Reis transmettait le son par variation de résistance d'un contact, pas par variation continue de courant comme le ferait celui de Bell~; la distinction technique est réelle mais n'enlève rien à l'antériorité du concept.}. Antonio Meucci, émigré italien installé à New York, déposa un \emph{caveat} en~1871 pour un dispositif de communication vocale, mais ne put payer les frais de renouvellement et perdit sa protection en~1874.

\dimmark{D10}{+}%
Ces échecs ne sont pas des accidents biographiques. Ce qu'ils révèlent est un fait structurel~: le concept de transmission de la voix existait depuis au moins vingt ans, dans trois pays et trois langues, mais il ne trouvait pas de véhicule institutionnel capable de le porter. L'industrie du télégraphe, seule infrastructure de communication à longue distance, n'en voulait pas~; les inventeurs isolés (Bourseul, Reis, Meucci) n'avaient ni les capitaux ni les réseaux d'influence nécessaires pour transformer un prototype en système commercial. Il fallut que deux hommes issus de cette industrie~--- mais chacun à sa manière en marge d'elle~--- arrivent au même point presque au même moment.


\subsection{Bell, Gray et le brevet~174,465}\label{subsec:bell}

Le 14~février~1876, deux demandes concurrentes arrivèrent au bureau des brevets de New York à quelques heures d'intervalle. Alexander Graham Bell, par l'entremise de son avocat Anthony Pollok, déposa une demande de brevet pour un appareil de transmission de la voix par courant électrique. Elisha Gray, par l'entremise de son avocat William Baldwin, déposa un \emph{caveat}~--- une déclaration d'intention de breveter~--- pour un dispositif de même nature\footnote{La chronologie exacte a fait l'objet de controverses durables. \textcite{evenson_telephone_2000} a montré que le \emph{caveat} de Gray, bien qu'arrivé physiquement au bureau des brevets avant la demande de Bell, fut enregistré dans le journal quotidien (\emph{cash blotter}) après celle-ci, en raison d'une procédure accélérée demandée par Pollok. La question de savoir si Bell avait eu connaissance du contenu du \emph{caveat} de Gray par l'intermédiaire de l'examinateur Zenas Fisk Wilber reste débattue~; voir \textcite{shulman_telephone_2008} pour l'enquête la plus récente.}. Le brevet de Bell, n\textsuperscript{o}~174,465, fut accordé le 7~mars~1876.

Gray était un professionnel du télégraphe, en poste chez Western Union depuis~1870, rompu aux problèmes du multiplexage. Bell était un outsider~: issu d'une famille d'émigrés anglais installée au Canada, spécialisée depuis trois générations dans l'éducation des sourds, il avait étudié les problèmes de l'audition et de la voix humaine avant de s'intéresser à l'électricité \parencite{bertho_histoire_1984}. S'il cherchait un télégraphe harmonique, c'est qu'il savait que les compagnies télégraphiques feraient la fortune de son inventeur~; mais il devait lutter contre ses propres commanditaires, Gardiner Hubbard et Thomas Sanders, pour se détourner du multiplexage et se consacrer au téléphone, dont aucun prototype ne fonctionnait encore au moment du dépôt du brevet. La tension entre ce que le financeur demande (le multiplexage, dont la valeur est immédiate et mesurable) et ce que l'inventeur poursuit (la voix, dont la valeur est incertaine) est une instance de la bidirectionnalité notée dans le préambule méthodologique~: la demande oriente la technique, mais la technique la plus transformatrice peut être celle qui ne répond à aucune demande formulée.

\dimmark{D5}{?}%
Ce qui importe pour l'analyse n'est pas la question de priorité, que les tribunaux tranchèrent en faveur de Bell après six cents procès étalés sur dix-huit ans, mais deux faits que l'épisode met en lumière. Le premier est un fait juridique. En~1881, dans l'affaire \emph{Spencer}, le juge Lowell statua que Bell avait créé un \og \emph{new art}\fg{}~--- une technique entièrement nouvelle~--- et qu'il avait droit à la revendication la plus large possible sur la transmission électrique de la parole\footnote{\og Bell had discovered a new art~--- that of transmitting speech by electricity~--- and has a right to hold the broadest claim for it which can be permitted in any case\fg{} (\textcite{beauchamp_who_2010}, p.~863).}. Cette qualification de \og nouvel art\fg{} transformait un brevet d'appareil en brevet de principe, couvrant non pas un dispositif particulier mais la totalité du champ technique. C'est ce brevet de principe qui, pendant dix-huit ans, donna à Bell et à ses successeurs un monopole légal sur la téléphonie aux États-Unis. Le mécanisme de monopole qui en résultait~--- une barrière juridique fondée sur la portée du brevet, non sur les rendements de réseau ni sur le verrouillage d'un consommable~--- était distinct de ceux qu'avaient fait apparaître le télégraphe et la tabulatrice.

\dimmark{D5}{?}%
Le second fait est un fait de perception. Gray, qui travaillait pour Western Union depuis~1870 et avait résolu plusieurs problèmes du multiplexage télégraphique, ne voyait dans le téléphone qu'une curiosité~: \og Le téléphone parlant, écrivait-il le lendemain à son associé, est un bel objet du point de vue scientifique\ldots{} mais si on le considère sous l'angle commercial, il n'a aucune valeur. Avec un fil on peut faire actuellement beaucoup plus, et la vitesse est la seule chose qui nous intéresse\fg{}\footnote{Cité par \textcite{bertho_histoire_1984}, p.~53. La formulation anglaise, rapportée par \textcite{hounshell_elisha_1975}, est~: \og \ldots talking telegraph is a beautiful thing in a scientific point of view\ldots{} But if you look at it in a business light it is of no importance.\fg{}}. La même année, William Orton, président de Western Union, refusa d'acheter les brevets de Bell pour cent mille dollars\footnote{\textcite{hounshell_elisha_1975}, p.~157. Orton voulait un télégraphe multiplex, pas un \og scientific curiosity\fg{}. Western Union changea d'avis l'année suivante et tenta de contester les brevets de Bell, sans succès.}. Ce que les professionnels du télégraphe ne parvenaient pas à voir, c'est que la valeur du téléphone résidait précisément dans ce que le télégraphe ne pouvait pas faire~: transmettre le non-codifiable. Leur grille d'évaluation~--- le nombre de messages par fil, la vitesse de transmission, le coût par mot~--- était celle d'un système fondé sur le codage, et elle rendait invisible la valeur d'un système fondé sur la voix. C'est un cas de ce que \textcite{rosenberg_perspectives_1976} appelle un \emph{focusing device} inversé~: le paradigme technique en place concentre l'attention sur certains problèmes (le multiplexage, la vitesse) au point d'aveugler les praticiens sur des possibilités qui relèvent d'un autre registre.

\dimmark{D6/D7}{+}%
Ce que le téléphone rend transmissible, c'est précisément ce que le codage télégraphique détruisait~: le ton d'une négociation, la crédibilité d'un interlocuteur, l'urgence d'une instruction. Mais la portée de cette rupture ne se manifesta pas immédiatement. \textcite{ploeckl_market_2023}, étudiant la diffusion du téléphone en Bavière entre~1883 et~1905 à travers 306~centraux téléphoniques, a montré que le téléphone était d'abord une commodité urbaine locale~: la taille de la population de la ville~--- et non les connexions avec d'autres villes~--- dominait largement la décision d'ouvrir un central et le nombre d'abonnés. Les premiers utilisateurs s'appelaient dans la même ville, pas à longue distance~; la présence d'une gare n'accélérait pas la diffusion et avait même un effet négatif sur l'adoption, ce qui signale que le téléphone, contrairement au télégraphe, naît hors de toute filiation avec le réseau ferroviaire. Ce n'est que dans les villes plus petites et plus tardives que l'accès aux réseaux extérieurs commença à peser sur l'adoption. Le téléphone ne déplaça donc pas d'emblée le canal par lequel la transmission agit sur l'économie~; il commença par quelque chose que le télégraphe ne faisait pas du tout~: la communication de proximité, intra-urbaine, quotidienne. Le déplacement vers la coordination organisationnelle à longue distance~--- l'AT\&T de Vail, la ligne transcontinentale de~1915~--- fut un processus de trente ans, pas un fait contemporain de l'invention\footnote{\textcite{juhasz_spinning_2018} ont proposé la notion de filtre de codifiabilité pour caractériser ce que les techniques de transmission rendent échangeable à distance. \textcite{steinwender_real_2018} a montré que le télégraphe déplaçait cette frontière en créant un marché unifié du coton entre Liverpool et New York. Le téléphone la déplace une seconde fois, mais la conséquence économique de ce déplacement~--- l'action sur les organisations plutôt que sur les marchés~--- ne se manifeste qu'une fois le réseau longue distance construit. Cette séquence sera reprise dans le bilan de cette section.}. Et c'est précisément ce décalage temporel qui éclaire le jugement de Gray~: en~1876, le téléphone était effectivement sans valeur dans les catégories du télégraphe, parce que sa valeur résidait dans un usage~--- la communication locale~--- que ces catégories ne savaient pas mesurer.


\subsection{Le Bell System et la construction du monopole}\label{subsec:bell_system}

Le brevet de Bell expirait en~1894. Les dix-huit années qui le séparaient de son dépôt furent employées par ses successeurs à construire, autour de la protection légale, un système intégré dont la logique dépassait de loin celle du brevet lui-même. La \emph{Bell Telephone Company}, créée en juillet~1877 par Bell et ses associés Gardiner Hubbard et Thomas Sanders, était une petite entreprise dont les ressources ne permettaient ni de fabriquer les appareils en série ni d'exploiter les réseaux. Le modèle choisi fut celui de la concession~: Bell accordait à des opérateurs locaux le droit d'exploiter le service téléphonique sur un territoire donné, en échange d'une part de leur capital et de l'obligation de louer (non d'acheter) les appareils Bell. En~1880, les trois cents opérateurs locaux ainsi créés desservaient environ 48\,000~abonnés, concentrés dans les grandes villes\footnote{Données de \textcite{joel_history_1975}, reprises par VdK. En~1895, 57\,\% des téléphones étaient dans des villes de plus de 50\,000~habitants, et 3\,\% seulement dans les zones rurales \parencite{mueller_universal_1993}.}.

\dimmark{D2}{*}%
En~1879, un accord mit fin à la première bataille entre Bell et Western Union~: Bell obtenait le marché du téléphone, Western Union conservait celui du télégraphe. Cet accord, négocié après que les tribunaux eurent reconnu la validité des brevets Bell, eut pour effet de séparer les deux industries de communication, qui partageaient pourtant les mêmes poteaux, les mêmes fils de cuivre et, dans une large mesure, les mêmes ingénieurs. La séparation n'était pas technique~; elle était juridico-commerciale, et elle allait durer un siècle.

C'est Theodore Vail qui donna au système sa structure durable. Vail n'était pas un technicien~: il venait du service postal américain, où il avait réorganisé l'acheminement du courrier par voie ferrée \parencite{bertho_histoire_1984}. C'était, selon la formule de Bertho, \og un stratège qui pense en termes de réseaux\fg{}. Sous sa direction, trois décisions architecturèrent le monopole. La première fut de déplacer la compagnie à New York et de lui donner des statuts permettant la prise de participations dans les compagnies locales~: les opérateurs cessèrent d'être des licenciés indépendants pour devenir des filiales dont la maison mère définissait la politique d'ensemble. La seconde fut de racheter, en~1882, la \emph{Western Electric Manufacturing Company}~--- alors la plus grande usine de matériel électrique des États-Unis, fondée par Elisha Gray et Enos Barton~--- pour intégrer la production du matériel. Aucun autre pays ne connaîtrait cette intégration de l'opérateur et du fabricant dans une même structure \parencite{bertho_histoire_1984}. La troisième fut de créer, en~1885, l'\emph{American Telephone and Telegraph Company} (AT\&T) pour les liaisons longue distance entre les réseaux locaux~: celui qui contrôlait l'interurbain contrôlait l'accès de chaque réseau local au reste du monde.

\dimmark{D2}{*}%
Le mécanisme de monopole qui en résultait procédait d'une logique distincte de celles observées dans les cas précédents. Le télégraphe avait convergé vers le monopole par les rendements de réseau~: chaque nouvel abonné augmentait la valeur du réseau pour tous les autres, et le premier opérateur à atteindre la masse critique éliminait ses concurrents (§\,\ref{sec:telegraph}). La tabulatrice avait convergé vers le monopole par le verrouillage du consommable~: les cartes perforées liaient le client à IBM indépendamment de la qualité de la machine (§\,\ref{sec:mechanography}). Le téléphone convergeait vers le monopole par l'intégration verticale et le contrôle de l'interconnexion~: Vail ne se contentait pas de posséder le réseau le plus étendu~; il possédait aussi l'usine qui fabriquait les appareils, le réseau longue distance qui connectait les villes entre elles, et le portefeuille de brevets qui interdisait aux concurrents d'utiliser les mêmes composants. Six cents procès en contrefaçon maintinrent cette barrière pendant dix-huit ans.

\emergence{D2~--- Le monopole du téléphone procède d'un troisième mécanisme, distinct des rendements de réseau (télégraphe) et du verrouillage du consommable (tabulatrice)~: l'intégration verticale (fabricant + opérateur + longue distance) combinée au contrôle de l'interconnexion. Trois technologies de communication, trois mécanismes de monopole, un même résultat~--- la concentration du marché. \questionch{2}{D2~--- Si le résultat (monopole) est stable mais le mécanisme varie, qu'est-ce qui détermine le choix du mécanisme~? La nature de la technologie, la structure institutionnelle, ou le hasard historique~?}}


\subsection{Quatre pays, quatre trajectoires, un résultat}\label{subsec:quatre_pays}

L'expérience américaine n'est qu'un des quatre modèles institutionnels par lesquels le téléphone fut organisé dans les pays industrialisés. Leur comparaison révèle que la convergence vers l'opérateur unique n'est pas un fait technique mais un fait politique dont la forme varie selon le contexte.

\dimmark{D2/D7}{*}%
En Prusse, le téléphone fut immédiatement placé sous l'autorité de l'administration des postes, sans phase de concession privée~: le monopole d'État du télégraphe s'étendit au téléphone par décision administrative. En Angleterre, fidèle à sa doctrine de libéralisme économique, le gouvernement accorda des concessions à des compagnies privées, chacune opérant un réseau urbain dans sa zone d'activité. Mais le \emph{Post Office} perçut rapidement le danger que le téléphone représentait pour les recettes du télégraphe, et se réserva d'abord les lignes interurbaines (1899), puis nationalisa la totalité des réseaux en~1912 \parencite{bertho_histoire_1984}.

Le cas français est le plus instructif, parce qu'il met en jeu un mécanisme institutionnel absent des cas américain et prussien. En~1879, le gouvernement de la IIIe~République accorda des concessions téléphoniques à des sociétés privées. Trois candidats fusionnèrent avant la fin de l'année pour former la \emph{Société générale des téléphones} (SGT), dotée d'un capital de 25~millions de francs, qui prit en charge les réseaux de Paris, Lyon, Marseille et de quelques villes de province. Les termes de la concession prévoyaient un partage des responsabilités entre la SGT (les abonnés et les centraux) et l'administration (les câbles), dont la cohabitation fut difficile dès le début. Lorsque, en~1884, l'incertitude sur le renouvellement de la concession devint palpable, la SGT cessa d'investir~: le service se dégrada, les listes d'attente s'allongèrent, les abonnés d'affaires se plaignirent des tarifs \parencite{bertho_histoire_1984}. La tentative de la SGT de négocier secrètement avec le ministre des Finances un accord de \og compagnie fermière\fg{} pour trente ans provoqua un scandale qui précipita la nationalisation.

\dimmark{D7}{+}%
Ce qui est remarquable dans l'épisode français n'est pas la nationalisation elle-même~--- on l'a vue aussi en Angleterre~--- mais la justification qui la sous-tend. En~1837, quand le Parlement français avait débattu du statut du télégraphe, l'argument dominant était celui de la sécurité de l'État~: le télégraphe devait être un monopole public parce que la transmission rapide d'informations était une affaire de souveraineté (§\,\ref{sec:telegraph}). En~1889, quand la nationalisation du téléphone fut décidée, l'argument avait changé de nature. La chambre de commerce de Valenciennes, dans une délibération de~1887, formula la position qui allait prévaloir~: \og L'État ne doit pas se dessaisir de l'organisation de grands services d'utilité publique qui n'ont aucunement le caractère d'une industrie particulière, mais dont le but est de mettre à la portée de tous les moyens d'action capables de faciliter les progrès du travail et de la production nationale\fg{}\footnote{Cité par \textcite{bertho_histoire_1984}, p.~62. La délibération de la chambre de commerce de Valenciennes se prononçait \og contre toute concession de ligne téléphonique à une compagnie privée\fg{}.}. Le monopole de sécurité du télégraphe avait cédé la place à un monopole de \emph{service public}~: même forme institutionnelle, fondement politique entièrement différent. Dans les deux cas, l'État prenait le contrôle du réseau~; dans le premier, il le faisait pour garder le pouvoir~; dans le second, pour fournir un service universel.

La nationalisation française eut cependant un effet que ses promoteurs n'avaient pas anticipé. Le financement des réseaux téléphoniques fut organisé par le système des \og avances remboursables\fg{}~: les collectivités locales avançaient les fonds, remboursés sur les recettes du réseau. Ce mode de financement calqua la géographie du réseau sur la géographie administrative plutôt que sur la rationalité économique~: chaque canton, parfois chaque commune, eut son propre réseau, mal relié aux autres. En~1921, la France comptait 14\,000~centraux téléphoniques pour 7\,000~en Allemagne, alors que la population allemande était plus nombreuse et plus urbanisée \parencite{bertho_histoire_1984}. Le défaut principal du réseau~--- l'insuffisance des circuits interurbains~--- découlait directement du mode de financement. À partir de~1910, techniciens et grand public convergèrent dans le même diagnostic~: le téléphone français était en crise. Ce que Bertho nomme des \og vices de structure\fg{} ne relevaient pas de la technique mais de l'architecture institutionnelle. Des travaux ultérieurs ont confirmé que cette fragmentation avait des conséquences macroéconomiques mesurables~: \textcite{roller_telecommunications_2001}, étudiant 21~pays de l'OCDE entre~1970 et~1990 avec un modèle simultané qui traite la causalité inverse, ont montré que l'effet de l'infrastructure téléphonique sur la croissance n'apparaît qu'au-delà d'un seuil de pénétration d'environ 40\,\%~--- un seuil que la fragmentation française, avec ses 14\,000~centraux mal reliés, rendait structurellement difficile à atteindre.

\dimmark{D8}{+}%
La période de monopole du brevet Bell (1879--1894) avait été, pour les actionnaires, la plus profitable de l'histoire de l'industrie~: \textcite{gabel_early_1969}, exploitant les données du Walker Report de la FCC (1938), a calculé un rendement de près de quarante-six pour cent sur la période, avec des dividendes moyens de quinze pour cent. Mais cette rentabilité avait pour contrepartie la stagnation~: des tarifs élevés, des installations insuffisantes, un service limité au centre-ville des grandes villes, et un développement résidentiel, suburbain et rural à peine esquissé. Aux États-Unis, la fin de la protection par brevet, en~1894, ouvrit une période de concurrence violente qui produisit le taux de croissance le plus rapide de l'histoire de l'industrie. En~1900, plus de six mille compagnies indépendantes opéraient sur le territoire américain, et le parc téléphonique avait été multiplié par quatre en cinq ans \parencite{mueller_universal_1993}. Soixante pour cent des villes américaines de plus de cinq mille habitants avaient deux réseaux concurrents en~1904, ce qui forçait les abonnés au \og \emph{dual service}\fg{}~: souscrire aux deux pour pouvoir joindre tout le monde. Les indépendants couvraient les zones rurales et les petites villes que le Bell System avait négligées par calcul~: tant que les brevets protégeaient le monopole, il n'y avait pas de raison de desservir les marchés peu rentables\footnote{\textcite{mueller_universal_1993} a montré que la notion de \og \emph{universal service}\fg{}, telle que Vail l'employa pour la première fois dans le rapport annuel d'AT\&T de~1907, ne signifiait pas un téléphone dans chaque foyer mais l'interconnexion des réseaux en un système unifié. C'est la concurrence, non le monopole, qui avait étendu le réseau géographiquement~; le slogan de Vail visait à légitimer la reconsolidation, pas à décrire une réalité.}. En~1907, la situation était critique~: les profits de Bell baissaient. Un syndicat de banquiers, présidé par J.~P.~Morgan, prit le contrôle de la compagnie et rappela Vail.

Le second Vail réorganisa l'entreprise autour de deux piliers. Le premier était la recherche fondamentale~: il assigna aux laboratoires d'ingénierie (qui deviendraient les Bell Labs en~1925) la mission de construire une ligne transcontinentale entre New York et San Francisco, soit six mille kilomètres, alors que la ligne la plus longue (Boston-Chicago, 1893) ne dépassait pas mille neuf cents kilomètres. La solution vint de la triode de Lee de Forest (1906), dont AT\&T acheta les brevets en~1914 pour neuf cent mille dollars~: l'amplification du signal par le tube à vide rendit la transmission longue distance possible, et la ligne transcontinentale fut ouverte au trafic commercial le 15~janvier~1915 \parencite{bertho_histoire_1984}. Le second pilier était la régulation négociée. En~1913, face à une enquête antitrust, Vail proposa ce qui est resté dans l'histoire sous le nom de \emph{Kingsbury Commitment}~: AT\&T revendrait ses participations dans Western Union, cesserait d'acquérir des compagnies indépendantes sans autorisation, et ouvrirait ses lignes longue distance aux concurrents, moyennant un droit de passage. L'accord semblait pro-concurrentiel~; dans les faits, il institutionnalisa le monopole. La clause qui limitait les acquisitions fut interprétée de manière permissive~: AT\&T pouvait acheter un opérateur à condition d'en revendre un autre de taille équivalente, ce qui revenait à un échange de territoires plutôt qu'à une ouverture du marché. En~1924, AT\&T avait absorbé 223~des 234~compagnies indépendantes \parencite{weiman_preying_1994}. Au milieu des années~1920, l'ère de la concurrence était terminée.

\dimmark{D2}{*}%
Quatre pays, quatre trajectoires institutionnelles~--- le monopole privé intégré (États-Unis), le monopole d'État immédiat (Prusse), la concession suivie de nationalisation (France, Angleterre)~--- et un même résultat~: un seul opérateur par territoire. Le pattern observé dans les sept trajectoires du télégraphe (§\,\ref{sec:telegraph}) se confirme, ce qui renforce l'hypothèse que la convergence vers le monopole est un fait structurel des réseaux de communication, pas un accident institutionnel. Mais le mécanisme est chaque fois différent~: rendements de réseau pour le télégraphe, intégration verticale pour le Bell System, décision administrative pour la Prusse, défaillance de la concession pour la France. Ce qui varie n'est pas le résultat mais le chemin~--- et le chemin, à son tour, détermine les pathologies du monopole résultant (fragmentation française, prédation américaine, rigidité prussienne).


\subsection{La commutation~: du calcul dans le réseau}\label{subsec:commutation}

Le téléphone, à la différence du télégraphe, exigeait une opération que le réseau télégraphique ne connaissait pas~: la commutation. Un message télégraphique était adressé à une station, où un opérateur le transcrivait et le remettait au destinataire~; le routage se faisait par la géographie des lignes, pas par un geste de connexion en temps réel. Un appel téléphonique supposait au contraire qu'un circuit physique fût établi entre deux abonnés, pour la durée de la conversation, parmi des milliers ou des millions de connexions possibles. Ce circuit devait être créé à la demande, maintenu le temps de l'échange, puis défait. La personne qui accomplissait ce geste~--- chercher le numéro demandé sur le tableau, enfoncer une fiche dans la prise correspondante, surveiller la fin de l'appel~--- était l'opératrice de commutation.

\dimmark{D4/D7}{+}%
Le geste paraissait simple~; sa complexité croissait quadratiquement avec le nombre d'abonnés. Pour $N$~abonnés, le nombre de connexions possibles est $N(N-1)/2$~: un réseau de mille abonnés suppose près d'un demi-million de circuits potentiels. À mesure que le nombre d'abonnés augmentait, les tableaux de commutation s'étendaient, le nombre d'opératrices croissait, et la qualité du service se dégradait, parce que chaque appel devait transiter par un nombre croissant d'intermédiaires. Le \emph{switching} n'était pas une tâche de transmission~; c'était une tâche de traitement de l'information~--- du routage, du calcul implicite~--- intégrée au c\oe{}ur du réseau de transmission. C'est la première fois dans le parcours du chapitre que les fonctions de TRANSMIT et de COMPUTE apparaissent inséparables dans un même système opérationnel.

L'automatisation du \emph{switching} fut tentée dès~1891, lorsque Almon Brown Strowger, un entrepreneur de pompes funèbres de Kansas City qui soupçonnait une opératrice de détourner ses appels vers un concurrent, breveta un commutateur mécanique à pas de progression (brevet US n\textsuperscript{o}~447,918). Le système de Strowger permettait à l'abonné de composer lui-même le numéro à l'aide d'un cadran rotatif, chaque impulsion actionnant un sélecteur électromécanique au central. La technologie fonctionnait~; quelques compagnies indépendantes l'adoptèrent dans les années~1900. Mais AT\&T, qui possédait le plus grand réseau et employait le plus grand nombre d'opératrices, n'adopta le commutateur mécanique qu'à partir de~1919, et la diffusion prit soixante ans supplémentaires pour couvrir l'ensemble du réseau.

\dimmark{D7/D10}{+}%
Ce délai est un fait surprenant au regard du protocole de ce chapitre~: une technologie disponible depuis les années~1890, manifestement capable de réduire les coûts de main-d'\oe{}uvre, et déployée par le concurrent le mieux équipé pour l'adopter, mit près d'un siècle à se généraliser. \textcite{feigenbaum_organizational_2024} ont montré que l'explication ne résidait ni dans un défaut de la technologie ni dans une résistance syndicale, mais dans les interdépendances organisationnelles entre la commutation et le reste du système Bell. Les opératrices n'étaient pas seulement des connectrices de circuits~; elles étaient le pivot d'un système de production complexe qui s'était construit autour d'elles. Elles comptaient les appels pour la facturation~; elles géraient les plaintes en temps réel~; elles fournissaient les renseignements et les services d'urgence~; elles connaissaient les numéros des abonnés réguliers. Automatiser le \emph{switching} exigeait de réinventer chacune de ces fonctions adjacentes~: créer des compteurs automatiques, des services de renseignement séparés, des procédures d'urgence sans opérateur, des annuaires imprimés plus complets. La commutation était ce que la théorie de l'architecture organisationnelle appelle une \emph{tâche intégrale} \parencite{ulrich_role_1995}~--- une tâche qui interagit avec toutes les autres et dont la modification implique la reconfiguration de l'ensemble.

\emergence{D7/D10~--- L'automatisation d'une tâche apparemment simple (la commutation) prend soixante ans parce que cette tâche est \emph{intégrale} au système de production. Le \emph{switching} portait de l'intelligence organisationnelle, pas seulement de la connexion physique. Hypothèse abductive~: la résistance à l'automatisation est proportionnelle au degré d'intégralité de la tâche dans le système. \questionch{2}{D7~--- Le concept de tâche intégrale s'applique-t-il aux vagues d'automatisation ultérieures (computing, robotique, IA)~?}}


\subsection{Les demoiselles du téléphone}\label{subsec:demoiselles}

\dimmark{D4}{*}%
Les premières opératrices de commutation furent, à New York en~1879, de jeunes garçons recrutés parmi les porteurs de dépêches télégraphiques. Ils se révélèrent, selon les termes de l'époque, \og trop remuants\fg{} pour se plier à la discipline physique du tableau de commutation~: rester assise, bridée par le casque d'écoute et le microphone, incapable de se lever ni de se rasseoir, tenue à une courtoisie sans faille envers des abonnés souvent irrités. Dès~1880, la \emph{Société générale des téléphones} à Paris n'embauchait plus que des femmes~; AT\&T fit de même aux États-Unis. En~1920, les opératrices de téléphone constituaient l'un des emplois féminins les plus répandus en Amérique~: environ 4\,\% des jeunes femmes blanches nées aux États-Unis qui travaillaient étaient opératrices, et AT\&T était le plus grand employeur de femmes du pays, avec des opératrices représentant plus de la moitié de sa main-d'\oe{}uvre \parencite{feigenbaum_answering_2024}.

Le profil de recrutement était précis~: diplômée du \emph{high school}, au moins dix-huit ans, vivant chez ses parents ou un proche, en bonne santé physique, dotée d'une voix agréable, d'une dextérité manuelle, d'une capacité à calculer rapidement et d'un tempérament stable, \og \emph{not easily ruffled by irritable customers}\fg{} \parencite{erickson_telephone_1946}. Les opératrices étaient presque exclusivement blanches, nées aux États-Unis~; AT\&T ne recruta pas de femmes noires avant~1940 \parencite{green_race_2001}. Le poste était mieux payé que le travail en usine, offrait une interaction humaine que le travail de bureau ne garantissait pas, mais imposait un rythme intense et un taux de rotation de l'ordre de 40\,\% par an \parencite{feigenbaum_answering_2024}.

Le parallèle avec le cas du calcul mécanisé (§\,\ref{subsec:zoom_tabulating}) est frappant. Au \emph{Nautical Almanac}, des calculateurs qualifiés, autonomes, travaillant à domicile, avaient été remplacés par des opératrices moins payées, recrutées dans les catégories inférieures de la fonction publique, supervisées de près. Au central téléphonique, le même schéma se reproduit~: une tâche exigeant des qualités \og spécifiquement féminines d'habileté, de patience, de résistance nerveuse\fg{} \parencite{bertho_histoire_1984} est confiée à des jeunes femmes dont la principale vertu, du point de vue du recruteur, est d'accepter un salaire inférieur à celui que des hommes auraient exigé pour le même travail. La féminisation et la compression salariale documentées par \textcite{davies_womans_1982} et \textcite{strom_beyond_1992} pour les machines de bureau se reproduisent dans le téléphone, ce qui renforce l'hypothèse que ce pattern est structurel et non contingent.

\dimmark{D4}{*}%
L'automatisation de la commutation, entre~1920 et~1940, élimina la majorité de ces emplois, une ville à la fois. \textcite{feigenbaum_answering_2024}, exploitant la variation dans le calendrier d'adoption du commutateur mécanique à travers près de trois mille villes américaines, ont montré que le nombre de jeunes opératrices chutait de 50 à 80\,\% dans les années qui suivaient l'installation du système automatique. L'effet sur les opératrices en poste était sévère~: une décennie après l'automatisation, celles qui travaillaient encore occupaient des emplois moins bien payés que leurs homologues dans les villes non automatisées, et les plus âgées d'entre elles avaient quitté le marché du travail. L'automatisation détruisait non seulement des emplois mais des carrières~: le \emph{switching} était l'un des rares métiers féminins de l'époque qui offrait une perspective de long terme, et sa disparition fermait cette possibilité.

\dimmark{D4}{+}%
Mais le résultat le plus significatif de l'étude concerne les cohortes suivantes. Les jeunes femmes qui entraient sur le marché du travail après l'automatisation ne connaissaient pas de baisse de leur taux d'emploi global. Les emplois perdus dans la commutation étaient compensés par une croissance de l'emploi dans le secrétariat (emplois de qualification comparable) et dans la restauration (emplois de qualification inférieure). Plus remarquable encore, la croissance de l'emploi secrétarial se concentrait dans des industries qui n'avaient pas employé de secrétaires auparavant~--- de nouveaux pairages occupation-industrie, ce que \textcite{autor_new_2024} appellent du \og \emph{new work}\fg{}. L'automatisation n'avait pas seulement déplacé les travailleuses d'un secteur vers un autre~; elle avait engendré, indirectement, la création d'usages nouveaux du travail féminin dans des secteurs qui n'en avaient pas eu besoin jusque-là. Ce mécanisme de réinstauration des tâches, théorisé par \textcite{acemoglu_race_2018}, était cependant conditionné~: dans les villes manufacturières, où les emplois féminins en col blanc étaient rares, et dans les villes frappées le plus durement par la Grande Dépression, l'automatisation réduisait l'emploi sans le compenser. La demande agrégée est une condition nécessaire pour que le mécanisme de réinstauration opère.

\emergence{D4~--- L'automatisation de la commutation (1920--1940) confirme et étend le pattern du \emph{Nautical Almanac}~: la mécanisation réorganise le travail plutôt qu'elle ne le supprime. Les incumbentes souffrent~; les cohortes suivantes se réallouent. Le mécanisme de compensation passe par la création de \og \emph{new work}\fg{} (Autor et al.~2023) dans des industries qui n'employaient pas de secrétaires auparavant. Mais ce mécanisme est conditionné par la demande agrégée~: en période de contraction, l'automatisation réduit l'emploi sans le compenser. \questionch{2}{D4~--- La conditionnalité du mécanisme de réinstauration à la demande agrégée a-t-elle une implication pour les vagues d'automatisation contemporaines (IA)~?}}


\subsection{La bobine Pupin et la matérialité du réseau}\label{subsec:pupin}

\dimmark{D3/0-TER}{*}%
Un dernier fait mérite d'être relevé avant de clore le cas, parce qu'il concerne la dimension matérielle que le récit du monopole et de l'automatisation a laissée en arrière-plan. En~1900, la communication longue distance par téléphone se heurtait à un obstacle physique~: l'affaiblissement du signal sur les lignes de cuivre. La ligne Boston-New York, inaugurée en~1884, représentait à peu près la limite pratique. Un quart du capital investi par les compagnies téléphoniques allait au fil de cuivre, et l'allongement des distances exigeait des conducteurs de section croissante, donc de coût croissant. Michael Pupin, professeur de physique à Columbia, breveta en~1900 une bobine d'inductance qui, insérée à intervalles réguliers sur la ligne, réduisait l'affaiblissement et permettait de doubler la distance de communication à section de cuivre égale, ou de réduire de moitié la section nécessaire pour une distance donnée. AT\&T acheta les brevets de Pupin pour un total de 455\,000~dollars entre~1901 et~1917 \parencite{brittain_turning_1970}. L'économie de cuivre réalisée a été estimée à 100~millions de dollars pour le premier quart du vingtième siècle \parencite{brittain_turning_1970}, soit un ratio de 1~pour~200 entre le coût du brevet et l'économie de matière qu'il permit.

Ce ratio illustre un mécanisme de substitution du capital intellectuel (le brevet, la connaissance mathématique de la propagation des ondes) au capital physique (le cuivre). C'est un fait de même nature que la substitution K/L documentée dans le cas du télégraphe (§\,\ref{sec:telegraph}), mais il opère sur un autre facteur~: non pas le travail mais la matière première. La bobine Pupin ne transmettait rien~; elle permettait de transmettre avec moins de cuivre. Le fait que le poste de dépense dominant du réseau téléphonique fût le cuivre, pas l'énergie ni le travail, signale que la matérialité du réseau de transmission pesait lourd dans l'économie du système, même si cette matérialité restait invisible pour l'abonné.

\dimmark{0-TER}{+}%
Quant à la dépendance énergétique, le cas du téléphone marque une transition qualitative par rapport au télégraphe, même si elle est discrète. Le télégraphe de Morse tirait son énergie de piles voltaïques, remplacées localement quand elles s'épuisaient~; la consommation était intermittente et décentralisée. Le téléphone, dans ses premières années, fonctionnait sur le même modèle~: chaque poste avait sa propre pile. Mais l'introduction de la \og batterie centrale\fg{}, où une pile humide unique, placée au central, alimentait l'ensemble des postes connectés, transforma le réseau téléphonique en un système à consommation permanente et centralisée\footnote{L'invention de la batterie centrale est attribuée à Pavel Golubitsky en Russie (années~1880). Elle fut rapidement adoptée par l'ensemble des réseaux urbains parce qu'elle simplifiait la maintenance et permettait l'extension à grande échelle.}. C'est un basculement modeste en valeur absolue~--- la consommation d'un central téléphonique de~1900 était négligeable en regard de celle d'une centrale électrique~--- mais c'est le premier pas dans la séquence d'escalade qui mènera, par le tube à vide (§\,\ref{sec:wireless}), le mainframe (§\,\ref{sec:computing}) et le datacenter (§\,\ref{sec:contemporain}), à une dépendance énergétique d'un tout autre ordre de grandeur.


\subsection*{Bilan~: ce que le téléphone ajoute à la traversée}\label{subsec:bilan_telephone}

Le cas du téléphone confirme certains patterns établis dans les cas précédents et en casse d'autres. La convergence vers le monopole se confirme~: quatre pays, quatre trajectoires institutionnelles, un même résultat. Mais le mécanisme est nouveau (intégration verticale et contrôle de l'interconnexion), ce qui renforce l'hypothèse que c'est le mode de contrôle du système, pas la technologie, qui détermine la structure de marché. La féminisation et la compression salariale se confirment (les demoiselles du téléphone prolongent les calculatrices du \emph{Nautical Almanac}), mais l'automatisation ajoute un résultat que le cas du calcul mécanisé n'avait pas produit~: les cohortes suivantes se réallouent, et la demande de travail se reconstitue par la création de tâches nouvelles dans des industries qui n'en avaient pas eu l'usage.

Ce qui est propre au téléphone, et qui ne pouvait pas émerger des cas précédents, tient en quatre observations. La première est la séquence temporelle du déploiement~: le téléphone commence comme commodité urbaine locale (\textcite{ploeckl_market_2023}), puis la concurrence post-brevets (1894--1907) étend le réseau géographiquement (\textcite{gabel_early_1969}), et c'est seulement après~1907 que la reconsolidation sous Vail en fait un outil de coordination organisationnelle à longue distance. Le déplacement du canal d'action~--- de la communication locale vers la coordination intra-organisationnelle~--- est un processus de trente ans, pas un fait contemporain de l'invention\footnote{\textcite{norton_transaction_1992}, utilisant le stock initial de téléphones en~1957 comme prédicteur de la croissance~1957--1977, a formalisé le canal par lequel le téléphone agit sur l'économie~: la réduction des coûts de transaction (au sens de \textcite{leff_externalities_1984}). \textcite{portes_determinants_2005} utilisent le trafic téléphonique bilatéral entre pays comme proxy de friction informationnelle dans un modèle de gravité des flux de capitaux~--- c'est l'un des rares usages économétriques du téléphone comme \emph{variable} d'information. \textcite{bloom_distinct_2014} ont montré que les technologies de communication centralisent les décisions dans la firme tandis que les technologies d'information les décentralisent~; cette distinction, qui recoupe celle que le présent chapitre construit entre le télégraphe et le téléphone, sera reprise au chapitre~2.}.

La deuxième observation est la communication domestique. Le téléphone est la première technique de l'information du parcours de ce chapitre dont l'utilisateur final n'est pas une entreprise, un État ou un militaire, mais un particulier chez lui. \textcite{hardy_role_1980}, dans une analyse cross-country couvrant 60~nations entre~1960 et~1973, a trouvé que les téléphones résidentiels contribuaient davantage à la croissance économique que les téléphones d'affaires (coefficient de~0,249 contre~0,097 avec la consommation d'énergie comme indicateur de développement)~--- un résultat qu'il qualifie lui-même de \og \emph{quite unexpected}\fg{}. Ce résultat s'éclaire sous Ploeckl~: si le téléphone est d'abord une commodité locale, c'est l'usage résidentiel qui capte l'essentiel de la valeur initiale. \textcite{martin_culture_1998} a montré que les femmes accédèrent au téléphone par l'intermédiaire des hommes d'affaires qui faisaient connecter leur domicile, et développèrent des usages~--- la sociabilité, la coordination domestique~--- que l'industrie considérait comme parasitaires avant de s'y adapter. La diffusion elle-même n'est pas monotone~: \textcite{fischer_technologys_1987} a documenté la perte d'un million de téléphones dans les fermes américaines entre~1920 et~1940, une dévolution technologique qui montre que l'adoption d'une technique de l'information peut reculer quand les conditions économiques se dégradent.

La troisième est la nature intégrale de la commutation~: le \emph{switching} n'était pas une tâche isolable mais le pivot d'un système de production complexe, et sa résistance à l'automatisation pendant soixante ans en est la preuve. La quatrième est la transition du monopole de sécurité au monopole de service public~: en~1837, l'État prenait le contrôle du réseau pour garder le pouvoir~; en~1889, pour fournir un service universel.

\emergence{D10~--- Le cas de Gray révèle un mécanisme plus précis que le \emph{focusing device} de \textcite{rosenberg_perspectives_1976}. Ce n'est pas seulement que le paradigme en place dirige l'attention vers certains problèmes~; c'est que le système métrique du paradigme~--- le nombre de messages par fil, la vitesse, le coût par mot~--- rend impossible de \emph{mesurer la valeur} d'innovations qui opèrent dans un autre registre. Gray ne manque pas d'intelligence~; il manque d'unités~--- et cette absence d'unités n'est pas une carence cognitive individuelle mais le résultat d'un cadre institutionnel (l'industrie télégraphique, son modèle d'affaires fondé sur la facturation au mot, ses contrats d'emploi qui rémunèrent l'optimisation du multiplexage) qui n'a pas besoin de telles unités pour fonctionner et qui en décourage activement la production. Que Bell, lui-même, ait dû lutter contre ses commanditaires Hubbard et Sanders pour se détourner du multiplexage le confirme indirectement~: ce ne sont pas seulement les esprits qui sont structurés par le paradigme, ce sont aussi les contrats qui les financent. Le téléphone ne produit pas de \og messages\fg{} au sens du télégraphe, et sa valeur~--- la communication de proximité, la sociabilité, le non-codifiable~--- n'existe pas dans la métrique du paradigme en place. \questionch{2}{D10~--- Ce mécanisme d'occultation de la valeur par le système métrique du paradigme se reproduit-il dans les cas ultérieurs (téléphone $\rightarrow$ internet, PC $\rightarrow$ smartphone)~?}}

Le tableau~\ref{tab:telephone-dimensions} récapitule les observations dimensionnelles du cas.

\begin{table}[p]
\caption{Le téléphone au prisme des dix dimensions~: observations, niveaux de certitude et questions ouvertes}\label{tab:telephone-dimensions}
\centering\footnotesize
\renewcommand{\arraystretch}{1.15}
\begin{tabularx}{\textwidth}{@{}cXcp{4.5cm}@{}}
\toprule
\textbf{Dim.} & \textbf{Observations du cas §\,\ref{sec:telephone}} & & \textbf{Question pour la suite} \\
\midrule
D1 & Le téléphone transmet un signal continu (la voix), non un signal discret (le code Morse). L'information transmise n'est pas réductible à un alphabet~; sa valeur réside dans ce que le codage détruirait (intonation, hésitation, crédibilité). & + & Cette distinction codifiable/non-codifiable détermine-t-elle le canal par lequel l'ICT affecte l'économie~? \\[3pt]
D2 & Convergence vers le monopole par un troisième mécanisme (intégration verticale + contrôle de l'interconnexion). Confirmé dans quatre contextes institutionnels distincts (US, FR, UK, DE). Prédation documentée par \textcite{weiman_preying_1994}. & * & Le mécanisme de monopole dépend-il de la technologie, de l'institution, ou du hasard historique~? \\[3pt]
D3 & Ratio Pupin~: \$455k de brevets pour \$100M de cuivre économisé (1:200). Substitution du capital intellectuel au capital physique (matière, pas travail). & * & Ce ratio de substitution se maintient-il avec les technologies ultérieures~? \\[3pt]
D4 & Féminisation et compression salariale~: même pattern que le calcul mécanisé. L'automatisation (1920--1940) élimine les emplois d'opératrices mais ne réduit pas l'emploi total des cohortes suivantes \parencite{feigenbaum_answering_2024}. Réallocation vers le secrétariat et la restauration. \og \emph{New work}\fg{} dans des industries nouvelles. & * & Le mécanisme de réinstauration est-il conditionné par la demande agrégée dans tous les cas~? \\[3pt]
D5 & 48\,000 abonnés (1880) $\rightarrow$ 801\,000 (1900) $\rightarrow$ 7~millions (1906). Transition tarifaire~: du tarif au message (comme le télégraphe) au forfait. Les téléphones résidentiels contribuent davantage à la croissance que les business phones (\textcite{hardy_role_1980}~: coeff.~0,249 vs 0,097). La diffusion n'est pas monotone~: un million de fermes perdent le téléphone entre~1920 et~1940 \parencite{fischer_technologys_1987}. & + & Le passage au forfait et la domestication modifient-ils la structure de la demande de communication~? \\[3pt]
D6 & Pas d'étude d'identification causale de l'effet du téléphone sur l'intégration des marchés. \textcite{ploeckl_market_2023} montre que le téléphone était d'abord une commodité urbaine locale, pas un outil d'intégration de marchés distants. L'absence reflète un déplacement progressif du canal d'action (local $\rightarrow$ longue distance $\rightarrow$ organisationnel). & * & Ce déplacement temporel du canal d'action se retrouve-t-il dans d'autres technologies de communication~? \\[3pt]
D7 & Le \emph{switching} est une tâche intégrale (Ulrich 1995)~: il interagit avec toutes les fonctions du système Bell. 60~ans pour automatiser. Transition justification du monopole~: sécurité (1837) $\rightarrow$ service public (1889). & + & Le concept de tâche intégrale s'applique-t-il aux vagues d'automatisation ultérieures~? \\[3pt]
D8 & Cycle~: monopole de brevet (1876--1894, stagnation, rendement 46\,\% pour les actionnaires, \textcite{gabel_early_1969}) $\rightarrow$ concurrence (1894--1907, 6\,000 compagnies, croissance la plus rapide de l'industrie) $\rightarrow$ reconsolidation (Morgan/Vail, 1907--1925). Le \emph{Kingsbury Commitment} (1913) institutionnalise le monopole sous couvert de régulation. Masse critique~: l'effet macro n'apparaît qu'au-delà de 40\,\% de pénétration \parencite{roller_telecommunications_2001}. & * & Ce cycle ouverture-concurrence-reconsolidation est-il intrinsèque aux réseaux de communication~? \\[3pt]
D9 & Le cuivre est le poste de dépense dominant du réseau (25\,\% du capital investi). La bobine Pupin réduit la dépendance au cuivre par substitution intellectuelle. & + & La géographie de la matière première (cuivre) contraint-elle la géographie du réseau~? \\[3pt]
D10 & Le système métrique du paradigme en place (messages/fil/temps) rend impossible de mesurer la valeur d'innovations d'un autre registre. Gray ne manque pas d'intelligence~; il manque d'unités. Le téléphone, contrairement au télégraphe, naît hors du réseau ferroviaire \parencite{ploeckl_market_2023}~: la séparation communication/transport est empiriquement confirmée. & * & L'occultation de la valeur par le système métrique du paradigme se reproduit-elle dans les vagues ultérieures~? \\[3pt]
\midrule
0-TER & Le cuivre domine la structure de coûts (matière). L'énergie de fonctionnement est faible mais en transition~: des piles individuelles (décentralisé, intermittent) à la batterie centrale (centralisé, permanent). Premier basculement dans la séquence d'électrification de TRANSMIT. & + & À quel moment le coût énergétique dépasse-t-il le coût matériel dans les réseaux de communication~? \\[3pt]
\bottomrule
\end{tabularx}
\renewcommand{\arraystretch}{1.0}
\end{table}



% =============================================================================

% =============================================================================
% §1.6 La communication sans fil et la radiodiffusion (~1895–1930)
% =============================================================================
\clearpage
% §1.6 — La communication sans fil et la radiodiffusion (~1895–1930)
% VERSION : v1 — 27 avril 2026
% STRUCTURE : quatre temps + bilan
%   (1) Ouverture : la libération du fil (Hertz → Marconi, cluster VdK, brevet 1897)
%   (2) Le monopole Marconi et le modèle maritime (3 piliers, Poldhu, Telefunken, WWI → RCA)
%   (3) La puissance dans le signal (Aitken : spark < arc, Navy $1.5M, câbles vulnérables)
%   (4) Le spectre et le broadcasting (Hazlett 1990, Adena et al. 2015, Gentzkow 2006)
%   Bilan : tableau D1–D10 + 0-TER, trois hypothèses abductives
% SOURCES PRINCIPALES : VdK wireless_engine, Aitken 1985 (ch. 2–3, 5),
%   Hazlett 1990 (JLE), Adena et al. 2015 (QJE/WZB), Gentzkow 2006 (QJE),
%   Jolly 1972, Raboy 2016, Douglas 1987
% =============================================================================

\section{La communication sans fil et la radiodiffusion (~1895--1930)}\label{sec:wireless}

Le télégraphe et le téléphone transmettaient de l'information à travers un fil de cuivre~; le wireless transmet à travers l'air. Le changement paraît simple, mais il modifie la structure même du couplage entre information et matière. Dans les cas précédents, la matérialité du réseau résidait dans le substrat~: le cuivre, les poteaux, la gutta-percha, les isolateurs. Le coût de transmission croissait avec la distance parce que le fil croissait avec la distance. Avec le wireless, la matérialité se déplace du substrat au signal~: ce n'est plus la longueur du conducteur qui détermine la portée, c'est la puissance de l'émetteur. L'infrastructure n'est plus un réseau qu'on pose~; c'est une énergie qu'on émet. Ce déplacement, qui paraît libérateur parce qu'il affranchit la transmission du câble, introduit en réalité un rapport qualitativement nouveau à l'énergie~: pour la première fois dans le parcours de ce chapitre, le coût d'exploitation énergétique de la transmission est continu et croissant avec l'ambition du service.

\dimmark{D10}{+}%
La démonstration par Heinrich Hertz, entre 1886 et 1888, de l'existence des ondes électromagnétiques prédites par Maxwell ne déboucha pas sur une application~: Hertz lui-même ne voyait aucun usage pratique à sa découverte \parencite{jolly_marconi_1972}. Ce qui suivit, entre 1890 et 1896, fut un cluster d'innovation au sens de \textcite{kooij_information_2022}~: Lodge en Angleterre, Branly en France, Righi en Italie, Popov en Russie, Tesla aux États-Unis, Bose en Inde travaillèrent indépendamment sur la détection des ondes hertziennes, sans qu'aucun d'entre eux ne construise un système de communication opérationnel \parencite{jolly_marconi_1972, raboy_marconi_2016}. Le passage du phénomène de laboratoire au système technique fut l'œuvre de Guglielmo Marconi, un jeune Italien de vingt et un ans sans formation universitaire, qui assembla les composants disponibles (oscillateur de Righi, cohéreur de Branly, antenne, terre) en un dispositif capable de transmettre un signal sur plusieurs kilomètres. Marconi n'inventa aucun composant~; il inventa le système, de la même manière qu'Edison n'avait pas inventé la lampe à incandescence mais le réseau électrique (§\,\ref{sec:electrification}). Le brevet britannique n\textsuperscript{o}~12,039, déposé en juin~1897, couvrait non pas un appareil mais un procédé de transmission de signaux électriques sans fil conducteur \parencite{jolly_marconi_1972}.


\subsection{Le monopole Marconi et le modèle maritime}\label{subsec:marconi}

\dimmark{D2/D7}{*}%
La stratégie commerciale de Marconi reposait sur trois piliers dont la combinaison produisit, pendant deux décennies, un monopole de fait sur la communication sans fil maritime. Le premier pilier était le portefeuille de brevets~: la \emph{Wireless Telegraph and Signal Company}, fondée en juillet~1897 avec un capital de 100\,000~livres sterling, investit environ 75\,000~livres dans l'acquisition et le dépôt de brevets avant~1900 \parencite{kooij_information_2022}. Le second était l'internationalisation~: plutôt que de vendre des licences à des opérateurs étrangers, Marconi créa des filiales nationales sous son contrôle (la \emph{Marconi International Marine Communication Company}, MIMCC, en~1900~; des filiales en Belgique, au Canada, en Argentine) et imposa à chacune l'obligation d'utiliser exclusivement le matériel Marconi \parencite{jolly_marconi_1972, kooij_information_2022}. Le troisième pilier était le leasing~: les équipements n'étaient pas vendus aux armateurs mais loués, accompagnés d'un opérateur Marconi embarqué, ce qui liait le client au fournisseur par un contrat de service continu\footnote{Le parallèle avec le modèle IBM (§\,\ref{sec:mechanography}) est frappant~: location de l'équipement, vente exclusive du consommable (les cartes chez IBM, le service d'opérateur chez Marconi), portefeuille de brevets comme barrière à l'entrée. Les deux monopoles procèdent du même triptyque, appliqué à des technologies sans rapport.}. Le refus d'interopérabilité~: Marconi interdit à ses opérateurs de communiquer avec les stations équipées de matériel concurrent, ce qui provoqua un conflit diplomatique à la conférence radiotélégraphique de Berlin en~1903 et conduisit à la Convention radiotélégraphique internationale de~1906, première tentative de régulation du spectre \parencite{jolly_marconi_1972, headrick_invisible_1991}.

\dimmark{0-TER}{+}%
La traversée transatlantique du 12~décembre~1901, lorsque Marconi reçut à St.~John's (Terre-Neuve) un signal émis depuis Poldhu (Cornouailles), marqua un seuil dans le rapport entre transmission et énergie. L'émetteur de Poldhu, un spark transmitter alimenté par un alternateur de 25~kilowatts, consommait une puissance comparable à celle d'une petite usine \parencite{jolly_marconi_1972, aitken_continuous_1985}. À seuil de réception donné, la portée variait comme la racine de la puissance d'émission~: doubler la portée exigeait de quadrupler la puissance, ce qui inscrivait la communication sans fil dans une logique de coût croissant radicalement différente de celle du télégraphe filaire.

\dimmark{D2/D8}{+}%
La concurrence vint de l'État. En Allemagne, l'administration impériale soutint la création de Telefunken (1903), société commune de Siemens et AEG, pour contrer le monopole Marconi. Telefunken bénéficia de contrats militaires et d'une politique de brevets agressive \parencite{kooij_information_2022, headrick_invisible_1991}. Le modèle institutionnel divergeait~: là où Marconi opérait comme un monopole privé international, Telefunken fonctionnait comme un instrument de politique industrielle. Les ventes cumulées de Marconi n'atteignaient que 6\,092~livres sterling en~1900 \parencite{kooij_information_2022}~; c'est la promesse stratégique, pas le marché, qui finançait l'expansion. La Première Guerre mondiale trancha le conflit~: les marines britannique et américaine prirent le contrôle des stations radio, et le monopole Marconi fut liquidé aux États-Unis en~1919 par la création de la Radio Corporation of America (RCA), orchestrée par la Navy comme \og \emph{chosen instrument}\fg{} de la politique de communication américaine \parencite{aitken_continuous_1985, kooij_information_2022}.

\dimmark{D5/D2}{*}%
\emergence{D2~--- Le monopole Marconi procède d'un quatrième mécanisme, distinct des rendements de réseau (télégraphe), du verrouillage du consommable (tabulatrice) et de l'intégration verticale (téléphone)~: le triptyque brevets + leasing + refus d'interopérabilité. Mais c'est l'État qui le défait (RCA)~: le wireless est la première technologie de communication où l'expropriation étatique en temps de guerre, et non la concurrence marchande, met fin au monopole privé. \questionch{2}{D2~--- L'État comme destructeur de monopole et créateur simultané d'un autre (RCA)~: ce mécanisme se retrouve-t-il dans les cas ultérieurs~?}}


\subsection{La puissance dans le signal}\label{subsec:puissance}

\dimmark{D8/0-TER}{*}%
L'exigence de puissance croissante posa un problème technique que les archives exploitées par \textcite{aitken_continuous_1985} permettent de documenter avec précision. En~1909, la Navy américaine publia un appel d'offres pour sa station d'Arlington (Virginie), premier élément d'un réseau de stations puissantes devant couvrir le Pacifique (Canal de Panama, Californie, Hawaï, Guam, Samoa, Philippines). Le Congrès alloua 1~million de dollars en~1911, puis augmenta l'enveloppe à 1,5~million. Le cahier des charges exigeait une portée de 3\,000~miles en toutes conditions.

La \emph{National Electric Signaling Company} (NESCO) de Fessenden remporta le contrat avec un spark transmitter synchrone rotatif de 100~kilowatts. La machine échoua à atteindre les spécifications. Ses quarante-huit électrodes en cuivre tournant à 1\,250~tours par minute dissipaient une part considérable de leur énergie en harmoniques inutiles \parencite{aitken_continuous_1985}. Un troisième émetteur, installé discrètement à Arlington par Cyril Elwell de la Federal Telegraph Company, était un arc de 30~kilowatts seulement, soit moins d'un tiers de la puissance du spark Fessenden. Il fit mieux.

\dimmark{0-TER}{*}%
Le fait que 100~kilowatts de spark produisent une portée inférieure à 30~kilowatts d'arc est le résultat 0-TER central de ce cas. L'explication est physique~: le spark dissipe son énergie en un spectre large d'harmoniques, tandis que l'arc concentre la sienne sur une seule fréquence, ce qui maximise le transfert d'énergie vers le récepteur\footnote{L'amiral Cone, chef du Bureau of Steam Engineering, objecta à Elwell qu'un ampère était un ampère et qu'on ne pouvait pas le changer (\og \emph{an ampere is an ampere and you cannot change it}\fg{}). Elwell rétorqua qu'un ampère d'onde continue valait deux ampères d'onde amortie. Cone avait tort~: c'est l'efficacité spectrale, pas la puissance brute, qui détermine la portée. L'épisode est rapporté par \textcite{aitken_continuous_1985}, ch.~3.}. La transition du spark à l'onde continue (arc, puis alternateur, puis tube à vide) est d'abord une innovation d'efficacité énergétique~: elle ne réduit pas la consommation totale mais augmente la portée par watt émis. La trajectoire de puissance qui suivit le confirme~: alternateur Alexanderson 2~kilowatts (1909), arc Federal 30~kilowatts (Arlington, 1913), arc 200~kilowatts et 500~kilowatts dans les années qui suivirent \parencite{aitken_continuous_1985}.

\dimmark{D5/0-GOV}{*}%
Le réseau de stations du Pacifique s'inscrivait dans une logique de souveraineté dont la Première Guerre mondiale allait révéler l'urgence. Le wireless répondait à une vulnérabilité stratégique que les câbles sous-marins avaient créée. \textcite{aitken_continuous_1985} documente cette vulnérabilité en détail~: l'Allemagne, coupée des câbles par la Grande-Bretagne dès~1914, ne disposait que de la station de Nauen pour communiquer avec ses diplomates et ses agents. Le télégramme Zimmerman de~1917, dans lequel l'Allemagne proposait une alliance au Mexique en échange de l'Arizona, du Nouveau-Mexique et du Texas, fut intercepté par le \emph{Naval Intelligence} britannique parce qu'il transitait par les câbles diplomatiques américains que Wilson avait mis à la disposition de l'Allemagne\footnote{L'affaire Zimmerman est rapportée par \textcite{aitken_continuous_1985}, ch.~5. Le président de Western Union, Newcomb Carlton, admit devant le Sénat en~1921 que tous les câbles entrant et sortant de Grande-Bretagne étaient remis quotidiennement au \emph{Naval Intelligence} britannique. Carlton fut assuré qu'ils n'étaient \og ni censurés ni même déchiffrés\fg{}~; les sacs de câbles étaient \og stockés dans un entrepôt de l'Amirauté et restitués le lendemain matin, non ouverts\fg{}.}. Au-delà de la diplomatie, le contrôle britannique des câbles s'étendait au commerce~: des négociants américains se plaignirent que leurs messages étaient retardés ou transmis à des concurrents britanniques par l'intermédiaire du \emph{Board of Trade} \parencite{aitken_continuous_1985}. Le wireless naissait comme technologie de souveraineté, pas de marché~: c'est la vulnérabilité des câbles qui finançait le développement de la radio, comme c'est la sécurité ferroviaire qui avait financé le déploiement du télégraphe en Angleterre (§\,\ref{sec:telegraph}). Le même pattern se reproduisait~: la demande étatique créait le marché avant que le marché n'existe.

\emergence{D5/0-GOV~--- La Navy finance le réseau Pacifique (1,5~million de dollars), exproprie le wireless pendant la guerre, crée RCA comme instrument national. La demande militaire crée le marché avant la demande commerciale~: même pattern que le Census pour Hollerith (§\,\ref{sec:mechanography}) et la Convention pour le télégraphe Chappe (§\,\ref{sec:telegraph}). \questionch{2}{D5~--- L'État comme premier client des technologies de l'information est-il un fait contingent (guerres, crises) ou structurel (les rendements croissants de réseau ne peuvent être financés que par un acteur qui ne calcule pas en termes de rentabilité immédiate)~?}}


\subsection{Le spectre et le broadcasting}\label{subsec:broadcasting}

Le troisième temps de l'histoire du wireless est celui où la technique change de nature~: de la communication point-à-point (comme le télégraphe) à la diffusion point-à-tous (le broadcasting). Deux faits nouveaux par rapport aux cas précédents émergent de cette transition.

\dimmark{D2}{*}%
Le premier fait est l'apparition d'un type de ressource sans précédent dans l'histoire des techniques de l'information. Le spectre électromagnétique est une ressource immatérielle et congestionnable~: deux émetteurs sur la même fréquence s'annulent mutuellement, comme deux trains sur la même voie se percutent. Mais, à la différence du cuivre ou du fil, le spectre n'est ni stockable, ni transportable, ni appropriable par la possession physique. C'est le premier bien commun de l'histoire de l'ICT~: ni bien privé rival (comme le cuivre du télégraphe), ni service (comme la commutation téléphonique), ni bien public pur (puisqu'il est congestionnable).

La manière dont ce bien commun fut régulé constitue l'un des épisodes les mieux documentés de l'économie de la régulation. \textcite{coase_federal_1959}, dans un article fondateur, avait soutenu que la régulation du spectre par la \emph{Federal Communications Commission} (FCC) était une erreur~: le problème de l'interférence aurait pu être résolu par la définition de droits de propriété privés sur les fréquences, comme on définit des droits de propriété sur la terre, et le marché aurait ensuite alloué ces droits à leurs usages les plus valorisés. \textcite{hazlett_rationality_1990}, dans une révision minutieuse de l'histoire législative, a montré que cette interprétation par l'erreur était elle-même erronée. Des droits de propriété privés sur le spectre existaient \emph{de facto} avant le Radio Act de~1927~: le Département du Commerce reconnaissait les transferts de licences, les stations de radio étaient achetées et vendues à des prix reflétant la valeur de la fréquence\footnote{La vente la plus spectaculaire eut lieu en septembre~1926~: la station WEAF de New York fut vendue par AT\&T à RCA pour 1~million de dollars, dont 800\,000~dollars pour le droit de fréquence seul \parencite{hazlett_rationality_1990}.}, et les stations partageant une même fréquence négociaient entre elles la répartition du temps d'antenne par des accords privés. En octobre~1926, le tribunal de l'Illinois statua, dans l'affaire \emph{Tribune Co.~v.~Oak Leaves Broadcasting Station}, que la priorité d'usage créait un droit de propriété sur la fréquence, en s'appuyant sur le droit commun des droits riverains et de la concurrence déloyale \parencite{hazlett_rationality_1990}.

\dimmark{D2/0-GOV}{*}%
Le Congrès réagit immédiatement~: en décembre~1926, les deux chambres votèrent une résolution déclarant qu'aucun droit privé sur le spectre ne serait reconnu, et exigèrent des broadcasters qu'ils signent des renonciations. Le Federal Radio Act de février~1927 créa la Federal Radio Commission et lui confia le pouvoir d'attribuer des licences selon le critère de l'\og intérêt public, de la commodité ou de la nécessité\fg{}. Ce que Hazlett montre, c'est que cette solution ne fut pas le produit d'une confusion entre la définition des droits et leur allocation~: elle fut un choix rationnel de distribution de rentes, dans lequel les grands broadcasters obtenaient la protection de leurs positions acquises (le gel de la bande à 550--1\,500~kHz, le refus d'élargir le spectre malgré la disponibilité technique), les régulateurs obtenaient un pouvoir discrétionnaire sur la vie et la mort des stations, et le Congrès obtenait un levier politique sur le contenu des émissions. La commission gela la bande AM à sa taille de~1924, utilisant moins de cinq pour cent de la capacité alors techniquement disponible \parencite{hazlett_rationality_1990}. Le \og chaos\fg{} de juillet~1926 à février~1927, invoqué par la Cour suprême comme justification de la régulation, avait été stratégiquement provoqué par le secrétaire au Commerce Herbert Hoover, qui refusa de faire appel du verdict \emph{Zenith} et de convoquer la conférence radio annuelle\footnote{Hoover avait déclaré qu'il \og accueillerait favorablement une affaire test\fg{} et utilisa la confusion qui suivit le verdict pour forcer le Congrès à légiférer \parencite{hazlett_rationality_1990}. Le \emph{New York Times} l'implora d'organiser un arrangement d'urgence avec l'industrie~; il s'y refusa.}. Le spectre ne fut pas régulé parce que le marché avait échoué~; il fut régulé parce qu'une solution de marché aurait privé les régulateurs et les incumbents de leurs rentes respectives.

\dimmark{D9}{*}%
Le second fait est la transformation du wireless en instrument politique de masse. La radiodiffusion, apparue aux États-Unis au début des années~1920 contre la volonté de l'industrie du point-à-point\footnote{\textcite{douglas_inventing_1987} a montré que le broadcasting naquit de la pratique des amateurs radio, pas d'une stratégie industrielle~: Marconi, les compagnies de câbles et le Post Office britannique étaient focalisés sur la communication point-à-point et ne voyaient pas la valeur de la diffusion à tous. La station KDKA de Pittsburgh (novembre~1920), propriété de Westinghouse, commença à émettre pour stimuler la vente de récepteurs, pas pour créer une industrie médiatique.}, créa le premier médium capable d'atteindre des millions de personnes simultanément. Les ventes de récepteurs radio aux États-Unis passèrent de 60~millions de dollars en~1922 à 850~millions en~1929 \parencite{kooij_information_2022}. RCA, sous la direction de David Sarnoff, vit son chiffre d'affaires bondir de 11~millions de dollars en~1922 à 60~millions en~1924 \parencite{kooij_information_2022}. La création de la National Broadcasting Company (NBC) en~1926, puis de la Columbia Broadcasting System (CBS) en~1928, organisa le broadcasting en réseaux nationaux financés par la publicité.

L'effet politique de ce nouveau médium a fait l'objet de deux études d'identification causale dont la convergence est remarquable. \textcite{adena_radio_2015}, exploitant la variation géographique et temporelle de la couverture radio en Allemagne entre~1929 et~1936 (signal calculé par le modèle de propagation ITM), ont montré que l'effet de la radio sur le soutien au parti nazi était bidirectionnel~: pendant la période démocratique (1929--1932), lorsque le contenu radiophonique était pro-Weimar, la radio ralentissait la progression du NSDAP dans les zones couvertes~; après la prise de contrôle du médium par les nazis en janvier~1933, la radio accélérait le soutien, le recrutement de membres et les actes antisémites. L'effet dépendait des prédispositions des auditeurs~: la propagande nazie était efficace dans les zones où l'antisémitisme historique était élevé, et contre-productive là où il était faible. Le même médium servait la démocratie ou la dictature selon qui le contrôlait~: le broadcasting est, dans la terminologie de la thèse, un pharmakon politique.

\dimmark{D9}{*}%
\textcite{gentzkow_television_2006}, par une stratégie d'identification similaire (variation géographique de l'introduction de la télévision aux États-Unis entre~1940 et~1960, exploitant la Seconde Guerre mondiale et le gel des licences FCC de~1948--1952 comme chocs exogènes), a montré que la télévision réduisait la participation électorale d'environ 2~points de pourcentage par décennie, expliquant la moitié du déclin de la participation aux élections non présidentielles depuis les années~1950. Le mécanisme identifié est le \emph{crowding out} de l'information politique~: la télévision, en tant que bien de divertissement supérieur, détournait les consommateurs des journaux et de la radio, qui contenaient davantage d'information politique locale. Le résultat est renforcé par un contraste frappant avec la radio~: \textcite{stromberg_radio_2004} a montré que la radio avait \emph{augmenté} la participation électorale dans les années~1930, probablement parce qu'elle causait moins de substitution à la lecture des journaux et apportait davantage d'information politique. Le broadcasting n'a donc pas un effet uniforme sur la participation~: la radio enrichit, la télévision appauvrit, et la différence tient à la structure de substitution entre les médias.

\emergence{D9~--- Le broadcasting crée le premier médium de masse. Son effet politique est bidirectionnel~: la radio renforce la démocratie (Weimar) ou la dictature (nazis) selon qui la contrôle \parencite{adena_radio_2015}~; la radio augmente la participation (Stromberg) là où la télévision la diminue (Gentzkow). Le même registre technique (TRANSMIT one-to-many) produit des effets opposés selon le contenu et la structure de substitution avec les autres médias. \questionch{2}{D9~--- Ce pharmakon politique est-il propre au broadcasting, ou se retrouve-t-il avec Internet et les réseaux sociaux~? Le mécanisme de crowding out de Gentzkow (la qualité du divertissement détourne de l'information) opère-t-il à nouveau avec le streaming et le smartphone~?}}


\subsection*{Bilan~: la communication sans fil au prisme des dix dimensions}\label{subsec:bilan_wireless}

Le tableau~\ref{tab:wireless-dimensions} récapitule ce que la traversée du cas wireless a permis d'établir. Trois niveaux de certitude sont distingués~: \textbf{*}~(résultat empirique solide, réplicable ou documenté indépendamment), \textbf{+}~(argument explicite mais fondation empirique incomplète), \textbf{?}~(piste ouverte).

\begin{table}[p]
\caption{La communication sans fil au prisme des dix dimensions~: observations, niveaux de certitude et questions ouvertes}\label{tab:wireless-dimensions}
\centering\footnotesize
\renewcommand{\arraystretch}{1.15}
\begin{tabularx}{\textwidth}{@{}cXcp{4.5cm}@{}}
\toprule
\textbf{Dim.} & \textbf{Observations du cas §\,\ref{sec:wireless}} & & \textbf{Question pour la suite} \\
\midrule
D2 & Trois mécanismes de monopole successifs~: brevet + leasing + refus d'interopérabilité (Marconi 1897--1914), licence étatique (Radio Act 1927), monopole créé par l'État (RCA 1919). Le spectre est le premier bien commun congestionnable immatériel de l'histoire de l'ICT. La régulation du spectre n'est pas une erreur mais un choix rationnel de distribution de rentes \parencite{hazlett_rationality_1990}. & * & La régulation du spectre préfigure-t-elle la régulation des plateformes numériques~? \\[3pt]
D5 & La Navy finance le réseau Pacifique (1,5~M\$), exproprie le wireless pendant la WWI, crée RCA. Même pattern que le Census pour Hollerith et la Convention pour le télégraphe Chappe~: l'État crée le marché avant que le marché n'existe. & * & L'État comme premier client est-il un fait constitutif de chaque vague d'ICT~? \\[3pt]
D8 & Données financières Marconi (ventes cumulées £6\,092 en~1900). RCA~: \$11M (1922) $\rightarrow$ \$60M (1924). Wireless bubble. Le pattern investissement $\rightarrow$ faillite $\rightarrow$ recapitalisation n'apparaît pas directement~; c'est l'expropriation étatique qui restructure le marché. & + & Le financement par la promesse stratégique (pas la rentabilité) est-il propre aux technologies de souveraineté~? \\[3pt]
D9 & Le broadcasting crée le premier médium de masse. Effet bidirectionnel~: radio pro-Weimar ralentit les nazis, radio nazie accélère le soutien \parencite{adena_radio_2015}. TV réduit la participation électorale \parencite{gentzkow_television_2006}, radio l'augmente \parencite{stromberg_radio_2004}. & * & Le pharmakon politique se reproduit-il avec Internet et les réseaux sociaux~? \\[3pt]
D10 & Le broadcasting naît contre la volonté de l'industrie du point-à-point. Les amateurs radio forcent le changement d'usage \parencite{douglas_inventing_1987}. Marconi, les câbles et le Post Office étaient focalisés sur le point-to-point. & + & L'user innovation comme mécanisme de changement de paradigme \\[3pt]
\midrule
0-TER & L'énergie passe du substrat (cuivre) au signal (onde). À efficacité spectrale fixée, la portée croît avec la racine carrée de la puissance d'émission~; doubler la portée exige de quadrupler la puissance. Mais cette loi est elle-même relâchée par l'efficacité de rayonnement~: 30~kW d'arc, concentrés sur une fréquence étroite, portent plus loin que 100~kW de spark dispersés en harmoniques \parencite{aitken_continuous_1985}. L'efficacité énergétique est la variable cachée de la trajectoire technique. C'est la première technologie ICT où le coût d'exploitation énergétique est continu et croissant avec l'usage. & * & Le déplacement de la matérialité du substrat au signal annonce-t-il le problème énergétique contemporain~? \\[3pt]
0-GOV & Les câbles sous-marins sont une vulnérabilité stratégique~: blocus câblé britannique, télégramme Zimmerman, interception commerciale \parencite{aitken_continuous_1985}. Le wireless naît comme technologie de souveraineté, pas de marché. La Navy orchestre la création de RCA. & * & La souveraineté numérique contemporaine prolonge-t-elle le même pattern~? \\[3pt]
\bottomrule
\end{tabularx}
\renewcommand{\arraystretch}{1.0}
\end{table}

La traversée du cas wireless suggère trois hypothèses que le chapitre~2 devra instruire. La première est que le spectre électromagnétique inaugure un type de ressource nouveau~: un bien commun congestionnable immatériel qui appelle une forme de régulation sans précédent dans l'histoire de l'ICT, et dont la régulation effective ne vise pas l'efficacité allocative mais la distribution de rentes. La seconde est que la transition du point-à-point au broadcasting est un changement de paradigme non anticipé par l'industrie en place, confirmant le pattern observé dans le cas du téléphone (§\,\ref{sec:telephone})~: les incumbents ne voient pas la valeur du nouveau registre parce que leur système métrique ne sait pas la mesurer. Marconi, comme Gray, ne manquait pas d'intelligence~; il manquait d'unités. La troisième hypothèse est propre à la dimension énergie-matière~: le wireless est la première technologie de l'information où le coût énergétique réside dans le signal lui-même, et non dans le substrat matériel du réseau. Le fil de cuivre coûtait cher à poser mais ne consommait rien une fois posé~; l'onde électromagnétique ne coûte rien à installer mais consomme en continu. Ce déplacement, du capital fixe matériel vers le flux énergétique d'exploitation, annonce le régime contemporain du datacenter~: une infrastructure dont le coût dominant n'est plus la construction mais l'alimentation.

\emergence{0-TER~--- Le wireless opère un basculement dans la structure de coûts de la transmission~: du capital fixe matériel (cuivre, gutta-percha) vers le flux énergétique d'exploitation (puissance d'émission). Ce basculement est masqué, dans les années~1900--1930, par les ordres de grandeur encore modestes~; mais la logique est en place, et elle se reproduira, amplifiée, à chaque étape de la numérisation. \questionch{3}{0-TER~--- Si le coût dominant de la transmission passe du substrat au signal, la trajectoire d'amélioration technique (spark $\rightarrow$ arc $\rightarrow$ tube $\rightarrow$ transistor) ne fait que repousser le moment où la consommation agrégée dépasse les gains d'efficacité unitaire. C'est le paradoxe de Jevons appliqué à la transmission.}}




% =============================================================================
% §1.7 L'électrification et le substrat commun (~1880–1930)
% =============================================================================
\clearpage
\section{L'électrification et le substrat commun (~1880--1930)}\label{sec:electrification}
% =============================================================================

% §1.7 — Ouverture avec analyse des mécanismes structurels
% DRAFT v3 — 22 avril 2026
% Cinq mécanismes : (1) information comme condition de viabilité du système de puissance,
% (2) load factor comme problème informationnel, (3) infrastructure partagée et path dependence,
% (4) monopole naturel par rendements d'échelle, (5) reverse salient.
% La bifurcation power/information n'est PAS posée ici — elle émerge dans le bilan.
% David/Devine → §1.8.
% Hughes cité via secondaires + VdK ; \wip{} là où le primaire est indispensable.

Les cas qui précèdent partageaient un trait commun que le récit n'avait pas eu besoin de nommer~: aucun d'entre eux ne dépendait du réseau électrique pour fonctionner. Le télégraphe de Morse tirait son énergie de piles voltaïques~; les calculatrices de Thomas de Colmar et de Burroughs étaient actionnées à la main~; les tabulatrices Hollerith consommaient un courant dont le coût représentait 1,5\,\% du budget d'exploitation. L'électricité traversait ces cas comme un détail technique, pas comme une condition de fonctionnement. La période qui s'ouvre ici change la donne~: la radiodiffusion, l'usage massif des tabulatrices dans les grandes organisations, et bientôt le calcul électronique supposent tous l'existence d'un réseau électrique fiable, centralisé et continu. Le réseau électrique n'est pas une technique de l'information au sens de ce chapitre~; mais la manière dont il s'est constitué, entre 1880 et 1930, révèle des mécanismes qui intéressent directement l'analyse, parce qu'ils mettent en jeu, pour la première fois, l'articulation entre des fonctions informationnelles et un système de puissance au sein d'un même artefact technique.

\dimmark{D3/D7}{+}%
Le 4~septembre 1882, la Pearl Street Central Generation Station ouvrit dans le sud de Manhattan. Ce que Thomas Edison avait conçu n'était pas une lampe à incandescence~; d'autres l'avaient précédé, de Swan en Angleterre à Sawyer aux États-Unis. C'était un système complet, calqué sur l'architecture de l'industrie du gaz~: une station centrale produisant l'électricité, un réseau de câbles souterrains la distribuant, un compteur kilowattheure la mesurant, un tarif à l'usage la facturant \parencite{kooij_electric_light_2015}. \textcite{hughes_networks_1983} a montré que cette conception systémique était la contribution décisive d'Edison, celle qui le distinguait de ses prédécesseurs~: non pas l'invention d'un composant, mais l'architecture d'un ensemble dont chaque élément était dimensionné par rapport aux autres\footnote{Hughes qualifie Edison de \emph{system builder} et organise l'évolution du système électrique en quatre phases~: invention et développement, transfert technologique, croissance du système (avec les \emph{reverse salients}), et \emph{technological momentum}. Le modèle est appliqué systématiquement à trois trajectoires nationales~: New York, Berlin et Londres \parencite{hughes_networks_1983}.}. Six générateurs de vingt-sept tonnes chacun, alimentés par des chaudières consommant cinq tonnes de charbon et 43\,500~litres d'eau par jour, fournissaient 600~kilowatts de courant continu à 110~volts. Quatre cents lampes brûlèrent le premier jour chez quatre-vingt-cinq clients, dont Drexel, Morgan \& Company et le \emph{New York Times}~; un an plus tard, 8\,573~lampes desservaient 513~clients \parencite{kooij_electric_light_2015}.

\dimmark{D6/D7}{+}%
Or le système de Pearl Street ne pouvait pas fonctionner sans un dispositif informationnel qui lui était constitutif. Avant le compteur kilowattheure (brevet US n\textsuperscript{o}~242,901, 14~juin 1881), Edison facturait au forfait par lampe installée~; le client payait le même montant qu'il consommât beaucoup ou peu, ce qui rendait impossible toute incitation à la sobriété et toute adaptation de la capacité à la demande réelle. Le compteur transformait un flux physique continu en un chiffre discret, lisible, facturable~: c'est un dispositif de traitement de l'information, logé au c\oe{}ur d'un système de puissance, sans lequel le modèle économique de la centrale n'existe pas. La chaîne causale est la suivante~: le compteur rend possible le tarif à l'usage~; le tarif à l'usage rend possible la diversification de la base de clients~; la diversification rend possible l'optimisation du taux d'utilisation de la capacité installée. C'est la première fois dans le parcours du chapitre que l'information n'agit ni sur un marché extérieur (comme le télégraphe créant le marché du coton chez \textcite{steinwender_real_2018}) ni sur la coordination interne d'une organisation (comme les six formes de McCallum sur l'Erie Railroad)~: elle rend un système de puissance \emph{viable}.

\dimmark{D2}{+}%
Le système d'Edison souffrait cependant d'une limitation physique qui allait déterminer la suite. Le courant continu ne pouvait être transporté qu'à basse tension, ce qui imposait la proximité entre la centrale et les clients~: Pearl Street ne desservait qu'un rayon d'environ un mile\footnote{La perte en ligne croît avec le carré du courant et la longueur du conducteur. En courant continu à 110~volts, doubler la distance de desserte exigeait de quadrupler la section des câbles de cuivre, rendant le coût prohibitif au-delà d'un mile environ.}. \textcite{hughes_networks_1983} nomme ce type de contrainte un \emph{reverse salient}~: le composant du système dont l'insuffisance empêche l'ensemble de progresser, par analogie avec le segment d'un front militaire qui n'avance pas au même rythme que le reste. La distance de transmission était le \emph{reverse salient} du système Edison~; sa résolution vint de l'extérieur du système, avec l'adoption du courant alternatif portée par Westinghouse et Tesla dans les années~1890. Le transformateur permettait de transporter l'électricité à haute tension sur de longues distances et de la distribuer à basse tension chez le client~: le \emph{reverse salient} était levé, mais au prix d'un remplacement complet de l'architecture du réseau\footnote{La \emph{War of Currents} entre Edison (DC) et Westinghouse (AC), qui opposa les deux systèmes pendant une décennie, est un cas où le choix technique n'avait pas de solution unique~: \textcite{hughes_networks_1983} montre que le système DC aurait pu survivre avec des améliorations incrémentales (le système à trois fils, les batteries de stockage), et que la victoire de l'AC ne fut ni inévitable ni purement technique. \textcite{granovetter_making_1998} soutiennent que la structure de l'industrie électrique fut construite socialement, pas déterminée techniquement.}. Le concept de \emph{reverse salient} est utile au-delà du cas Edison~: il sera repris en fin de chapitre pour formuler l'hypothèse que l'énergie est, au vingt et unième siècle, le \emph{reverse salient} du système informationnel lui-même.

\dimmark{D2/D10}{+}%
Samuel Insull, ancien secrétaire d'Edison devenu directeur de la \emph{Chicago Edison Company}, comprit que la clé du réseau n'était pas la lampe mais le \emph{load factor}~: le taux d'utilisation de la capacité installée. Une centrale qui ne tourne qu'aux heures de pointe produit de l'électricité coûteuse~; une centrale qui tourne en continu produit de l'électricité bon marché. Le problème était de nature informationnelle~: il fallait savoir qui consommait quand, comment la demande se répartissait dans la journée, quelles catégories de clients pouvaient lisser la charge. En diversifiant la base (résidentiel le soir, industriel le jour, tramways en continu), Insull fit baisser le coût unitaire~; en négociant des tarifs dégressifs, il augmenta la consommation et justifia des centrales de plus en plus grandes\footnote{\wip{Développer avec Hughes 1983, ch.~8 (\emph{Chicago~: The Dominance of Technology})~: le modèle Insull, les rendements d'échelle, le monopole naturel régulé par les \emph{public utility commissions} créées à partir de 1907. La comparaison avec Berlin (ch.~7, Rathenau, coordination technologie-politique) et Londres (ch.~9, primauté du politique) produirait trois trajectoires d'électrification, comme les sept trajectoires du télégraphe en §1.2, convergeant vers un résultat commun~: le système intégré à grande échelle.}}. Le mécanisme de monopole qui en résultait était distinct des deux précédents~: ni les rendements de réseau du télégraphe, ni le verrouillage par le consommable d'IBM, mais les rendements d'échelle dans la génération, combinés à une régulation publique qui garantissait au monopoleur son territoire en échange d'un contrôle des tarifs. \textcite{hughes_networks_1983} nomme \emph{technological momentum} la résistance au changement d'un système dans lequel des investissements massifs, des institutions éducatives, des sociétés professionnelles et des cadres réglementaires se sont co-construits autour d'une architecture technique~: une fois le réseau en place, ses composants co-évoluent et résistent au remplacement. Le réseau électrique acquit ce momentum avant 1914~; les techniques de l'information, elles, allaient désormais évoluer dans le milieu qu'il avait créé.

\dimmark{D10/0-TER}{+}%
Un dernier trait mérite d'être relevé, parce qu'il n'est visible qu'avec le recul que donne la traversée des cas précédents. Edison n'était pas électricien de formation~: il était opérateur de télégraphe. Sa carrière commença par le multiplexage télégraphique~; le quadruplex (1874), qui permit de transmettre quatre messages simultanément sur un seul fil, fit sa réputation et sa fortune initiale \parencite{kooij_electric_light_2015}. C'est la maîtrise du circuit électrique, acquise dans le domaine de la communication, qu'il transposa dans le domaine de la puissance. Les poteaux du télégraphe portèrent les premiers fils de distribution électrique. Western Electric, qui fabriquait du matériel télégraphique, fabriqua ensuite du matériel téléphonique, puis du matériel électrique. Les mêmes chaînes d'approvisionnement en cuivre servaient les deux réseaux. Cette parenté n'est pas une coïncidence~: elle signale que les deux trajectoires de l'électricité~--- la puissance et l'information~--- naissent dans les mêmes ateliers, avec les mêmes matériaux, par les mêmes ingénieurs, avant de diverger institutionnellement en industries séparées\footnote{La chronologie renforce l'observation~: le télégraphe électrique (1837) précède la dynamo auto-excitatrice (1866) de trente ans. La trajectoire de l'information précède celle de la puissance. Cette antériorité et cette parenté seront reprises dans le bilan de cette section.}. Au tournant du siècle, la divergence était consommée~: l'industrie électrique et l'industrie des télécommunications avaient des acteurs, des régulateurs et des revues savantes distincts. Mais le substrat restait commun, et la dépendance des techniques de l'information à ce substrat allait s'intensifier dans chacun des cas qui suivent.

\dimmark{0-TER}{+}%
Cette intensification restera pourtant invisible pendant quatre décennies supplémentaires, grâce à un mécanisme que \textcite{dennard_design_1974} ont formalisé sans en mesurer les conséquences énergétiques. Le \emph{MOSFET scaling} de Dennard établit que si les dimensions d'un transistor diminuent d'un facteur $k$, la tension et le courant diminuent du même facteur, de sorte que la densité de puissance~--- watts par centimètre carré~--- reste constante. En pratique, le voltage des transistors passe de 5~volts au début des années~1990 à 1,2~volt vers~2005 \parencite{roussilhe_phase_2025}. Tant que cette loi tient, la puissance de calcul peut croître exponentiellement (selon la trajectoire de \textcite{moore_cramming_1965}) sans que la consommation électrique des composants n'augmente proportionnellement. C'est le bouclier qui permet au discours de dématérialisation de se maintenir~: l'ordinateur est perçu comme de plus en plus puissant et de moins en moins énergétivore par opération, ce qui est vrai~--- mais la consommation agrégée croît parce que le volume de calcul croît plus vite que l'efficacité. Vers~2005, Dennard atteint sa limite physique aux alentours de 3~GHz~; l'industrie contourne l'obstacle par la multiplication des c\oe{}urs (\emph{multi-core}), mais la digue électrique a cédé. La puissance maximale (TDP) des processeurs, stabilisée entre 75 et 150~watts pendant la décennie précédente, recommence à croître~; avec les GPU dédiés à l'IA générative, elle atteint 700~watts en~2024 et les feuilles de route projettent 1\,400~watts d'ici~2026 \parencite{roussilhe_phase_2025, sun_chip_2024}. Le \emph{scaling} de Dennard aura ainsi joué le rôle d'un facteur d'occultation au sens de \textcite{ensmenger_environmental_2018}~: pendant quarante ans, il a maintenu le substrat électrique hors du champ de vision des économistes de l'ICT, non par dissimulation délibérée mais parce que le mécanisme physique rendait la question énergétique littéralement sans objet à l'échelle du composant. La section~\ref{sec:computing} reprend là où ce bouclier cède.

\subsection{Berlin, Londres et la divergence des trajectoires}\label{subsec:trajectoires_europe}

\dimmark{D2/0-GOV}{+}%
Le système Edison que la section précédente a décrit ne s'est pas déployé à l'identique dans les trois capitales où il fut transféré. \textcite{hughes_networks_1983} a documenté les trajectoires divergentes de New York, Berlin et Londres entre 1882 et 1914~; leur comparaison révèle que la forme institutionnelle, et non la supériorité technique, détermine le rythme d'électrification. À Berlin, Emil Rathenau, qui avait acquis les droits Edison pour l'Allemagne après avoir visité l'exposition de Paris en 1881, fonde la \emph{Deutsche Edison Gesellschaft} (1883), devenue \emph{Allgemeine Elektricitäts-Gesellschaft} (AEG) en 1887. La concession municipale de 1884 lui accorde le monopole de la distribution dans un périmètre défini, en échange du partage des revenus avec la ville~; AEG contrôle simultanément la fabrication des équipements et l'exploitation du réseau, une intégration verticale que ni Edison à New York ni aucun concurrent à Londres n'atteindra\footnote{Siemens \& Halske, le concurrent berlinois d'AEG, adopte une stratégie complémentaire plutôt que frontale~: les deux firmes se partagent le marché allemand et coopèrent dans les projets d'infrastructure à grande échelle. \textcite{hughes_networks_1983}, ch.~7, documente cette coordination comme un trait distinctif du modèle allemand.}. En 1914, la Berliner Elektrizitäts-Werke (BEW), filiale d'AEG, dessert un rayon de trente kilomètres avec six centrales et une capacité installée de 155\,000~kW~; le taux de charge (\emph{load factor}) y est supérieur à celui de Chicago et bien au-dessus de celui de Londres.

\dimmark{0-GOV}{+}%
À Londres, la trajectoire fut inverse. L'\emph{Electric Lighting Act} de 1882, promu par Joseph Chamberlain au Board of Trade, imposa aux compagnies privées une clause de rachat municipal à vingt et un ans au prix des actifs isolés (\emph{scrap iron clause})~: la municipalité pouvait racheter câbles, dynamos et bâtiments à leur valeur séparée, sans prime pour la valeur du système en fonctionnement\footnote{L'expression \emph{scrap iron} provient des débats parlementaires de 1882~: le secrétaire du Board of Trade expliqua que le rachat se ferait à la valeur des «~briques comme briques et du mortier comme mortier~» (\emph{bricks as bricks and mortar as mortar}). Voir \textcite{hughes_networks_1983}, p.~57.}. Cette clause gela l'investissement privé~: pourquoi immobiliser du capital dans une infrastructure dont la valeur serait amputée au moment du rachat~? En 1883, soixante-neuf autorisations provisoires furent accordées~; en 1884, quatre~; en 1885, aucune. La station Edison de Holborn Viaduct, ouverte en 1882, fut abandonnée par Edison \& Swan en 1886. Quand le Parlement amenda la loi en 1888 pour porter la durée de concession à quarante-deux ans, il était trop tard~: les villes avaient acquis le droit de gérer elles-mêmes l'éclairage, et la fragmentation s'enracina. En 1918, Londres comptait encore soixante-cinq compagnies d'électricité, dix fréquences différentes et vingt-quatre tensions~; l'unification ne viendrait qu'avec le \emph{Central Electricity Board} de 1926.

\dimmark{D2/D10}{*}%
La comparaison Hughes établit un résultat qui dépasse le cas de l'électricité. La trajectoire technique du réseau n'explique pas sa forme institutionnelle~: Berlin et Londres partaient du même artefact (le système Edison), disposaient d'ingénieurs compétents et de capitaux disponibles~; c'est la configuration politique qui discrimina les trajectoires. À Berlin, la coalition Rathenau--banques d'investissement--municipalité avait le pouvoir d'imposer une concession stable~; à Londres, la coalition aristocratie--science--industrie se heurta à des intérêts contradictoires au Parlement. Hughes résume~: \og~The most penetrating explanation for the failure in London and the success in Berlin is neither technological nor economic; it is political.\fg{}\footnote{\textcite{hughes_networks_1983}, p.~78.} Ce résultat est exactement parallèle à celui que la section~\ref{sec:telegraph} avait établi pour le télégraphe~: sept trajectoires nationales convergeant vers le monopole par sept mécanismes différents, dont la variété tient aux configurations politiques locales et non aux propriétés techniques du réseau. La dimension \textbf{0-GOV} n'est pas un résidu~; elle est, dans les deux cas, le facteur déterminant de la forme institutionnelle.

\emergence{La divergence Berlin/Londres confirme que les réseaux électriques, comme les réseaux télégraphiques, convergent vers des formes institutionnelles comparables (monopole régulé, intégration verticale, contrôle étatique) par des trajectoires politiquement déterminées. La supériorité technique n'explique ni le succès de Berlin ni l'échec de Londres~: c'est la capacité des coalitions d'acteurs à stabiliser un cadre institutionnel favorable à l'investissement long qui discrimine les trajectoires. Le mécanisme vaut pour 0-TER~: là où l'électrification est retardée par la fragmentation, les techniques de l'information restent plus longtemps sur substrat mécanique ou manuel, ce qui retarde aussi leur couplage au réseau de puissance.}


% --- Les subsections continuent ici : phonographe, IBM tabulatrices, Shannon, Enigma ---


\subsection{Le phonographe et les supports de mémoire}\label{subsec:phonographe}

\dimmark{0-TER}{*}%
Edison invente le phonographe en 1877, un an après le téléphone de Bell. Il est frappant que le même homme ait conçu, en l'espace d'une décennie, le compteur kilowattheure (information \emph{pour} la puissance), le quadruplex (information \emph{pour} la transmission) et le phonographe (information \emph{pour} la mémoire)~: trois fonctions en un seul atelier, non comme programme délibéré mais comme convergence d'un même génie du circuit électrique vers ses trois usages possibles. Le gramophone de Berliner (1887) et le Télégraphone de Poulsen (1898) prolongent la même logique~: des informations peuvent être fixées sur un support matériel et relues à volonté, sans consommer de courant entre les deux actes.

Ce trait est précisément ce qui distingue le régime pré-numérique du stockage de tous les régimes qui suivront. La cire, le disque de gomme-laque, le fil magnétique, le film photographique retiennent l'information sans énergie~: l'inscription est une déformation physique permanente du support, et la lecture en est la traduction mécanique ou optique. Le profil énergétique de STORE dans toute cette période est celui d'une ressource à constitution ponctuelle et à mobilisation ponctuelle~; entre les deux, l'énergie est nulle. Cette passivité n'est pas une limitation technique en attente de correction~: c'est une propriété structurelle des supports analogiques, qui tient à leur nature physique. Elle cessera d'être la norme avec la mémoire à accès dynamique (DRAM, Intel 1103, 1970), dont les cellules doivent être rafraîchies des centaines de fois par seconde sous peine de perdre leur contenu. La frontière entre le régime passif et le régime actif-permanent de STORE correspond à cette transition, qui sera traitée dans la section suivante.


\subsection{La radio et la télévision de masse}\label{subsec:broadcasting_mass}

\dimmark{D1/D5}{+}%
La radiodiffusion de masse ~--- NBC en 1926, BBC en 1927, puis les systèmes nationaux européens tout au long des années~1930~--- accomplit quelque chose que le téléphone, le télégraphe et même la presse n'avaient pas réalisé sous cette forme~: la transmission simultanée d'un même signal à des millions de récepteurs sans que le nombre de récepteurs affecte ni la qualité du signal ni son coût de production. C'est la première réalisation technologique d'un bien strictement non rival~: la énième auditrice d'un concert radiodiffusé ne diminue pas ce que reçoit la première. Les ventes de récepteurs radio aux États-Unis illustrent l'ampleur du phénomène~: soixante millions de dollars en 1922, huit cent cinquante millions en 1929, en l'espace de sept ans seulement.

\dimmark{D8}{+}%
Cette croissance pose immédiatement la question du modèle économique. Un signal radiodiffusé ne peut pas être rendu excluable par un mécanisme technique simple~: quiconque possède un récepteur peut capter l'émission sans payer. Deux solutions émergent~: la redevance, adoptée en Europe dans le cadre d'opérateurs publics (BBC, ORTF, ARD), et la publicité, choisie aux États-Unis par les réseaux commerciaux. Ces deux solutions impliquent deux régimes de propriété et deux rapports à la demande~: la redevance découple la production du contenu du consentement à payer des auditeurs individuels~; la publicité découple le coût du signal de sa valeur pour l'auditeur, au profit d'un tiers payeur dont l'intérêt n'est pas le programme mais l'attention. Le rapport ICT-énergie dans les deux cas est identique~: le réseau électrique est nécessaire à la réception domestique de masse, mais ce coût est diffus et invisible~--- il est noyé dans la facture d'électricité du ménage, pas dans le coût d'exploitation du diffuseur.


\subsection{IBM et les tabulatrices dans les grandes organisations}\label{subsec:ibm_tab}

\dimmark{D2/D5}{*}%
La période 1924--1935 voit se refermer le cycle ouvert par Hollerith~: la «~Computing-Tabulating-Recording Company~», rebaptisée IBM («~International Business Machines~») par Thomas Watson en 1924, consolide une position de quasi-monopole que ni la concurrence de Remington Rand ni l'antitrust de 1936 ne parviendront à défaire. Deux événements accélèrent ce mouvement. Le premier est technique~: en 1928, IBM introduit la carte à quatre-vingt colonnes à trous rectangulaires, brevetée, qui brise l'interopérabilité entre les systèmes Hollerith et Powers et verrouille les clients existants sur les matériels IBM. Le second est institutionnel~: en 1935, le Social Security Act impose à l'administration américaine de tenir les comptes individuels de cotisation de plus de vingt-six millions de travailleurs, un volume de données qui excède la capacité de tout système manuel~; la tabulatrice, dont IBM est le fournisseur dominant, devient dès lors une infrastructure d'État, financée par les impôts et déployée à une échelle que le marché privé n'aurait pas atteinte aussi rapidement. Cette articulation entre crise administrative de la grande organisation et déploiement de la mécanographie comme infrastructure de coordination s'inscrit dans ce que \textcite{beniger_control_1986} a appelé la \emph{control revolution}~: la révolution informationnelle n'y serait pas une rupture post-industrielle mais une réponse continue à la crise de contrôle déclenchée par la révolution industrielle elle-même, dont les chemins de fer du XIX\textsuperscript{e}~siècle, l'organisation chandlérienne et la mécanographie hollerithienne constituent les phases successives. Le cadre est retenu ici comme lecture disponible~; la décomposition par fonctions TRANSMIT/STORE/COMPUTE et le suivi du couplage énergétique~0-TER en restent les apports propres au présent chapitre.

\dimmark{D4}{+}%
Le cas du \emph{Nautical Almanac Office} que étudie \textcite{daston_calculations_2018} éclaire ce que la mécanisation fait réellement au travail de calcul. Entre 1928 et 1937, le bureau mécanise progressivement le calcul des éphémérides~: le budget annuel passe de cinq cents livres sterling à six cent sept, soit une hausse de vingt et un pour cent, tandis que le nombre de collaborateurs augmente de même. La tabulatrice ne remplace pas les calculateurs humains~: elle déplace le goulot d'étranglement. Ce qui manquait avant la mécanisation était le calcul arithmétique brut~; ce qui manque après est ce que les superviseurs de l'époque appellent \og~analytical intelligence~\fg{}~: la capacité à formuler les problèmes, à vérifier les résultats, à concevoir les séquences d'opérations. Vingt à trente pour cent du temps de travail est consacré à la supervision et au contrôle~; la mécanisation déplace la pénurie vers ces tâches. Le mécanisme est déjà celui qu'\textcite{autor_skills_2003} formalisera soixante ans plus tard~: la mécanisation des tâches routinières codifiables augmente la demande de tâches cognitives non routinières, elle ne la supprime pas.

\dimmark{0-TER}{*}%
Le profil énergétique de ce régime reste négligeable. Les tabulatrices IBM des années~1930 consomment entre cinq cents watts et deux kilowatts par unité~; l'électricité représente, comme noté plus haut, moins de deux pour cent des budgets d'exploitation. Ce qui est crucial pour l'argument de la thèse n'est pas la consommation absolue, mais la structure du coût~: ce coût est discret (lié à la durée d'utilisation), pas continu. Quand les tabulatrices sont éteintes, elles ne consomment pas~; quand les cartes perforées sont dans leurs boîtes, elles n'exigent rien. STORE reste passif, COMPUTE est actif seulement pendant le calcul. La dissémination de l'ICT dans les organisations des années~1930 ne crée pas de consommation continue~: elle crée seulement de l'usage.


\subsection{Boole, Shannon et la logique devenue circuit}\label{subsec:shannon}

\dimmark{D10}{+}%
George Boole publie en 1847 ses \emph{Mathematical Analysis of Logic}, dans lequel il montre que les opérations logiques (et, ou, non, implication) peuvent être représentées par des équations algébriques dont les seules valeurs sont 0 et~1. Le geste est purement mathématique~: Boole ne pense pas au circuit électrique. Il faut attendre quatre-vingt-onze ans pour que Claude Shannon, alors étudiant au MIT, réalise dans sa thèse de master (1938) que les circuits de commutation électrique~--- interrupteurs ouverts ou fermés, courant passant ou non~--- réalisent physiquement les opérations de l'algèbre de Boole. Le~1 logique est le circuit fermé~; le~0 logique est le circuit ouvert. Cette identification transforme l'algèbre de Boole d'un outil mathématique en un programme d'ingénierie~: toute opération logique peut être câblée, et toute suite d'opérations logiques peut être automatisée par des circuits électriques convenablement agencés. Vannevar Bush, directeur du MIT, avait construit dès 1930 un analyseur différentiel électromécanique~; c'est dans ce contexte que Shannon réalise son travail, sous sa direction.

\dimmark{D10}{+}%
Dix ans plus tard, Shannon publie \emph{A Mathematical Theory of Communication} (1948), qui accomplit un second geste~: définir l'information indépendamment de son support physique et de son contenu sémantique. L'information se mesure en bits~; l'entropie informationnelle est formellement analogue à l'entropie thermodynamique de Boltzmann, dont Shannon reprend le terme. Ce second geste fonde la séparation analytique entre l'information et son support~--- séparation qui rendra possible la conception de réseaux universels, indépendants du type de contenu qu'ils transportent, et qui contribuera en même temps à occulter la matérialité et le coût énergétique du signal. La convergence numérique des années~1980--2000 hérite des deux gestes de Shannon~: l'universalité du bit (tout contenu est encodable) et la neutralité du substrat (tout substrat peut transporter n'importe quel bit). Ce que Shannon n'a pas résolu~--- et que \textcite{landauer_irreversibility_1961} tranchera en 1961~--- est précisément ce que cette neutralité cache~: transmettre, stocker et calculer des bits a un coût physique irréductible, que le formalisme de la théorie de l'information rend invisible par construction.

\dimmark{D10}{+}%
Trois ans avant ce papier fondateur, \textcite{wiener_cybernetics_1948} avait publié \emph{Cybernetics, or Control and Communication in the Animal and the Machine}, qui articule l'entropie shannonienne au contrôle par rétroaction et nomme \emph{cybernétique} le champ qui en résulte. Douze ans plus tard, dans \emph{Science}, Wiener met en garde contre la subordination de l'humain à la machine~: les automates qu'il avait contribué à fonder pouvaient, écrit-il, imposer leur logique aux humains qui les emploient si l'on n'y prenait pas garde~\parencite{wiener_consequences_1960}. Le geste qui retournera cette mise en garde en programme d'innovation intervient dès les premières années de la décennie suivante~: \textcite{minsky_steps_1961}, dans \emph{Proceedings of the IRE}, formule la promesse que les machines pourraient, dans le temps d'une vie humaine, surpasser l'humain en intelligence générale. La trajectoire qui mène de l'inquiétude wienerienne à la promesse minskyenne, puis aux investissements massifs des décennies suivantes dans la recherche en intelligence artificielle, illustre une dynamique récurrente dans le déploiement des techniques de l'information~: les imaginaires des fondateurs sont retournés par les acteurs qui les suivent en arguments d'investissement, indépendamment de l'intention initiale. \textcite{johnstone_deep_2026} ont récemment identifié cette dynamique comme structurelle dans la trajectoire de la digitalisation~; le présent chapitre la retrouvera à plusieurs reprises dans les sections qui suivent.


\subsection{Enigma, Bletchley Park et Colossus}\label{subsec:enigma}

\dimmark{D5/D8}{+}%
La Seconde Guerre mondiale accomplit dans le calcul ce que la Première avait accompli dans l'aviation~: elle mobilise des ressources sans limites ordinaires et produit en quelques années des sauts technologiques que le marché aurait mis des décennies à réaliser. En~1939, la Wehrmacht chiffre ses communications avec la machine Enigma, dont le nombre de configurations possibles excède $10^{23}$~; briser Enigma à la main est hors de portée humaine. À Bletchley Park, dix mille personnes~--- mathématiciens, linguistes, ingénieurs, dont Alan Turing~--- travaillent au déchiffrement. La \og~bombe~\fg{} électromécanique que Turing finalise en 1940, fondée sur la logique des contradictions plutôt que sur la force brute, réduit l'espace des configurations à explorer à un volume tractable\footnote{La \og~bombe~\fg{} de Bletchley Park dérive elle-même d'une machine polonaise, la \emph{bomba kryptologiczna}, conçue par Marian Rejewski, Jerzy Różyki et Henryk Zygalski à partir de~1938. Le transfert de technologie de Varsovie à Londres, réalisé quelques semaines avant l'invasion allemande de la Pologne, est l'un des transferts technologiques les plus conséquents de la Seconde Guerre mondiale~; il reste peu connu en dehors de l'histoire des services de renseignement.}. Colossus (1943), conçu par Tommy Flowers de la British Post Office pour décrypter le code Lorenz utilisé par le haut commandement allemand, franchit une étape supplémentaire~: dix tubes à vide, cinq cent brins de papébande perforée lus à cinq mille caractères par seconde, et le premier ordinateur électronique programmable de grande taille qui fonctionne en conditions opérationnelles.

\dimmark{D8/0-TER}{+}%
Colossus fonctionne en continu~: les opérations de déchiffrement ne s'arrêtent pas la nuit, parce que les communications ennemies non plus. C'est la première occurrence dans ce parcours historique d'un système de calcul conçu pour une disponibilité permanente, non par choix architectural mais par nécessité opérationnelle. Le pattern se reproduira~: SAGE tournera sans interruption parce qu'une attaque nucléaire n'attend pas les heures ouvrables (section~\ref{sec:computing})~; les serveurs web fonctionneront sans interruption parce qu'un client potentiel se connecte à n'importe quelle heure (section~\ref{sec:convergence}). La disponibilité permanente, qui produira le couplage énergétique continu décrit dans la section suivante, s'esquisse ici dans le contexte militaire où elle est la moins surprenante~--- et la moins comptabilisée.

\dimmark{D10}{+}%
Bletchley Park illustre enfin une dynamique de financement qui traversera toute l'histoire du calcul~: c'est l'État militaire qui crée le marché technologique que le secteur privé exploitera ensuite, une fois le risque socialisé et les preuves de concept établies. Le secret militaire maintient les résultats hors du domaine public pendant plusieurs décennies~; les ingénieurs de Bletchley Park ne peuvent pas breveter leurs découvertes. Les bénéficiaires du transfert sont les grandes firmes d'électronique~--- IBM, Ferranti, GEC~--- qui recruteront ces ingénieurs après la guerre et s'appuieront sur des savoirs dont le coût de développement a été supporté par le contribuable. Ce mécanisme, que \textcite{mazzucato_entrepreneurial_2013} a systématisé pour la période contemporaine, est structurel dans l'histoire du calcul bien avant qu'elle ne le nomme.


\subsection{Bilan~: la séquence d'électrification des trois fonctions}\label{subsec:bilan_electrification}

\dimmark{0-TER}{*}%
Le parcours de cette section a fait apparaître une séquence qui n'avait pas été nommée jusque-là. Les trois fonctions informationnelles~--- transmettre, stocker, calculer~--- ne s'électrifient pas au même moment~:

\begin{itemize}
\item \textbf{TRANSMIT} s'électrifie dès l'origine~: le télégraphe électrique (1837) \emph{est} électrique par définition. La pile voltaïque, puis le réseau à courant continu et alternatif, portent le signal dès le premier jour. L'électrification de la transmission n'est pas une transition~; c'est une condition initiale.

\item \textbf{COMPUTE} reste mécanique ou manuel pendant plus d'un siècle. La règle à calcul, la calculatrice de Thomas de Colmar, les tabulatrices Hollerith ne tirent pas leur énergie du réseau électrique~: elles sont actionnées à la main ou par des moteurs électriques dont la consommation est marginale dans le bilan d'exploitation. L'électronique de puissance n'entre dans le calcul qu'avec l'ENIAC (1946), soit cent neuf ans après Morse.

\item \textbf{STORE} bascule en dernier. Le papier, le carton perforé, la cire du phonographe, le film photographique, le disque de gomme-laque stockent l'information sans énergie~: l'inscription est une déformation physique permanente du support, et la lecture en est la traduction mécanique ou optique. La mémoire DRAM (Intel 1103, 1970) est la première technologie de stockage qui \emph{exige} un flux d'énergie continu~--- plusieurs centaines de rafraîchissements par seconde~--- pour maintenir son contenu. Avant elle, stocker était gratuit entre l'écriture et la lecture~; après elle, stocker consomme en permanence.
\end{itemize}

\dimmark{0-TER}{*}%
Cette séquence~--- T~(1837) $\rightarrow$ C~(1946) $\rightarrow$ S~(1970)~--- a une conséquence directe pour l'argument de la thèse. Tant que STORE reste passif, le coût énergétique de l'information est ponctuel~: il se produit à la transmission et au calcul, pas entre les deux. L'information peut dormir sur une étagère sans consommer de watt. Le régime énergétique est additif~: chaque usage ajoute une dépense discrète, mais il n'y a pas de fond continu. C'est ce régime que les sections §\ref{sec:telegraph} à §\ref{sec:wireless} ont traversé sans avoir besoin de le nommer, parce qu'il ne posait pas de problème~: l'énergie de l'information était invisible parce qu'elle était négligeable.

\dimmark{0-TER}{+}%
Le basculement vers un régime permanent s'opère en deux temps. Le premier temps est la DRAM (1970)~: la mémoire de travail des ordinateurs devient un flux continu de rafraîchissement. Le second temps est le serveur web (1991) et le datacenter (années~2000)~: l'exigence de disponibilité permanente (\emph{always-on}) transforme STORE d'une ressource ponctuelle en une infrastructure continue. Un fichier stocké sur un serveur web n'existe que si le serveur est allumé~; une page web accessible à tout moment exige une alimentation à tout moment. La section~\ref{sec:convergence} documentera ce basculement~; la section~\ref{sec:contemporain} en mesurera les conséquences. Ce que la présente section établit est le point de départ~: jusqu'en 1970, le substrat électrique des techniques de l'information est une condition d'existence, pas un coût d'exploitation continu. L'occultation de la dimension 0-TER dans les cas du XIX\textsuperscript{e}~siècle n'est pas une erreur d'observation~: elle reflète un régime physique où l'énergie de l'information était réellement négligeable.

\questionch{3}{Comment quantifier la consommation d'énergie de chaque fonction (TRANSMIT, STORE, COMPUTE) à travers les périodes identifiées par ce chapitre~? Le croisement des séries de capacités informationnelles de \textcite{hilbert_worlds_2011} avec les séries d'efficacité énergétique de \textcite{koomey_implications_2011} permettrait-il de reconstituer la consommation agrégée par fonction et par décennie~? Ce croisement n'a jamais été publié. Le chapitre~3 de cette thèse le tentera.}


\begin{table}[htbp]
\caption{Dimensions économiques de l'électrification (§\ref{sec:electrification}).}\label{tab:electrification-dimensions}
\centering\small
\renewcommand{\arraystretch}{1.15}
\begin{tabularx}{\textwidth}{@{}cXc@{}}
\toprule
\textbf{Dim.} & \textbf{Ce que §\ref{sec:electrification} établit} & \textbf{Niveau} \\
\midrule
D2 & Monopole naturel par rendements d'échelle (Insull)~; fragmentation = retard (Londres vs Berlin). & ★ \\
D3 & Capital de réseau irréversible (centrale + distribution)~; intégration verticale AEG. & + \\
D4 & Création d'emplois dans la maintenance et l'exploitation~; déplacement du travail vers la supervision (Nautical Almanac). & + \\
D5 & Demande dérivée (éclairage, puis moteur, puis information)~; diversification des usages pour optimiser le \emph{load factor}. & ★ \\
D6 & Le compteur kilowattheure transforme un flux continu en signal discret facturable~: l'information rend le système de puissance viable. & + \\
D7 & Standardisation des fréquences et tensions comme condition d'interconnexion~; fragmentation londonienne (24 tensions, 10 fréquences). & + \\
D8 & Cycle investissement lourd $\rightarrow$ consolidation (Chicago Edison $\rightarrow$ Commonwealth Edison)~; financement par les banques d'investissement (Deutsche Bank, Drexel Morgan). & ★ \\
D10 & Lock-in DC vs AC (guerre des courants)~; \emph{technological momentum} de Hughes~; \emph{scaling} de Dennard comme bouclier d'occultation. & ★ \\
\midrule
0-TER & Séquence T\,(1837) $\rightarrow$ C\,(1946) $\rightarrow$ S\,(1970)~; régime passif du stockage pré-DRAM~; STORE permanent = consommation continue. & ★ \\
0-GOV & Régulation des monopoles naturels (\emph{public utility commissions})~; divergence Berlin/Londres par la configuration politique, non technique. & ★ \\
\bottomrule
\end{tabularx}
\renewcommand{\arraystretch}{1.0}
\end{table}


% =============================================================================

% =============================================================================
% §1.8 Le calcul programmable et la révolution électronique (1945–1975)
% =============================================================================
\clearpage
% =============================================================================
% §1.8 Le calcul programmable et la révolution électronique (1945–1975)
% =============================================================================
% Résultats dimensionnels visés :
%   0-TER  — le computing comme problème énergétique visible
%   D3     — séparation hardware/software, naissance du capital intangible
%   D2     — consent decree comme anti-monopole (négatif du télégraphe)
%   D5/D8  — État comme créateur de marché
%   D10    — Moore's Law comme dispositif de coordination sectorielle
% Sources empiriques principales :
%   Nordhaus 2007, Koomey et al. 2011, Dennard et al. 1974,
%   Flamm 1987/1988, Moore 1965, Mody 2017, Merges & Nelson 1990
% =============================================================================

\section{Le calcul programmable et la révolution électronique (1945--1975)}\label{sec:computing}

Les deux cas précédents ont mis en lumière un contraste de fonction et de trajectoire. Le télégraphe transmet de l'information entre des lieux distants, à un coût énergétique si faible par rapport au transport physique qu'il en est invisible ; la mécanisation du calcul traite de l'information sur place, à un coût en travail humain que la machine redistribue sans le supprimer. Dans les deux cas, l'énergie consommée par le traitement de l'information reste un poste négligeable : 160 watts pour un bureau télégraphique contre 370 kilowatts pour la locomotive qui acheminait les mêmes dépêches ; neuf livres sterling d'électricité sur un budget annuel de six cent sept livres au Nautical Almanac Office mécanisé.

Avant que le calcul électronique n'existe matériellement, son concept est posé. \textcite{turing_computable_1936} formule, dans \emph{On Computable Numbers} paru aux \emph{Proceedings of the London Mathematical Society} en 1937, le concept de la machine universelle~--- une machine unique capable d'exécuter tout calcul reproductible pourvu qu'on lui fournisse la description du calcul à effectuer\footnote{Le mot «~ordinateur~» étant postérieur à cette formulation et propre à la langue française (1955, sur proposition de Jacques Perret pour IBM France), Turing parle de \emph{computing machine}, dont sa version universelle est la matrice conceptuelle. La projection rétrospective du mot ordinateur sur le texte de 1937 est un anachronisme léger, sans conséquence pour l'argument.}. Le résultat est purement mathématique~: aucun dispositif matériel n'en est tiré à ce moment. Neuf ans séparent cette formulation théorique de l'ENIAC, et quinze ans la séparent du premier ordinateur commercial~; l'écart entre la formulation conceptuelle et la disponibilité industrielle, caractéristique des techniques de l'information qui se reproduira à plusieurs reprises dans ce chapitre, ne tient pas à une difficulté technique surmontée par les ingénieurs, mais à l'absence préalable d'une demande qui aurait justifié la mobilisation des ressources nécessaires à la matérialisation du concept. La généalogie qui mène d'ENIAC à EDVAC puis à l'\emph{IAS Machine} peut donc se lire à deux niveaux indépendants~: une généalogie matérielle dont les jalons sont des artefacts physiques, et une généalogie conceptuelle dont les jalons sont des publications théoriques. Le présent chapitre suit les deux.

\dimmark{0-TER}{*}%
Le cas qui s'ouvre ici rompt avec cette invisibilité. En 1946, l'Electronic Numerical Integrator and Computer consomme à lui seul 150 kilowatts, soit la puissance d'une petite centrale, pour exécuter cinq mille additions par seconde. Pour la première fois dans l'histoire du traitement de l'information, la consommation d'énergie de la machine est du même ordre de grandeur que celle des activités qu'elle sert. Ce n'est pas un accident de jeunesse technologique ; c'est le début d'une course de quatre-vingts ans entre la puissance de calcul et la consommation électrique qui la rend possible, course dont les termes se modifient à chaque génération technologique sans que le couplage structurel se desserre.

Cette visibilité de l'énergie est cependant d'un type particulier. Elle est visible comme \emph{contrainte d'ingénierie}~: l'ENIAC nécessite une alimentation dédiée, un refroidissement actif par ventilateurs, une tension régulée pour éviter la corruption des calculs par les fluctuations du réseau\footnote{L'ENIAC opère sur dix-huit mille tubes à vide dont chacun peut tomber en panne --- en moyenne un tube toutes les deux heures en fonctionnement continu. La stabilité de l'alimentation électrique conditionne directement la fiabilité du calcul. Ce couplage entre qualité de l'énergie et intégrité de l'information est une propriété structurelle des premiers ordinateurs électroniques, non un détail opérationnel.}. Elle est en revanche \emph{absente des comptabilités nationales et des modèles économiques de la productivité}~: la \emph{Survey of Current Business} n'isole pas de poste «~énergie du traitement de l'information~»; les travaux de \textcite{flamm_creating_1988} sur le financement de l'informatique ne mentionnent pas le coût électrique d'exploitation. La visibilité technique de l'énergie coexiste avec son invisibilité économique et statistique. C'est cette dissymétrie~--- et non l'ignorance brute d'une contrainte~--- que la thèse cherche à documenter.

\dimmark{0-TER}{+}%
Cette dissymétrie a une justification physique que les travaux de Landauer ont établie dès 1961. \textcite{landauer_irreversibility_1961} montre que l'effacement d'un bit dans une mémoire physique dissipe nécessairement une quantité minimale d'énergie de l'ordre de $kT\ln 2$ joules, où $k$ est la constante de Boltzmann et $T$ la température du système~--- environ $3 \times 10^{-21}$~joules à température ambiante\footnote{La borne de Landauer a été confirmée expérimentalement par \textcite{berut_experimental_2012} et par \textcite{jun_high_2014}. \textcite{bennett_thermodynamics_1982} a montré que le calcul réversible (sans effacement d'information) échappe en théorie à cette borne~; en pratique, tout ordinateur réel efface de l'information et la borne s'applique. Le ratio entre la consommation énergétique effective des processeurs modernes et la borne de Landauer est de l'ordre de $10^{6}$~--- autant de marge d'amélioration technologique encore disponible, et autant de preuve que la consommation actuelle tient à l'architecture des composants, pas aux lois de la physique.}. Ce résultat est fondamental pour la thèse~: il établit que le coût énergétique du calcul n'est pas un accident technologique corrigible par l'ingénierie~--- il est irréductible au niveau du bit. La convergence des trois fonctions sur un substrat numérique après 1983 ne produit pas seulement un coût énergétique économiquement visible~; elle produit un couplage physiquement nécessaire entre traitement de l'information et dissipation d'énergie.

Ce cas couvre trente ans, de l'ENIAC (1946)\footnote{La paternité du \emph{stored-program concept} --- le principe selon lequel le programme est stocké dans la même mémoire que les données --- est disputée dans l'historiographie. Konrad Zuse présente le Z3 en 1941, premier ordinateur programmable fonctionnel à relais électromécaniques, dont les autorités nazies sous-estimeront l'importance. Eckert et Mauchly (ENIAC/EDVAC) revendiquent la priorité sur Von Neumann, dont le \emph{First Draft of a Report on the EDVAC} (1945) formalise l'architecture à programme enregistré de façon si complète que son nom reste attaché au concept. Ce qui importe analytiquement pour ce cas n'est pas la priorité d'invention mais la conséquence architecturale~: l'espace mémoire programme et l'espace mémoire données sont unifiés dans la même structure. C'est cette fusion~--- et non l'existence d'un ordinateur programmable~--- qui marque le pivot de la convergence STORE/COMPUTE.} au microprocesseur Intel 4004 (1971) et au papier fondateur du \emph{MOSFET scaling} de Dennard (1974). Il traverse au moins quatre artefacts majeurs : l'ordinateur programmable à tubes, le mainframe commercial à transistors, le circuit intégré, le microprocesseur. L'unité du cas ne tient pas à un objet technique unique, comme le télégraphe de Morse ou la tabulatrice de Hollerith, mais à une dynamique commune : chaque bond de performance computationnelle ouvre un nouveau front de consommation énergétique, tandis que chaque gain d'efficacité par composant est réinvesti en densité de calcul. \textcite{nordhaus_two_2007} mesure la chute du coût par calcul sur deux siècles ; \textcite{koomey_implications_2011} mesure le doublement de l'efficacité énergétique par calcul tous les dix-huit mois entre 1946 et 2000. Personne ne croise les deux séries. C'est dans cet écart que la thèse se loge.


% -----------------------------------------------------------------------------
\subsection{Du tube au transistor~: la transition matérielle (1947--1965)}\label{subsec:transistor}
% -----------------------------------------------------------------------------

\dimmark{D10}{+}%
Le 23~décembre 1947, aux Bell Telephone Laboratories, John Bardeen, Walter Brattain et William Shockley réalisent le premier transistor à pointe de contact. Le dispositif, encore fragile et coûteux, remplace le tube à vide dans sa fonction d'amplification et de commutation~: plus petit (quelques millimètres contre plusieurs centimètres), plus fiable (pas de filament à chauffer), et surtout beaucoup moins gourmand en énergie (milliwatts contre watts). Le brevet est déposé en 1948~; le prix Nobel couronne les trois inventeurs en 1956. La trajectoire qui va du premier transistor au circuit intégré couvre moins de quinze ans, une compression temporelle sans équivalent dans les cas précédents.

\dimmark{D5/D8}{+}%
Cette compression n'est pas le fait du marché civil. \textcite{flamm_creating_1988} documente que la demande militaire absorbe la quasi-totalité de la production de transistors dans les premières années. Le programme de miniaturisation des missiles balistiques intercontinentaux (Minuteman) exige des composants légers et fiables~; la Navy et l'Air Force achètent 100\,\% de la production américaine de circuits intégrés en 1962. La courbe d'apprentissage qui en résulte fait chuter le prix du transistor de 45~dollars en 1954 à 2~dollars en 1960~--- une division par vingt en six ans. Le marché civil hérite d'une technologie dont le coût de développement a été socialisé par le budget de la défense.

\dimmark{D10}{+}%
Le circuit intégré (1958--1965) franchit une étape supplémentaire. Jack Kilby chez Texas Instruments et Robert Noyce chez Fairchild Semiconductor déposent des brevets indépendants la même année (1958-1959) pour l'intégration de plusieurs transistors sur un même substrat de silicium. La querelle de priorité qui s'ensuit~--- tranchée par un accord de licences croisées en 1966~--- importe moins que la conséquence architecturale~: le circuit intégré permet de graver des milliers, puis des millions de transistors sur une puce unique, réduisant les connexions externes et les points de défaillance. Noyce fonde Intel en 1968 avec Gordon Moore~; le premier microprocesseur commercial (Intel 4004) sort en 1971. La trajectoire de la loi de Moore (§\,\ref{subsec:moore_law}) commence ici.

\dimmark{0-TER}{+}%
Le profil énergétique de cette transition est paradoxal. L'efficacité par opération s'améliore de façon spectaculaire~: le ratio joules/opération chute de six ordres de grandeur entre l'ENIAC (1946) et l'IBM~360 (1964)\footnote{L'ENIAC consommait environ 150\,kW pour 5\,000 additions par seconde, soit 30~joules par opération. L'IBM~360/50 (1964) consommait environ 20\,kW pour 500\,000 opérations par seconde, soit 0,04~joule par opération. Le gain est de l'ordre de $10^{3}$. Entre l'IBM~360 et l'Intel~4004 (1971), le gain est encore de trois ordres de grandeur. Sources~: \textcite{nordhaus_two_2007}, \textcite{koomey_implications_2011}.}. Mais la puissance totale installée augmente~: les mainframes des années 1960 consomment 20 à 50~kW chacun, et leur nombre croît. Le paradoxe est déjà en place~: l'efficacité par unité s'améliore, la consommation agrégée croît parce que le volume de calcul croît plus vite que l'efficacité. C'est le mécanisme que \textcite{jevons_coal_1865} avait identifié pour le charbon et que la littérature contemporaine nomme \emph{rebound effect}~; il s'appliquera au calcul pendant les cinq décennies suivantes.


% -----------------------------------------------------------------------------
\subsection{La demande de calcul comme problème militaire}\label{subsec:demand_pull_computing}
% -----------------------------------------------------------------------------

\dimmark{D5}{+}%

La demande qui fait naître le calcul électronique n'est ni commerciale ni statistique. Elle est militaire. En 1942, le Ballistic Research Laboratory d'Aberdeen, dans le Maryland, doit produire les tables de tir de l'artillerie américaine : chaque table exige le calcul de plusieurs centaines de trajectoires, chaque trajectoire requiert la résolution numérique d'équations différentielles non linéaires tenant compte de la gravité, de la résistance de l'air, de la rotation terrestre et de la température\footnote{Le calcul d'une seule trajectoire demandait environ 750 multiplications. Avec un \emph{differential analyzer} mécanique, le temps de calcul était d'environ 20 minutes par trajectoire ; à la main, une calculatrice humaine (\emph{computer}, au sens d'alors) y consacrait plusieurs jours.}. Le volume de travail dépasse la capacité de toute organisation humaine. Le laboratoire emploie alors plus d'une centaine de calculateurs, pour l'essentiel des femmes titulaires de diplômes en mathématiques, dont le travail ne suffit pas à répondre au rythme des demandes du front.

\dimmark{D8}{+}%
C'est dans ce contexte que John Mauchly et J.~Presper Eckert, à la Moore School of Electrical Engineering de l'université de Pennsylvanie, proposent en 1942 une machine entièrement électronique capable d'exécuter les mêmes calculs. L'armée finance le projet en avril 1943, pour un coût final de 400\,000 dollars, soit l'équivalent d'environ 7 millions de dollars de 2024. La machine, achevée en février 1946, arrive trop tard pour la guerre qu'elle devait servir. Le point analytique est ailleurs : seul un acteur non marchand, confronté à un besoin opérationnel sans substitut et disposant de ressources sans contrainte de rentabilité, peut financer un prototype dont le coût unitaire exclut tout marché civil. Le pattern est identique à celui que l'on a observé pour le télégraphe, lorsque le Congrès américain vote en 1843 les 30\,000 dollars de la ligne Washington-Baltimore, et pour la tabulatrice, lorsque le Census Bureau finance le concours de 1888 qui donne naissance à la machine de Hollerith. L'État crée le marché avant que le marché n'existe.

\dimmark{D5}{*}%
\textcite{flamm_creating_1988} documente cette structure avec une précision quantitative que la littérature d'histoire des techniques n'atteint pas toujours. Dans les années 1950, le gouvernement fédéral finance environ les deux tiers de la recherche et du développement en informatique aux États-Unis. Lorsque IBM décide de lancer la production en série de l'IBM 650 au milieu de la décennie, la vente projetée de cinquante machines au gouvernement fédéral, sur un total prévisionnel de deux cent cinquante unités, pèse de façon déterminante dans la décision. La demande militaire ne se limite pas au financement direct : la Guerre froide, l'explosion de la première bombe atomique soviétique en 1949, puis la guerre de Corée accélèrent les achats fédéraux de matériel informatique et poussent l'administration Truman à investir massivement dans la technologie de calcul. L'Atomic Energy Commission, la National Security Agency et le Department of Defense soutiennent chacun le développement de systèmes de calcul avancés pour la conception d'armes nucléaires, la cryptanalyse et la défense aérienne.

Le recul de la part publique dans le financement est lui-même une mesure de la transition vers le marché civil. \textcite{flamm_targeting_1987} montre que cette part passe d'environ deux tiers dans les années 1950 à moins d'un cinquième dans les années 1980, à mesure que les ventes privées d'IBM et de ses concurrents prennent le relais. La trajectoire n'est pas linéaire ; elle est marquée par des paliers qui correspondent aux grandes commandes fédérales (le projet SAGE de défense aérienne au milieu des années 1950, les programmes spatiaux de la NASA dans les années 1960). Chaque palier de financement public crée un seuil de capacité industrielle que le marché civil exploite ensuite, sans jamais avoir eu à supporter le risque initial.

\dimmark{0-TER}{+}%
Le projet SAGE (\emph{Semi-Automatic Ground Environment}, 1952--1963) est le cas le plus éloquent de ce couplage militaire entre informatique et énergie. Destiné à détecter et intercepter les bombardiers soviétiques via un réseau de vingt-trois centres reliés à des stations radar à travers l'Amérique du Nord, chaque centre SAGE est équipé d'un ordinateur AN/FSQ-7 de deux cent cinquante tonnes qui consomme trois mégawatts~--- vingt fois la consommation de l'ENIAC\footnote{Les données de puissance sont documentées dans les sources primaires~: \og~It weighed half a million pounds and consumed three megawatts of power\fg{} (interactions.acm.org, 2010, \og Project SAGE, a half-century on\fg{}). Deux machines AN/FSQ-7 étaient construites par centre, l'une en fonctionnement et l'autre en veille chaude permanente~--- une contrainte de disponibilité continue qui fixe la consommation à six mégawatts par centre, indépendamment de la charge de calcul.}. L'opération exige des générateurs spécialisés, des systèmes de refroidissement dédiés, et des bâtiments conçus autour des contraintes électriques de la machine. Western Electric est retenu comme contractant principal pour les bâtiments et l'alimentation électrique interne~--- au même niveau qu'IBM pour le hardware et la RAND Corporation pour le software. Pour la première fois dans l'histoire de l'informatique, l'ingénierie électrique n'est pas un service externe~: elle est une composante d'architecture à part entière, définie simultanément avec la machine et non après elle. Le couplage ICT/énergie cesse d'être un problème opérationnel pour devenir un problème de conception.

\dimmark{D8/D10}{+}%
Deux moments conceptuels prolongent SAGE et préparent la décennie suivante. En avril 1963, J.~C.~R. Licklider, alors directeur du \emph{Information Processing Techniques Office} de l'ARPA, distribue à ses collaborateurs un mémorandum dans lequel il propose le développement d'un \emph{intergalactic computer network} permettant l'opération intégrée de ressources de calcul géographiquement réparties~\parencite{licklider_memorandum_1963}. La même année, Paul Baran publie pour la RAND Corporation un rapport décrivant un réseau de communication distribué fondé sur la commutation par paquets, conçu pour rester opérationnel après une frappe nucléaire qui détruirait les nœuds centraux d'un réseau classique~\parencite{baran_distributed_1963}. Les deux propositions partagent un contexte~--- la défense Cold War et la conviction que la fragilité du réseau commuté par circuit imposait une architecture nouvelle~--- mais leur articulation effective demandait des moyens institutionnels dont aucun des auteurs ne disposait à ce moment. Ce que l'ARPA construira à partir de 1965~--- l'ARPANET, opérationnel en 1969 entre quatre nœuds universitaires~--- combinera la vision distribuée de Baran et l'ambition d'intégration de Licklider~\parencite{abbate_inventing_1999}. Cette généalogie est analytiquement importante pour la suite de la thèse~: le réseau dont naîtra Internet n'est pas un héritage civil contraint par la défense, c'est une commande de la défense dont la séquence étatique de mise en œuvre suit la même logique que celle de SAGE, du télégraphe Morse de 1843 et de la tabulatrice Hollerith de 1888. \textcite{edwards_closed_1996} a proposé de cette période une lecture canonique~: les ordinateurs de la Cold War ne sont pas seulement des artefacts techniques mais des dispositifs discursifs qui structurent la pensée du contrôle, du commandement et de l'ennemi~--- ce que l'analyse présente prolonge en y ajoutant le compte des flux matériels et financiers que la lecture discursive d'Edwards laisse hors champ.

\emergence{Le pattern demande militaire $\rightarrow$ prototype financé par l'État $\rightarrow$ marché civil n'est pas propre au computing. Il reproduit la séquence observée pour le télégraphe (§\,\ref{sec:telegraph}) et pour la tabulatrice (§\ref{subsec:hollerith}), mais à une échelle de financement incomparable. Flamm documente que le gouvernement fédéral finance les deux tiers de la R\&D informatique dans les années 1950, contre un simple vote de crédits ponctuels pour Morse en 1843. L'ampleur du financement reflète l'ampleur de la demande : le calcul balistique, la cryptanalyse et la simulation nucléaire posent des problèmes de taille $n$ tels qu'aucune organisation de travail humain, si ingénieuse soit-elle, ne peut les résoudre dans les délais opérationnels. La machine n'est pas ici un substitut au travail qualifié, comme le calculateur Burroughs remplaçant le comptable ; elle est la condition de possibilité d'un calcul qui n'existait pas avant elle.}

\questionch{2}{Comment la transition de la demande militaire à la demande civile modifie-t-elle la structure des coûts et la trajectoire technologique du computing~? Le chapitre~2 devra formaliser le mécanisme par lequel le financement public sans contrainte de rentabilité crée un seuil technologique que le marché civil ne pourrait pas atteindre seul (D8), et par lequel la demande dérivée de la défense oriente la trajectoire de l'innovation dans des directions qui ne sont pas celles que le marché aurait choisies (D5, D10).}


% -----------------------------------------------------------------------------
\subsection{Les architectures alternatives bloquées (1959--1973)}\label{subsec:alternatives_bloquees}
% -----------------------------------------------------------------------------

\dimmark{D8/D10}{+}%
Le récit qui précède pourrait laisser entendre que la trajectoire ARPANET~--- décentralisation par commutation par paquets, financement militaire américain, ouverture progressive aux universités puis au marché~--- était la seule architecture computationnelle viable au tournant des années 1960. Deux cas contemporains et un troisième, moins connu, montrent qu'il n'en était rien. Le présent chapitre les traite avec la même attention que les cas d'émergence réussie, parce qu'ils révèlent les conditions nécessaires par leur absence, en symétrie avec ce qui a été établi au XIX\textsuperscript{e}~siècle pour la \emph{Difference Engine} de Babbage et pour le télégraphe électrique français entre 1837 et 1851 (§\ref{sec:telegraph}).

\dimmark{D8}{+}%
En Union soviétique, le mathématicien Viktor Glouchkov propose à partir de 1959 le projet OGAS (\emph{Общегосударственная автоматизированная система}, Système national automatisé), conçu pour articuler par un réseau d'ordinateurs les flux d'information de planification entre les ministères, les régions et les usines de l'économie planifiée. \textcite{peters_how_2016} documente le trajet du projet sur trois décennies (1959--1989) et son échec final, qu'il n'attribue pas à des limitations techniques~--- l'URSS dispose alors d'une informatique générale compétente, dont l'ordinateur Oural (1955--1990) documente la diffusion industrielle dans une lecture complémentaire~\parencite{koretsky_ural_2022}~--- mais à la concurrence bureaucratique entre ministères, dont chacun craignait la perte de contrôle informationnel qu'aurait imposé un système intégré. Le titre choisi par Peters pour son ouvrage~--- \emph{How Not to Network a Nation}~--- circonscrit l'objet de la démonstration~: ce n'est pas que la nation n'a pas pu être mise en réseau, c'est que les conditions politiques et institutionnelles d'une telle mise en réseau n'étaient pas réunies, et l'explication des conditions absentes est aussi informative que celle des conditions réunies aux États-Unis pour ARPANET.

\dimmark{D8/0-GOV}{+}%
Au Chili, le gouvernement Allende lance en 1971 le projet Cybersyn (\emph{Synco}), conçu par le cybernéticien britannique Stafford Beer et soutenu directement par Fernando Flores au ministère de l'Économie~: un réseau de Telex reliant les usines nationalisées à un centre de coordination équipé d'un IBM 360/50, permettant le pilotage de l'économie socialiste en temps quasi réel par un tableau de bord cybernétique. \textcite{medina_cybernetic_2011} a reconstitué l'architecture du projet, sa trajectoire courte et son contexte intellectuel. Le projet est interrompu par le coup d'État du 11~septembre 1973~: la salle de contrôle est démantelée, les archives dispersées, les ingénieurs exfiltrés ou tués. La frontière entre non-émergence et destruction politique est ici prise au mot~: Cybersyn n'a pas échoué par défaut de demande ni par manque de financement, il a été supprimé par une intervention politique externe. La différence avec OGAS, dont l'échec est lent et endogène, mérite d'être préservée analytiquement~: les mécanismes de non-émergence ne sont pas tous du même type.

\dimmark{D10}{+}%
Ces cas ne sont pas des contre-factuels~: ils ne disent pas ce qui se serait passé si OGAS ou Cybersyn avaient été déployés, ni ne prétendent que ces architectures auraient été socialement préférables à celle qui a effectivement émergé. Ils disent que la trajectoire ARPANET~--- décentralisée, militaire, américaine~--- n'épuise pas l'espace des configurations computationnelles disponibles au tournant des années 1970, et que d'autres architectures, fondées sur d'autres rapports entre l'État et le calcul, ont été techniquement portées puis institutionnellement bloquées. \textcite{johnstone_deep_2026} font de cette observation un argument structurel du déploiement de la digitalisation~: la trajectoire dominante s'impose moins par sa supériorité intrinsèque que par l'absence des conditions politiques qui auraient permis l'émergence d'alternatives, et le mécanisme se reproduira lors de bifurcations ultérieures que la section~\ref{sec:convergence} examinera (refus d'AT\&T en 1971 de racheter l'ARPANET~; rejet en 1994 du projet de loi Inouye réservant vingt pour cent des capacités télécoms américaines à des usages publics). La symétrie avec la non-émergence du télégraphe électrique français entre 1837 et 1851 et avec la \emph{Difference Engine} de Babbage est analytique~: dans les trois époques, c'est la dimension institutionnelle (\textbf{0-GOV}) et la dimension financement (D8) qui déterminent la réalisation effective ou bloquée des architectures techniquement disponibles, non la supériorité intrinsèque d'une option sur les autres.

\questionch{2}{Quelles conditions institutionnelles (\textbf{0-GOV}) et financières (D8) doivent être réunies pour qu'une architecture computationnelle techniquement portée atteigne effectivement le déploiement~? Et symétriquement~: quels mécanismes bloquent une architecture alternative quand même ses conditions techniques sont réunies~? Le chapitre~2 devra formaliser ces mécanismes en distinguant les non-émergences endogènes (OGAS, fragmentation bureaucratique) des non-émergences par destruction externe (Cybersyn, coup d'État).}


% -----------------------------------------------------------------------------
\subsection{La loi de Moore et l'économie politique du semi-conducteur}\label{subsec:moore_law}
% -----------------------------------------------------------------------------

\dimmark{D10}{*}%
En~1965, Gordon Moore, cofondateur d'Intel, publie dans \emph{Electronics} un article de quatre pages qui deviendra le texte le plus cité de l'histoire de l'industrie électronique \parencite{moore_cramming_1965}. L'observation est empirique~: le nombre de composants par circuit intégré, à coût unitaire minimal, double approximativement tous les dix-huit mois. Mais la force du texte ne réside pas dans l'observation~; elle réside dans sa fonction performative. Moore lui-même l'a formulé sans ambiguïté~: \og \emph{My prediction was about the future direction of the semiconductor industry, and I have found that the industry is best understood through some of its underlying economics}\fg{}\footnote{Cité par \textcite{roussilhe_phase_2025}, qui s'appuie sur Babbage (2024), \og Moore on Moore\fg{}, \emph{The Chip Letter}. \textcite{mody2017moores} documente en détail la construction sociale de la \og loi\fg{} et son rôle dans l'organisation de la R\&D du secteur.}. La \og loi\fg{} de Moore n'est pas une loi physique~; c'est un dispositif de coordination sectorielle qui aligne les anticipations des fabricants de puces, des éditeurs de logiciel et des investisseurs sur une feuille de route commune. Tant que l'industrie croit en sa capacité à maintenir le rythme, elle investit les sommes nécessaires pour le maintenir~; si elle cessait d'y croire, les investissements tariraient et la \og loi\fg{} s'arrêterait. Comme le formule \textcite{roussilhe_phase_2025}, c'est un \og baromètre de la croyance de l'industrie en elle-même\fg{}.

\dimmark{D8}{+}%
Cette croyance a un coût. La \og seconde loi de Moore\fg{}, ou loi de Rock, stipule que le prix d'une fonderie de semi-conducteurs de nouvelle génération double tous les quatre ans \parencite{ross_moores_1995}. Le coût de construction d'une fonderie capable de graver en dessous de 10~nanomètres dépasse désormais vingt milliards de dollars, ce qui exclut de fait toutes les entreprises sauf trois ou quatre dans le monde \parencite{gruber_semiconductor_1996, roussilhe_phase_2025}. L'écart entre les deux \og lois\fg{} définit la condition de survie de l'industrie~: il faut que les économies d'échelle obtenues par la miniaturisation (première loi) croissent plus vite que les coûts de production (seconde loi). Quand cet écart se resserre, la pression à \og consommer\fg{} la puissance de calcul mise sur le marché s'intensifie~--- un mécanisme dont les conséquences sont traitées dans la section~\ref{sec:convergence}.

\dimmark{0-TER}{*}%
L'évolution de la puissance électrique maximale des composants (TDP, \emph{Thermal Design Power}) fournit un proxy factuel de la trajectoire énergétique. L'Intel 4004, premier microprocesseur commercialisé en~1971, affichait un TDP de 0,63~watt. Le Motorola~68000 (1979) atteignait 1,5~watt~; l'Intel Pentium (1994), 11,7~watts~; le Pentium~4 Extreme Edition (2003), 92~watts \parencite{sun_chip_2024, roussilhe_phase_2025}. Entre~2005 et~2018, le TDP des processeurs grand public se stabilisa entre 75 et 150~watts~--- la période où le \emph{scaling} de Dennard, décrit dans la section précédente (§\,\ref{sec:electrification}), contenait encore la consommation. La suite de cette série~--- l'explosion des TDP avec les GPU dédiés à l'IA générative~--- relève de la section~\ref{sec:contemporain}.


% -----------------------------------------------------------------------------
\subsection{Le \emph{consent decree} et la naissance du logiciel comme marchandise}\label{subsec:consent_decree}
% -----------------------------------------------------------------------------

\dimmark{D2/0-GOV}{*}%
Le 14~janvier 1949, le Department of Justice dépose une plainte antitrust contre AT\&T, accusant le monopole téléphonique d'avoir utilisé sa position dominante pour exclure les concurrents du marché des équipements. L'affaire se conclut le 24~janvier 1956 par un \emph{consent decree} dont les conséquences dépassent largement le secteur des télécommunications. AT\&T conserve son monopole sur le service téléphonique, mais doit licencier l'ensemble de ses brevets~--- y compris le transistor~--- à tout demandeur, moyennant des redevances raisonnables. Texas Instruments, Fairchild Semiconductor et la future Intel bénéficient directement de cette diffusion forcée~: la technologie du transistor, développée aux frais des abonnés téléphoniques, devient accessible à l'industrie électronique tout entière\footnote{\textcite{temin_fall_1987} documente l'histoire du \emph{consent decree} de 1956 et ses effets sur la structure de l'industrie américaine des télécommunications. La clause de licence obligatoire est l'un des mécanismes par lesquels l'antitrust américain a involontairement accéléré la diffusion des technologies de l'information.}.

\dimmark{D2/D3}{*}%
Treize ans plus tard, le mécanisme se reproduit pour IBM. En janvier 1969, le dernier jour de l'administration Johnson, le Department of Justice ouvre une enquête antitrust contre IBM, accusant la firme de monopoliser le marché des ordinateurs en liant (\emph{bundling}) matériel, logiciel et services dans une offre indissociable. Le 23~juin 1969, six mois après l'ouverture de l'enquête et avant même toute décision judiciaire, IBM annonce l'\emph{unbundling}~: désormais, le matériel, le logiciel et les services seront facturés séparément. La décision est présentée comme volontaire~; elle anticipe une contrainte judiciaire que la firme préfère devancer. L'affaire elle-même traînera jusqu'en 1982, date à laquelle l'administration Reagan l'abandonnera~; mais le mal est fait, du point de vue d'IBM. Le logiciel est devenu une marchandise autonome.

\dimmark{D3}{*}%
Cette séparation a une conséquence comptable qui prépare la section~\ref{sec:convergence}. Avant 1969, le logiciel n'existe pas comme catégorie économique distincte~: il est inclus dans le prix de la machine, sans valeur comptable propre. Après 1969, le logiciel devient un bien séparé, appropriable, valorisable. Les premières entreprises de logiciel indépendant~--- Computer Associates (1976), SAP (1972), Oracle (1977)~--- naissent dans le sillage de l'\emph{unbundling}. \textcite{campbell-kelly_software_2003} documente cette naissance comme la création d'une industrie qui n'existait pas dix ans plus tôt. Le capital intangible de \textcite{corrado_intangible_2005}, dont la mesure structure la comptabilité nationale depuis les années 2000, trouve son origine institutionnelle dans une décision antitrust de 1969.

\dimmark{D3/0-TER}{+}%
Le point analytique pour la thèse est le suivant. L'\emph{unbundling} crée une asymétrie entre la visibilité comptable du logiciel et l'invisibilité de son substrat énergétique. Le logiciel, une fois séparé du matériel, peut être valorisé comme actif immatériel~: il figure dans les bilans, il est amorti, il génère des revenus de licence. Mais le logiciel n'existe que s'il est exécuté, et l'exécution consomme de l'électricité. Cette consommation n'est pas attachée au logiciel dans les comptes~: elle est noyée dans les charges d'exploitation de l'utilisateur, pas dans le coût de production du vendeur. L'\emph{unbundling} de 1969 inaugure ainsi une dissociation comptable entre la valeur du logiciel (visible, appropriable) et le coût énergétique de son usage (diffus, non imputable), dissociation qui structure encore aujourd'hui la comptabilité du secteur numérique.


% -----------------------------------------------------------------------------
\subsection{Bilan~: trois ruptures et un paradoxe}\label{subsec:bilan_computing}
% -----------------------------------------------------------------------------

\dimmark{0-TER}{*}%
La période 1945--1975 voit se produire trois ruptures simultanées dont la conjonction prépare les développements des sections suivantes.

\paragraph{Première rupture~: l'architecture à programme enregistré.} Le \emph{stored-program concept}, formalisé par von Neumann en 1945 et matérialisé dans l'EDVAC puis l'IAS Machine, unifie l'espace mémoire du programme et l'espace mémoire des données. Le programme devient donnée~: il peut être copié, modifié, transmis comme n'importe quel fichier. Cette fusion est la condition technique de l'industrie du logiciel~: sans elle, le logiciel resterait câblé dans la machine et ne pourrait pas être vendu séparément. Mais elle a aussi une conséquence énergétique~: la mémoire qui contient le programme consomme de l'énergie en permanence (rafraîchissement DRAM). Stocker un programme n'est plus gratuit énergétiquement, contrairement au carton perforé ou au câblage fixe.

\paragraph{Deuxième rupture~: la transition tube-transistor-IC.} L'efficacité énergétique par opération s'améliore de neuf ordres de grandeur entre l'ENIAC (1946) et l'Intel 4004 (1971). Cette amélioration résulte de la miniaturisation (moins de matière à commuter, moins d'énergie dissipée) et du passage à l'électronique à semi-conducteurs (pas de filament à chauffer). Le tableau~\ref{tab:efficacite-calcul} documente les jalons de cette trajectoire.

\begin{table}[htbp]
\caption{Efficacité énergétique du calcul~: jalons 1946--1971.}\label{tab:efficacite-calcul}
\centering\small
\renewcommand{\arraystretch}{1.15}
\begin{tabular}{@{}llrrr@{}}
\toprule
\textbf{Année} & \textbf{Machine} & \textbf{Puissance (kW)} & \textbf{Opérations/s} & \textbf{Ops/kWh} \\
\midrule
1946 & ENIAC & 150 & 5\,000 & $1{,}2 \times 10^{2}$ \\
1954 & IBM 704 & 30 & 40\,000 & $4{,}8 \times 10^{3}$ \\
1964 & IBM 360/50 & 20 & 500\,000 & $9{,}0 \times 10^{4}$ \\
1971 & Intel 4004 & 0,0005 & 60\,000 & $4{,}3 \times 10^{11}$ \\
\bottomrule
\end{tabular}
\renewcommand{\arraystretch}{1.0}
\par\smallskip
\footnotesize Sources~: \textcite{nordhaus_two_2007}, \textcite{koomey_implications_2011}. Ops/kWh = opérations par kilowattheure.
\end{table}

\paragraph{Troisième rupture~: l'\emph{unbundling} et la naissance du capital intangible.} L'\emph{unbundling} d'IBM (1969) crée le logiciel comme marchandise séparée du matériel. Cette séparation institutionnelle prépare la mesure du capital intangible de \textcite{corrado_intangible_2005}, qui structure aujourd'hui la comptabilité nationale. Elle prépare aussi la dissociation comptable entre la valeur du logiciel (visible, appropriable) et le coût énergétique de son usage (diffus, non imputable).

\dimmark{0-TER}{*}%
\paragraph{Le paradoxe.} Ces trois ruptures ont un effet combiné paradoxal. L'efficacité par opération s'améliore spectaculairement~; la consommation agrégée croît parce que le volume de calcul croît plus vite que l'efficacité. Le même mécanisme de rebond (\emph{rebound effect}) que \textcite{jevons_coal_1865} avait identifié pour le charbon s'applique au calcul~: chaque gain d'efficacité libère des ressources qui sont réinvesties en volume, de sorte que la consommation totale augmente. La section~\ref{sec:convergence} documentera comment ce paradoxe se déploie entre 1975 et 2007~; la section~\ref{sec:contemporain} montrera qu'il s'intensifie après 2007 avec l'IA générative.

\emergence{La période 1945--1975 voit trois ruptures simultanées~: l'architecture à programme enregistré qui fusionne STORE et COMPUTE~; la transition tube-transistor-IC qui améliore l'efficacité énergétique de neuf ordres de grandeur~; l'\emph{unbundling} qui crée le logiciel comme capital intangible séparé du matériel. Ces ruptures préparent la convergence numérique de §\,\ref{sec:convergence} et l'explosion de §\,\ref{sec:contemporain}. Le paradoxe qui les traverse~--- amélioration de l'efficacité unitaire, croissance de la consommation agrégée~--- est le mécanisme central de l'occultation de la dimension 0-TER dans les décennies qui suivent.}

\questionch{3}{Comment quantifier le rebond énergétique du calcul sur la période 1945--2020~? Le croisement des séries d'efficacité de \textcite{koomey_implications_2011} avec les séries de capacité de \textcite{hilbert_worlds_2011} et les statistiques de consommation électrique des datacenters de \textcite{masanet_recalibrating_2020} permettrait de tester l'hypothèse Jevons sur le calcul. Ce croisement n'a jamais été publié sous cette forme.}


\begin{table}[htbp]
\caption{Dimensions économiques du calcul programmable (§\ref{sec:computing}).}\label{tab:computing-dimensions}
\centering\small
\renewcommand{\arraystretch}{1.15}
\begin{tabularx}{\textwidth}{@{}cXc@{}}
\toprule
\textbf{Dim.} & \textbf{Ce que §\ref{sec:computing} établit} & \textbf{Niveau} \\
\midrule
D2 & \emph{Consent decree} 1956 (AT\&T) et \emph{unbundling} 1969 (IBM)~: l'antitrust comme mécanisme de diffusion technologique. & ★ \\
D3 & \emph{Unbundling} 1969~: naissance du logiciel comme capital intangible séparé du matériel. & ★ \\
D5 & Demande militaire (tables de tir, cryptanalyse, simulation nucléaire, SAGE)~: l'État crée le marché que le privé ne pourrait financer. & ★ \\
D8 & L'État finance 2/3 de la R\&D dans les années 1950~; le marché civil prend le relais dans les années 1970--1980. & ★ \\
D10 & Loi de Moore comme dispositif de coordination performatif~; \emph{scaling} de Dennard comme bouclier d'occultation énergétique. Architectures alternatives bloquées (OGAS, Cybersyn). & ★ \\
\midrule
0-TER & Efficacité $\times 10^{9}$, mais consommation agrégée croît (rebond Jevons). Architecture von Neumann couple STORE et énergie (rafraîchissement DRAM). ENIAC = 150\,kW~; SAGE = 3\,MW par centre. & ★ \\
0-GOV & Antitrust (consent decree, unbundling)~; financement fédéral massif~; trajectoires alternatives bloquées par la configuration politique (OGAS, Cybersyn). & ★ \\
\bottomrule
\end{tabularx}
\renewcommand{\arraystretch}{1.0}
\end{table}

% =============================================================================
% §1.9 La convergence sur le réseau (1975–2007)
% =============================================================================
\clearpage
% =============================================================================
% §1.9 La convergence sur le réseau (1975–2007)
% =============================================================================
% Résultats dimensionnels visés :
%   0-TER  — la numérisation crée une consommation permanente (STORE actif)
%   D2     — plateformes et marchés winner-take-all
%   D3     — capital intangible ≥ tangible (Corrado et al. 2005)
%   D7     — coûts de transaction quasi nuls, modularité (Langlois 2003)
%   D8     — cycle Perez (dotcom = railway mania)
%   D10    — open standards comme dispositif de coordination
% Sources empiriques principales :
%   Hilbert & López 2011, David 1990, Corrado/Hulten/Sichel 2005/2009,
%   Perez 2002, Koomey 2011, Flamm 1988, Shapiro & Varian 1999
% =============================================================================
 
\section{La convergence sur le réseau (1975--2007)}\label{sec:convergence}
 
La section précédente s'achevait sur une concomitance ignorée : en 1974, le microprocesseur, le \emph{Dennard scaling} et le premier choc pétrolier coexistent sans que les trajectoires du calcul, de l'énergie et de l'économie se croisent. La période qui s'ouvre ici est celle où les trois fonctions --- transmettre, stocker, calculer --- se numérisent sur le même substrat électrique, mais pas au même rythme, pas selon les mêmes logiques, et pas avec les mêmes conséquences.
 
Ce que la mesure empirique va confirmer, la théorie des systèmes sociotechniques l'avait anticipé. \textcite{ropohl_systemtheorie_1979}, dans son~\emph{Allgemeine Technologie} dont la première édition paraît en 1979, observe~: \og~Le téléphone, le télex, la radio, la télévision et l'ordinateur ont jusqu'à présent constitué des mondes séparés, aussi bien dans leurs réseaux que dans leurs terminaux. Mais ces mondes distincts convergent désormais vers un cosmos unifié~--- un réseau de communication mondial intégré, avec des terminaux universels capables de recevoir, traiter et envoyer indifféremment du texte, du son, de l'image et du film.~\fg\footnote{\og~Bislang sind Telefon, Fernschreiben, Rundfunk, Fernsehen und Computer sowohl mit ihren Netzen wie auch mit ihren Endgeräten getrennte Welten gewesen. Inzwischen aber wachsen diese verschiedenen Welten zu einem umfassenden Kosmos zusammen, zu einem globalen integrierten Kommunikationsnetz mit universellen Endgeräten, die gleichermassen Schrift, Ton, Bild und Film empfangen, verarbeiten und versenden können.~\fg~\parencite[Fallbeispiel \og~Computer~\fg, 3.~Aufl.]{ropohl_systemtheorie_1979}. L'observation est formulée en~1979, au moment précis où TCP/IP est en cours de développement (il sera adopté en 1983) et où le Web n'existe pas encore. Ropohl la confirme dans la troisième édition de 2009, ajoutant uniquement la référence à la \og~loi de Moore~\fg{} et à l'intégration ordinateur-téléphone mobile, qu'il décrit comme \og~en cours~\fg{} à la date de révision.} L'observation est formulée à un moment où la convergence n'est pas encore réalisée~--- TCP/IP sera adopté en 1983, le Web inventé en 1989~--- mais où elle est déjà inscrite dans la logique des systèmes sociotechniques existants. Ce que Ropohl appelle \emph{informationstechnische Integration} est exactement ce que la période 1975--2007 produit empiriquement.

\dimmark{0-TER}
L'étude de référence pour saisir cette période dans sa globalité est celle de \textcite{hilbert_worlds_2011}, publiée dans \emph{Science}, qui mesure les capacités technologiques mondiales de stockage, de communication et de calcul en suivant soixante technologies analogiques et numériques entre 1986 et 2007. L'étude repose sur plus de mille cent sources primaires, documentées dans un supplément méthodologique de trois cent deux pages \parencite{lopez_methodological_2012}. Trois résultats structurent la période.
 
Le premier est l'asymétrie des taux de croissance. La capacité de calcul à usage général croît de 58~\% par an entre 1986 et 2007, soit un doublement tous les dix-huit mois et une vitesse dix fois supérieure à celle du PIB mondial. La capacité de télécommunication bidirectionnelle croît de 28~\% par an, cinq fois plus vite que le PIB. La capacité de stockage installée croît de 23~\% par an, quatre fois plus vite. La diffusion de l'information par broadcasting, en revanche, ne croît que de 6~\% par an, au rythme de l'économie. Le calcul croît donc dix fois plus vite que le stockage, qui croît lui-même quatre fois plus vite que l'économie : les trois fonctions ne sont pas simplement en expansion ; elles divergent.
 
Le deuxième résultat est l'asynchronie de la numérisation. Les télécommunications sont numériques dès 1990 (99,9~\% en 2007). Le stockage ne bascule qu'au début des années 2000 (25~\% numérique en 2000, 94~\% en 2007 ; les cassettes VHS analogiques dominaient encore en 1986). Le broadcasting reste à 75~\% analogique en 2007 (la télévision hertzienne terrestre n'est éteinte aux États-Unis qu'en 2009, en France en 2011). La numérisation n'est pas un événement unique ; c'est un processus de trente ans dont le rythme diffère par fonction.
 
Le troisième résultat est la composition des fonctions dans le substrat numérique. \textcite{hilbert_worlds_2011} définissent le calcul, à la suite de \textcite{turing_computable_1936}, comme « la transmission répétée d'information à travers l'espace (communication) et le temps (stockage), guidée par une procédure algorithmique ». Autrement dit, le calcul numérique contient structurellement de la transmission et du stockage. Lorsque les trois fonctions partagent le même substrat, leurs demandes énergétiques se composent au lieu de rester séparées. C'est dans cette composition que se loge le couplage structurel entre information et énergie que les cas précédents n'avaient fait qu'entrevoir.
 
 
% -----------------------------------------------------------------------------
\subsection{Le calcul distribué sur infrastructure analogique (1975--1995)}\label{subsec:pc}
% -----------------------------------------------------------------------------
 
\dimmark{D10}
En 1975, le microprocesseur Intel 8080 permet à Ed Roberts de construire l'Altair 8800, vendu en kit à 439 dollars par correspondance. Deux ans plus tard, l'Apple~II et le TRS-80 rendent l'ordinateur personnel accessible à un public non spécialiste. En 1981, IBM entre sur le marché avec un ordinateur personnel fondé sur des standards ouverts : le processeur Intel 8088, le système d'exploitation MS-DOS de Microsoft, et une architecture documentée que n'importe quel fabricant peut reproduire. La décision d'IBM de publier les spécifications de sa machine, prise sous la pression d'un développement accéléré (le Project Chess ne dure qu'un an), crée les conditions d'un marché concurrentiel par clonage. En trois ans, Compaq, Dell et des dizaines de fabricants asiatiques produisent des compatibles IBM-PC à des prix décroissants. \textcite{vanderkooij_information_2023} identifie cet épisode comme le déclencheur de la Third Computer Revolution (1980--2000).

\dimmark{D5/D10}{+}%
La pénétration du PC dans les foyers américains, négligeable au début de la décennie, atteint quinze pour cent en 1989~; sur la même période, le nombre de \emph{network hosts} connectés aux réseaux interuniversitaires passe d'environ deux mille en 1985 à cent soixante mille à la fin de la décennie. Cette diffusion ne tient pas qu'à la trajectoire technique de l'Altair vers l'IBM~PC~; elle est médiatisée par une culture qui se forme dans la baie de San Francisco entre 1968 et 1985 et qui infléchit durablement l'architecture des systèmes. Le \emph{Whole Earth Catalog} de Stewart Brand, publié pour la première fois en 1968, articule une critique technologique de contre-culture avec une attention aux outils «~de petite échelle~» dont les ordinateurs personnels feront partie~\parencite{turner_counterculture_2008, kirk_whole_2001}. Ted Nelson publie en 1974 \emph{Computer Lib~/ Dream Machines}, manifeste auto-édité qui appelle les utilisateurs à défier la «~\emph{computer priesthood}~» d'IBM et du gouvernement et à concevoir les ordinateurs comme \emph{creativity systems} plutôt que comme machines de calcul d'entreprise~\parencite{nelson_computer_1974}. La trajectoire des fondateurs d'Apple, Steve Jobs et Steve Wozniak, croise celle du \emph{Homebrew Computer Club} de Menlo Park où se diffusent ces valeurs~\parencite{markoff_dormouse_2005, levy_hackers_1984}.

\dimmark{D2/0-GOV}{*}%
La tension entre la culture du partage gratuit et la nécessité de monétisation de l'industrie naissante se cristallise au tournant de la décennie. Le 3~février 1976, Bill Gates~--- alors âgé de vingt ans et cofondateur d'une jeune entreprise nommée Micro-Soft~--- publie dans le bulletin MITS \emph{Altair Computer Notes} une «~lettre ouverte aux hobbyists~» dans laquelle il reproche aux utilisateurs de copier sans autorisation l'interpréteur BASIC qu'il a écrit pour l'Altair~\parencite{gates_hobbyists_1976}. L'argument est économique~: si le logiciel ne se vend pas, personne ne pourra investir dans son développement. La réponse de la communauté est largement hostile~: la pratique du partage est tenue pour intrinsèque à la pratique informatique elle-même, et l'idée de \emph{vendre} du logiciel apparaît comme une trahison des valeurs partagées. La résolution institutionnelle de cette tension se fera dans deux directions concurrentes que la suite du chapitre suivra~: le mouvement \emph{Free Software} fondé par Richard Stallman en 1984~\parencite{stallman_gnu_1984}, qui codifie la \emph{copyleft license} comme architecture juridique alternative et donnera GNU puis Linux~; et le modèle de licensing massif que Gates et Allen construisent à partir de 1980 et qui fera de Microsoft l'entreprise la plus valorisée du monde en 1998~\parencite{campbell_kelly_computer_2014}. Les deux architectures coexistent encore aujourd'hui dans un partage fonctionnel~: les systèmes dérivés d'Unix libre dominent les serveurs et les infrastructures cloud, Microsoft domine les postes de travail et le logiciel d'entreprise. \textcite{johnstone_deep_2026} identifient cette persistance comme une trace structurelle de la phase \emph{counterculture} dans l'architecture même du système informationnel contemporain, qui contraint encore la trajectoire un demi-siècle après sa cristallisation.

\dimmark{D7}
Le tableur transforme la fonction économique du micro-ordinateur. VisiCalc (1979), conçu par Dan Bricklin et Bob Frankston pour l'Apple~II, est la première application qui justifie l'achat d'un ordinateur par une entreprise pour des raisons de productivité et non de prestige technologique. Lotus 1-2-3 (1983) prend le relais sur l'IBM PC. Le tableur ne remplace pas un employé ; il modifie la manière dont les cadres interagissent avec les données financières, permettant la simulation instantanée de scénarios. Le parallèle avec la tabulatrice au Nautical Almanac Office (§\ref{subsec:hollerith}) est exact : la machine ne supprime pas le travail, elle en change la nature et la localisation.
 
\dimmark{D4}
Pendant cette phase, le calcul est distribué mais isolé. Le micro-ordinateur ne communique que par le réseau téléphonique analogique, via un modem dont le débit ne dépasse pas 14\,400 bits par seconde en 1995. Le stockage est local : disquettes, puis disques durs dont la capacité passe de 10 mégaoctets (IBM PC XT, 1983) à quelques gigaoctets au milieu des années 1990. L'infrastructure de communication reste analogique ; le réseau téléphonique commuté (PSTN) porte des données numériques sous forme de signaux modulés. \textcite{hilbert_worlds_2011} montrent que la capacité de télécommunication par habitant ne passe que de 0,16 à 0,23 mégaoctet compressé par jour entre 1986 et 1993, un doublement bien plus lent que celui de la capacité de calcul (de 0,06 à 0,8 MIPS par habitant, soit un facteur treize sur la même période).
 
\dimmark{D1}
C'est dans ce décalage entre la puissance de calcul locale et la bande passante de communication que s'installe le paradoxe de la productivité. En 1987, Robert Solow observe que « l'on voit l'ère de l'ordinateur partout sauf dans les statistiques de productivité ». \textcite{david_dynamo_1990}, dans un article court mais décisif publié dans l'\emph{American Economic Review}, propose une explication par analogie historique : l'électrification avait suivi le même schéma. L'ampoule est inventée en 1879, mais en 1900, seuls 3~\% des foyers américains disposent de l'éclairage électrique et les moteurs électriques ne représentent que 5~\% de la force motrice industrielle. Il faut attendre les années 1920 pour que la diffusion atteigne 50~\% et que les gains de productivité deviennent mesurables. David montre que le délai s'explique par le coût de remplacement d'installations encore fonctionnelles et par la nécessité de réinventer l'organisation du travail (le passage du moteur central aux moteurs individuels par machine).
 
Le parallèle dynamo/computer est le point de départ méthodologique de cette thèse (§\ref{sec:electrification}). Ce que David ne pose pas, et que les données de \textcite{hilbert_worlds_2011} permettent de voir, c'est que le décalage entre calcul et communication n'est pas seulement un problème de diffusion organisationnelle ; c'est un décalage physique entre deux taux de numérisation. Le calcul est numérique depuis 1945 ; la communication ne le devient qu'en 1990. Pendant quarante-cinq ans, le calcul numérique circule sur un réseau analogique. Le paradoxe de la productivité est, entre autres, le symptôme de cette hétérogénéité de substrat.
 
\emergence{Le PC distribue le COMPUTE chez l'utilisateur, mais le TRANSMIT reste analogique et le STORE reste local. La consommation d'énergie du traitement de l'information est encore ponctuelle : l'ordinateur peut être éteint, la disquette ne consomme rien sur l'étagère. Le profil énergétique de cette phase est \emph{additif} : le PC s'ajoute au papier sans le remplacer (la consommation de papier de bureau augmente dans les années 1980--1990).}
 
 
% -----------------------------------------------------------------------------
\subsection{La numérisation du réseau et la double infrastructure (1983--2000)}\label{subsec:www}
% -----------------------------------------------------------------------------
 
\dimmark{D10}
L'adoption de TCP/IP comme standard d'Internet en janvier 1983 est le moment technique où la transmission d'information bascule dans le numérique. La commutation par paquets transforme la nature du TRANSMIT : dans le réseau téléphonique commuté, chaque appel monopolise un circuit physique de bout en bout ; dans le réseau à paquets, l'information est découpée, routée par des algorithmes (ce qui est du calcul), mise en cache dans des mémoires tampons (ce qui est du stockage), et réassemblée à destination. Le TRANSMIT numérique n'est plus une fonction pure ; il contient structurellement du COMPUTE et du STORE. Les données de \textcite{hilbert_worlds_2011} confirment cette transformation : la capacité de télécommunication effective passe de 0,23 mégaoctet compressé par habitant et par jour en 1993 à 1,01 en 2000, puis à 27 en 2007 --- soit une multiplication par cent dix-sept en quatorze ans, tirée par l'introduction du haut débit.
 
\dimmark{0-TER}
La transition vers le numérique s'effectue cependant \emph{sur} l'infrastructure analogique existante, pas à sa place. Le modem dial-up des années 1990 transmet des paquets TCP/IP sur la boucle locale en cuivre du réseau téléphonique. Le DSL, introduit à la fin de la décennie, réutilise la même paire de cuivre en exploitant des fréquences que la voix n'utilise pas. La fibre optique, techniquement mature depuis les travaux de Corning en 1970 et le premier câble transatlantique TAT-8 en 1988, ne remplace la boucle locale cuivre que progressivement, parce que le cuivre est déjà posé et que la fibre exige de creuser de nouvelles tranchées. Pendant toute la période de transition, les opérateurs maintiennent la double infrastructure : les anciens commutateurs de circuits \emph{et} les nouveaux routeurs IP. La transition numérique est un double coût d'infrastructure, pas une substitution.
 
\dimmark{D6}
En 1989, Tim Berners-Lee propose au CERN un système d'hypertexte pour faciliter l'accès à l'information interne de l'organisation. Les protocoles qu'il développe --- HTTP, HTML, URL --- sont rendus publics en 1993. L'apparition du navigateur graphique Mosaic la même année, puis de Netscape en 1994, ouvre le Web au grand public. En 1998, Google introduit un moteur de recherche fondé sur l'analyse des liens hypertextes (le PageRank). Le nombre d'utilisateurs d'Internet dans le monde passe de quelques millions en 1993 à 280 millions en 1999 \parencite{abbate_inventing_1999}.
 
Le Web ajoute une dimension nouvelle au traitement de l'information : le stockage distribué permanent. Chaque page web est un fichier hébergé sur un serveur allumé vingt-quatre heures sur vingt-quatre, accessible à tout moment par n'importe quel utilisateur connecté. C'est la première infrastructure de STORE permanent à l'échelle mondiale. Un livre peut être fermé sans coût ; un serveur web ne peut pas être éteint sans que les pages qu'il héberge disparaissent. La mise en réseau du stockage transforme un acte ponctuel (stocker un fichier sur un disque local) en une obligation continue (alimenter le serveur qui rend le fichier accessible).
 
\dimmark{0-TER}
C'est dans cette période que l'acronyme \emph{ICT} (\emph{Information and Communication Technology}) se stabilise. L'OCDE et l'UNESCO l'adoptent à la fin des années 1990 pour désigner la convergence des télécommunications, de l'informatique et de l'audiovisuel. L'acronyme consacre l'unification de trajectoires techniques jusque-là distinctes. Il entérine aussi, dans son silence sur l'énergie, l'occultation que cette thèse cherche à lever : \emph{ICT} ne contient ni \emph{E} pour \emph{energy} ni \emph{M} pour \emph{materials}. La dimension 0-TER est absente du nom même du secteur.
 
La matérialité de l'infrastructure de réseau reproduit les structures observées au XIXe~siècle pour le télégraphe. Les câbles sous-marins en fibre optique suivent largement les routes des câbles télégraphiques posés un siècle plus tôt. La silice ultrapure nécessaire à la fabrication de la fibre, les navires câbliers spécialisés, les stations de répétition sous-marines sont invisibles pour l'utilisateur, tout comme l'étaient la gutta-percha malaise et les poteaux télégraphiques du réseau Chappe. Le parallèle matériel avec §\ref{subsec:materialite} est structurel, pas métaphorique : même fonction (TRANSMIT longue distance), même occultation du substrat physique, même concentration géographique de la production.

\dimmark{0-GOV/D10}{+}%
La commercialisation d'Internet entre 1993 et 2000 n'efface pas l'existence de bifurcations institutionnellement portées et politiquement écartées, en symétrie avec les architectures alternatives bloquées de la section précédente (§\ref{subsec:alternatives_bloquees}). Trois moments le rappellent. En 1971, AT\&T se voit offrir par l'ARPA la possibilité de racheter l'infrastructure ARPANET~; après examen, l'entreprise refuse, jugeant que la commutation par paquets ne sera jamais commercialement viable~\parencite{rose_att_2012, abbate_inventing_1999}. La décision laisse l'ARPANET sous tutelle militaire puis universitaire jusqu'à la privatisation de NSFNET en 1995, et ouvre l'espace dans lequel les fournisseurs d'accès commerciaux des années 1990 viendront se loger~; \textcite{russell_open_2014} a documenté plus largement le scepticisme des opérateurs télécommunications historiques à l'égard de la commutation par paquets pendant la majeure partie des années 1970, scepticisme qui a paradoxalement permis l'émergence d'Internet en dehors de leur tutelle. En 1994, le sénateur démocrate Daniel Inouye propose un projet de loi qui aurait obligé les opérateurs de télécommunications à réserver vingt pour cent de leurs capacités à des usages publics, éducatifs et non commerciaux~\parencite{tarnoff_internet_2022}~; le projet ne passe pas, faute d'un mouvement social organisé pour le porter. La privatisation intégrale de NSFNET suit en 1995, et le \emph{Framework for Global Electronic Commerce} promulgué par l'administration Clinton en juillet 1997 consacre cette orientation en posant l'autorégulation industrielle comme doctrine internationale du numérique~\parencite{clinton_framework_1997}. Ces trois moments ne sont pas des accidents~: ils marquent les points où une autre trajectoire d'Internet, plus régulée et plus partagée, aurait été institutionnellement possible mais a été écartée, et leur cumul fixe pour les vingt années suivantes un cadre où l'autorégulation marchande paraît la seule architecture envisageable.


% -----------------------------------------------------------------------------
\subsection{Le capital intangible et la mesure de l'invisible}\label{subsec:logiciel}
% -----------------------------------------------------------------------------
 
\dimmark{D3}
La séparation du hardware et du software imposée à IBM en 1969\footnote{Sous la pression de l'enquête antitrust ouverte par le Department of Justice (\emph{United States v.~IBM}, lancée en janvier 1969), IBM annonce dès le 23~juin 1969 sa nouvelle politique de \emph{unbundling}~: facturer séparément le hardware, le software et les services, ce qui ouvre la voie à une industrie indépendante du logiciel. L'enquête antitrust elle-même s'étend jusqu'en janvier 1982, date à laquelle elle est abandonnée sans condamnation par l'administration Reagan.} a créé les conditions d'une industrie du logiciel indépendante. En trente ans, cette industrie passe d'une activité artisanale de programmation sur mesure à un marché mondial dominé par quelques acteurs (Microsoft pour les systèmes d'exploitation, Oracle pour les bases de données, SAP pour la gestion intégrée). Le logiciel cesse d'être un poste de R\&D interne pour devenir un bien marchand, puis une catégorie de capital dans les comptes nationaux.
 
\textcite{corrado_measuring_2005} et \textcite{corrado_intangible_2009} fournissent la mesure de cette transformation. Leur estimation montre que l'investissement total des entreprises américaines dans le capital intangible --- logiciel, R\&D, design, formation, capital organisationnel, image de marque --- atteignait environ mille milliards de dollars en 1999, un montant approximativement égal à l'investissement en capital tangible la même année. Le basculement se situe au milieu des années 1990. En 2003, environ huit cent milliards de dollars d'investissement intangible restaient exclus des comptes nationaux publiés, ce qui conduit à l'exclusion de plus de trois mille milliards de dollars de stock de capital intangible.

\dimmark{D8/D6}{+}%
La croissance du capital intangible n'est pas uniformément distribuée à travers les secteurs. Un acteur peu visible dans la chronologie populaire de l'internet en a pourtant été l'un des premiers déployeurs massifs~: le secteur financier. \textcite{schiller_digital_2014} documente avec une précision quantitative rare que l'\emph{American Bankers Association} reconnaît en 1970 les réseaux de données comme \emph{operational basis for finance}, quinze ans avant la diffusion grand public de la téléinformatique. Le budget de traitement de données de Morgan Stanley est multiplié par cinq entre 1979 et 1985~; les premières formes d'\emph{algorithmic trading} apparaissent au début des années 1980 et se généralisent dans les deux décennies suivantes~\parencite{mackenzie_trading_2021}. La financiarisation des économies occidentales au sens que lui donne \textcite{epstein_financialization_2005} est techniquement portée par le déploiement précoce des ICT dans la banque d'investissement, le trading et la gestion d'actifs, bien avant que la commercialisation d'internet ne rende ce déploiement visible. Cette antériorité est analytiquement importante~: une part du capital intangible que mesurera \textcite{corrado_intangible_2009} dans les années 1990 a déjà été investie par le secteur financier dans la décennie précédente, et l'angle économique sous lequel le secteur ICT sera lu à partir de 2000 est en grande partie un héritage de l'angle financier sous lequel il a été déployé entre 1970 et 1990.

\dimmark{D1}
L'inclusion du capital intangible dans le cadre de la comptabilité de la croissance modifie les résultats observés. \textcite{corrado_intangible_2009} montrent que le taux de croissance de la production par travailleur augmente plus rapidement lorsque les actifs intangibles sont comptabilisés, et que l'approfondissement du capital (\emph{capital deepening}) devient la source dominante sans ambiguïté de la croissance de la productivité du travail. Le capital intangible n'est donc pas un détail comptable ; c'est la catégorie qui détermine l'interprétation des sources de la croissance économique dans la période 1995--2007.
 
\dimmark{0-TER}
Le paradoxe que ces mesures ne posent pas est le suivant : le capital intangible est traité dans les comptes nationaux comme s'il n'avait pas de substrat matériel. Le logiciel ne s'use pas physiquement (pas d'amortissement par usure) ; il n'a pas de localisation contrainte (il peut être copié) ; il ne consomme pas d'énergie en tant que tel (une ligne de code sur un disque n'a pas de coût énergétique propre). Mais le logiciel n'existe que s'il est exécuté par un processeur alimenté en électricité, et les données qu'il traite n'existent que sur des serveurs alimentés en permanence. Le capital intangible a un coût énergétique indirect que la catégorie comptable rend invisible. La séparation entre « intangible » et « tangible » dans la comptabilité nationale reproduit, au niveau des comptes, la séparation entre économie de l'information et économie de l'énergie que cette thèse identifie au niveau disciplinaire.

\dimmark{D3/0-TER}{*}%
Ce découplage comptable a une conséquence pratique~: les entreprises qui investissent massivement dans le capital intangible (logiciels, données, algorithmes) externalisent la consommation énergétique correspondante vers les fournisseurs de cloud et les opérateurs de datacenters, dont les coûts apparaissent comme des charges d'exploitation et non comme l'amortissement d'un actif. La structure comptable rend donc structurellement difficile l'imputation du coût énergétique de l'information à ses bénéficiaires économiques. Le chapitre~2 reprendra ce mécanisme sous la dimension D3~; le chapitre~3 tentera d'en estimer l'ampleur.

\questionch{La mesure du capital intangible de \textcite{corrado_intangible_2009} peut-elle être complétée par une estimation du coût énergétique indirect de ce capital ? Le chapitre~3 posera cette question en croisant les séries de capital intangible avec les estimations de consommation d'énergie du secteur ICT de \textcite{malmodin_energy_2018}.}
 
 
% -----------------------------------------------------------------------------
\subsection{La bulle spéculative et le cycle Perez}\label{subsec:dotcom}
% -----------------------------------------------------------------------------
 
\dimmark{D8}
Le NASDAQ passe d'environ mille points en 1995 à 5\,048 le 10 mars 2000, puis retombe à 1\,114 en octobre 2002. La capitalisation détruite entre le sommet et le creux dépasse cinq mille milliards de dollars. L'épisode n'est pas un accident ; c'est une régularité structurelle que \textcite{perez_technological_2002} formalise dans un modèle qui distingue cinq phases dans la relation entre capital financier et révolution technologique : irruption, ferveur spéculative (\emph{frenzy}), point de retournement (\emph{turning point}), synergie et maturité.
 
Le même schéma se rejoue pour chaque GPT. La canal mania des années 1790 en Grande-Bretagne, la railway mania des années 1840, la corporate mania électrique des années 1880--1890 et le radio boom des années 1920 précèdent chacun un krach suivi d'une phase de déploiement systématique de l'infrastructure. Le krach ne détruit pas l'infrastructure ; il la transfère à bas coût des spéculateurs vers les utilisateurs. Les chemins de fer construits pendant la bulle des années 1840 sont les mêmes qui transportent les marchandises dans les années 1850 ; les câbles de fibre optique posés pendant la bulle dotcom sont les mêmes qui portent le trafic Internet dans les années 2010.
 
L'épisode Iridium condense le même cycle à l'échelle d'un projet. Motorola investit cinq milliards de dollars dans une constellation de soixante-six satellites de communication en orbite basse, lancée commercialement en novembre 1998. Neuf mois plus tard, la société dépose le bilan, faute de clients (le prix du terminal est de 3\,000 dollars, celui de la minute de communication de 3 à 8 dollars, à un moment où le téléphone cellulaire terrestre devient bon marché et couvre l'essentiel du marché solvable). La constellation est rachetée en décembre 2000 pour 25 millions de dollars, soit 0,5~\% de l'investissement initial. L'infrastructure survit à la faillite de son promoteur et fonctionne encore aujourd'hui pour des usages de niche (maritime, aviation, zones non couvertes). Le parallèle avec les faillites des premiers câbles transatlantiques de 1858 est exact.

\dimmark{D5/D6}{+}%
La bulle laisse cependant une tension irrésolue que le seul cycle Perez ne couvre pas. La phase \emph{frenzy} a vu se déployer une infrastructure et une base d'utilisateurs sans qu'un modèle commercial stable n'émerge~: les abonnements payants sont massivement rejetés par les utilisateurs hérités de la culture du partage gratuit (§\ref{subsec:pc}), et la publicité bannière des premiers sites reste peu ciblée et peu rentable. Une formulation conceptuelle anticipe la résolution sans la produire~: \textcite{goldhaber_attention_1997} publie en avril 1997 dans \emph{First Monday} un article intitulé «~The Attention Economy and the Net~» dans lequel il propose que la ressource désormais rare et donc objet de la concurrence économique n'est pas l'information~--- abondante et reproductible à coût marginal nul~--- mais l'attention de l'utilisateur. La formulation reste prospective en 1997~: aucun modèle commercial ne l'opérationnalise au moment où elle est écrite. Mais elle nomme avec précision ce que les acteurs de l'industrie publicitaire et des moteurs de recherche cherchent sans le formuler~--- un dispositif qui capture et valorise l'attention plutôt que l'information. La résolution arrivera trois ans plus tard avec l'invention par Google d'un dispositif qui ne se contente pas de capturer l'attention mais l'inscrit dans un marché de prédictions comportementales~: mécanisme dont la section~\ref{sec:contemporain} examinera la portée et les conséquences matérielles.

\emergence{Le cycle Perez montre que le financement de l'infrastructure informationnelle obéit aux mêmes dynamiques que celui des infrastructures physiques antérieures. Le capital financier afflue vers une technologie dont la valeur perçue dépasse la valeur d'usage immédiate, crée une bulle, s'effondre, et laisse derrière lui une infrastructure que la phase de déploiement exploite. Ce résultat confirme que l'ICT n'est pas immatérielle : elle exige des investissements en capital fixe (câbles, serveurs, antennes, satellites) dont la logique financière est celle de toute infrastructure lourde.}
 
 
% -----------------------------------------------------------------------------
\subsection{Le téléphone intelligent et la convergence des fonctions}\label{subsec:smartphone}
% -----------------------------------------------------------------------------
 
\dimmark{0-TER}
Les données de \textcite{hilbert_worlds_2011} montrent qu'en 2007, la part numérique atteint 99,9~\% pour la télécommunication, 94~\% pour le stockage, mais seulement 25~\% pour le broadcasting. La convergence est donc achevée pour deux fonctions sur trois, et en cours pour la dernière. C'est dans cette configuration que l'iPhone d'Apple, lancé en juin 2007, rend la convergence visible dans un objet unique.
 
Un seul appareil réunit un processeur ARM (COMPUTE), une mémoire flash (STORE), un réseau cellulaire et Wi-Fi (TRANSMIT bidirectionnel), un écran tactile (interface), un récepteur GPS, un accéléromètre et une caméra (capteurs). Le routage IP est du calcul ; le cache du navigateur est du stockage ; le protocole HTTP est du contrôle. Les trois fonctions sont techniquement inséparables dans l'architecture de l'appareil. Ce que la tripartition TRANSMIT/STORE/COMPUTE distingue analytiquement, le smartphone fusionne physiquement.
 
\dimmark{D2}
Le smartphone n'invente pas la convergence ; il la rend portable et permanente. L'appareil est allumé en continu, connecté en permanence, et génère un flux constant de données (géolocalisation, consultation de pages, échanges de messages) qui transite par une infrastructure de réseau elle aussi permanente. Le résultat structurel est un déplacement du COMPUTE et du STORE de l'appareil vers le serveur distant : ce que le PC des années 1980 avait distribué chez l'utilisateur, le smartphone le recentralise dans le datacenter. L'architecture \emph{client léger / serveur lourd} qui en résulte est le point de départ de la section suivante (§\ref{sec:datacenter}).
 
Les données de la Banque mondiale, issues de l'UIT, confirment que le basculement se situe autour de 2005 : les abonnements mobiles dépassent les lignes fixes cette année-là (33,9 pour cent habitants contre 19,1), et le nombre de lignes fixes commence à décliner, passant de 19,1 en 2005 à 10,6 en 2023. Le mobile atteint 109,8 abonnements pour cent habitants en 2023, soit plus d'un abonnement par personne sur Terre. Le téléphone fixe, technologie dominante pendant un siècle (§\ref{subsec:telephone}), devient en vingt ans un résidu.
 
\emergence{La convergence T/S/C sur un substrat numérique unique est le résultat structurel de la période 1975--2007. Elle produit trois conséquences pour la thèse. Premièrement, la consommation d'énergie de l'information cesse d'être ponctuelle pour devenir permanente : le réseau est allumé en continu, les serveurs fonctionnent en permanence, le smartphone est toujours connecté. Deuxièmement, les demandes énergétiques des trois fonctions se composent au lieu de rester séparées : le TRANSMIT numérique contient du COMPUTE et du STORE, le COMPUTE contient du STORE et du TRANSMIT. Troisièmement, la mesure économique du secteur (le capital intangible de \textcite{corrado_intangible_2009}, le PIB du secteur ICT de l'OCDE, les indices de prix hédoniques du BEA) reste séparée de la mesure énergétique (\textcite{koomey_implications_2011}, \textcite{malmodin_energy_2018}, IEA). Personne ne croise les deux. C'est dans cet écart que s'ouvre la section suivante.}
 
\questionch{Comment formaliser le couplage entre capacités informationnelles (mesurées par \textcite{hilbert_worlds_2011} en bits et MIPS) et consommation d'énergie (mesurée par \textcite{koomey_implications_2011} en kWh par calcul) ? Le croisement des deux séries donnerait une estimation de la consommation d'énergie totale du traitement de l'information mondiale. Ce croisement n'a jamais été publié. Le chapitre~3 de cette thèse le tentera.}


\begin{table}[htbp]
\caption{Dimensions économiques de la convergence (§\ref{sec:convergence}).}\label{tab:convergence-dimensions}
\centering\small
\renewcommand{\arraystretch}{1.15}
\begin{tabularx}{\textwidth}{@{}cXc@{}}
\toprule
\textbf{Dim.} & \textbf{Ce que §\ref{sec:convergence} établit} & \textbf{Niveau} \\
\midrule
D1 & L'information numérique est non rivale~; le circuit qui la transporte est rival. Le paradoxe de la productivité (Solow 1987, David 1990) reflète un décalage de substrat, pas seulement d'organisation. & + \\
D2 & Plateformes \emph{winner-take-all} (Microsoft, Google)~; architecture ouverte IBM-PC $\rightarrow$ clonage $\rightarrow$ concentration sur l'OS et les applications. & + \\
D3 & Capital intangible $\geq$ tangible dès 1995 (\textcite{corrado_intangible_2009})~; \emph{unbundling} IBM (1969) comme origine institutionnelle~; découplage comptable entre valeur du logiciel et coût énergétique de son usage. & ★ \\
D5 & Paradoxe de la productivité (Solow 1987, David 1990)~; le gain de productivité n'apparaît qu'après recomposition organisationnelle complète. & ★ \\
D6 & Le Web (1993) crée un signal de prix instantané à l'échelle mondiale~; le moteur de recherche (Google, 1998) devient l'interface d'accès à l'information. & + \\
D7 & Standards ouverts (TCP/IP 1983, HTTP 1993) comme condition de la convergence~; modularité au sens de \textcite{langlois_modularity_2003}. & ★ \\
D8 & Cycle Perez~: dotcom = railway mania~; l'infrastructure (câbles, serveurs) survit au krach et sert la phase de déploiement. Secteur financier comme premier déployeur massif (1970--1990). & ★ \\
D10 & Open source (GNU 1984, Linux 1991) vs licensing massif (Microsoft)~: deux trajectoires coexistantes~; culture \emph{Whole Earth Catalog} comme trace structurelle. Bifurcations écartées (refus AT\&T 1971, échec Inouye 1994). & + \\
\midrule
0-TER & STORE devient permanent (serveur web \emph{always-on})~; premier régime de consommation continue~; la numérisation du réseau (1983--2000) est un double coût d'infrastructure, pas une substitution. & ★ \\
0-GOV & Privatisation NSFNET (1995)~; \emph{Framework for Global Electronic Commerce} (Clinton 1997) consacre l'autorégulation~; trajectoires alternatives (Inouye 1994) écartées. & + \\
\bottomrule
\end{tabularx}
\renewcommand{\arraystretch}{1.0}
\end{table}


% =============================================================================
% §1.10 L'information qui consomme (2007–)
% =============================================================================
\clearpage
\section{L'information qui consomme~: convergence numérique et substrat énergétique (2007--)}\label{sec:contemporain}
% =============================================================================

\subsection{Le \emph{turning point} et l'invention du surplus comportemental (2001--2008)}\label{subsec:turning_point}

\dimmark{D5/D6}{+}%
La section précédente s'est achevée sur une tension irrésolue~: la phase \emph{frenzy} de la fin des années 1990 avait déployé une infrastructure et capté une base d'utilisateurs sans qu'un modèle commercial stable ne stabilise les flux financiers. \textcite{goldhaber_attention_1997} avait nommé le problème sans le résoudre. La résolution intervient entre 2001 et 2003 dans une firme alors périphérique à la bulle qui venait d'éclater~: Google découvre que les traces de comportement des utilisateurs~--- clics, requêtes, durées de session, abandons~--- ne sont pas seulement des données nécessaires à l'amélioration de son moteur de recherche, mais des actifs prédictifs valorisables sur un marché publicitaire qui se constitue par leur agrégation. La traduction opérationnelle est \emph{AdWords} (lancé en octobre 2000, repensé en 2002) puis \emph{AdSense} (2003)~: un mécanisme d'enchères sur l'attention des utilisateurs ciblés par leur historique comportemental. \textcite{zuboff_surveillance_2019} a proposé de cette séquence le terme de \emph{surplus comportemental}~: une catégorie d'actifs informationnels dérivés du comportement humain et capturables sans contrepartie marchande directe, dont l'agrégation forme la matière première d'un nouveau type de marché (\emph{behavioural futures markets}).

\dimmark{D5/D10}{+}%
L'année 2007--2008 marque la généralisation du mécanisme à l'échelle de l'écosystème. Facebook, fondé en 2004 et longtemps réticent à déployer la publicité par crainte de cannibaliser sa croissance utilisateur, adopte en 2007 un \emph{Ad Platform} comparable à celui de Google, puis publie en 2008 \emph{Facebook Beacon} qui tente d'étendre la capture aux comportements hors plateforme~--- une extension trop visible qui provoque un recul public mais inscrit le principe~\parencite{west_data_2019}. La même année, le lancement de l'iPhone en juin 2007 puis de l'App Store en juillet 2008 transforme l'équipement utilisateur en plateforme permanente de capture comportementale~: chaque application installée devient un capteur de données, chaque séquence d'interaction une enregistrement valorisable. \textcite{johnstone_deep_2026} identifient cette séquence 2001--2008 comme le \emph{turning point} de la trajectoire numérique au sens de \textcite{perez_technological_2002}, par lequel la phase \emph{frenzy} bascule vers la phase \emph{synergy}~: ce qui était dispersion spéculative devient configuration stable, ce qui était promesse devient infrastructure. La nuance que la présente analyse veut introduire concerne le contenu même de la stabilisation~: la résolution \emph{commerciale} de la tension fin 1990 (l'invention du modèle économique manquant) est aussi le moment où le couplage \emph{énergétique} cesse d'être ponctuel pour devenir continu à l'échelle planétaire, puisque la capture comportementale exige des serveurs en fonctionnement permanent. Les sections suivantes documentent cette continuité.


\subsection{L'informatique en nuage et le centre de données à très grande échelle}\label{subsec:cloud}

\dimmark{D2/D3}{+}%
Amazon Web Services ouvre ses premiers services en 2006, Google Compute Engine suit en 2012, Microsoft Azure s'impose dans les années suivantes~: en une décennie, le calcul cesse d'être un service local que l'entreprise produit pour elle-même et devient un flux acheté à des opérateurs d'infrastructure spécialisés. Le mouvement est structurellement analogue à celui qu'Edison et Insull avaient accompli entre 1882 et 1910~: remplacer des générateurs locaux dispersifs par une centrale mutualisée exploitant les économies d'échelle. L'entreprise qui maintenait sa propre salle informatique comme Edison avait ses lampes à arc propres cède la place à l'abonné du réseau~; le capital informatique fixe se transforme en coût variable d'exploitation. Le rendement d'échelle est ici triple~: le taux d'utilisation des serveurs monte de dix à soixante-dix pour cent par mutualisation des pic de charge~; le coût d'achat des composants baisse par l'effet du volume~; le coût d'opération (refroidissement, alimentation, maintenance) est optimisé à une échelle qu'aucune entreprise individuelle ne peut atteindre.

\dimmark{D8/0-TER}{+}%
La puissance électrique installée par ces centres de données devient un indicateur économique significatif. La consommation électrique mondiale des centres de données atteignait environ 415~TWh en 2024~--- soit 1,5~\% de la demande électrique mondiale~--- et l'Agence internationale de l'énergie projette une quasi-multiplication par deux d'ici 2030, à environ 945~TWh \parencite{iea_energy_ai_2025}. Pour situer l'ordre de grandeur~: c'est l'équivalent de la consommation électrique annuelle du Japon ajoutée en six ans. Dans certains états américains, la concentration géographique de l'infrastructure rend le coût sectoriel visible à l'échelle régionale~: en Virginie du Nord, les centres de données représentent plus de vingt-six pour cent de la consommation électrique régionale~; les files d'attente de raccordement au réseau de transport à haute tension atteignent sept ans pour les projets les plus importants.

\dimmark{D10}{+}%
Cette concentration crée un effet de dépendance de sentier du même type que le \emph{technological momentum} que \textcite{hughes_networks_1983} avait identifié pour le système électrique d'Insull~: des investissements massifs, des standards d'interconnexion, des institutions de régulation et des contrats d'approvisionnement énergétique s'accumulent autour d'une architecture technique et résistent collectivement au remplacement. Les quatre plus grands opérateurs~--- Amazon, Google, Microsoft, Meta~--- ont multiplié par trois leur consommation électrique entre 2018 et 2023 (de trente-cinq à plus de cent dix~TWh) et ont engagé plus de trois cent vingt milliards de dollars de dépenses d'investissement pour 2025 seul, un rythme qui n'a aucun précédent dans l'histoire du secteur informationnel.

\dimmark{0-TER}{+}%
L'empreinte agrégée du secteur déborde le seul poste électrique des centres de données. Sur la consommation énergétique opérationnelle, \textcite{bieser_review_2023} synthétisent les estimations disponibles et situent la contribution du secteur ICT aux émissions mondiales de gaz à effet de serre entre un et demi et quatre pour cent en~2023~; l'ampleur de l'intervalle, plus que la valeur centrale, est elle-même un résultat~: aucune méthodologie partagée ne stabilise le périmètre du secteur ni l'allocation de la consommation électrique entre équipements terminaux, réseaux et centres de données. Sur les flux de matière, \textcite{forti_global_2020} documentent que les déchets d'équipements électriques et électroniques (DEEE) atteignaient 53,6 millions de tonnes en 2019 et sont projetés à 74,7 millions de tonnes en 2030, dont les terminaux informatiques et de télécommunication constituent une part croissante. Sur l'angle climatique, \textcite{mazzucato_ai_2024} situent les émissions cumulées du \emph{cloud carbon} au même ordre de grandeur que celles du transport aérien commercial mondial, ordre de grandeur qui n'est pas une mesure consolidée mais une comparaison heuristique destinée à fixer l'échelle. Ces trois séries ne convergent pas vers un compte unifié~; elles convergent vers un constat~: l'empreinte matérielle du système informationnel contemporain est documentable, son ampleur est significative à l'échelle des grandes industries lourdes, et l'absence d'un cadre statistique partagé pour la mesurer est elle-même un fait analytique sur lequel le chapitre~\ref{ch:biophysique} aura à revenir.


\subsection{La photolithographie et le coût de la miniaturisation}\label{subsec:lithographie}

\dimmark{D4/D10}{+}%
Les sections précédentes ont documenté la consommation d'énergie du système informationnel du côté de l'usage~: centres de données, réseaux, terminaux. Elles n'ont pas examiné le substrat de \emph{production} qui rend possible l'existence même de ces équipements. Or la fabrication des semi-conducteurs, et en particulier la photolithographie qui en constitue l'étape la plus critique, est elle-même le lieu d'une intensification énergétique et matérielle considérable dont l'histoire éclaire la configuration actuelle du marché des puces.

La loi de Moore, formulée en 1965 comme observation empirique sur le doublement du nombre de transistors par circuit intégré, est devenue au cours des décennies suivantes une institution industrielle au sens fort du terme. \textcite{mody_long_2017} a montré que cette régularité n'est pas le résultat passif d'une trajectoire technique~; c'est une prophétie auto-réalisatrice, coordonnée par l'\emph{International Technology Roadmap for Semiconductors} (ITRS, 1991--2016, puis IRDS), autour de laquelle l'ensemble de l'industrie aligne ses cycles d'investissement, ses calendriers de R\&D et ses décisions d'équipement. La loi de Moore n'est pas une loi de la nature~; c'est un dispositif de coordination industrielle dont la force repose sur le fait que chaque acteur a intérêt à ce que les autres y croient.

\dimmark{0-TER}{*}%
La contrainte physique fondamentale de la miniaturisation est la résolution de la gravure optique. Les motifs du circuit sont projetés par un faisceau lumineux à travers un masque sur la surface du wafer de silicium~; la finesse des motifs gravables est bornée par la longueur d'onde de la lumière utilisée (limite de diffraction). Chaque génération de gravure a donc exigé de descendre en longueur d'onde~: lumière visible dans les années 1970, ultraviolet profond (DUV) à 248~nm puis 193~nm dans les années 1990 et 2000, et finalement ultraviolet extrême (EUV) à 13,5~nm à partir de la fin des années 2010. Le passage au DUV à immersion (193~nm avec lentille immergée dans l'eau) a permis de prolonger la trajectoire jusqu'au n{\oe}ud de 7~nm, au prix de techniques de \emph{multi-patterning} (gravure en plusieurs passes) qui multipliaient le nombre d'étapes de fabrication et donc le coût et la consommation d'énergie par couche gravée. Lorsque le multi-patterning est devenu économiquement prohibitif, seul le passage à l'EUV pouvait prolonger la trajectoire de miniaturisation.

\dimmark{D4/D8}{*}%
Ce passage a exigé trente ans de développement (1994--2024) et des investissements cumulés de l'ordre de plusieurs dizaines de milliards de dollars. L'industrialisation de l'EUV a été lancée en 1994 par une coalition de fabricants de semi-conducteurs, parmi lesquels Intel, TSMC et Samsung, associés à ASML, entreprise néerlandaise issue en 1984 d'une coentreprise entre ASM International et Philips. La source lumineuse EUV repose sur un mécanisme sans équivalent dans les générations précédentes~: un laser CO$_2$ de haute puissance vaporise des gouttelettes d'étain en plasma, qui émet un rayonnement à 13,5~nm~; ce rayonnement est collecté et focalisé par un système de miroirs multicouches (fabriqués par Carl Zeiss SMT en Allemagne) dont la surface doit être polie à une précision inférieure à l'atome. Le rendement de conversion énergétique de ce processus est de l'ordre de quatre à six pour cent~: plus de quatre-vingt-dix pour cent de l'énergie injectée est dissipée en chaleur\footnote{Les estimations de rendement varient selon les sources~; le chiffre de quatre à six pour cent correspond aux données publiées par les analystes de l'industrie pour les systèmes NXE~3400 en production vers 2020--2023. La génération suivante (High-NA, EXE~5000) utilise une optique à ouverture numérique de 0,55 au lieu de 0,33, ce qui améliore la résolution mais ne modifie pas fondamentalement le rendement de la source.}. Le refroidissement de la machine et de son environnement représente à lui seul jusqu'à trente pour cent de la consommation totale de l'installation.

\dimmark{D4}{*}%
Le résultat industriel de cette trajectoire est un monopole mondial sans équivalent dans l'histoire des technologies de production. ASML est aujourd'hui le seul fabricant de machines de lithographie EUV~; ses concurrents japonais Nikon et Canon, qui partageaient le marché DUV dans les années 1990, n'ont pas accompli la transition. La machine High-NA EUV (EXE~5200, livrée à partir de 2024) contient plus de cent mille composants, pèse l'équivalent de deux avions de ligne, et coûte environ 350~millions d'euros l'unité. Ses trois principaux clients~--- TSMC, Samsung et Intel~--- sont aussi les entreprises qui ont co-financé son développement par des prises de participation directes au capital d'ASML et par des engagements d'achat anticipé, reproduisant un schéma de financement consortial que l'industrie des semi-conducteurs avait inauguré avec SEMATECH (consortium public-privé américain créé en 1987) et que l'institut IMEC (Louvain, fondé en 1984) a prolongé en Europe. Ce mode de financement, dans lequel des concurrents sur le marché aval co-investissent dans leur fournisseur commun en amont, n'a pas d'équivalent dans les autres GPT documentées par ce chapitre~; il constitue une forme de coordination par la \emph{roadmap} qui se distingue tant de la coordination par le marché que de la coordination par l'État.

\dimmark{D9}{*}%
La division géographique du travail qui résulte de cette trajectoire est elle-même un fait historique sans précédent dans l'industrie informationnelle. Le design des puces les plus avancées est concentré aux États-Unis (Apple, Nvidia, Qualcomm, AMD)~; la fabrication est dominée par TSMC à Taïwan et Samsung en Corée du Sud~; l'équipement critique de lithographie est produit aux Pays-Bas (ASML) avec des optiques allemandes (Carl Zeiss SMT) et des lasers allemands (Trumpf). Aucun pays ne maîtrise la chaîne complète. Cette configuration n'est pas une nécessité technique~; c'est le résultat contingent de quarante ans de décisions d'investissement, de politiques industrielles et de spécialisations cumulatives. Elle structure la géopolitique contemporaine des semi-conducteurs~: les restrictions américaines à l'exportation de machines EUV vers la Chine (à partir de 2019 pour les machines les plus avancées), le \emph{CHIPS and Science Act} américain (2022), le \emph{European Chips Act} (2023) et le plan \emph{Made in China 2025} sont autant de réponses à une vulnérabilité que cette division du travail a créée. Le parallèle avec la domination britannique sur la gutta-percha et les câbles sous-marins au XIXe~siècle (§\ref{subsec:materialite}) est structurel, non métaphorique~: dans les deux cas, le contrôle d'un input matériel critique confère un pouvoir de veto sur l'infrastructure informationnelle mondiale. Mais les conditions de production du monopole diffèrent~: le monopole britannique reposait sur l'accès géographique à une matière première extractive~; le monopole d'ASML repose sur l'accumulation de compétences et de coûts fixes irrécouvrables dans un processus de fabrication d'une complexité sans équivalent.

\dimmark{0-TER}{*}%
La conséquence énergétique de cette trajectoire est la plus directement pertinente pour l'argument de la thèse. Chaque machine EUV en fonctionnement consomme de l'ordre d'un mégawatt, soit environ dix fois la consommation d'une machine DUV de génération précédente. Une usine de fabrication de semi-conducteurs avancés (un \emph{fab} aux n{\oe}uds 5~nm ou 3~nm) opère plusieurs dizaines de ces machines simultanément, vingt-quatre heures sur vingt-quatre. La consommation électrique de TSMC représentait en 2023 environ six pour cent de la consommation totale d'électricité de Taïwan, une proportion projetée à 12,5~pour cent en 2025\footnote{Estimation citée par Bloomberg en août 2022 et reprise dans la littérature de suivi sectoriel. La proportion exacte dépend de la trajectoire de déploiement des n{\oe}uds avancés et de l'évolution de la demande électrique totale de Taïwan.}. La poursuite de la miniaturisation, loin de réduire l'empreinte matérielle de l'ICT, déplace l'intensification du produit vers le procédé~: le coût énergétique unitaire de chaque transistor gravé continue de baisser, mais le coût énergétique du système de production qui le rend possible augmente à chaque génération technologique. C'est une externalisation de la contrainte matérielle, du même type que celle que le chapitre a documentée pour le passage du cuivre à la fibre optique (§\ref{subsec:www})~: la contrainte ne disparaît pas~; elle se déplace vers un niveau d'infrastructure moins visible.

\emergence{La trajectoire de la photolithographie révèle que la loi de Moore n'est pas une loi naturelle de dématérialisation~; c'est un dispositif de coordination industrielle dont la prolongation exige une intensification croissante du substrat énergétique et matériel de production. Le monopole d'ASML, la division géographique du travail et la consommation électrique des fabs avancés sont les résultats historiques de cette trajectoire. Pour la thèse, le cas établit que l'intensification énergétique du système informationnel ne se joue pas seulement côté usage (centres de données, terminaux, réseaux) mais aussi côté production (fabrication des puces elles-mêmes), et que les deux dynamiques se composent.}

\footnote{Pour le récit géopolitique de la concentration du secteur, voir \textcite{miller_chip_2022}. Les données techniques sur les machines EUV proviennent des rapports annuels ASML (2020--2024) et des présentations investisseurs. Sur le financement consortial SEMATECH et ses effets sur la structure de l'industrie américaine, voir \textcite{irwin_learning_1994} et \textcite{flamm_targeting_1987}.}


\subsection{Le stockage numérique et la mémoire permanente}\label{subsec:store}

\dimmark{0-TER}{*}%
La section~\ref{subsec:phonographe} avait établi que le régime pré-numérique du stockage est passif~: la cire, le disque, le film fixent l'information sans énergie, et cette information peut être relue des décennies plus tard sans qu'un seul watt ait été consommé entre les deux actes. Ce régime prend fin avec la mémoire à accès dynamique~: la DRAM (Intel~1103, 1970) est la première mémoire dont les cellules doivent être rafraîchies plusieurs centaines de fois par seconde pour maintenir leur contenu~; sans alimentation continue, l'information disparaît. La DRAM n'est pas une amélioration du stockage passif~: c'est une rupture de régime, un changement dans la nature du couplage entre information et énergie. La première s'applique dès lors à la mémoire de travail des ordinateurs. L'étape suivante est la généralisation de ce régime au stockage persistant lui-même.

\dimmark{0-TER}{*}%
Le Web (Berners-Lee, 1991) est l'événement qui accomplit cette généralisation. Un document stocké sur un serveur web n'est pas seulement un fichier~: c'est un fichier auquel n'importe qui peut accéder à n'importe quel moment, depuis n'importe où, sans préavis. Pour que cet accès soit possible, le serveur doit être allumé en permanence. La consommation d'énergie du serveur ne dépend pas du nombre de lectures~: un fichier consulté une fois par an et un fichier consulté un million de fois par jour consomment la même énergie sur le serveur. Ce qui a changé n'est pas la nature de l'information stockée~: c'est l'exigence de disponibilité permanente qui lui est associée, et cette exigence impose un flux d'énergie continu indépendamment de l'usage. C'est précisément le mécanisme identifié dans la section~\ref{sec:convergence}~: STORE couplé à TRANSMIT sur infrastructure \emph{always-on} produit une consommation continue, là où tous les substrats précédents (papier, carton, vinyle, film) ne consommaient qu'à la fabrication et à la lecture.

\dimmark{0-TER}{*}%
\textcite{hilbert_worlds_2011} mesurent la montée en puissance de ce régime~: en 2002, pour la première fois, la capacité de stockage numérique mondiale dépasse la capacité analogique~; en 2007, quatre-vingt-quatorze pour cent de l'information mondiale est numérique en termes de capacité de stockage, et quatre-vingt-dix-neuf virgule neuf pour cent des télécommunications sont numérisées\footnote{Le chiffre de 94~\% concerne le stockage~; le 99,9~\% concerne les télécommunications. Ces deux grandeurs ne doivent pas être confondues. L'étude de \textcite{hilbert_worlds_2011} couvre la période 1986--2007 et ne fournit pas de données post-2007.}. Le coût énergétique de ce basculement est invisible dans leurs données~: Hilbert et López mesurent les capacités fonctionnelles (bits par seconde de transmission, bits de stockage, MIPS de calcul) mais pas les kilowattheures associés. L'angle mort est structurel~: leur cadre analytique, comme celui de Shannon (1948) dont il procède, traite l'information indépendamment de son substrat énergétique.


\subsection{L'intelligence artificielle générative et les lois d'échelle}\label{subsec:genai}

\dimmark{0-TER}{*}%
Le basculement de régime de la couche compute s'opère avec la diffusion des lois d'échelle (\emph{scaling laws}) dans l'entraînement des grands modèles de langage. \textcite{kaplan_scaling_2020} établissent empiriquement que la performance prédictive d'un modèle de langage évolue comme une loi de puissance du compute investi, du nombre de paramètres et de la taille du corpus d'entraînement. \textcite{hoffmann_training_2022} précisent les rapports optimaux entre ces trois variables (les «~lois Chinchilla~\fg). Ce que ces résultats impliquent pour l'économie du système informatique est considérable~: améliorer la frontière des performances nécessite structurellement \emph{plus} de compute, pas moins. Il n'existe pas de substitut organisationnel, algorithmique ou architectural qui permettrait d'atteindre le même niveau de performance sans augmenter le compute~--- du moins dans le paradigme Transformer qui domine l'industrie depuis 2017. Le rapport improvement/énergie, que tous les cas précédents avaient vu évoluer dans le même sens (plus de performance pour la même énergie ou moins), s'inverse pour cette couche spécifique.

\dimmark{0-TER}{*}%
La traduction en puissance de rack est documentable. Les serveurs équipés de GPU Ampere (2020) consommaient environ treize kilowatts par rack~; les serveurs Blackwell (2024) en consomment environ cent trente~; les feuilles de route prévoient un mégawatt par rack pour les générations Rubin (2026 et au-delà)\footnote{Ces chiffres sont issus des rapports techniques et des annonces constructeurs (NVIDIA). \textcite{sun_chip_2024} et \textcite{roussilhe_phase_2025} les documentent dans le contexte de la rupture de la loi de Dennard.}. La traduction en coût énergétique par usage est tout aussi significative~: une requête ChatGPT exige environ dix fois plus d'énergie qu'une recherche Google de génération précédente. Pour la première fois dans l'histoire traversée par ce chapitre, l'énergie est un coût marginal de production du service informationnel~: chaque inférence a un coût électrique directement proportionnel à sa longueur et à la taille du modèle qui la produit.

\dimmark{D2/D8}{+}%
La stratification empirique crée une asymétrie analytique essentielle. Les couches réseau et stockage classique restent en régime Koomey~: \textcite{mytton_network_2024} documentent que la consommation énergétique des réseaux n'est pas proportionnelle au volume de données transmises~; chez Telefónica, le trafic a augmenté de cinq cent soixante-cinq pour cent entre 2015 et 2021 tandis que la consommation baissait de sept virgule deux pour cent. Le coût énergétique par bit transmis continue de décroître d'environ trente pour cent par an. Il n'y a pas de rupture dans ces couches. La rupture est localisée dans le compute IA-intensif~--- et c'est cette localisation qui la rend identifiable, mesurable, et modélisable dans le cadre IMACLIM. Parce que la couche qui a basculé est celle qui concentre l'essentiel de la croissance de l'investissement~--- plus de sept cent quinze milliards de dollars de dépenses d'investissement annoncées par les grands opérateurs pour 2026~--- elle tire l'ensemble du système malgré sa localisation dans un segment technique spécifique\footnote{Pour une analogie structurelle~: dans les années 1900, l'automobile représentait un segment étroit de l'économie de la mobilité, dominée par le cheval et le rail. Parce qu'elle concentrait l'investissement et l'innovation, elle a reconfiguré l'ensemble du système de transport, de l'urbanisme et de l'infrastructure pétrolière. Le compute IA occupe une position structurellement analogue.}.

\dimmark{D8}{*}%
La traduction la plus spectaculaire de cette rupture localisée est l'émergence des contrats d'approvisionnement en énergie nucléaire. Microsoft signe un accord d'achat d'électricité sur vingt ans avec Constellation Energy pour le redémarrage de la centrale Three Mile Island (835~MW)~; Amazon contractualise 1,92~GW depuis la centrale de Susquehanna sur dix-sept ans~; Meta signe un accord de 1,100~MW depuis la centrale Clinton. Plus de dix gigawatts de capacité nucléaire ont été contractés par les grandes firmes numériques en 2024--2025 aux seuls États-Unis \parencite{iea_energy_ai_2025}. Pour la première fois dans l'histoire de l'industrie tech, des entreprises dont le métier est le traitement de l'information signent des contrats de \emph{production} d'énergie sur quinze à vingt ans~: les horizons d'investissement des deux secteurs se synchronisent, et le couplage financier-temporel entre compute et énergie acquiert une matérialité contractuelle qu'aucune phase précédente de l'histoire informationnelle n'avait produite.


\subsection{La blockchain et les cryptomonnaies}\label{subsec:blockchain}

\dimmark{0-TER}{*}%
Bitcoin (Nakamoto, 2009) est le cas extrême du couplage calcul-énergie~: son protocole de preuve de travail (\emph{proof-of-work}) exige que la sécurité du réseau soit achetée en kilowattheures. Le mécanisme est délibéré~: pour qu'une transaction soit validée, des machines doivent résoudre un problème cryptographique dont la difficulté est ajustée de façon à ce que la résolution dure environ dix minutes, quelle que soit la puissance de calcul disponible sur le réseau. Augmenter la puissance de calcul ne réduit pas la consommation d'énergie par transaction~: cela augmente la difficulté et maintient la consommation énergétique totale. L'indice de Cambridge sur la consommation électrique du Bitcoin (\emph{Cambridge Bitcoin Electricity Consumption Index}) estime la consommation annuelle du réseau Bitcoin à cent cinquante à deux cents~TWh~--- un ordre de grandeur comparable à la consommation d'un pays européen de taille moyenne. Ethereum a migré vers un protocole de preuve d'enjeu (\emph{proof-of-stake}) en 2022, réduisant sa consommation de plus de quatre-vingt-dix-neuf pour cent~: ce contraste entre les deux protocoles illustre que le couplage compute-énergie n'est pas une nécessité physique pour toute blockchain, mais un choix d'architecture de gouvernance.

\dimmark{D10/0-TER}{+}%
Le cas mérite d'être pris pour ce qu'il est. La preuve de travail ne consomme pas de l'électricité par accident d'ingénierie~; elle transforme délibérément une dépense computationnelle vérifiable en mécanisme de sécurité, c'est-à-dire qu'elle remplace une autorité centrale par un coût physique imposé aux participants. La dépense d'énergie n'est pas un résidu du protocole~: elle est la condition par laquelle le protocole rend coûteuse l'attaque, et c'est ce qui distingue le registre distribué d'un simple fichier répliqué. Mais cette propriété fait de Bitcoin un bon révélateur et un mauvais modèle général. Elle révèle qu'une architecture informationnelle peut convertir de l'énergie en confiance distribuée, c'est-à-dire acheter de la rareté dans un monde numérique où la copie est presque gratuite~; elle ne dit pas que tout système informationnel doit s'organiser ainsi. Le cas doit rester dans le chapitre, mais comme cas limite~: il montre une modalité du couplage information-énergie inscrit dans la règle de gouvernance elle-même, non la loi générale du numérique contemporain.


\subsection{L'IA agentique et l'automatisation cognitive}\label{subsec:agentique}

\dimmark{D5/D10}{+}%
La génération de texte, d'image ou de code ne clôt pas la trajectoire de l'intelligence artificielle contemporaine~; elle en déplace le centre de gravité. Tant que le modèle produit une réponse à une requête, le calcul reste attaché à un acte humain discret~: l'utilisateur demande, la machine infère, le résultat revient dans l'interface. Avec ce qu'on appelle l'IA agentique, le calcul change de statut. Il ne sert plus seulement à produire une sortie informationnelle~; il sert à enchaîner des opérations, à choisir des outils, à interroger des bases, à écrire du code, à vérifier partiellement ses propres résultats, et à relancer le processus tant que l'objectif n'est pas atteint. Le modèle devient alors moins un générateur qu'un opérateur.

Ce déplacement ne doit pas être décrit comme une autonomie au sens fort. Les agents contemporains restent fragiles, dépendants de leur environnement logiciel, exposés aux erreurs d'exécution, aux hallucinations et aux objectifs mal spécifiés\footnote{La littérature récente sur les agents fondés sur les grands modèles de langage reste instable~: elle documente des progrès dans l'usage d'outils, la planification et la décomposition de tâches, mais aussi des taux d'échec élevés dès que les environnements deviennent longs, ouverts ou faiblement spécifiés. Le point retenu ici n'est donc pas l'autonomie des agents, qui demeure une promesse technique, mais la modification de l'unité économique pertinente~: une tâche humaine peut être transformée en séquence d'appels de modèles, d'outils et de vérifications.}. Mais la rupture n'a pas besoin d'être parfaite pour avoir une portée économique~: il suffit que l'agent transforme une requête humaine en trajet computationnel composé de plusieurs inférences enchaînées pour que le rapport entre acte humain et budget de calcul cesse d'être de un pour un. Le rapprochement avec les techniques antérieures de l'information~--- le télégraphe qui n'avait pas supprimé tous les délais mais certains délais pertinents, la tabulatrice qui n'avait pas supprimé le travail administratif mais déplacé son goulot d'étranglement~--- ne porte pas la preuve~; il sert seulement à indiquer que la structure du déplacement n'est pas neuve.

\dimmark{D5/0-TER}{+}%
C'est ici que le changement devient analytique. Le numérique de plateforme restait borné par le temps humain~: on regardait une vidéo, on écrivait un message, on cherchait une information. Le numérique agentique relâche partiellement cette borne, car une tâche humaine peut devenir un processus computationnel composé d'opérations intermédiaires. La contrainte ne disparaît pas~; elle change de lieu. Elle passe du temps de l'utilisateur au budget de calcul, de l'attention humaine à la capacité d'orchestration, de la rareté cognitive à la rareté infrastructurelle. En termes contemporains, on dirait que la demande devient partiellement \emph{machine-to-machine}~; cette formulation reste conjecturale tant que les mesures empiriques manquent. Le chapitre~\ref{ch:structures} aura à déterminer si cette modification de l'unité économique pertinente constitue un effet de rebond d'une nature inédite (D5), ou simplement une croissance ordinaire de l'usage telle que la littérature contemporaine sur le rebond la modélise déjà.


\subsection{L'Internet des objets et l'\emph{edge computing}}\label{subsec:iot}

\dimmark{D9/0-TER}{+}%
Le déploiement des objets connectés à partir des années 2010~--- compteurs intelligents, capteurs industriels, dispositifs domestiques, terminaux logistiques~--- prolonge un mouvement déjà engagé par la téléphonie mobile~: l'informatique sort des centres de calcul pour se diffuser dans l'environnement matériel ordinaire. Le décompte précis du parc d'équipements reste un objet de communication industrielle plus que de mesure scientifique~; les rapports des équipementiers réseau (Cisco, Ericsson) avancent des chiffres de plusieurs dizaines de milliards d'objets connectés à la fin de la décennie 2020, mais leur fiabilité est limitée par la définition retenue du périmètre et par l'absence de méthodologie partagée\footnote{Les rapports industriels (Cisco \emph{Annual Internet Report}, Ericsson \emph{Mobility Report}, ITU \emph{Measuring the Information Society}) doivent être traités comme des documents de marché plus que comme des sources de mesure. Leur fonction est de coordonner les anticipations d'investissement du secteur, et les ordres de grandeur qu'ils annoncent doivent être tenus comme des indicateurs de tendance, non comme des résultats consolidés.}. Ce qui importe pour le présent chapitre n'est pas le nombre exact de capteurs mais la transformation qualitative qu'ils induisent~: le calcul migre vers la périphérie du réseau parce que la latence, la bande passante et la confidentialité ne permettent plus toujours d'acheminer chaque donnée jusqu'à un centre de données distant.

\dimmark{0-TER}{+}%
La promesse énergétique de cette migration vers l'\emph{edge} est ambiguë. Réduire le volume de données transmises diminue la consommation des réseaux longue distance~; calculer localement consomme moins par opération que d'attendre l'aller-retour vers le cloud~; chiffrer et préfiltrer à la source économise du stockage en aval. Mais la multiplication des nœuds de calcul périphériques crée elle-même une consommation distribuée, dont l'agrégation est moins visible que celle des centres de données mais n'est pas pour autant négligeable~: chaque équipement maintient une alimentation, une mémoire, une connectivité, et une part de calcul. Surtout, les usages les plus dépendants de l'IoT massif~--- véhicule autonome, ville intelligente, jumeau industriel~--- réagrègent les flux périphériques dans des modèles entraînés et exécutés en cloud, où la rupture identifiée pour la couche IA générative s'applique de plein droit. La consommation totale n'est donc pas réduite par l'\emph{edge}~; elle est redistribuée géographiquement entre le terminal, l'antenne, le nœud de bord et le centre de données. Le Ch.~\ref{ch:biophysique} aura à mesurer cette redistribution.


\subsection{L'IA physique, la robotique et les véhicules autonomes}\label{subsec:physical_ai}

\dimmark{0-TER}{+}%
Une frontière supplémentaire se déplace lorsque le calcul sort de l'écran et entre dans l'action physique. Les véhicules autonomes (Waymo en exploitation commerciale limitée à San Francisco et Phoenix, Tesla Full Self-Driving en bêta sur flotte client), les robots humanoïdes (Tesla Optimus, Figure, Apptronik), les robots mobiles industriels (Boston Dynamics, Agility Robotics) et les systèmes d'assistance physique constituent une famille hétérogène mais structurellement comparable~: ce sont des plateformes informationnelles dont la sortie n'est plus un fichier ou un signal, mais un mouvement dans un environnement matériel. Cette caractéristique modifie le profil de consommation. La machine embarque un budget de calcul significatif pour percevoir, planifier et agir en temps quasi réel~; elle dépend en parallèle d'un entraînement en cloud qui ingère des volumes massifs de données capturées sur le terrain.

\dimmark{0-TER}{+}%
Le statut de ce résultat doit rester prudent. La plupart de ces plateformes sont en déploiement limité ou en démonstration~: les promesses industrielles sont nombreuses, les résultats publiés sont fragmentés, et l'écart entre prototype et industrialisation est documenté de longue date dans la robotique. Ce que l'on peut tenir, en revanche, est la structure du couplage énergétique. L'IA physique combine trois sources de consommation~: l'inférence embarquée, dont la puissance dépend du compromis entre autonomie d'action et coût du composant~; l'entraînement en cloud, qui relève du régime décrit à la section~\ref{subsec:genai}~; et l'entraînement en simulation, où des moteurs physiques distribués (NVIDIA Omniverse et plateformes équivalentes) génèrent des environnements virtuels destinés à compléter les données réelles. Cette troisième source est peu mesurée mais non négligeable~: simuler le monde physique pour entraîner des machines physiques transfère une partie de l'apprentissage de l'extraction de données vers la génération de scénarios, déplaçant ainsi le coût énergétique sans nécessairement le réduire.


\subsection{Les jumeaux numériques et la simulation à grande échelle}\label{subsec:digital_twins}

\dimmark{0-TER}{+}%
Le jumeau numérique désigne une copie informationnelle d'un système physique, tenue synchronisée avec lui par un flux continu de données issues de capteurs et par un modèle qui en simule le comportement. Le concept a été formulé au début des années 2000 dans le domaine de la gestion du cycle de vie produit, puis étendu à des objets plus vastes~: usines (Siemens Xcelerator), réseaux électriques, infrastructures urbaines, systèmes climatiques (NVIDIA Earth-2, programme \emph{Destination Earth} de l'Union européenne). Son intérêt analytique pour le présent chapitre tient moins à la technologie elle-même qu'à son régime de fonctionnement~: à la différence d'un modèle numérique classique, qu'on exécute ponctuellement pour répondre à une question, le jumeau numérique est censé tourner en continu pour maintenir sa correspondance avec le réel.

\dimmark{0-TER}{+}%
Cette continuité a une conséquence énergétique de principe~: un jumeau opérationnel consomme du calcul même lorsque le système physique qu'il représente n'évolue pas, parce que la synchronisation impose un cycle permanent d'ingestion, de simulation et de vérification. Pour les jumeaux d'objets industriels limités, l'ordre de grandeur reste maîtrisable~; pour les jumeaux climatiques ou terrestres à haute résolution, il devient un sujet en soi, puisque la simulation d'un système physique à grande échelle pour le suivre en temps quasi réel mobilise des infrastructures de calcul de classe supercalculateur. Le déploiement de ces objets est encore largement industriel et expérimental, et les mesures consolidées de leur empreinte sont rares~; le présent chapitre ne peut donc qu'enregistrer la trajectoire, pas la quantifier. Le Ch.~\ref{ch:biophysique} aura à déterminer si les jumeaux numériques constituent une catégorie d'usage suffisamment volumineuse pour modifier la trajectoire agrégée de consommation, ou s'ils restent un cas particulier dont la portée est plus symbolique qu'énergétique.


\subsection{Les réseaux 5G/6G et l'infrastructure de transmission}\label{subsec:5g}

\dimmark{D9/0-TER}{+}%
La cinquième génération de réseaux mobiles, déployée à partir de 2019, se distingue de la quatrième par une densification des stations de base, l'usage de bandes de fréquence millimétriques (entre 24 et 39~GHz) et une intégration plus poussée des fonctions de calcul dans le réseau lui-même. Le coût énergétique de ce déploiement est souvent présenté comme une rupture~: une station de base 5G consomme deux à trois fois plus qu'une station 4G équivalente. La formulation est exacte mais incomplète, et c'est précisément l'incomplétude qui doit être tenue~: la consommation par station ne dit rien de la consommation par bit transmis. Les bandes millimétriques sont plus directionnelles et requièrent des cellules plus petites, mais elles transportent aussi un volume de données beaucoup plus important par unité de spectre. La conséquence agrégée n'est donc pas mécaniquement une augmentation de la consommation~: elle dépend du rapport entre la croissance des stations et la croissance des données transportées.

\dimmark{0-TER}{+}%
C'est exactement le mécanisme que \textcite{mytton_network_2024} ont identifié pour les générations précédentes~: la consommation énergétique des réseaux n'est pas proportionnelle au trafic, parce que les améliorations d'efficacité par bit absorbent une part substantielle de la croissance des usages. Tant que ce mécanisme tient, les réseaux mobiles continuent d'opérer dans le régime Koomey décrit à la section~\ref{subsec:genai}, et leur empreinte énergétique reste un sujet de gestion plutôt qu'un sujet de rupture. Les questions ouvertes pour le Ch.~\ref{ch:biophysique} portent sur les configurations qui pourraient le faire céder~: la densification radicale anticipée pour la 6G, l'intégration calcul-réseau qui ferait remonter une part du compute dans l'infrastructure de transmission, ou le déploiement parallèle de constellations satellitaires en orbite basse (Starlink, Kuiper) dont le coût énergétique par bit reste mal documenté et qui ajoutent une couche d'infrastructure dont la matérialité~--- lancements, durée de vie courte, remplacement permanent~--- est d'un autre ordre que celle des antennes terrestres.


\subsection{La fenêtre 2008 et les trois cadres régulatoires post-crise}\label{subsec:fenetre_2008}

\dimmark{D8/0-GOV}{+}%
La crise financière de 2008 mérite d'être nommée dans l'histoire de la trajectoire numérique non pour ses causes proximales, qui relèvent de l'analyse macro-financière, mais pour ce qu'elle a ouvert et refermé à propos du couplage ICT-économie. La crise survient au moment exact où la phase \emph{synergy} commence à se stabiliser (§\ref{subsec:turning_point})~: le \emph{cloud computing} démarre commercialement, le modèle \emph{surplus comportemental} se généralise, l'iPhone et l'App Store sont lancés. Plusieurs propositions de politique publique formulées à ce moment auraient pu réorienter cette stabilisation. \textcite{hines_green_2008} et la \emph{Green New Deal Group} britannique proposent en juillet 2008 un \emph{Green New Deal} qui associe décarbonation rapide de l'économie, régulation financière renforcée et investissement public massif dans les infrastructures bas-carbone. La proposition ne sera retenue ni par les administrations américaine et britannique post-crise, ni par les institutions européennes, dont la réponse principale combine plans de relance limités, politique monétaire ultra-accommodante (\emph{quantitative easing} à partir de 2009) et maintien de l'autorégulation industrielle dans le numérique. \textcite{varoufakis_techno_2023} a proposé du résultat agrégé l'expression de \emph{techno-féodalisme}~: la combinaison de \emph{cheap credit} (taux quasi nuls pendant une décennie) et de captation rentale par les plateformes numériques produit une expansion sans précédent des centres de données et des firmes \emph{big tech}, dont la capitalisation atteint en 2020 un niveau que la «~révolution dot-com~» de 1999 n'avait jamais approché. Le lien causal direct entre crise 2008 et expansion data centers reste à quantifier~; ce que documente l'historiographie disponible est que la fenêtre d'une bifurcation régulatoire et infrastructurelle s'est ouverte en 2008--2009 et s'est refermée sans qu'aucune des options alternatives portées à ce moment ne soit retenue~\parencite{lundvall_techno_2017}.

\dimmark{0-GOV/D10}{+}%
C'est dans le cadre laissé par cette fermeture que se constituent, entre 2010 et 2024, les trois grands cadres régulatoires du numérique contemporain. \textcite{bradford_digital_2023} les caractérise par leur ancrage normatif différencié~: le modèle américain reste \emph{market-driven}, prolongeant l'autorégulation issue du \emph{Framework} de 1997 (§\ref{subsec:www})~; le modèle européen adopte une orientation \emph{rights-driven} dont les jalons sont le Règlement général sur la protection des données (RGPD, 2018), le \emph{Digital Services Act} et le \emph{Digital Markets Act} (2022), et l'\emph{AI Act} (2024)~; le modèle chinois est \emph{state-led} et combine contrôle politique du contenu, soutien stratégique à l'industrie nationale et régulation antimonopolistique à géométrie variable. \textcite{moldicz_techno_2022} documente la course à la \emph{souveraineté technologique} qui se déploie entre ces trois pôles après 2018~: contrôles américains à l'exportation des semi-conducteurs avancés, plan \emph{Made in China 2025}, plans européens \emph{Chips Act} et \emph{Cloud sovérain}. Ces trois cadres ne convergent pas vers une régulation mondiale unifiée du couplage ICT-énergie~; ils convergent vers une fragmentation régulatoire qui rend plus difficile, plutôt que plus facile, la formulation de politiques visant le couplage lui-même.


\subsection{La critique des \emph{twin transitions} et les niches alternatives}\label{subsec:critique_twin}

\dimmark{0-TER/D10}{+}%
Le discours dominant de la décennie 2020 sur la digitalisation est celui des \emph{twin transitions}~: l'idée portée par la Commission européenne et par les organisations internationales (OCDE, AIE) que la transition numérique et la transition écologique sont structurellement complémentaires et se renforcent mutuellement~\parencite{european_commission_twin_2022}. La littérature académique récente est venue questionner cette hypothèse de complémentarité. \textcite{santarius_digital_2022} ont proposé le concept de \emph{digital sufficiency} pour interroger l'hypothèse que la digitalisation contribuerait nécessairement à la décarbonation. \textcite{creutzig_digitalization_2022} dans la \emph{Annual Review of Environment and Resources} produisent un état de l'art qui documente la coexistence d'effets décarbonants (gestion fine de l'énergie, optimisation logistique) et d'effets rebond significatifs (croissance des usages absorbant les gains d'efficacité). \textcite{bennich_untangling_2025} et \textcite{mouthaan_systemic_2023} démontrent par revue systématique que l'hypothèse de complémentarité mécanique n'est pas étayée par les données disponibles, et que les conditions sous lesquelles digitalisation et décarbonation coopèrent effectivement relèvent de configurations politiques précises, non d'une dynamique structurelle automatique. \textcite{makitie_digital_2023} prolongent en montrant que la conception même d'une \emph{digital green transition} comme objet politique unifié pose des problèmes analytiques mal résolus. Le chapitre~\ref{ch:biophysique} aura à reprendre ce dossier en évaluant si, sous quelles conditions et dans quelles proportions, le déploiement numérique observé depuis 2007 peut effectivement contribuer à la décarbonation des autres secteurs ou, au contraire, pèse sur les budgets carbone disponibles.

\dimmark{D10}{+}%
La critique des \emph{twin transitions} ne dit pas qu'aucune alternative au déploiement numérique dominant n'existe~; elle dit que les alternatives existantes restent confinées à des niches dont la portée systémique reste à établir. Plusieurs configurations méritent d'être nommées, sans qu'on en surestime la portée ni qu'on les présente comme un programme de substitution. Le projet \emph{Fairphone} (lancé en 2013 aux Pays-Bas) propose un téléphone modulaire conçu pour la réparation et l'allongement de la durée d'usage~; \textcite{haucke_smartphone_2018} documente sa trajectoire économique limitée mais analytiquement instructive. Le projet \emph{Te Hiku Media} en Nouvelle-Zélande développe depuis 2018 des modèles de reconnaissance vocale en langue maori sous gouvernance communautaire stricte, et le collectif \emph{DAIR} (\emph{Distributed AI Research Institute}) fondé par Timnit Gebru porte une alternative au modèle \emph{scale-at-all-costs} de l'IA contemporaine~\parencite{hao_empire_2025}. La municipalité de Barcelone a développé entre 2015 et 2023 le projet \emph{DECODE} de souveraineté numérique citoyenne~\parencite{smith_smart_2021}, et la plateforme \emph{Energy Community Platform} de l'Union européenne (2024) soutient le déploiement de coopératives énergétiques utilisant des outils numériques sous gouvernance partagée~\parencite{energy_community_2024}. Ces configurations partagent un trait analytique~: elles découplent l'usage du numérique du couplage \emph{surplus comportemental + scaling + concentration} qui structure la trajectoire dominante depuis 2007. Le chapitre~\ref{ch:biophysique} et le chapitre~\ref{ch:experimentation} auront à déterminer si ces niches peuvent constituer des configurations de référence pour l'analyse de trajectoires alternatives, ou si leur marginalité même est un résultat de structure qu'aucune dynamique endogène du système actuel ne paraît susceptible de modifier.


\subsection{La limite thermodynamique~: le principe de Landauer}\label{subsec:landauer}

\dimmark{0-TER}{+}%
L'histoire que ce chapitre a suivie pourrait donner l'impression que l'information échappe toujours à ses supports. Un code remplace un signal optique~; un fil remplace un messager~; une carte perforée remplace une colonne de chiffres~; un transistor remplace une lampe~; un modèle statistique remplace une règle explicite. À chaque étape, le support change, et l'on pourrait croire que la trajectoire conduit vers une abstraction croissante. Le principe de Landauer rappelle que cette abstraction a un plancher.

\textcite{landauer_irreversibility_1961} établit que l'effacement irréversible d'un bit d'information entraîne une dissipation minimale d'énergie proportionnelle à $kT\ln 2$. La formule est minuscule à l'échelle d'un bit~; elle est décisive à l'échelle du principe. Elle ne dit pas que les ordinateurs contemporains sont proches de cette limite, ni que l'énergie des centres de données serait réductible à l'effacement logique. Elle dit quelque chose de plus profond pour l'histoire économique de l'information~: le calcul n'est pas seulement une opération formelle, car toute opération physiquement réalisée modifie un support matériel et dissipe de l'énergie lorsqu'elle est irréversible. \textcite{bennett_logical_1973,bennett_thermodynamics_1982} a ensuite montré qu'un calcul logiquement réversible peut contourner cette dissipation minimale, mais au prix d'une autre contrainte~: conserver l'information au lieu de l'effacer, c'est déplacer la charge vers la mémoire, la complexité et la gestion des états.

\dimmark{0-TER}{+}%
Ce détour thermodynamique ne doit pas être transformé en causalité historique directe\footnote{La littérature sur Landauer est traversée par une difficulté terminologique~: l'entropie informationnelle de Shannon, l'entropie thermodynamique et l'entropie algorithmique ne désignent pas la même grandeur, même si les analogies entre elles ont été fécondes. Bennett a précisément contribué à clarifier le paradoxe du démon de Maxwell en distinguant l'acquisition d'information de l'effacement de mémoire. Le chapitre n'utilise donc pas Landauer pour expliquer la consommation actuelle des centres de données~; il l'utilise pour borner conceptuellement l'idée d'un calcul immatériel.}. Landauer n'explique ni la loi de Moore, ni l'économie des semi-conducteurs, ni la croissance contemporaine des infrastructures de calcul. Il fournit une borne conceptuelle. Or cette borne suffit à rompre avec l'un des gestes fondateurs de la théorie de l'information~: la séparation entre information et énergie. Shannon avait raison de séparer le contenu sémantique du signal pour construire une théorie mathématique de la communication~; il n'en résultait pas que l'information réalisée puisse être séparée indéfiniment de son support. La fécondité de la séparation formelle a produit, par retour, une occultation matérielle.

Cette conclusion ne signifie pas que toute augmentation de calcul conduit mécaniquement à une augmentation proportionnelle d'énergie. L'histoire récente de l'informatique est aussi celle d'améliorations d'efficacité, d'architectures spécialisées et de gains matériels qui déplacent sans cesse le rapport entre opération logique et dissipation effective. Mais elle interdit de traiter le calcul comme une activité purement immatérielle. Le bit n'a pas de joule dans la théorie de Shannon~; le bit effacé en a un dans le monde de Landauer. Le Ch.~\ref{ch:structures} pourra alors reprendre la question sous une autre forme~: si le calcul possède un plancher physique, quelles structures économiques déterminent la vitesse à laquelle les sociétés s'en approchent, s'en éloignent, ou l'oublient~?


\begin{table}[htbp]
\caption{Dimensions économiques de la période contemporaine (§\ref{sec:contemporain}).}\label{tab:contemporain-dimensions}
\centering\small
\renewcommand{\arraystretch}{1.15}
\begin{tabularx}{\textwidth}{@{}cXc@{}}
\toprule
\textbf{Dim.} & \textbf{Ce que §\ref{sec:contemporain} établit} & \textbf{Niveau} \\
\midrule
D2 & Concentration cloud (AWS, Azure, GCP)~; monopole ASML sur la lithographie EUV~; modèle Insull reproduit à l'échelle mondiale. & ★ \\
D3 & Externalisation du compute vers le cloud~: capital fixe $\rightarrow$ coût variable~; découplage comptable entre valeur logicielle et coût énergétique d'usage. & + \\
D4 & Déplacement du travail vers l'orchestration et la supervision (IA agentique)~; division internationale du travail semi-conducteur (design US, fab Taïwan, équipement NL). & + \\
D5 & Surplus comportemental (\textcite{zuboff_surveillance_2019})~; rebond potentiel de l'IA agentique (tâche $\rightarrow$ séquence d'inférences)~; scaling laws = amélioration $\propto$ compute. & ★ \\
D6 & AdWords/AdSense (2002--2003)~: marché de l'attention opérationnalisé~; requête ChatGPT $\approx$ 10$\times$ énergie recherche Google. & + \\
D8 & Contrats nucléaires 15--20 ans (Microsoft/Three Mile Island, Amazon/Susquehanna, Meta/Clinton)~: horizons d'investissement ICT et énergie synchronisés. Capex 2025 $>$ 320~Mrd\$. & ★ \\
D9 & Division géographique du travail semi-conducteur (US/TW/KR/NL/DE)~; restrictions export EUV vers Chine~; CHIPS Act, European Chips Act, Made in China 2025. & ★ \\
D10 & Loi de Moore comme dispositif de coordination (ITRS/IRDS)~; scaling laws IA comme nouvelle feuille de route~; niches alternatives (Fairphone, Te Hiku, DECODE). & ★ \\
\midrule
0-TER & Rupture Dennard ($\sim$2005)~; datacenters 415~TWh (2024) $\rightarrow$ 945~TWh (2030)~; TDP GPU 700~W (2024) $\rightarrow$ 1\,400~W (2026)~; EUV $\approx$ 1~MW/machine~; TSMC = 6~\% élec. Taïwan~; Bitcoin $\approx$ 150--200~TWh/an~; Landauer = plancher thermodynamique. & ★ \\
0-GOV & Fenêtre 2008 refermée sans bifurcation~; trois cadres régulatoires (\textcite{bradford_digital_2023})~: US \emph{market-driven}, EU \emph{rights-driven}, CN \emph{state-led}~; critique des \emph{twin transitions}. & ★ \\
\bottomrule
\end{tabularx}
\renewcommand{\arraystretch}{1.0}
\end{table}


% =============================================================================


% =============================================================================
% §1.11 Des cas aux archétypes
% =============================================================================
\clearpage
\section{Des cas aux archétypes~: ce que l'histoire révèle et ce qu'elle ne peut pas expliquer}\label{sec:reduction}
% =============================================================================

Le parcours qui précède a traversé dix cas sur deux siècles, distribués entre les sections §\,\ref{sec:era0} (avant l'électricité) et §\,\ref{sec:contemporain} (systèmes informationnels contemporains). Du télégraphe optique de Chappe au datacenter d'IA générative, du calcul à la plume de Clairaut et Lalande à l'inférence par GPU à 1\,400~watts, de la gutta-percha du câble transatlantique au bouquet élémentaire des fonderies de semi-conducteurs, le matériau est hétérogène par nature et par époque. La méthode a été inductive~: chaque cas a été observé pour lui-même, les catégories analytiques ont émergé de la comparaison entre les cas, et les patterns provisoires qu'un cas faisait apparaître ont été modifiés ou infirmés par le suivant. Le moment est venu de faire le geste inverse~: non plus observer des cas, mais rassembler ce que les observations révèlent de commun, et nommer ce qu'elles ne parviennent pas à expliquer.

La répartition des cas dans la synthèse n'est pas uniforme. Les régularités qui apparaissent dans la première partie de ce bilan s'établissent sur le matériau du XIX\textsuperscript{e} et du début du XX\textsuperscript{e} siècle, où plusieurs cas comparables permettent l'induction. Les cas contemporains (§\,\ref{sec:computing}, §\,\ref{sec:convergence}, §\,\ref{sec:contemporain}) interviennent plus tard dans le bilan~: ce sont eux qui font apparaître l'\emph{anomalie} énergétique et les deux postulats qui cessent de tenir. Les premiers servent à identifier ce qui se répète~; les seconds, à identifier ce qui rompt. La double fonction est délibérée.

\subsection{Trois régularités et une anomalie}\label{subsec:regularites}

La première régularité est institutionnelle. Chaque réseau de communication observé dans ce chapitre converge vers une forme ou une autre de monopole~: sept trajectoires nationales pour le télégraphe (§\,\ref{sec:telegraph}), quatre pour le téléphone (§\,\ref{sec:telephone}), trois mécanismes successifs pour le wireless (§\,\ref{sec:wireless}), le monopole régulé d'Édison et d'Insull pour l'électrification (§\,\ref{sec:electrification}), la domination d'IBM à 85,7\,\% du parc mécanographique (§\,\ref{sec:mechanography}). Les conditions initiales varient radicalement~--- sécurité ferroviaire, souveraineté militaire, contrôle policier, modernisation autoritaire, concession coloniale~--- mais le résultat institutionnel est le même. Ce qui varie, c'est le mécanisme~: rendements de réseau pour le télégraphe, verrouillage du consommable pour la tabulatrice, intégration verticale et contrôle de l'interconnexion pour le téléphone, brevets et refus d'interopérabilité pour Marconi, rendements d'échelle dans la génération pour l'électricité. La forme apparente est récurrente~; les mécanismes qui la produisent ne le sont pas. La concentration des infrastructures informationnelles est donc une régularité morphologique sans mécanisme commun, et non un résultat causalement unifié~: ce constat oblige le Ch.~\ref{ch:structures} à la décomposer en dimensions séparées plutôt qu'à la traiter comme un mécanisme à part entière, dans les rendements croissants (D2), la nature du capital (D3), la coordination par les standards (D7) et l'allocation politique des ressources rares comme le spectre (D6).

L'envers de cette régularité institutionnelle est tout aussi instructif. Babbage occupe la place du cas qui ne devient pas infrastructure malgré la conjonction d'éléments qui auraient pu en faire un. Une demande de calcul vérifiable existe~: les tables du Nautical Almanac contiennent des erreurs identifiées, et Prony a démontré que la division du travail intellectuel pouvait produire des résultats reproductibles. Une capacité technique existe~: les plans de la Difference Engine sont défendables, ceux de l'Analytical Engine plus exigeants mais conceptuellement cohérents. Pourtant aucune des deux ne deviendra machine de production. Ce qui manque est identifiable~: une organisation capable d'absorber les coûts d'une fabrication mécanique de précision sans débouché commercial avéré, un véhicule financier adapté à un projet dont les bénéfices ne sont pas appropriables par le financeur, une promesse stabilisée auprès des utilisateurs administratifs qui auraient pu être les premiers clients. Babbage n'est pas un génie en avance sur son temps~; c'est un cas de désajustement entre projet technique, demande solvable, administration du financement et trajectoire d'usage\footnote{La formule \og en avance sur son temps\fg{} déplace le problème vers une métaphysique chronologique sans intérêt analytique. Le matériau historique suggère plutôt que chaque configuration d'émergence est singulière, et que ce qui manque à Babbage est identifiable au cas par cas.}. La leçon est plus générale~: les non-émergences ne sont pas des accidents extérieurs à la trajectoire technique~; elles rendent visible, en creux, la configuration d'éléments dont la rencontre permet à une capacité informationnelle de devenir infrastructure. Demande de coordination, support matériel stabilisé, organisation capable de porter l'usage, véhicule financier adapté, promesse socialement crédible~: cette configuration n'est pas une théorie causale (chacun de ces éléments dépend des autres) mais une grille heuristique que le matériau permet d'inférer.

Une seconde irrégularité interne au matériau force à préciser ce que la convergence morphologique cache. Calculatrice mécanique et tabulatrice mécanisent des opérations cognitives proches pour des organisations dont la fonction comptable, statistique ou administrative est comparable, et appartiennent à la même grande famille du calcul instrumenté. Pourtant l'une produit un marché fragmenté de plusieurs dizaines de fabricants, où Burroughs, Felt, Brunsviga, Comptometer, Monroe, Marchant et d'autres se partagent une clientèle qui possède ses machines, peut changer de fournisseur et opère hors d'un service vendeur~; l'autre produit une structure concentrée que la fusion CTR de 1911 puis l'IBM de Watson transforment en monopole effectif pendant un demi-siècle, fondé sur la location, le format propriétaire des cartes, la formation des opérateurs et le service de maintenance. L'écart interdit d'expliquer la structure de marché par la seule fonction technique. Il oblige à reconnaître que le mode juridique et économique d'appropriation du capital informationnel~--- vente ou location, possession ou abonnement, format ouvert ou propriétaire, machine séparée ou machine inséparable de ses consommables~--- peut être un déterminant de la structure de marché au même titre que les coûts fixes, les rendements croissants ou les externalités. Le mécanisme ne s'applique pas à toutes les configurations~: pour le câble sous-marin, l'ampleur du capital fixe et la spécificité des actifs ferment le choix, et l'on ne loue pas un câble transatlantique à des utilisateurs finaux. Mais lorsque la technique laisse plusieurs régimes d'usage techniquement disponibles, ce choix devient un déterminant économique indépendant des propriétés techniques, et c'est dans cette borne que le Ch.~\ref{ch:structures} aura à reprendre la dimension D3.

La seconde régularité concerne le travail. Dans chaque cas où une technique de l'information est introduite dans une organisation existante~--- la calculatrice mécanique dans les bureaux d'assurance (§\,\ref{sec:calculating}), la tabulatrice au \emph{Nautical Almanac} et au PLM (§\,\ref{sec:mechanography}), le commutateur automatique dans le réseau Bell (§\,\ref{sec:telephone})~---, le nombre de travailleurs se maintient ou augmente~; leur profil change~; le coût unitaire baisse par compression salariale plutôt que par suppression de postes. La mécanisation réorganise le travail au lieu de le remplacer~--- un résultat que \textcite{reinstaller_market_2001}, \textcite{davies_womans_1982}, \textcite{feigenbaum_answering_2024} et \textcite{daston_calculation_2018} documentent dans quatre registres différents (quantitatif, genre, emploi, organisation). Le seul cas qui produit une identification causale de ce mécanisme est celui de l'automatisation de la commutation téléphonique~: \textcite{feigenbaum_answering_2024}, exploitant la variation géographique du calendrier d'adoption dans près de trois mille villes américaines, montrent que les cohortes suivantes se réallouent vers des tâches nouvelles dans des industries qui n'employaient pas de secrétaires auparavant~--- du \og \emph{new work}\fg{} au sens de \textcite{autor_new_2024}. Mais ce mécanisme de réinstauration est conditionnel~: dans les villes frappées par la Grande Dépression, l'automatisation réduit l'emploi sans le compenser. La demande agrégée est une condition nécessaire pour que la restructuration l'emporte sur la suppression.

Cette régularité est le résultat le plus robuste du chapitre sur la dimension du travail. Mais elle ne dit pas tout. En relisant les cas traversés, trois observations apparaissent qui ne se réduisent pas à la question \og que fait la technique au travail~?\fg{} et qui suggèrent que le rapport entre les techniques de l'information et le travail est plus riche que ce que la formulation \og restructuration sans suppression\fg{} laisse entendre.

La première observation est que la technique informationnelle ne fonctionne qu'insérée dans un dispositif de travail intellectuel qu'elle ne peut pas remplacer. McCallum, sur l'\emph{Erie Railroad}, déploie six formes informationnelles dont le télégraphe n'est que la quatrième~; les cinq autres~--- le \emph{rule book}, la hiérarchie managériale, le reporting, la comptabilité analytique, la standardisation~--- sont du travail de conception codifié (§\,\ref{sec:telegraph}). Au \emph{Nautical Almanac}, Comrie estime que vingt à trente pour cent du temps du surintendant sont désormais consacrés à la conception des séquences dans lesquelles humains et machines doivent s'articuler~--- ce que \textcite{daston_calculation_2018} nomme l'\emph{analytical intelligence}. Bolle, au PLM, formule la même leçon dans une devise~: \og \emph{d'abord organiser, ensuite mécaniser}\fg{} (§\,\ref{sec:mechanography}). Dans ces trois cas, la machine est un composant d'un système dont le reste est du travail qualifié de conception~--- un travail que la mécanisation \emph{crée} au lieu de le réduire. Le chapitre~2 devra déterminer si ce coût de conception est un coût transitoire d'ajustement, comme le suggère \textcite{david_dynamo_1990} pour le parallèle dynamo/ordinateur, ou un coût permanent inhérent à toute mécanisation du traitement de l'information.

La seconde observation est que le stock de compétences existant ne subit pas passivement la technique~: il en conditionne la forme. L'arbitrage entre le système Morse (un fil, un opérateur formé au code) et le système Wheatstone (cinq fils, un opérateur non qualifié) n'est pas un choix purement technique~: c'est un arbitrage entre capital physique et capital humain, dans lequel la disponibilité de compétences détermine quel design s'impose (§\,\ref{sec:telegraph}). L'administration française du télégraphe fait construire par Bréguet un appareil électrique qui reproduit les mouvements du sémaphore optique pour ne pas avoir à requalifier les opérateurs~: la technologie adoptée est déformée par le stock de compétences en place. Et le métier Jacquard (§\,\ref{sec:era0}), en substituant du capital informationnel (les cartes perforées) au travail qualifié des tireurs de lacs, provoqua la résistance la plus célèbre de l'histoire des techniques~: les tisserands luddites (1811--1816) ne détruisaient pas des machines à vapeur mais des machines à information, des métiers dont les cartes perforées encodaient le savoir-faire qu'ils avaient passé des années à acquérir. Ce n'est pas la technique qui transforme le travail~: c'est le travail~--- sous forme de compétences disponibles ou indisponibles, défendues ou dévaluées~--- qui sélectionne, déforme ou rejette la technique. Le chapitre~2 devra instruire cette direction causale inverse avec les outils appropriés et déterminer si le stock de capital humain spécifique peut constituer une contrainte structurelle sur l'émergence d'un régime technologique~--- et pas seulement un facteur d'ajustement.

La troisième observation est que certaines tâches résistent à l'automatisation non par défaut technique mais parce qu'elles portent un savoir organisationnel implicite dont l'extraction exige la refonte du système entier. L'automatisation de la commutation téléphonique a pris soixante ans, de Strowger (1891) à la généralisation~; \textcite{feigenbaum_organizational_2024} ont montré que l'explication ne résidait ni dans la résistance syndicale ni dans un défaut de la technologie, mais dans le fait que la commutation interagissait avec la facturation, les réclamations, les renseignements et les urgences~--- une tâche \emph{intégrale} au sens de \textcite{ulrich_role_1995}, dont la modification impliquait la reconfiguration de l'ensemble du système Bell. Le chapitre~2 devra déterminer si ce type de résistance structurelle se retrouve dans les vagues d'automatisation ultérieures, et si la transgression de la frontière du savoir tacite par les modèles de langage~--- le \og paradoxe de Polanyi\fg{} de \textcite{autor_polanyi_2014}~--- modifie ou non cette donnée.

Le chapitre a enfin fait apparaître, sans le thématiser, un pattern biographique récurrent~: des individus formés dans un paradigme technique transfèrent leurs compétences vers le paradigme suivant par leur mobilité. Edison, opérateur de télégraphe, transpose la maîtrise des circuits électriques vers l'électrification (§\,\ref{sec:electrification}). Siemens, officier d'artillerie, transpose la connaissance des signaux électromagnétiques vers le télégraphe prussien (§\,\ref{sec:telegraph}). Shannon connecte l'algèbre de Boole au circuit de commutation (§\,\ref{sec:electrification}). La récurrence suggère que la diffusion des paradigmes techniques passe en partie par la mobilité des personnes qui en maîtrisent les composants~--- mais c'est une observation biographique, pas un résultat que le chapitre ait pu formaliser.

Ces quatre observations ne se substituent pas à la régularité centrale. Elles la complètent en montrant que le rapport entre les techniques de l'information et le travail est multidirectionnel~: la technique restructure le travail, mais elle crée aussi du travail de conception qu'elle ne remplace pas~; le travail, sous forme de compétences, sélectionne la technique~; et le savoir tacite porté par certaines tâches résiste à l'automatisation pendant des décennies. La littérature contemporaine sur l'automatisation raisonne presque exclusivement dans la première direction. Le chapitre~2, en mobilisant le cadre \emph{task-based} de \textcite{autor_skills_2003} et le modèle de \textcite{acemoglu_race_2018}, devra déterminer lesquelles de ces directions sont des phénomènes propres au dix-neuvième siècle et lesquelles sont des traits structurels du rapport ICT-travail qui persistent~--- ou se transforment~--- à l'ère de l'intelligence artificielle.

La troisième régularité concerne le rapport entre la technique de l'information et la demande qui l'appelle. La bidirectionnalité posée dans le préambule méthodologique~--- la technique transforme l'économie, mais l'économie appelle la technique~--- n'est pas symétrique à chaque période. Au dix-neuvième siècle, c'est surtout la demande qui tire~: la coordination militaire pour le télégraphe Chappe, la sécurité ferroviaire pour le \emph{block system}, le commerce pour le câble transatlantique, le recensement pour la tabulatrice. Au vingtième siècle, le basculement s'opère~: le modèle IBM (location et cartes exclusives) crée une dépendance au fournisseur~; Vail construit le réseau longue distance avant que la demande n'existe~; la Navy finance le wireless et le computing avant toute demande commerciale. Au vingt et unième siècle, \textcite{roussilhe_phase_2025} montre que l'offre ne se contente plus de précéder la demande~: elle doit la produire activement et continument, parce que la structure de coûts de l'industrie des semi-conducteurs exige que la puissance de calcul mise sur le marché chaque année soit \og consommée\fg{} pour maintenir les économies d'échelle. La flèche causale s'est inversée~--- et cette inversion a une conséquence énergétique directe que le chapitre a documentée mais ne peut pas modéliser.

L'anomalie est l'énergie. De §\,\ref{sec:era0} à §\,\ref{sec:convergence}, la dimension 0-TER apparaît dans chaque cas mais n'y constitue jamais un résultat mesurable. Le ratio entre la puissance du circuit informationnel et la puissance du système qu'il coordonne est de l'ordre de 1\,:\,2\,000 dans le cas ferroviaire. L'électricité représente 1,5\,\% du coût d'exploitation des tabulatrices au \emph{Nautical Almanac}. Le signal télégraphique tire son énergie de piles dont le coût n'apparaît nulle part dans les comptes de Western Union. La matérialité extractive~--- la gutta-percha, le cuivre, les poteaux~--- est réelle mais n'est documentée que rétrospectivement \parencite{tully_victorian_2009, fitzmaurice_material_2023}. L'énergie traverse le chapitre comme un fait sans observateur.


\subsection{Deux postulats qui cessent de tenir}\label{subsec:postulats}

Cette anomalie n'est pas un accident d'observation. Elle est produite par deux postulats implicites que le matériau historique n'a jamais contestés~--- jusqu'au cas contemporain.

Le premier postulat est que le coût marginal de l'information tend vers zéro. Le télégramme coûte à envoyer, mais l'information qu'il porte~--- le cours du coton à Liverpool~--- se diffuse gratuitement une fois connue. La calculatrice mécanique a un coût fixe d'achat, mais chaque addition supplémentaire ne coûte rien. Le logiciel Windows a un coût de développement, mais chaque PC supplémentaire qui l'exécute ne coûte rien à Microsoft en puissance de calcul~: il n'y a pas de coût marginal pour l'éditeur de logiciel \parencite{roussilhe_computing_2025}. Ce postulat tient de §\,\ref{sec:telegraph} à §\,\ref{sec:convergence}. Il cesse de tenir avec l'IA générative~: chaque inférence a un coût marginal positif, directement proportionnel à la consommation électrique par token. Pour la première fois dans l'histoire des techniques de l'information traversée par ce chapitre, l'énergie est un coût marginal de production du service informationnel.

Le second postulat est que l'efficacité par unité fonctionnelle s'améliore continument. Le coût par opération arithmétique baisse d'environ trois pour cent par an de 1850 à 1945 \parencite{nordhaus_two_2007}. L'efficacité énergétique par calcul double tous les dix-huit mois entre 1946 et 2000 \parencite{koomey_implications_2011}. Le \emph{scaling} de Dennard (§\,\ref{sec:electrification}) maintient la densité de puissance constante pendant que la miniaturisation progresse. Ce postulat cesse de tenir sur deux fronts~: le \emph{scaling} de Dennard atteint sa limite physique vers~2005, et l'empreinte environnementale par centimètre carré de fabrication de circuits intégrés, qui baissait depuis les années~1990, remonte à partir du n\oe{}ud de 7~nm \parencite{pirson_environmental_2022}. L'amélioration relative a cessé~; les impacts absolus croissent.

Ce qui est remarquable n'est pas que ces deux postulats cassent~--- toute trajectoire technique finit par rencontrer ses limites~--- mais qu'ils cassent simultanément. Tant que le coût marginal était nul, la question de l'efficacité énergétique ne se posait pas à l'utilisateur final~: Microsoft pouvait \og gaspiller les transistors\fg{} \parencite{gilder_coming_1995} sans que cela ne lui coûte quoi que ce soit. Tant que l'efficacité par composant s'améliorait, la consommation agrégée pouvait croître sans que le couplage ne devienne visible~: le \emph{scaling} de Dennard absorbait la croissance. Quand les deux cessent de tenir en même temps, le couplage entre information et énergie, invisible pendant deux siècles, apparaît. Le fait que personne n'ait croisé les séries de \textcite{nordhaus_two_2007} (coût par calcul) et de \textcite{koomey_implications_2011} (efficacité énergétique par calcul) n'est pas un oubli~; c'est le signe que les catégories d'analyse en vigueur ne contenaient pas la case.


\subsection{L'occultation comme pattern structurel}\label{subsec:occultation}

\textcite{ensmenger_environmental_2018} est le premier à constituer cette absence en objet d'étude. Dans son programme pour une histoire environnementale du \emph{computing}, il identifie trois axes~: les liens entre production de puissance de calcul et processus traditionnels d'extraction et de pollution~; la production de discours d'invisibilisation (la \og dématérialisation\fg{}, le \og cloud\fg{})~; l'asymétrie géographique entre zones d'extraction et zones de consommation. Mais ce qui fait la force de sa proposition, c'est que le deuxième axe~--- le discours~--- n'est pas un bruit à traverser pour atteindre les \og vrais faits\fg{} du premier axe. Il est constitutif~: c'est parce que le \emph{computing} est perçu comme immatériel que ses coûts matériels ne sont pas comptés, et c'est parce qu'ils ne sont pas comptés qu'ils peuvent croître sans rencontrer de résistance.

Le chapitre qui s'achève a documenté ce pattern bien au-delà du \emph{computing}. \textcite{belloc_telegraphie_1894} constate que la télégraphie est \og surtout remarquable par ses conséquences économiques\fg{} sans mentionner ses conséquences matérielles. D'Alembert juge le calcul \og plus laborieux que profond\fg{}, effaçant le travail du calcul comme objet digne d'attention. Gray ne voit pas la valeur commerciale du téléphone parce que ni la métrique du paradigme télégraphique~--- messages par fil, coût par mot~---, ni le modèle d'affaires de son employeur (Western Union, monopole télégraphique en place), ne contiennent la case (§\,\ref{sec:telephone}). \textcite{nordhaus_two_2007} mesure le coût du calcul sur deux siècles sans inclure l'électricité. Ce ne sont pas des erreurs~; ce sont des effets de cadrage. Chaque époque produit une grille qui capture certaines dimensions~--- le monopole naturel pour les régulateurs du télégraphe, la productivité du travail pour les historiens du calcul, l'intégration des marchés pour les économistes du commerce~--- et qui, structurellement, ne contient pas la case énergie-matière. L'occultation n'est pas produite par l'ignorance mais par les catégories d'analyse elles-mêmes.

\textcite{roussilhe_phase_2025} montre que ce pattern se renforce dans la période contemporaine. Les \og lois\fg{} qui organisent l'industrie des semi-conducteurs~--- la loi de Moore, les \emph{scaling laws} d'OpenAI, la \og loi de Huang\fg{}~--- ne sont pas des descriptions physiques~; ce sont des dispositifs de coordination sectorielle qui alignent les anticipations des fabricants, des éditeurs de logiciel et des investisseurs sur une feuille de route commune. Leur fonction est performative~: tant que l'industrie y croit, elle investit les sommes nécessaires pour les réaliser. Comme le formule Roussilhe, la loi de Moore est un \og baromètre de la croyance de l'industrie en elle-même\fg{}. La régulation du spectre électromagnétique (§\,\ref{sec:wireless}) procédait de la même structure~: \textcite{hazlett_rationality_1990} a montré que le Radio Act de~1927 n'était pas une correction de défaillance de marché mais un choix de distribution de rentes, organisé comme une réponse technique à un problème qui était en réalité politique. Les \og lois\fg{} technologiques et les régulations qui les accompagnent ne décrivent pas le monde~; elles le produisent~--- et ce qu'elles produisent inclut systématiquement l'invisibilité de la dimension énergétique.


\subsection{Ce que le chapitre a révélé et ce qu'il ne peut pas faire}\label{subsec:limites}

Le tableau~\ref{tab:ch1-questions} rassemble, pour chacune des dix dimensions et des deux forces transversales, ce que la traversée historique a permis d'établir et les questions qu'elle ouvre pour le chapitre~2. Trois niveaux sont distingués~: \textbf{★}~fait documenté indépendamment, \textbf{+}~argument explicite mais fondation incomplète, \textbf{?}~piste ouverte.

\begin{table}[p]
\caption{Les dix dimensions après la traversée historique~: acquis et questions ouvertes pour le chapitre~2.}\label{tab:ch1-questions}
\centering\scriptsize
\renewcommand{\arraystretch}{1.18}
\begin{tabularx}{\textwidth}{@{}cp{1.8cm}Xp{5.2cm}@{}}
\toprule
 & \textbf{Dimension} & \textbf{Ce que le ch.~1 établit} & \textbf{Questions pour le ch.~2} \\
\midrule
1 & Propriétés du bien & Le circuit télégraphique est rival~; l'information transportée ne l'est pas. La codifiabilité filtre les effets du commerce (Juhász~★). & ★ La codifiabilité module-t-elle \emph{toutes} les dimensions~? \newline ? Si l'IA rend codifiable presque tout, est-elle une ICT qui élimine le filtre~? \\[4pt]
2 & Rendements et concentration & Sept trajectoires institutionnelles convergent vers le monopole, par cinq mécanismes différents (★). Le résultat est stable~; le mécanisme varie. & ★ Qu'est-ce qui détermine le choix du mécanisme~? \newline + Trait des réseaux de communication, ou des infrastructures lourdes en général~? \\[4pt]
3 & Nature du capital & Le gradient isolé/intégré détermine irréversibilité et structure de marché (★ calculatrice vs tabulatrice). La forme du capital, pas la technologie, discrimine. & ★ Le mode d'appropriation (vente vs location+consommable) est-il le déterminant endogène de la structure de marché~? \newline + L'irréversibilité est-elle financière (Dixit-Pindyck) ou systémique (Hughes-Unruh)~? \\[4pt]
4 & Travail et nature du travail & La mécanisation réorganise le travail au lieu de le remplacer~--- quatre cas sur 150~ans (★). Conditionnalité~: la demande agrégée est nécessaire pour que la compensation opère (★ Feigenbaum-Gross). & ★ La conditionnalité à la demande agrégée tient-elle quand l'automatisation touche les tâches cognitives non-routinières~? \newline + Le travail de conception (\emph{analytical intelligence}) est-il un coût transitoire ou permanent~? \newline + Le stock de compétences conditionne-t-il la forme de la technologie adoptée~? \newline + Le savoir tacite des tâches intégrales constitue-t-il une résistance mesurable~? \\[4pt]
5 & Demande et préférences & La demande de calcul est dérivée et inélastique au prix de la machine (★). Le modèle «~service gratuit financé par la publicité~» naît avec la radio, pas avec Internet. & ? Effet de rebond par les tarifs~? \newline + La dissociation prix perçu / coût de production est-elle un trait constitutif de l'économie numérique~? \\[4pt]
6 & Signal de prix & L'information synchrone \emph{crée} le marché~--- Steinwender ★, quatre réplications. Effet conditionnel~: sans transport physique, pas d'arbitrage (Andrabi ★). & ★ La codifiabilité est-elle le seul filtre, ou la fonction ICT (TRANSMIT vs COMPUTE) en est-elle un autre~? \newline + Le seuil de McCallum (l'information comme condition de fonctionnement, pas simple facilitateur) se généralise-t-il~? \\[4pt]
7 & Coûts de transaction & La hiérarchie informationnelle émerge quand le prix ne coordonne pas (★ McCallum, six formes). La commutation est une tâche intégrale~: 60~ans pour automatiser (★ Feigenbaum). & + La frontière firme-marché est-elle déterminée par les techniques d'information disponibles~? \newline + La devise de Bolle (\og d'abord organiser, ensuite mécaniser~\fg{}) a-t-elle un équivalent contemporain~? \\[4pt]
8 & Financement & Pattern investissement $\rightarrow$ faillite $\rightarrow$ recapitalisation documenté pour le câble (★), la railway mania, le dot-com. L'État crée le marché avant le marché (★ Census, Navy, Convention). & + Le cycle Perez est-il propre aux réseaux de communication ou commun aux infrastructures lourdes~? \newline ? Le financement par la promesse stratégique (pas la rentabilité) est-il propre aux technologies de souveraineté~? \\[4pt]
9 & Géographie & La géographie du réseau reflète les rapports de puissance (★ câbles coloniaux, concessions). Domination britannique sur les matériaux du câble (★). & + La contrainte matérielle (cuivre, gutta-percha) préfigure-t-elle les terres rares et le lithium~? \newline ? La localisation des DC suit-elle l'énergie, le pouvoir, ou les rendements croissants~? \\[4pt]
10 & Innovation & Le système métrique du paradigme en place rend invisible la valeur des innovations d'un autre registre~--- Gray, Marconi (★). Verrouillage de standard sur des millénaires (sexagésimal, VN, carte 80~col.). & ★ L'occultation par le système métrique se reproduit-elle à chaque vague~? \newline + Le facteur limitant de l'innovation migre-t-il du talent vers l'énergie dans la période contemporaine~? \\[4pt]
\midrule
0-TER & Énergie-matière & L'énergie traverse le chapitre comme un fait sans observateur. Ratio 1:2\,000 (circuit informationnel / système coordonné). Deux postulats (coût marginal nul, efficacité croissante) cassent simultanément. & ★ Le ratio évolue-t-il~? À quel seuil le couplage devient-il visible~? \newline + L'occultation est-elle produite par les catégories d'analyse ou par les ordres de grandeur~? \\[4pt]
0-GOV & Souveraineté & La forme du réseau reflète l'intérêt de l'acteur qui le finance~--- de Chappe aux PPAs nucléaires (★). & + La souveraineté numérique prolonge-t-elle le même pattern~? \\
\bottomrule
\end{tabularx}
\renewcommand{\arraystretch}{1.0}
\end{table}

\begin{table}[p]
\caption{Matrice dimensions $\times$ périodes~: trajectoire de l'occultation énergétique.}\label{tab:synoptique-ch1}
\centering\scriptsize
\renewcommand{\arraystretch}{1.12}
\setlength{\tabcolsep}{3pt}
\begin{tabular}{@{}llcccccccccccc@{}}
\toprule
 & \textbf{Période} & \rotatebox{70}{\textbf{D1}} & \rotatebox{70}{\textbf{D2}} & \rotatebox{70}{\textbf{D3}} & \rotatebox{70}{\textbf{D4}} & \rotatebox{70}{\textbf{D5}} & \rotatebox{70}{\textbf{D6}} & \rotatebox{70}{\textbf{D7}} & \rotatebox{70}{\textbf{D8}} & \rotatebox{70}{\textbf{D9}} & \rotatebox{70}{\textbf{D10}} & \rotatebox{70}{\textbf{0-TER}} & \rotatebox{70}{\textbf{0-GOV}} \\
\midrule
§\ref{sec:era0} & Avant 1794 & $\bullet$ & --- & $\circ$ & $\bullet$ & $\circ$ & --- & $\circ$ & --- & $\circ$ & $\bullet$ & --- & $\circ$ \\
§\ref{sec:telegraph} & 1794--1880 & $\bullet$ & $\bullet$ & $\circ$ & $\bullet$ & $\bullet$ & $\bullet$ & $\circ$ & $\bullet$ & $\bullet$ & $\bullet$ & --- & $\bullet$ \\
§\ref{sec:calculating} & 1820--1914 & $\circ$ & $\circ$ & $\bullet$ & $\bullet$ & $\circ$ & --- & --- & $\circ$ & --- & $\circ$ & --- & --- \\
§\ref{sec:mechanography} & 1890--1940 & $\circ$ & $\bullet$ & $\bullet$ & $\bullet$ & $\circ$ & --- & $\circ$ & $\bullet$ & --- & $\bullet$ & --- & $\circ$ \\
§\ref{sec:telephone} & 1876--1920 & $\circ$ & $\bullet$ & $\circ$ & $\bullet$ & $\bullet$ & $\circ$ & $\bullet$ & $\circ$ & $\circ$ & $\bullet$ & --- & $\bullet$ \\
§\ref{sec:wireless} & 1895--1930 & $\circ$ & $\bullet$ & --- & $\circ$ & $\bullet$ & $\bullet$ & --- & $\circ$ & --- & $\bullet$ & --- & $\bullet$ \\
§\ref{sec:electrification} & 1880--1930 & --- & $\bullet$ & $\bullet$ & $\circ$ & $\bullet$ & --- & $\bullet$ & $\bullet$ & --- & $\bullet$ & $\nearrow$ & $\bullet$ \\
§\ref{sec:computing} & 1945--1975 & --- & $\bullet$ & $\bullet$ & $\circ$ & $\bullet$ & --- & --- & $\bullet$ & --- & $\bullet$ & $\circ$ & $\bullet$ \\
§\ref{sec:convergence} & 1975--2007 & $\circ$ & $\bullet$ & $\bullet$ & $\circ$ & $\bullet$ & $\bullet$ & $\bullet$ & $\bullet$ & $\circ$ & $\bullet$ & $\bullet$ & $\circ$ \\
§\ref{sec:contemporain} & 2007-- & $\circ$ & $\bullet$ & $\bullet$ & $\nearrow$ & $\bullet$ & $\circ$ & $\circ$ & $\bullet$ & $\bullet$ & $\bullet$ & $\bullet$ & $\bullet$ \\
\midrule
\multicolumn{2}{@{}l}{\textbf{Synthèse}} & \multicolumn{10}{c}{\emph{Dimensions stables~: D2, D5, D8, D10}} & $\boldsymbol{\nearrow\nearrow}$ & \\
\bottomrule
\end{tabular}
\renewcommand{\arraystretch}{1.0}
\par\medskip
\footnotesize
\textbf{Légende}~: $\bullet$~=~dimension centrale dans cette période~; $\circ$~=~dimension présente mais secondaire~; ---~=~dimension non pertinente ou non documentée~; $\nearrow$~=~dimension en émergence.
\par\smallskip
\textbf{Lecture de la colonne 0-TER}~: La dimension énergie-matière passe de «~---~» (non pertinente) dans les sections §\ref{sec:era0}--§\ref{sec:wireless}, à «~$\nearrow$~» (émergente) en §\ref{sec:electrification} avec l'électrification du substrat, à «~$\circ$~» (présente mais secondaire) en §\ref{sec:computing} où l'énergie est une contrainte d'ingénierie mais pas un coût économique visible, à «~$\bullet$~» (centrale) à partir de §\ref{sec:convergence}. Cette trajectoire illustre l'\emph{occultation progressive puis la réémergence} de la contrainte énergétique dans l'analyse des techniques de l'information.
\end{table}

Le chapitre a révélé une structure. Dix dimensions économiques (D1--D10) et deux forces transversales (0-TER, 0-GOV) ont émergé de la confrontation séquentielle des cas. Des hypothèses abductives ont été formulées~: le mode d'appropriation du capital, et non la technologie, détermine la structure de marché~; la mécanisation réorganise le travail au lieu de le remplacer~; la technique ne se déploie pas sans un véhicule de financement adapté à la structure des coûts et des bénéfices. Un gradient d'occultation de la dimension énergie-matière a été identifié, produit non par l'ignorance mais par les catégories d'analyse en vigueur à chaque époque. Et deux postulats implicites~--- le coût marginal nul de l'information, l'amélioration continue de l'efficacité~--- ont été nommés parce que le cas contemporain les a fait voler en éclats.

\dimmark{D10}{*}%
Un parallèle disciplinaire récent mérite d'être nommé ici par souci de positionnement bibliographique. \textcite{johnstone_deep_2026}, dans la revue \emph{Technological Forecasting and Social Change}, ont mobilisé le cadre des \emph{Deep Transitions} développé par \textcite{schot_deep_2018} pour analyser l'évolution du «~méta-régime numérique~» sur la même période 1937--2024 que celle parcourue par le présent chapitre, et la séquence de phases qu'ils en dérivent~--- gestation, irruption, frenzy, turning point, synergy, maturity~--- recouvre une partie des mêmes faits empiriques. Les deux approches ne se substituent pas l'une à l'autre. Leur cadre est évolutionniste et historique au sens de la sociologie des transitions sociotechniques (STS), construit autour des notions de règles dominantes, de règles précurseurs et de méta-régime~; le présent chapitre est dans le registre de l'induction analytique au sens de \textcite{znaniecki_method_1934}, de la causalité par cas nécessaire ou suffisante au sens de \textcite{trannoy_causalite_2025}, et de l'occultation matérielle au sens de \textcite{ensmenger_environmental_2018}. Le chapitre mobilise ponctuellement le récit de Johnstone et collègues là où il éclaire des points précis de l'observation par cas (§\ref{sec:electrification}, §\ref{sec:computing}, §\ref{sec:convergence}, §\ref{sec:contemporain}), sans transposer leur cadre théorique~: l'objet du présent chapitre est de produire un matériau brut sur lequel le cadre disciplinaire de la suite de la thèse (économie de la croissance et des techniques génériques, comptabilité énergétique, modélisation IMACLIM) puisse opérer, non d'opérationnaliser une approche STS-évolutionniste sur le même matériau. La distinction des registres importe aussi pour éviter une redondance interprétative~: ce que \textcite{johnstone_deep_2026} décrivent comme \emph{turning point} (2001--2008) apparaît dans le présent chapitre comme une rupture de couplage énergétique permanent (§\ref{subsec:turning_point})~; ce qu'ils nomment \emph{frenzy} (1996--2000) apparaît comme une résolution différée du modèle économique de la commercialisation d'Internet (§\ref{subsec:dotcom})~; ce qu'ils caractérisent comme phase de \emph{synergy} contemporaine (2008--) est décomposé dans le présent chapitre en cloud, store, GenAI, blockchain, IoT et régulation fragmentée (§\ref{sec:contemporain}). Les éclairages se complètent sans se recouvrir.

Mais le chapitre ne peut pas aller plus loin. L'induction identifie des régularités~; elle ne les explique pas. L'abduction génère des hypothèses~; elle ne les teste pas. Le constat que sept pays convergent vers le monopole télégraphique n'est pas un mécanisme~; le constat que l'occultation de 0-TER se reproduit à chaque cas n'est pas une cause. Pour formaliser les mécanismes qui produisent ces régularités~--- et pour déterminer lesquels sont propres aux techniques de l'information et lesquels sont communs à toutes les infrastructures lourdes~---, il faut changer d'instrument. L'analyse économique systématique du chapitre~2 prend les dix dimensions que ce chapitre a fait émerger et les déploie comme des canaux causaux, chacun avec ses mécanismes, ses conditions de validité et ses équations. L'analyse énergie-matière du chapitre~3 prend la dimension transversale 0-TER et la quantifie fonction par fonction~--- transmission, stockage, calcul~--- pour mesurer ce que l'occultation cachait. Le chapitre qui s'achève a rendu les deux suivants nécessaires~; il ne pouvait pas les remplacer.


% =============================================================================
% NOTE MÉTHODOLOGIQUE — Rapport Ch.1 / évaluation des IAMs
% Ajoutée le 28 avril 2026, à intégrer dans la rédaction finale.
% =============================================================================
%
% Ce que le Ch.1 produit pour le dispositif expérimental (Ch.4-5) :
%
% 1. Des FAITS STYLISÉS au sens de Wilson et al. (2021, Climatic Change).
%    La co-saturation transport-communication (Grübler), la périodisation par
%    le rapport improvement/énergie, la diversité des configurations push/pull,
%    le gradient isolé/intégré du capital (D3) sont des patterns historiques
%    généralisables que le modèle IMACLIM doit reproduire. Si une trajectoire
%    projetée viole un fait stylisé du Ch.1, c'est un signal d'erreur dans le
%    modèle, pas une prédiction.
%
% 2. Des CONDITIONS D'ÉMERGENCE ET DE NON-ÉMERGENCE.
%    Chappe émerge par commande étatique-militaire, mais le télégraphe
%    commercial n'émerge que quand le rail crée la demande. Babbage existe
%    techniquement mais l'industrie n'émerge pas (absence de demande et de
%    financement). Le wireless émerge comme outil militaire, mais la
%    radiodiffusion n'émerge que quand le récepteur domestique et le modèle
%    publicitaire apparaissent. Ces cas de non-émergence alimentent les
%    scénarios du Ch.4 (conditions sous lesquelles les trajectoires
%    alternatives n'atteignent pas l'échelle).
%
% 3. Des TESTS DE COHÉRENCE pour le Ch.5.
%    Chaque fait stylisé est une contrainte que les résultats du modèle
%    doivent respecter. Si le module numérique projette un découplage
%    compute/énergie sans précédent historique, ou une concentration de
%    marché incompatible avec le gradient D3, le résultat est suspect.
%
